% ======================================================================
%  Remainder Terms of the Riemann Zeta Function  --  rewrite (v7)
%  Result-first restructure. Minimal, BasicTeX-compatible preamble so
%  LaTeX Workshop previews it with no extra package installs.
%  Add packages only as a section needs one.
% ======================================================================
\documentclass[11pt]{article}

\usepackage[margin=1in]{geometry}
\usepackage{amsmath,amssymb,amsthm}
\usepackage{graphicx}
\usepackage{float}
\usepackage{booktabs}
\usepackage{xcolor}
\usepackage[colorlinks=true,linkcolor=blue,citecolor=blue,urlcolor=blue]{hyperref}

% ---- Lean source listings (Appendix A) ------------------------------------
\usepackage{listings}
\definecolor{leanblue}{RGB}{0,0,180}
\definecolor{leangreen}{RGB}{0,128,0}
\lstdefinelanguage{Lean}{
  morekeywords={
    import, open, variable, variables, section, end, axiom,
    def, theorem, lemma, example, abbrev, set, unfold, conv_lhs,
    noncomputable, namespace, where, by, intro, intros,
    apply, fun, match, with, if, then, else, do,
    let, in, exact, rfl, simp, dsimp, rw, have, from, ring, field_simp,
    push_cast, ring_nf, linear_combination, congr, exact_mod_cast, simpa, only
  },
  sensitive=true,
  morecomment=[l]{--},
  morestring=[b]"
}
\lstdefinestyle{leanstyle}{
  language=Lean,
  basicstyle=\ttfamily\footnotesize,
  keywordstyle=\color{leanblue}\bfseries,
  commentstyle=\color{leangreen},
  stringstyle=\color{purple},
  identifierstyle=\color{black},
  columns=fullflexible,
  keepspaces=true,
  showstringspaces=false,
  breaklines=true,
  breakindent=2em,
  numbers=left,
  numberstyle=\tiny\color{gray},
  stepnumber=1,
  literate={ℝ}{{$\mathbb{R}$}}1
           {ℂ}{{$\mathbb{C}$}}1
           {ℤ}{{$\mathbb{Z}$}}1
           {ω}{{$\omega$}}1
           {φ}{{$\varphi$}}1
           {ψ}{{$\psi$}}1
           {χ}{{$\chi$}}1
           {σ}{{$\sigma$}}1
           {θ}{{$\theta$}}1
           {→}{{$\to$}}1
           {←}{{$\leftarrow$}}1
           {‖}{{$\|$}}1
           {≠}{{$\neq$}}1
           {¬}{{$\neg$}}1
           {∃}{{$\exists$}}1
           {∀}{{$\forall$}}1
           {⌈}{{$\lceil$}}1
           {⌉}{{$\rceil$}}1
           {₀}{{$_0$}}1
           {₁}{{$_1$}}1
           {₂}{{$_2$}}1
           {∓}{{$\mp$}}1
           {·}{{$\cdot$}}1
           {ℕ}{{$\mathbb{N}$}}1
           {Σ}{{$\Sigma$}}1
           {∑}{{$\sum$}}1
           {∈}{{$\in$}}1
           {ζ}{{$\zeta$}}1
           {π}{{$\pi$}}1
           {⊢}{{$\vdash$}}1
}

% New, regenerated figures live in ./figures/. The original paper folder is
% kept as a fallback path only for reference while we rebuild.
\graphicspath{{figures/}{../Riemann_zeta_function_remainders_and_other_observations_6/}}

% Theorem environments (match the original).
\newtheorem{theorem}{Theorem}[section]
\newtheorem{proposition}[theorem]{Proposition}
\newtheorem{conjecture}[theorem]{Conjecture}
\newtheorem{lemma}[theorem]{Lemma}
\newtheorem{corollary}[theorem]{Corollary}
\theoremstyle{definition}
\newtheorem{definition}[theorem]{Definition}
\newtheorem{remark}[theorem]{Remark}

% ----------------------------------------------------------------------
% Link matrix (2-D rotation + translation, 3x3 homogeneous). Left-aligned
% columns with a phantom minus so cos/sin line up and the real minus sign
% overhangs to the left. Usage: \linkmat{<angle>}{<length>}
% ----------------------------------------------------------------------
% Link matrix with the step written as a fraction, so the factor grows in
% height rather than width: third column is (cos t)/#2, (sin t)/#2.
\newcommand{\linkmatd}[2]{%
\left(\begin{array}{@{}l@{\,}l@{\,}l@{}}
\cos #1 & -\sin #1           & \frac{\cos #1}{#2}\\[1pt]
\sin #1 & \phantom{-}\cos #1 & \frac{\sin #1}{#2}\\[1pt]
\multicolumn{1}{@{}c@{\,}}{0} & \multicolumn{1}{c@{\,}}{0} & \multicolumn{1}{c@{}}{1}
\end{array}\right)}
% Chain-2 link matrix: step $-|\chi|/#2$, again written as a fraction.
\newcommand{\linkmatx}[2]{%
\left(\begin{array}{@{}l@{\,}l@{\,}l@{}}
\cos #1 & -\sin #1           & -\frac{|\chi|\cos #1}{#2}\\[1pt]
\sin #1 & \phantom{-}\cos #1 & -\frac{|\chi|\sin #1}{#2}\\[1pt]
\multicolumn{1}{@{}c@{\,}}{0} & \multicolumn{1}{c@{\,}}{0} & \multicolumn{1}{c@{}}{1}
\end{array}\right)}
\newcommand{\linkmat}[2]{%
\left(\begin{array}{@{}l@{\ \,}l@{\ \,}l@{}}
\cos #1 & -\sin #1           & #2\cos #1\\
\sin #1 & \phantom{-}\cos #1 & #2\sin #1\\
\multicolumn{1}{@{}c@{\ \,}}{0} & \multicolumn{1}{c@{\ \,}}{0} & \multicolumn{1}{c@{}}{1}
\end{array}\right)}

\title{\textbf{Remainder Terms of the Riemann Zeta Function}}
\author{Paul Stahura\\[0.3em]\small\texttt{paul+zeta@stahura.net}}
\date{\today}

\begin{document}
\maketitle

\begin{center}
{\color{red}\textbf{DRAFT, rewrite in progress (v7).}}
\end{center}

\begin{abstract}
Starting from the Riemann--Siegel decomposition of $\zeta(s)$ given in Siegel's
1932 paper, we introduce a single change of variable, $t=I(T)$, that replaces
Siegel's coupled pair ``imaginary part $t$ / summation index $m$'' with one real
index $T$. Under this reparameterization the formula is \emph{identically}
Riemann--Siegel. We then split Siegel's single remainder integral $R$ into two
exact pieces, $R_{1ps}$ and $R_{2ps}$, and prove that
$R = R_{1ps}+R_{2ps}$, so that
\[
\zeta \;=\; \Sigma_{1} + R_{1ps} + \Sigma_{2} + R_{2ps}.
\]
Our central observation is that each of these remainders is nothing more than
\emph{one additional, fractional partial summand} appended to its Dirichlet sum:
\[
\zeta(s)=\sum_{n=1}^{m} n^{-s} + \hat d_1\,(m+1)^{-s}
        + \chi(s)\sum_{n=1}^{m} n^{s-1} + \hat d_2\,\chi(s)\,(m+1)^{s-1},
\]
with $\hat d_1,\hat d_2$ real numbers (always positive on the critical line),
the fractions of those two
summands that are used. As a corollary, when $\sigma=\tfrac12$ one
has $d_1=d_2$ (equivalently $\hat d_1=\hat d_2$); the fact is formally
verified in Lean. The decomposition was
discovered through experimental mathematics using a spiral visualization of the
partial sums, described later in the paper.
\end{abstract}

\tableofcontents
\newpage

% ======================================================================
\section{Introduction}
% ======================================================================

This paper makes three claims, in increasing order of importance.

\begin{enumerate}
\item \textbf{A reparameterization.} We introduce a change of variable
$t=I(T)$ that collapses Siegel's two roles for $t$ (the imaginary part of the
input and the index of summation) into a single real ``index'' $T$. With this
substitution the Riemann--Siegel formula is unchanged in value; it is the
\emph{same} formula, viewed through a different variable, $T$, instead of $t$.

\item \textbf{A decomposition of the remainder.} Siegel's formula writes
$\zeta = \Sigma_1 + \Sigma_2 + R$, where $\Sigma_1,\Sigma_2$ are finite Dirichlet
sums and $R$ is a remainder integral. We define two exact remainder terms,
$R_{1ps}$ and $R_{2ps}$, and prove
\begin{equation}\label{eq:main}
R \;=\; R_{1ps} + R_{2ps},
\qquad\text{hence}\qquad
\zeta \;=\; \Sigma_{1} + R_{1ps} + \Sigma_{2} + R_{2ps}.
\end{equation}
Equation~\eqref{eq:main} is the main result of the paper.

\item \textbf{An interpretation.} The two new remainders are not opaque
integrals: each is exactly \emph{one more summand} of its Dirichlet sum,
namely the $(m+1)$st term, scaled by a real ``fractional'' weight
$\hat d_1$ or $\hat d_2$, always positive on the critical line and positive
elsewhere outside narrow windows around the poles of
\S\ref{sec:pole-locations}, as proved in \S\ref{sec:d1-positive}
(\S\ref{sec:summand} distinguishes these
dimensionless \emph{fractional amounts} from the closely related absolute
lengths $d_1,d_2$). In this sense the remainder is simply the partial sum
carried one fractional step further.
\end{enumerate}

A corollary of the interpretation is that on the critical line $\sigma=\tfrac12$
the two weights coincide, $d_1=d_2$; we return to this at the end, where we also
record that the fact has been formally verified in Lean. We emphasize that this
critical-line symmetry, while pleasing, is \emph{not} the focus of the paper.
The decomposition~\eqref{eq:main} and its ``one more partial summand'' reading
are.

\paragraph{Origin of the result.} The decomposition was not found analytically.
It emerged from experimental mathematics: plotting the partial sums of $\zeta$ as
a Euler-like spiral in the complex plane and studying its geometry with a
visualization tool we have developed (and published) over several years. The
spiral itself has precedent. Erickson~\cite{Erickson2005} plots the partial
sums, identifies the shapes as Cornu spirals, labels their centers $C_n$, finds
$C_0=\zeta(s)$, and traces the symmetry he sees among the $C_n$ back to the
approximate functional equation. Nickel~\cite{Nickel2013,Nickel2015} studies the
same Argand
diagram of the steps $n^{-s}$ geometrically, pairing each early step with a
conjugate region of linked Euler (Cornu) spirals, recovering the functional
equation from that symmetry and locating zeros where two vectors point in
opposite directions with equal magnitudes, which happens exactly when
$\sigma=\tfrac12$. Kapitonets~\cite{Kapitonets2019} names the polyline of steps
the Riemann spiral and works with its signed radius of curvature, its reverse
points and its inflection points. That
geometric story (the spiral, a local coordinate frame, and the ``yin'' and
``yang'' curves that actually produce $R_{1ps}$ and $R_{2ps}$) is deferred to
\S\ref{sec:geometry} and \S\ref{sec:yinyang}, so that the reader sees the clean
algebraic result first and its experimental origin second.

\paragraph{Roadmap.} \S\ref{sec:siegel} recalls Siegel's 1932 decomposition.
\S\ref{sec:IT} introduces the mapping $I(T)$ and shows the reparameterized
formula is identical to Riemann--Siegel. \S\ref{sec:decomp} defines
$R_{1ps},R_{2ps}$ and proves~\eqref{eq:main}. \S\ref{sec:summand} establishes the
``one more partial summand'' interpretation, and \S\ref{sec:matrix-product}
recasts each partial sum together with its remainder as a product of
homogeneous transformation matrices, one per link, the remainder entering as
one extra fractional link.
\S\ref{sec:other-remainders} discusses other remainders: \S\ref{sec:kuznetsov}
contrasts the Kuznetsov / Siegel~$f$ remainders, and
\S\ref{sec:remainders-summary} names the three splits $R_{rs}$, $R_{ps}$, and
$R_{ak}$ side by side, and \S\ref{sec:remainder-scaling} records how the
remainders scale when $\{T\}$ is held fixed.
\S\ref{sec:d1-function} studies the weight $d_1$ as a function of the index:
its handoff jumps, two exactly continuous coordinates built from it, and
zeta-free closed-form approximations on and off the critical line.
\S\ref{sec:geometry} tells the experimental/geometric story
behind the result. \S\ref{sec:IT-functions} collects the $I(T)$ functions,
including the bisector-frame origin of $I(T)$ and variants for
$L$-functions. \S\ref{sec:yinyang} develops the yin and yang
curves from which the weights $d_1$ and $d_2$ were derived, and proves
their common limit curve as $T\to\infty$.
\S\ref{sec:consequences} collects further observations, including the
critical-line corollary for zeta zeros, the link-number connection, and the
crossing walks \(\Sigma_{1x}\) and \(\Sigma_{2x}\) of \S\ref{sec:sum-x}.
\S\ref{sec:prior-lit} sets the construction beside the prior literature on the
geometry of these spirals.
A glossary of the geometric vocabulary, with a short dictionary to names
used by other authors, follows as \S\ref{sec:glossary}.

% ======================================================================
\section{Siegel's 1932 decomposition}\label{sec:siegel}
% ======================================================================

We begin with the Riemann--Siegel decomposition of $\zeta(s)$ into two finite
Dirichlet sums plus a remainder integral, as it appears in Siegel's 1932
paper~\cite{Siegel1932} (his equation 13):
\begin{equation}\label{eq:siegel13}
\zeta(s)
=
\sum_{n=1}^{m} n^{-s}
\;+\;
\frac{(2\pi)^s}{2\,\Gamma(s)\,\cos\!\bigl(\tfrac{\pi s}{2}\bigr)}
\sum_{n=1}^{m} n^{s-1}
\;+\;
\frac{(2\pi)^s e^{\tfrac{\pi i s}{2}}}{\Gamma(s)\,\bigl(e^{2\pi i s}-1\bigr)}
\int_{C_2} g(x)\,\mathrm{d}x,
\end{equation}
where the path $C_2$ consists of two half-lines meeting at the point
$\eta\bigl(1-\tfrac{\epsilon}{2}\bigr)$ for a small fixed $\epsilon>0$, one
half-line passing through the saddle point $\eta=+\sqrt{t/2\pi}$ and the
other through $-(m+\tfrac12)$, and $g$ is defined below.

It is convenient to name the three pieces. Recall the completion factor
\begin{equation}
\chi(s) \;:=\; 2^{s}\pi^{s-1}\sin\!\Bigl(\tfrac{\pi s}{2}\Bigr)\Gamma(1-s)
\;=\;
\frac{(2\pi)^s}{2\,\Gamma(s)\,\cos\!\bigl(\tfrac{\pi s}{2}\bigr)},
\end{equation}
so that, writing $s=\sigma+it$ and taking the first $m$ terms,
\begin{align}
\Sigma_{1}(s) &= \sum_{n=1}^{m} n^{-s}, \\[2pt]
\Sigma_{2}(s) &= \frac{(2\pi)^s}{2\,\Gamma(s)\,\cos\!\bigl(\tfrac{\pi s}{2}\bigr)}
\sum_{n=1}^{m} n^{s-1}
= \chi(s)\sum_{n=1}^{m} n^{s-1}, \\[2pt]
R(s) &= \frac{(2\pi)^s e^{\tfrac{\pi i s}{2}}}{\Gamma(s)\,\bigl(e^{2\pi i s}-1\bigr)}
\int_{C_2} g(x)\,\mathrm{d}x,
\end{align}
where
\begin{equation}
g(x)=x^{s-1}\,\frac{e^{-2\pi i m x}}{e^{2\pi i x}-1}.
\end{equation}
With these names, \eqref{eq:siegel13} reads
\begin{equation}\label{eq:siegel-named}
\zeta(s) = \Sigma_{1} + \Sigma_{2} + R.
\end{equation}

In Siegel's formulation $m$ is the index of summation, and to minimize the
remainder he takes
\begin{equation}\label{eq:siegel-m}
m=\left\lfloor\sqrt{\tfrac{t}{2\pi}}\right\rfloor .
\end{equation}
Thus $t$ plays two roles at once: it is the imaginary part of the input $s$,
and, through \eqref{eq:siegel-m}, it also fixes the number of terms $m$. In
this way $m$ depends on $t$. The next section reverses the dependency.

% ======================================================================
\section{The same formula, reparameterized: The \texorpdfstring{$I(T)$}{I(T)} mapping}\label{sec:IT}
% ======================================================================

We introduce a change-in-variable mapping for $t$:
\begin{equation}\label{eq:IT}
t
\;=\;
\frac{\pi\,\bigl(2\,T + 1\bigr)}{\log\!\bigl(\tfrac{1}{T} + 1\bigr)}
\;=\; I(T),
\qquad
m \;=\; \lfloor T\rfloor .
\end{equation}
Here $T$ is a real number, while $m=\lfloor T\rfloor$ is an integer. So now
$t$ depends on $m$, the reverse of Siegel's \eqref{eq:siegel-m}. The single
function $I$ therefore determines both the imaginary part of the input and the
integer index of summation; for this reason we call $T$ ``the index'' even
though it is real. Both $t$ and $m$ are functions of $T$.

For the remainder of the paper we replace Siegel's $s=\sigma+it$ by
\begin{equation}
s \;=\; \sigma + i\,I(T),
\end{equation}
and we carry $\sigma$ and $T$ as two separate arguments of $\zeta$. Nothing about
the \emph{value} of $\zeta$ changes; we have only renamed the variable that
generates $t$ and $m$. In this notation Siegel's decomposition
\eqref{eq:siegel-named} becomes
\begin{equation}\label{eq:zeta-T}
\zeta(\sigma,T)
=
\Sigma_{1} + \Sigma_{2} + R,
\end{equation}
where, with $m=\lfloor T\rfloor$,
\begin{align}
\Sigma_{1}(\sigma,T) &= \sum_{n=1}^{m}
\frac{1}{n^{\,\sigma + i\,I(T)}}, \\[4pt]
\Sigma_{2}(\sigma,T) &= \chi\!\bigl(\sigma + i\,I(T)\bigr)
\sum_{n=1}^{m}
\frac{1}{n^{\,1 - \sigma - i\,I(T)}} .
\end{align}
Using the identity
\begin{equation}
\frac{(2\pi)^s e^{\tfrac{\pi i s}{2}}}{\Gamma(s)\bigl(e^{2\pi i s}-1\bigr)}
\;=\;
\frac{\chi(s)}{\,e^{i\pi s}-1\,},
\end{equation}
the remainder becomes
\begin{equation}\label{eq:R-T}
R(\sigma,T)
=
\frac{\chi\!\bigl(\sigma + i\,I(T)\bigr)}{\,e^{\,i\pi (\sigma + i\,I(T))}-1\,}
\int_{C_2} g(x)\,\mathrm{d}x .
\end{equation}

The point of this section is only that \eqref{eq:zeta-T} \emph{is}
\eqref{eq:siegel-named}: the reparameterization is an exact identity, not an
approximation. Everything that follows takes place inside this reindexed but
otherwise unchanged Riemann--Siegel formula.

\begin{remark}
The derivation of $I(T)$ and its approximations and inverse are developed in
\S\ref{sec:IT-origin}. The context \emph{behind} the mapping is not needed to
state or use~\eqref{eq:zeta-T}.
\end{remark}

% ======================================================================
\section{Decomposing the remainder: \texorpdfstring{$R=R_{1ps}+R_{2ps}$}{R = R1ps + R2ps}}\label{sec:decomp}
% ======================================================================

We now split Siegel's single remainder $R$ into two exact pieces, named
$R_{1ps}$ and $R_{2ps}$ (the subscript ``$ps$'' distinguishes them from the
approximate remainders of other authors). Define
\begin{equation}\label{eq:R1ps-def}
R_{1ps}(\sigma, T) = d_{1}(\sigma, T)\,\cdot\, \lceil T \rceil^{-i\,I(T)},
\end{equation}
\begin{equation}\label{eq:R2ps-def}
R_{2ps}(\sigma, T) = d_{2}(\sigma, T)\,\cdot\, \lceil T \rceil^{\,i\,I(T)}
\cdot \frac{\chi(\sigma + i\,I(T))}{|\chi(\sigma + i\,I(T))|},
\end{equation}
with real weights
\begin{align}
d_{1}(\sigma,T) &= |R|\,
\frac{\sin\!\bigl(\omega - \arg R + \arg\chi\bigr)}{\sin\!\bigl(2\omega + \arg\chi\bigr)},
\label{eq:d1}\\[4pt]
d_{2}(\sigma,T) &= |R|\,
\frac{\sin\!\bigl(\omega + \arg R\bigr)}{\sin\!\bigl(2\omega + \arg\chi\bigr)},
\label{eq:d2}
\end{align}
where $\arg R=\arg R(\sigma,T)$, $\arg\chi=\arg\chi(\sigma+iI(T))$, and
\begin{equation}
\omega(T) \;=\; I(T)\,\log\!\bigl(\lfloor T+1\rfloor\bigr).
\end{equation}
Geometrically, $\omega$ is the \emph{absolute} direction angle, measured from
the positive real axis, of the next summand of the first Dirichlet sum: with
$t=I(T)$ and $m=\lfloor T\rfloor$, the summand $n^{-s}=n^{-\sigma}e^{-it\ln n}$
rotates by $\theta_n=-t\ln\tfrac{n+1}{n}$ relative to its predecessor, and these
relative rotations telescope,
\begin{equation}\label{eq:omega-telescope}
\theta_1+\theta_2+\cdots+\theta_m
\;=\;
-t\ln\tfrac{2}{1}-t\ln\tfrac{3}{2}-\cdots-t\ln\tfrac{m+1}{m}
\;=\;
-t\ln(m+1)
\;=\;
-\omega,
\end{equation}
so the $(m+1)$st summand points along $e^{-i\omega}$.

The two quotients are one intersection seen from two sides. In the frame that
lays the bisector link of \S\ref{sec:bisector} along the real axis, the segment
joining the two ends of the reverse bisector link, the yin and yang points,
crosses that axis at distance $d_1$ from the joint, and the same crossing read in
the reverse frame is $d_2$; \S\ref{sec:derive-r1ps}--\S\ref{sec:derive-r2ps}
derive \eqref{eq:d1}--\eqref{eq:d2} that way, which is how they were found.

Note: both $d_1$ and $d_2$ are real, and on the critical line they are always
positive. Writing $\psi=\arg\chi(\sigma+iI(T))$, the
definitions \eqref{eq:R1ps-def}--\eqref{eq:R2ps-def} express $R$ as a combination
of two fixed \emph{unit} directions,
\begin{equation}\label{eq:R-as-cone}
R \;=\; d_1\,e^{-i\omega} \;+\; d_2\,e^{\,i(\omega+\psi)},
\end{equation}
namely the phases of the next summand of each Dirichlet sum. Requiring
$d_1,d_2\in\mathbb{R}$ turns \eqref{eq:R-as-cone} into a $2\times2$ real linear
system (its real and imaginary parts), whose solution by Cramer's rule is exactly
the pair of sine expressions \eqref{eq:d1}--\eqref{eq:d2}; hence $d_1,d_2$ are
real. They are positive exactly when $R$ lies in the cone spanned by the two
directions $e^{-i\omega}$ and $e^{i(\omega+\psi)}$, so that both numerators carry the
same sign as the common denominator $\sin(2\omega+\psi)$. On the critical line
$R$ never leaves the cone: there $d_1=d_2$ (Corollary~\ref{cor:equal}) and the
common value stays strictly inside the link, the normalized fraction
$\lceil T\rceil^{\sigma}d_1$ sweeping ${\approx}\,[0.23,\,0.78]$ in every unit
interval (\S\ref{sec:d1-function}); we verify $d_1=d_2>0$ at every one of
$30{,}562$ samples of a dense grid over $1\le T\le30$, including magnified
windows around the near-parallel instants at fractional parts
${\approx}\,\tfrac14,\tfrac34$, and in unit windows at $T=100$, $300$, and
$1000$ (\texttt{check\_d1\_positive\_critical.py}). Off the critical line the
weights are positive outside a narrow window around each pole of
\S\ref{sec:pole-locations}; inside such a window $R$ exits the cone and one of
the two weights turns negative (Remark~\ref{rem:d-ratio}, where $d_1=-d_2$
exactly at a pole). Both statements, the unconditional positivity on the
critical line and the $O(1/T)$ width of the off-line windows, are proved in
\S\ref{sec:d1-positive}. Geometrically, positivity is the statement that each
fractional summand points \emph{forward} along its
partial sum, a positive fraction of the next term.

\begin{theorem}\label{thm:decomp}
For all $\sigma$ and $T$,
\[
R_{1ps} + R_{2ps} = R.
\]
Consequently,
\[
\zeta = \Sigma_{1} + R_{1ps} + \Sigma_{2} + R_{2ps}.
\]
\end{theorem}

\begin{proof}
Write $\phi = \arg R$, $\psi = \arg\chi$, and $\omega=\omega(T)$, so that
$R=|R|e^{i\phi}$ and, for non-integer $T$, $\lceil T\rceil^{\,iI(T)}=e^{i\omega}$
and $\chi/|\chi|=e^{i\psi}$.

\smallskip
\noindent\textbf{Step 1 (substitute).} From \eqref{eq:R1ps-def}--\eqref{eq:d2},
\[
R_{1ps} = |R|\,\frac{\sin(\omega - \phi + \psi)}{\sin(2\omega + \psi)}\, e^{-i\omega},
\qquad
R_{2ps} = |R|\,\frac{\sin(\omega + \phi)}{\sin(2\omega + \psi)}\, e^{i\omega}\, e^{i\psi}.
\]

\noindent\textbf{Step 2 (add).}
\[
R_{1ps} + R_{2ps}
= \frac{|R|}{\sin(2\omega + \psi)}
\Bigl[\, e^{-i\omega}\sin(\omega - \phi + \psi) + e^{i(\omega + \psi)}\sin(\omega + \phi)\,\Bigr].
\]

\noindent\textbf{Step 3 (expand via $\sin\theta = \tfrac{e^{i\theta}-e^{-i\theta}}{2i}$).}
\begin{align*}
e^{-i\omega}\sin(\omega - \phi + \psi)
&= \frac{e^{i(\psi - \phi)} - e^{i(-2\omega + \phi - \psi)}}{2i}, \\
e^{i(\omega + \psi)}\sin(\omega + \phi)
&= \frac{e^{i(2\omega + \psi + \phi)} - e^{i(\psi - \phi)}}{2i}.
\end{align*}

\noindent\textbf{Step 4 (add the expansions).} The $e^{i(\psi-\phi)}$ terms
cancel, leaving
\[
\frac{1}{2i}\Bigl[\, e^{i(2\omega + \psi + \phi)} - e^{i(\phi - \psi - 2\omega)}\,\Bigr].
\]

\noindent\textbf{Step 5 (factor $e^{i\phi}$).}
\[
= \frac{e^{i\phi}}{2i}\Bigl[\, e^{i(2\omega + \psi)} - e^{-i(2\omega + \psi)}\,\Bigr]
= e^{i\phi}\sin(2\omega + \psi).
\]

\noindent\textbf{Step 6 (substitute back).}
\[
R_{1ps} + R_{2ps}
= \frac{|R|}{\sin(2\omega + \psi)}\cdot e^{i\phi}\sin(2\omega + \psi)
= |R|e^{i\phi} = R.
\]

\smallskip
\noindent\textbf{Integer $T$.} At integer $T$ one has
$\lceil T\rceil^{\,iI(T)} = -e^{i\omega}$, and the factor $(-1)^{\lfloor T\rfloor}$
in $R$ shifts $\phi=\arg R$ by $\pi$. That $\pi$ shift flips the signs of both
sine terms in $d_1,d_2$, while the $-1$ contributes a further sign; the two sign
flips cancel, so $R = R_{1ps}+R_{2ps}$ holds for all $T$.
\end{proof}

A formalization of the algebra of Theorem~\ref{thm:decomp} in Lean is given in
Appendix~\ref{app:lean-decomp}. It assumes no axioms of its own; the appendix
records what its three hypotheses are and what is left outside its scope.
Since $R=R_{1ps}+R_{2ps}$, we have in particular proved the exact formula
\begin{equation}\label{eq:zeta-ps}
\zeta(\sigma,T)
=
\Sigma_{1} + R_{1ps} + R_{2ps} + \Sigma_{2}.
\end{equation}

\begin{figure}[htbp]
\centering
\includegraphics[width=0.82\textwidth]{fig_remainder_average}
\caption{$R_{1ps}(\sigma,T)-1$ (green, left), $\tfrac12 R(\sigma,T)$ (violet,
middle), $R_{2ps}(\sigma,T)+1$ (red, right), for $\sigma=\tfrac12$ and
$2<T<3$. The middle trajectory is the average of the left and right
trajectories (the unit shifts cancel), since
$R_{1ps}+R_{2ps}=R$. Generated by \texttt{fig\_remainder\_average.py}.}
\label{fig:remainder-average}
\end{figure}

% ======================================================================
\section{The remainders are ``one more partial summand''}\label{sec:summand}
% ======================================================================

Theorem~\ref{thm:decomp} is more than an algebraic identity. Unlike Siegel's
integral $R$ (or the approximate remainders $R_{ak}$ of Kuznetsov), which bear no
resemblance to the terms of the Dirichlet sums, each of $R_{1ps}$ and $R_{2ps}$
is \emph{exactly the next term} of its partial sum, scaled by a real
number (positive in the sense of the note after \eqref{eq:R-as-cone}). We call
such a term a \emph{fractional summand}.

Concretely, the last term actually included in $\Sigma_1$ is the $\lfloor
T\rfloor$-th term, and the very next term would be the $\lceil T\rceil$-th. The
remainder $R_{1ps}$ lies along that next term, shortened by a real factor; and
similarly for $R_{2ps}$ along the next term of $\Sigma_2$. Table~\ref{tab:frac}
makes this precise.

\begin{table}[htbp]
\centering
\resizebox{\textwidth}{!}{%
\begin{tabular}{ccccc}
\toprule
 & Last summand & Next summand & Fractional amount & Partial summand \\
\midrule
$R_{1ps}$
& $\dfrac{1}{\lfloor T\rfloor^{\,\sigma + iI(T)}}$
& $\dfrac{1}{\lceil T\rceil^{\,\sigma + iI(T)}}$
& $\lceil T \rceil^{\sigma}\, d_{1}$
& $d_{1}\,\lceil T\rceil^{-iI(T)}$ \\[10pt]
$R_{2ps}$
& $\dfrac{\chi(\sigma + iI(T))}{\lfloor T\rfloor^{\,1 - \sigma - iI(T)}}$
& $\dfrac{\chi(\sigma + iI(T))}{\lceil T\rceil^{\,1 - \sigma - iI(T)}}$
& $\dfrac{\lceil T \rceil^{\,1-\sigma}}{|\chi(\sigma + iI(T))|}\,d_{2}$
& $d_{2}\,\lceil T\rceil^{\,iI(T)}\,\dfrac{\chi(\sigma + iI(T))}{|\chi(\sigma + iI(T))|}$ \\
\bottomrule
\end{tabular}%
}
\caption{Each remainder is the $(m+1)$st (``next'') summand of its Dirichlet sum,
scaled by a real fractional amount. The final column is the resulting
\emph{partial summand} $=$ (fractional amount) $\times$ (next summand), which
simplifies to $R_{1ps}$ and $R_{2ps}$ respectively.}
\label{tab:frac}
\end{table}

Reading \eqref{eq:zeta-ps} in this light, and writing $s=\sigma+iI(T)$ and
$m=\lfloor T\rfloor$, the zeta function takes the clean ``one more summand'' form
\begin{equation}\label{eq:zeta-clean}
\boxed{\;
\zeta(s)=\sum_{n=1}^{m} n^{-s} + \hat d_1\,(m+1)^{-s}
        + \chi(s)\sum_{n=1}^{m} n^{s-1} + \hat d_2\,\chi(s)\,(m+1)^{s-1}
\;}
\end{equation}
with $\hat d_1,\hat d_2$ reals, positive on the critical line (and off it
outside the pole windows of \S\ref{sec:pole-locations}). In words: each finite Dirichlet sum,
carried exactly one fractional term past its cutoff, reproduces $\zeta$
with no remainder terms.

The coefficients here are the \emph{fractional amounts} of
Table~\ref{tab:frac}, not the weights $d_1,d_2$ of
\eqref{eq:R1ps-def}--\eqref{eq:R2ps-def}. The hat is what distinguishes them,
and the two differ by the length of the link they sit on:
\begin{equation}\label{eq:frac-vs-weight}
\hat d_1=\lceil T\rceil^{\sigma}\,d_1,
\qquad
\hat d_2=\frac{\lceil T\rceil^{\,1-\sigma}}{|\chi|}\,d_2.
\end{equation}
The unhatted weight multiplies a \emph{unit} vector, so $|R_{1ps}|=d_1$ is an
absolute length, whereas $\hat d_1$ multiplies the full summand $(m+1)^{-s}$
and is therefore the dimensionless fraction of that summand which is used. At
$\sigma=\tfrac12$, $T=6.18$ for instance, $d_1=0.0875$ while
$\hat d_1=\sqrt7\,d_1=0.2315$.

\begin{figure}[htbp]
\centering
\includegraphics[width=0.86\textwidth]{fig1_spiral_summands}
\caption{The fractional-summand picture: pictorial (spiral) representation of the
four pieces of \eqref{eq:zeta-clean} in the complex plane for $\sigma=0.5$ and
$T\approx 6.18$ (equivalently $t=I(T)\approx 279.85$, $m=6$). The blue path is the
partial-sum spiral $\Sigma_1=\sum_{n\le m}n^{-s}$; the green path,
laid tip-to-tail from the point $\Sigma_1+R_{1ps}$, is
$\Sigma_2=\chi\sum_{n\le m}n^{\,s-1}$. The short red and orange segments $R_{1ps}=d_1e^{-i\omega}$ and
$R_{2ps}=d_2e^{\,i(\omega+\arg\chi)}$ lie along link $m$ of each spiral
(the faint dashed vectors show the full next summand, of which each remainder is
a ``fractional'' part), so that the four pieces add up to the point $\zeta(s)$.
Underneath, the thin faint blue path continues the $\Sigma_1$ chain past $m$
out to link $\lfloor I(T)/\pi\rfloor=89$, winding into its spiral near
$\zeta$.
On the critical line $d_1=d_2$, so the two remainders have equal length.
The fractional links are short here ($d_1=d_2\approx0.087$), so the inset
zooms in on them. Generated
by \texttt{fig1\_spiral\_summands.py}.}
\label{fig:spiral-summands}
\end{figure}

\begin{figure}[htbp]
\centering
\includegraphics[width=0.72\textwidth]{fig2_remainder_zoom}
\caption{Zoom into the joint region of Figure~\ref{fig:spiral-summands}
(imaginary part $\approx 1.5$--$2.6$), showing the theorem
$R_{1ps}+R_{2ps}=R$ as a small triangle. Starting from the joint $\Sigma_1$
(end of the sum-1 partial sum), the red link $R_{1ps}$ reaches
$\Sigma_1+R_{1ps}$ and the orange link $R_{2ps}$ reaches
$\Sigma_1+R_{1ps}+R_{2ps}$; the straight purple vector from $\Sigma_1$ to that
last joint is Siegel's remainder $R=R_{1ps}+R_{2ps}$ (see the legend). The
curly braces labelled $d_1$ and $d_2$ span the two fractional links,
marking their lengths, i.e.\ how far each reaches from its own joint.
The thin blue and green lines show the entire extent each fractional summand
would have if it were not shortened to the lengths $d_1$ and $d_2$: each
extends its fractional link to the full $(m+1)$st summand. The faint blue and
green chains are the neighboring links of the two partial sums, arriving at
and departing from the joint region. Generated by
\texttt{fig2\_remainder\_zoom.py}.}
\label{fig:remainder-zoom}
\end{figure}

% ======================================================================
\section{Summands as links and joints}\label{sec:matrix-product}
% ======================================================================

Equation~\eqref{eq:zeta-ps} presents $\zeta$ as the sum of the four terms
$\Sigma_1+R_{1ps}+\Sigma_2+R_{2ps}$, two of which are partial sums. Read geometrically, each summand of a
partial sum is one \emph{link} of a planar chain, and each remainder
$R_{1ps},R_{2ps}$ is one additional \emph{fractional} link. One can think
of $\hat d_1$ (and likewise $\hat d_2$) as the percentage of the way along
that additional link at which the chain stops: the link taken at its full
length would be the whole $(m+1)$st summand, and $R_{1ps}$ is that
fraction $\hat d_1$ of it. Write $s=\sigma+iI(T)$,
$t=I(T)$, $m=\lfloor T\rfloor$, and $\omega=t\ln(m+1)$, and denote two complex numbers by
\begin{equation}
B_1=\Sigma_1+R_{1ps},\qquad B_2=\Sigma_2+R_{2ps},\qquad \zeta=B_1+B_2 .
\end{equation}

Each remainder is a single extra link appended to its partial sum,
pointing in the direction of the $(m+1)$-th summand with a real length $d_1$ or
$d_2$ (the fractional summands of \S\ref{sec:summand}):
\begin{equation}\label{eq:R-phasors}
R_{1ps}=d_1\,\lceil T\rceil^{-iI(T)}=d_1\,e^{-i\omega},\qquad
R_{2ps}=d_2\,\lceil T\rceil^{\,iI(T)}\,\frac{\chi}{|\chi|}=d_2\,e^{\,i(\omega+\arg\chi)},
\end{equation}
with $d_1,d_2$ as in \eqref{eq:d1}--\eqref{eq:d2}. The first equality in each is
the definition \eqref{eq:R1ps-def}--\eqref{eq:R2ps-def}; the second uses that, for
non-integer $T$, $\lceil T\rceil=\lfloor T+1\rfloor$, so
$\lceil T\rceil^{\pm iI(T)}=e^{\pm iI(T)\log\lceil T\rceil}=e^{\pm i\omega}$ by the
definition $\omega=I(T)\log\lfloor T+1\rfloor$, together with
$\chi/|\chi|=e^{\,i\arg\chi}$.

\subsection{Sum form}\label{sec:sum-form}

Splitting $\zeta$ into its partial sum plus this last
fractional link,
\begin{alignat}{2}
B_1 &= \phantom{\chi(s)}\sum_{n=1}^{m} n^{-s}    &&{}+ d_1\,e^{-i\omega}, \label{eq:B1-sum}\\
B_2 &= \chi(s)\sum_{n=1}^{m} n^{\,s-1}           &&{}+ d_2\,e^{\,i(\omega+\arg\chi)}. \label{eq:B2-sum}
\end{alignat}

\paragraph{Componentwise.} Writing out real and imaginary parts,
\begin{align}
\Re B_1 &= \sum_{n=1}^{m}\frac{\cos(t\ln n)}{n^{\sigma}} + d_1\cos\bigl(t\ln(m+1)\bigr),\\
\Im B_1 &= -\sum_{n=1}^{m}\frac{\sin(t\ln n)}{n^{\sigma}} - d_1\sin\bigl(t\ln(m+1)\bigr),\\
\Re B_2 &= \sum_{n=1}^{m}\frac{|\chi|\cos(\arg\chi+t\ln n)}{n^{\,1-\sigma}} + d_2\cos\bigl(\arg\chi+t\ln(m+1)\bigr),\\
\Im B_2 &= \sum_{n=1}^{m}\frac{|\chi|\sin(\arg\chi+t\ln n)}{n^{\,1-\sigma}} + d_2\sin\bigl(\arg\chi+t\ln(m+1)\bigr).
\end{align}

\subsection{Product form}\label{sec:product-form}

The goal of this subsection is to write each link
as a homogeneous coordinate transformation matrix, and then to group the
links into a single transformation matrix, so that $\zeta$ takes the form
of a matrix product with four components: the two sums $\Sigma_1,\Sigma_2$
and the two partial-summand remainders $R_{1ps},R_{2ps}$.

Each summand is a link in the plane. We number the
links of each chain from zero: link $n$ runs from joint $n$ to joint $n+1$
(joint $n$ being the value of the partial sum after $n$ summands), so link $n$
represents the $(n+1)$st summand. The $m$ summands of a partial sum are thus
links $0,\dots,m-1$, and each remainder is the \emph{fractional link} $m$. We
encode one link by the $3\times3$ matrix with link length $l$ and rotation
$\theta$ from the previous link,\footnote{This is the planar case of the
Denavit--Hartenberg convention of robot kinematics~\cite{DenavitHartenberg1955},
in which each link of a serial chain is encoded by the $4\times4$ homogeneous
transformation
$\mathrm{Rot}_z(\theta_i)\,\mathrm{Trans}_z(d_i)\,\mathrm{Trans}_x(a_i)\,
\mathrm{Rot}_x(\alpha_i)$
built from four parameters: the joint angle $\theta_i$, the joint offset $d_i$,
the link length $a_i$ and the link twist $\alpha_i$. A planar chain has
$d_i=\alpha_i=0$, which leaves the two parameters $(\theta_i,a_i)$ and the
$3\times3$ matrix \eqref{eq:link-matrix}; \eqref{eq:zeta-serial} is then the
forward-kinematic product of the chain. One collision of names to watch: the
Denavit--Hartenberg offset $d_i$ has nothing to do with the weights $d_1,d_2$ of
\eqref{eq:d1}--\eqref{eq:d2}.}
\begin{equation}\label{eq:link-matrix}
M(\theta,l)=\linkmat{\theta}{l},
\end{equation}
where the accumulated position is the top two entries of the last column.
The link matrix rotates by $\theta$ and then translates forward $l$
along the new heading. Its inverse is again a link-shaped matrix, of a
second type that translates \emph{back} $l$ along the current heading and
then undoes the rotation:
\begin{equation}\label{eq:link-matrix-2}
M(\theta,l)^{-1}=
\begin{pmatrix}
\cos\theta & \phantom{-}\sin\theta & -l\\
-\sin\theta & \phantom{-}\cos\theta & \phantom{-}0\\
\phantom{-}0 & \phantom{-}0 & \phantom{-}1
\end{pmatrix}.
\end{equation}
Forward links build chain 1; inverses are what let chain 2 be walked
backwards.

\emph{Chain 1.} Collect the $m$ summand links of $\Sigma_1$ into a single
product and give the fractional link its own name:
\begin{equation}\label{eq:P1-TR1-defs}
P_{\Sigma_1}=\prod_{n=0}^{m-1}M\bigl(\theta_n,\,(n+1)^{-\sigma}\bigr),
\qquad
M_{R_{1ps}}=M\bigl(-t\ln\tfrac{m+1}{m},\,d_1\bigr),
\end{equation}
where link $n$ carries the two parameters
$\theta_n$, its rotation from the previous link, and
$l_n=(n+1)^{-\sigma}$, its length: as in \eqref{eq:omega-telescope},
$\theta_0=0$ and $\theta_n=-t\ln\tfrac{n+1}{n}$ for $n\ge1$. Written out, with
the first two factors and the last two shown, the chain that
reaches $B_1$ is
\begin{equation}\label{eq:P1-written}
\footnotesize
P_{\Sigma_1}M_{R_{1ps}}=
\underbrace{\left(\begin{array}{@{}c@{\ \,}c@{\ \,}c@{}}
1 & 0 & 1\\
0 & 1 & 0\\
0 & 0 & 1
\end{array}\right)}_{\text{link }0}
\underbrace{\linkmatd{\theta_1}{2^{\sigma}}}_{\text{link }1}\;\cdots\;
\underbrace{\linkmatd{\theta_{m-1}}{m^{\sigma}}}_{\text{link }m-1}
\underbrace{\linkmat{\theta_m}{d_1}}_{\text{fractional link }m},
\end{equation}
the leading factor being the unit step $1^{-s}=1$ along the real axis, where
$\theta_0=0$ and $l_0=1$ collapse \eqref{eq:link-matrix} to a pure translation,
and the trailing factor being $M_{R_{1ps}}$, whose rotation
$\theta_m=-t\ln\tfrac{m+1}{m}$ continues the same pattern while its length is
$d_1$ rather than a power of $m+1$. The product carries the frame from the origin
to the bisector point $B_1=\Sigma_1+R_{1ps}$, whose coordinates are the top two
entries of the last column, and it arrives with accumulated rotation
$-t\ln(m+1)=-\omega$ by the telescoping \eqref{eq:omega-telescope}.

\emph{Chain 2.} The second chain is traversed tip-first, exactly as
Figure~\ref{fig:spiral-summands} draws leg 2: the fractional link
$R_{2ps}$ leaves $B_1$, and the summand links follow in descending order.
Walking a link backwards is exactly its inverse \eqref{eq:link-matrix-2};
regrouping each rotation with the backward translation that follows it
returns everything to the forward type \eqref{eq:link-matrix}, with
negated parameters and each rotation paired with the preceding link's
length. Two bookkeeping factors appear at the ends of the backward walk: a
leading pure translation $M(0,-d_2)$, and a trailing pure rotation
$M(-\arg\chi,0)$, the functional-equation factor $\chi$ peeled off on its
own. The leading translation absorbs the zero-length joint
$M(2\omega+\arg\chi+\pi,\,0)$ that must sit between the two chains (a
rotation followed by a translation is a single link), defining the
backward fractional link
\begin{equation}\label{eq:TR2-def}
M_{R_{2ps}}
\;:=\;
M\bigl(2\omega+\arg\chi+\pi,\,0\bigr)\,M(0,\,-d_2)
\;=\;
M\bigl(2\omega+\arg\chi+\pi,\;-d_2\bigr);
\end{equation}
the backward summand links collect into
\begin{equation}\label{eq:P2-def}
P_{\Sigma_2}
\;:=\;
\prod_{n=m}^{1}M\bigl(\theta_n,\;-|\chi|\,n^{\sigma-1}\bigr)
\quad\text{($n$ descending, $\theta_n=-t\ln\tfrac{n+1}{n}$)};
\end{equation}
and the trailing rotation is simply dropped, since a post-multiplied
zero-length joint never moves the position column, only the final
orientation.

Written out as \eqref{eq:P1-written} was, with the first two factors and the last
two shown and with $\psi=\arg\chi$ as before, the second chain is
\begin{equation}\label{eq:P2-written}
\footnotesize
\begin{aligned}
M_{R_{2ps}}P_{\Sigma_2}
&=\underbrace{\linkmat{(2\omega+\psi+\pi)}{-d_2}}_{\text{fractional link }m}
\underbrace{\linkmatx{\theta_m}{m^{1-\sigma}}}_{\text{link }m-1}\;\cdots\\[3pt]
&\hphantom{{}={}}\cdots\;
\underbrace{\linkmatx{\theta_2}{2^{1-\sigma}}}_{\text{link }1}
\underbrace{\linkmat{\theta_1}{-|\chi|}}_{\text{link }0},
\end{aligned}
\end{equation}
the leading rotation being the merged joint of \eqref{eq:TR2-def}. The links run
backwards, from the fractional link $m$ down to link $0$, each brace naming the
link whose length that factor carries; the rotation it carries is the $\theta_n$
of the link after it, which is the regrouping. Its
last column is not $B_2$ but $e^{i\omega}B_2$, leg 2 read in the frame chain 1
leaves behind, whose orientation is $-\omega$; multiplying on the left by
$P_{\Sigma_1}M_{R_{1ps}}$ turns that back onto world axes and adds $B_1$, which
is \eqref{eq:zeta-serial}. All three columns agree to $28$ digits at
$\sigma=\tfrac12$, $T=6.18$ (\texttt{check\_chain2\_matrix.py}).

\paragraph{The full $3\times3$ equation for $\zeta$.}
With these four names the two chains combine into one serial product, in a
single link type:
\begin{equation}\label{eq:zeta-serial}
Z \;=\; P_{\Sigma_1}\,M_{R_{1ps}}\,M_{R_{2ps}}\,P_{\Sigma_2}
\;=\;
\begin{pmatrix}
\cos\Phi & -\sin\Phi & \Re\,\zeta\\
\sin\Phi & \phantom{-}\cos\Phi & \Im\,\zeta\\
0 & 0 & 1
\end{pmatrix},
\qquad \Phi=\pi+\arg\chi.
\end{equation}
The factors line up one to one with the terms of \eqref{eq:zeta-ps},
repeated here for comparison:
\begin{equation*}\tag{\ref{eq:zeta-ps}}
\zeta(\sigma,T)
=
\Sigma_{1} + R_{1ps} + R_{2ps} + \Sigma_{2}.
\end{equation*}
Figure~\ref{fig:zeta-chain} draws \eqref{eq:zeta-serial} with one arrow
per factor.

\begin{figure}[H]
\centering
\includegraphics[width=0.8\textwidth]{fig_zeta_chain}
\caption{The serial chain of \eqref{eq:zeta-serial} at $\sigma=\tfrac12$,
$T=6.18$ ($t=I(T)\approx279.85$, $m=6$), one arrow for the net displacement
of each matrix factor: $P_{\Sigma_1}$ (blue, $O$ to $\Sigma_1$),
$M_{R_{1ps}}$ (red, to
$B_1$), $M_{R_{2ps}}$ (orange, to $B_1+R_{2ps}$; its rotation parameter
carries the merged joint $2\omega+\arg\chi+\pi$), and $P_{\Sigma_2}$ (green,
to $\zeta$). The four arrows are the four terms of
$\zeta=\Sigma_1+R_{1ps}+R_{2ps}+\Sigma_2$; the two fractional links are
short here ($d_1=d_2\approx0.087$), so the inset zooms in on them. Colors
match
Figure~\ref{fig:spiral-summands}. Generated by
\texttt{fig\_zeta\_chain.py}.}
\label{fig:zeta-chain}
\end{figure}

Three features of \eqref{eq:zeta-serial}
deserve comment. First, the rotation $2\omega+\arg\chi+\pi$ carried by
$M_{R_{2ps}}$ has three jobs at once: $+\omega$ cancels chain 1's
accumulated rotation $-t\ln(m+1)=-\omega$, a
further $+(\omega+\arg\chi)$ pre-rotates against the counter-rotation of
the backward-walked chain 2, and the final $+\pi$ rotates the backward
translations to point forward again; together
$\omega+(\omega+\arg\chi)+\pi=2\omega+\arg\chi+\pi$. Second, the links of
$P_{\Sigma_2}$ carry chain 1's very rotation parameters $\theta_n$:
in this form the two chains use identical rotations and differ only in
their lengths, $(n+1)^{-\sigma}$ forward versus $-|\chi|\,n^{\sigma-1}$
backward. Third, the net displacements of the four factors $P_{\Sigma_1}$,
$M_{R_{1ps}}$, $M_{R_{2ps}}$, $P_{\Sigma_2}$ are exactly the four terms
$\Sigma_1$, $R_{1ps}=d_1e^{-i\omega}$,
$R_{2ps}=d_2e^{\,i(\omega+\arg\chi)}$, $\Sigma_2$ of \eqref{eq:zeta-ps},
and the final orientation $\Phi=\pi+\arg\chi$ is the direction of
$-\chi/|\chi|$: the trailing rotation $M(-\arg\chi,0)$ (the
functional-equation factor $\chi$, peeled off automatically when the base
link of chain 2 is regrouped) was discarded, so $\chi$'s direction is
left showing in the orientation block, and $Z$ packages the pair
$(\zeta,\arg\chi)$ in one matrix. (Appending the discarded
$M(-\arg\chi,0)$ would leave the position column untouched and make the
orientation the constant $\pi$.)

% ======================================================================
\section{Other remainders}\label{sec:other-remainders}
% ======================================================================

\subsection{Kuznetsov's remainders}\label{sec:kuznetsov}

Our remainders $R_{1ps},R_{2ps}$ are not the only way to split Siegel's $R$. It is worth
contrasting them with the remainder of Alexey Kuznetsov~\cite{Kuznetsov2025},
whose approximation to the zeta function we also use for numerical work. The
essential difference is one of \emph{kind}: Kuznetsov's remainder is
an \emph{approximation} to $R$, whereas $R_{1ps}+R_{2ps}=R$ is an \emph{exact}
identity (Theorem~\ref{thm:decomp}).

Writing Kuznetsov's approximate remainder as $R_{ak}$ (the subscript ``$ak$''
for Alexey Kuznetsov, and to distinguish it from Siegel's exact, unsubscripted
$R$), his method gives
\begin{equation}\label{eq:Rak}
R \;\approx\; R_{ak}
\;=\; -\tfrac12(-1)^{\lfloor T\rfloor}
\Bigl(I_{1}(\sigma,T) + \chi\bigl(\sigma+iI(T)\bigr)\,I_{2}(\sigma,T)\Bigr),
\end{equation}
where $I_1,I_2$ are the two integrals of Kuznetsov's construction, each
evaluated as a short finite series in precomputed coefficients
$\omega_0,\omega_1[n],\lambda[n]$ (his notation; this $\omega_0$ is
unrelated to the direction angle $\omega(T)$ of \S\ref{sec:decomp}); in
some of our calculations we use $8$ such
coefficients, which already gives several digits of accuracy across the range of
interest. Just as we split $R$ into two pieces, the two terms of \eqref{eq:Rak}
define Kuznetsov's two half-remainders
\begin{equation}\label{eq:Rak-split}
R_{ak}=R_{1ak}+R_{2ak},\qquad
R_{1ak}=-\tfrac12(-1)^{\lfloor T\rfloor}I_{1},\qquad
R_{2ak}=-\tfrac12(-1)^{\lfloor T\rfloor}\chi\,I_{2},
\end{equation}
so that, in parallel with \eqref{eq:zeta-ps},
\[
\zeta \;\approx\; \Sigma_1 + R_{1ak} + \Sigma_2 + R_{2ak}.
\]
The split \eqref{eq:Rak-split} mirrors the functional equation: $R_{1ak}$ pairs
with the first Dirichlet sum $\Sigma_1$ and $R_{2ak}$ carries the
functional-equation factor $\chi$ alongside $\Sigma_2$.
Figure~\ref{fig:kuznetsov-zoom} compares the two routes from $\Sigma_1$ to
$\Sigma_1+R$: the exact fractional links on one side and Kuznetsov's dashed
pair on the other.

\begin{figure}[!htb]
\centering
\includegraphics[width=0.72\textwidth]{fig4_kuznetsov_zoom}
\caption{Exact versus Kuznetsov remainders, in the same zoom region as
Figure~\ref{fig:remainder-zoom}. The exact fractional summands
$R_{1ps}$ (red) and $R_{2ps}$ (orange) run from $\Sigma_1$ to $\Sigma_1+R$ on
one side of the resultant $R$ (purple); Kuznetsov's approximate half-remainders
$R_{1ak},R_{2ak}$ are drawn as the orange dashed pair. Since $R_{1ak}+R_{2ak}\approx R$
(here to $\sim10^{-14}$), the dashed path lands on $\Sigma_1+R$; unlike the
exact links, the $ak$ links do not lie along a summand direction. Computed with
$8$ Kuznetsov coefficients and generated by \texttt{fig4\_kuznetsov\_zoom.py}.}
\label{fig:kuznetsov-zoom}
\end{figure}

This pairing is exactly Siegel's: in his 1932 paper~\cite{Siegel1932} Siegel
writes $\zeta$ as
$f(s)$ plus its reflected counterpart, and the first side is
\begin{equation}\label{eq:siegel-f}
f(s) \;\approx\; \Sigma_1 + R_{1ak},
\end{equation}
i.e.\ the first partial sum completed by an approximation to the first half of
the remainder. The sign is $\approx$ rather than $=$ because $R_{1ak}$ is the
quantity $-\tfrac12(-1)^{\lfloor T\rfloor}I_1$ of \eqref{eq:Rak-split}, with
$I_1$ evaluated as a truncated series; read the other way round, as the
\emph{definition} $R_{1ak}:=f(s)-\Sigma_1$, the relation is an identity and the
approximation moves into the numerical evaluation of $R_{1ak}$.

The form of \eqref{eq:siegel-f} is Siegel's own, and it is worth recording where
in his paper to look for it. He defines $f$ in his \S3, equation~58, as the
contour integral
\begin{equation}\label{eq:siegel-58}
f(s)
\;=\;
\int_{0\swarrow 1}
\frac{x^{-s}\,e^{\pi i x^{2}}}{e^{\pi i x}-e^{-\pi i x}}\,\mathrm{d}x ,
\end{equation}
where his symbol $0\swarrow 1$ marks a straight path crossing the real axis
between $0$ and $1$, running from upper right to lower left parallel to the
bisector of the first and third quadrants; it is the reflection in the real axis
of the path $0\nwarrow 1$ carrying the integral $\Phi$ of his \S1. The Gaussian
factor $e^{\pi i x^{2}}$ decays in both directions along that line, so
\eqref{eq:siegel-58} converges. He then develops $f$ asymptotically in his \S4,
where his equation~84 reads
\begin{equation}\label{eq:siegel-84}
f(s)
\;=\;
\sum_{n=1}^{m_{1}} n^{-s}
\;+\;
O\!\bigl((\lvert s\rvert/2\pi e)^{-\sigma/2}\bigr)
\qquad (\sigma\ge 0,\; t>0),
\end{equation}
with $m_{1}=[\Re\eta_{1}-\Im\eta_{1}]$ and $\eta_{1}=\sqrt{s/2\pi i}$, which to
leading order is Siegel's cutoff $\lfloor\sqrt{t/2\pi}\rfloor$
of~\eqref{eq:siegel-m}. Earlier in the same section he states it in words as
well: developing $f$ asymptotically yields ``as a principal term a sum of
$[\sqrt{t/2\pi}\,]$ summands, that is $\sum_{n=1}^{m}n^{-s}$''. So ``first
partial sum plus a remainder'' is Siegel's own description of $f$, and it
confirms that $\approx$ is the right sign in \eqref{eq:siegel-f}: at
$\sigma=\tfrac12$ the error term of \eqref{eq:siegel-84} is only
$O(t^{-1/4})$. One caution on the cutoff: because
$\sqrt{t/2\pi}=T+\tfrac12-\tfrac1{24T}+O(T^{-2})$ by
\eqref{eq:siegel-cutoff-index}, Siegel's $m_{1}$ is one more than our
$m=\lfloor T\rfloor$ whenever $\{T\}>\tfrac12+\tfrac1{24T}$. This does not disturb
\eqref{eq:siegel-f}, since $R_{1ak}$ carries the cutoff as a parameter and
re-centers with it; it is Siegel's choice, not ours, that makes the remainder
smallest. Equation numbers cited here are those of Siegel's paper as reprinted
in his collected works, the numbering carried over by the Barkan--Sklar
translation~\cite{BarkanSklar2018}.\footnote{One incidental finding while
checking \eqref{eq:siegel-f} against the source: the translation's
equation~83 prints the sum as $\sum n^{s-1}$, while equation~84, derived from
it, has $\sum n^{-s}$. The numerics say $n^{-s}$ is correct, so equation~83
carries a typo in the Barkan--Sklar translation.}

Neither Siegel
nor Kuznetsov names the other side, but it is $\Sigma_2+R_{2ak}$, and adding the
two recovers $\zeta$ to the same accuracy. The contrast with our result is that
$\Sigma_1+R_{1ps}$ and $\Sigma_2+R_{2ps}$ are exact however they are read, and,
moreover, each remainder there is a single fractional summand
(\S\ref{sec:summand}).

\subsection{The three remainders naming convention summary}\label{sec:remainders-summary}

There are a number of remainders in play, so to keep them straight we have named
them with subscripts. Siegel's remainder integral is written $R$ when unsplit; each
named split below recovers that same $R$ (exactly, or approximately for the
Kuznetsov form):
\begin{enumerate}
\item \textbf{Riemann--Siegel half-split.}
\[
R \;=\; R_{rs} \;=\; R_{1rs}+R_{2rs},
\qquad
R_{1rs}=R_{2rs}=\tfrac12 R.
\]
Here $R_{1rs}$ is simply half the remainder integral.
\item \textbf{Partial-summand split} (\S\ref{sec:decomp}--\S\ref{sec:summand}).
\[
R \;=\; R_{ps} \;=\; R_{1ps}+R_{2ps}.
\]
Here $R_{1ps}$ is a partial next summand: the $(\lceil T\rceil)$-th Dirichlet
term of $\Sigma_1$, shortened by the real weight $d_1$.
\item \textbf{Kuznetsov / Siegel $f$ split} (\S\ref{sec:kuznetsov}).
\[
R \;\approx\; R_{ak} \;=\; R_{1ak}+R_{2ak}.
\]
Here $R_{1ak}$ is Siegel's $f(s)$ minus the first partial sum:
$R_{1ak}=f(s)-\Sigma_1$ (approximately, once $R_{1ak}$ is evaluated
numerically).
\end{enumerate}

All three are measured from the same pair of joints, both at index
$m=\lfloor T\rfloor$: each $R_{1\bullet}$ leaves joint $m$ of the $\Sigma_1$
chain, which is the partial sum $\Sigma_1$ itself, and each $R_{2\bullet}$
arrives at joint $m$ of the $\Sigma_2$ chain, which is the point $\Sigma_1+R$.
A split therefore chooses nothing but the point $\Sigma_1+R_{1\bullet}$ at which
the two halves meet, its own bisector point in the sense of
\S\ref{sec:bisector}; the anchoring joints are common to all three.

\subsection{Scaling the remainders with the fractional part of \texorpdfstring{$T$}{T} unchanged}\label{sec:remainder-scaling}

Write $T=\lfloor T\rfloor+\{T\}$ for the integer and fractional parts of the
index. Empirically, the \emph{shape} of the remainder figure depends almost
entirely on the fractional part $\{T\}$. By shape we mean the relative angles
among $R$, $R_{1ps}$, $R_{2ps}$, $R_{1ak}$, and $R_{2ak}$ once the
configuration is rotated into the frame in which $R$ lies along the positive
real axis. The integer part $\lfloor T\rfloor$ mainly retunes the overall
size: with $n=\lceil T\rceil$ the lengths scale like $n^{-\sigma}$ times a
slowly varying amplitude. Figure~\ref{fig:remainder-scale-grid} makes this
visible. Each panel shows only the remainders, drawn in the $R$-frame
(rotated so $R$ lies on the real axis and translated so the midpoint of $R$
sits at the origin), with a common axis scale. The top row holds
$\{T\}=0.18$ at $T=6.18$, $50.18$, and $100.18$; the bottom row holds
$\{T\}=0.72$ at $T=6.72$, $50.72$, and $100.72$. Across a row the figure
shrinks but keeps its shape; across a column (same $\lfloor T\rfloor$,
different $\{T\}$) the shape changes.

\begin{figure}[H]
\centering
\includegraphics[width=\textwidth]{fig_remainder_scale_grid}
\caption{Remainders in the $R$-frame at fixed fractional part $\{T\}$,
varying $\lfloor T\rfloor$. The frame is centered on the midpoint of $R$
(so $R$ runs from $-\lvert R\rvert/2$ to $+\lvert R\rvert/2$ on the real
axis). Top row: $\{T\}=0.18$ at $T=6.18$, $50.18$, $100.18$. Bottom row:
$\{T\}=0.72$ at $T=6.72$, $50.72$, $100.72$. Colors match
Figure~\ref{fig:kuznetsov-zoom}: $R_{1ps}$ red, $R_{2ps}$ orange, $R$
purple, $R_{1ak}$ and $R_{2ak}$ orange dashed. The large open circle marks the
left endpoint of $R$, which is the anchoring joint $\Sigma_1$ at $m=\lfloor
T\rfloor$ that all three splits leave from. All six panels share the
same axis scale. Generated by \texttt{fig\_remainder\_scale\_panels.py}.}
\label{fig:remainder-scale-grid}
\end{figure}

That figure holds $\{T\}$ fixed and climbs the strip. The complementary view
holds the chain fixed and lets $T$ sweep a whole unit, which is
Figure~\ref{fig:remainder-frame-sweep}: the same $R$-frame in sixteen snapshots
from $T=4$ to $T=5$ at $m=4$, the sweep that
Figure~\ref{fig:bisector-frame} will draw for the neighboring links. Two things
show up there that a fixed $\{T\}$ cannot show. First, on the critical line both
splits are symmetric, $|R_{1ps}|=|R_{2ps}|$ exactly and likewise for the $ak$
pair, so both apexes ride the perpendicular bisector of $R$ and the state of a
split is the single number giving the height of its apex above the chord.
Second, the two splits move quite differently: the $ak$ apex marches steadily
from below $R$ to well above it, while the $ps$ apex swings down onto the chord,
crosses to the far side, and comes back, touching the chord exactly at the two
heights where $\sin(2\omega+\arg\chi)$ vanishes
(\S\ref{sec:pole-locations}). There the cone of \eqref{eq:R-as-cone} has
collapsed to a line, and on the critical line the split collapses with it into
the plain halving $R_{1ps}=R_{2ps}=\tfrac12R$: the $d_1$ pole of
\eqref{eq:d1} is removable at $\sigma=\tfrac12$, its numerator vanishing with
its denominator, which is why the sweep passes through those heights smoothly.
Off the line it is a genuine pole --- at $\sigma=0.6$ and $\sigma=0.75$, $d_1$
runs through infinity and changes sign there --- and that is what the bands of
\S\ref{sec:pole-locations} are.

\begin{figure}[H]
\centering
\includegraphics[width=\textwidth]{fig_remainder_frame_sweep}
\caption{The remainder splits in the $R$-frame as $T$ sweeps one unit,
$T=4\to5$ at $\sigma=\tfrac12$, sixteen snapshots read left to right from the
top left, the companion of Figure~\ref{fig:bisector-frame} for the remainders.
The frame and the colors are those of
Figure~\ref{fig:remainder-scale-grid}: $R$ purple along the real axis with its
midpoint at the origin, $R_{1ps}$ red and $R_{2ps}$ orange from the $\Sigma_1$
joint (open circle) to the far end of $R$, $R_{1ak}$ and $R_{2ak}$ orange
dashed, all sixteen panels to one scale, with $|R|$ given in each. The chain
length is held at $m=4$ throughout, as
Figure~\ref{fig:bisector-frame} holds the bisector link, so $T=5$ is the handoff
instant seen from below. The dotted arrow on $R_{1ps}$ and on $R_{2ps}$ runs
from the middle of that segment to the middle of it as the next panel depicts
it, tracing the motion between panels. Both apexes stay on the perpendicular
bisector of $R$, the two halves of each split being equal in length on the
critical line. The $ps$ apex crosses the chord between the fourth and fifth
panels and again between the twelfth and thirteenth, at the two heights of
\eqref{eq:pole-locations} where the split is the plain halving; the $ak$ apex
crosses once, between the tenth and eleventh.
Generated by \texttt{fig\_remainder\_frame\_sweep.py}.}
\label{fig:remainder-frame-sweep}
\end{figure}

To make the scaling precise, factor out the leading $n^{-\sigma}$ from each
half-remainder length by writing
\[
|R_{1\bullet}(\sigma,T)| \;=\; \kappa_{1\bullet}(\sigma,T)\,n^{-\sigma},
\qquad n=\lceil T\rceil,
\]
for the three first halves
$R_{1\bullet}\in\{R_{1rs},\,R_{1ps},\,R_{1ak}\}$ of
\S\ref{sec:remainders-summary}. Each split is then on the same footing:
$R=R_{1rs}+R_{2rs}=R_{1ps}+R_{2ps}\approx R_{1ak}+R_{2ak}$, and for the
$rs$ half-split one has $R_{1rs}=R_{2rs}=\tfrac12 R$ so
$\kappa_{2rs}=\kappa_{1rs}$. The three amplitudes admit the closed forms
\begin{equation}\label{eq:kappa}
\begin{aligned}
\kappa_{1rs}
&=
\tfrac12\,n^{\sigma}\,|R|
\;\approx\;
\tfrac14\,n^{\sigma}\,\bigl|I_{1}(s)+\chi(s)\,I_{2}(s)\bigr|,
\\[4pt]
\kappa_{1ps}
&=
d_{1}\,n^{\sigma}
=
2\kappa_{1rs}\,
\frac{\sin(\omega-\arg R+\psi)}{\sin(2\omega+\psi)},
\\[4pt]
\kappa_{1ak}
&=
\tfrac12\,n^{\sigma}\,|I_{1}(s)|,
\end{aligned}
\end{equation}
with $s=\sigma+iI(T)$, $\omega=I(T)\log n$, $\psi=\arg\chi(s)$, and
$I_1,I_2$ the Kuznetsov integrals of \S\ref{sec:kuznetsov}; the
``$\approx$'' in the first line holds when $R$ is evaluated by Kuznetsov's
method \eqref{eq:Rak}, and the second equality in the $\kappa_{1ps}$ line is
exact with $\kappa_{1rs}$ read as $\tfrac12 n^{\sigma}|R|$. Consequently, at
fixed $\sigma$ and fixed $\{T\}=f$, the passage from $T_1=N_1+f$ to
$T_2=N_2+f$ multiplies every half-remainder length by
\begin{equation}\label{eq:remainder-scale}
\frac{|R_{1\bullet}(\sigma,T_2)|}{|R_{1\bullet}(\sigma,T_1)|}
=
\Bigl(\frac{N_1+1}{N_2+1}\Bigr)^{\sigma}
\frac{\kappa_{1\bullet}(\sigma,T_2)}{\kappa_{1\bullet}(\sigma,T_1)}.
\end{equation}
When the fractional part is held fixed the ratio of $\kappa$'s is nearly $1$,
so the leading factor $\bigl((N_1+1)/(N_2+1)\bigr)^{\sigma}$ already accounts
for almost all of the size change. Figure~\ref{fig:remainder-scale-match}
checks this in the $R$-frame: the left panel is $T=6.18$; the right panel is
$T=50.18$ with all remainder vectors scaled by the exact factor
$\lambda=|R|_{6.18}/|R|_{50.18}$ from \eqref{eq:remainder-scale}
(equivalently $\kappa_{1rs,1}/\kappa_{1rs,2}$ times the $n^{-\sigma}$ ratio).
The two panels nearly coincide.

\begin{figure}[H]
\centering
\includegraphics[width=0.92\textwidth]{fig_remainder_scale_match}
\caption{Same fractional part $\{T\}=0.18$: $R$-frame remainders at
$T=6.18$ (left) and at $T=50.18$ after scaling by
$\lambda=\bigl(n_2/n_1\bigr)^{\sigma}(\kappa_{1rs,1}/\kappa_{1rs,2})$
(right). The scaled figure matches the reference to high accuracy
($|R_{1ps}|$ agrees to about $0.3\%$). Generated by
\texttt{fig\_remainder\_scale\_panels.py}.}
\label{fig:remainder-scale-match}
\end{figure}

% ======================================================================
\section{The positive real function \texorpdfstring{$d_1$}{d1} and its periodicity}\label{sec:d1-function}
% ======================================================================

The weights $d_1$ and $d_2$ of \S\ref{sec:decomp} are the quantitative heart
of the two remainders: $R_{1ps}=d_1e^{-i\omega}$ and
$R_{2ps}=d_2e^{\,i(\omega+\arg\chi)}$ are fixed in direction, so everything
the remainders do is carried by the two real functions
$d_1(\sigma,T)$ and $d_2(\sigma,T)$, positive on the critical line
(\S\ref{sec:decomp}). This section studies $d_1$ as a function
of the index $T$: its critical-line symmetry $d_1=d_2$, its behavior across
the integer-$T$ handoffs, and an asymptotic closed form. It then turns the
numerics into theorems: positivity, on the critical line without exception
and off it outside narrow windows (\S\ref{sec:d1-positive}); the exact
limiting waveform of the normalized fraction (\S\ref{sec:d1-limit}); and
uniform two-sided bounds valid at every height (\S\ref{sec:d1-bounds}).

\subsection{The exact formulas}\label{sec:d1-exact}

\begin{corollary}\label{cor:equal}
When $\sigma=\tfrac12$ the two fractional weights coincide, $d_1=d_2$, and
hence $|R_{1ps}| = |R_{2ps}|$: the two fractional summands are always equal
in length on the critical line.
\end{corollary}

This follows from the definitions \eqref{eq:d1}--\eqref{eq:d2} together with the
functional symmetry of $\chi$ on the critical line. Explicitly, when
$\sigma=\tfrac12$ the functional equation forces $|\chi|=1$ and
$\arg R = \tfrac12\arg\chi$ modulo $\pi$ (proved in
Appendix~\ref{app:lean-critical}). Writing $\psi=\arg\chi$ and
substituting $\arg R = \tfrac{\psi}{2}$ into \eqref{eq:d1}--\eqref{eq:d2}, the two
numerators collapse to the \emph{same} value:
\begin{equation}
d_1 = |R|\,\frac{\sin\!\bigl(\omega - \tfrac{\psi}{2} + \psi\bigr)}{\sin(2\omega+\psi)}
    = |R|\,\frac{\sin\!\bigl(\omega + \tfrac{\psi}{2}\bigr)}{\sin(2\omega+\psi)},
\qquad
d_2 = |R|\,\frac{\sin\!\bigl(\omega + \tfrac{\psi}{2}\bigr)}{\sin(2\omega+\psi)},
\end{equation}
so $d_1 = d_2$. Equivalently, setting $d_1=d_2$ in \eqref{eq:R-as-cone} gives
$R = d_1\bigl(e^{-i\omega}+e^{i(\omega+\psi)}\bigr)
   = 2d_1\cos\!\bigl(\omega+\tfrac{\psi}{2}\bigr)\,e^{\,i\psi/2}$,
which indeed has argument $\tfrac{\psi}{2}$ (or $\tfrac{\psi}{2}+\pi$, where the
cosine is negative), consistent with $\arg R=\tfrac12\arg\chi$ modulo $\pi$. We
have verified this in Lean; the proof is recorded in
Appendix~\ref{app:lean-critical}. We stress that this critical-line equality is a pleasant
consequence of the decomposition, not its motivation.

Since $d_1=d_2$ on the critical line, a single curve there tells the whole
story of both remainders, and it is worth drawing. Figure~\ref{fig:d1-critical}
plots $d_1(\tfrac12,T)$ for $1\le T\le7$, and the graph is unexpectedly
tame: as a function of the height $t$ the zeta landscape is famously erratic,
but viewed through the index $T$ the weight $d_1$ oscillates once per unit of
$T$, every unit interval the same shape, with only two blemishes: the
amplitude decays slowly, and the curve jumps at each integer $T$, where
$m=\lfloor T\rfloor$ increments and the fractional summands hand off to the
next link. The bottom panel shows that a single normalization (by the length
of the link that carries $d_1$) removes the decay and most of the jumps at
once, leaving a curve that looks nearly \emph{periodic} in $T$. That is a
strong hint that $d_1$ is, at heart, a fixed waveform in the fractional part
of $T$ dressed up by the link geometry. The rest of this section chases that
hint down, first by removing the jumps exactly, then by identifying the
waveform in closed form.

\begin{figure}[H]
\centering
\includegraphics[width=0.92\textwidth]{fig_d1_critical}
\caption{Top: the common fractional weight $d_1=d_2$ on the critical line,
plotted against the index $T$ for $1\le T\le 7$. The curve oscillates once per
unit of $T$ and jumps at each integer $T$, where $m=\lfloor T\rfloor$
increments and the fractional summands hand off to the next link; the
amplitude decays slowly as $T$ grows. The grey step is the length
$\lceil T\rceil^{-\sigma}$ of the link that carries the fractional summand,
which $d_1$ never exceeds; the panel keeps its original top, so the steps for
$1\le T<3$ lie at or above the frame. The double-headed arrow on $6<T<7$ spans
that link length. Bottom: the same formula for $d_1$
normalized by the length $\lceil T\rceil^{-\sigma}$ of link $m$, which
carries it: the \emph{normalized
distance from joint} $\lceil T\rceil^{\sigma}d_1$, the crossing fraction
along link $m$ at which the point $B_1=\Sigma_1+R_{1ps}$ sits (cf.\
\S\ref{sec:yinyang}). The panel is drawn on the full scale $0$ to $1$ of the
normalized link, so $0$ is joint $m$ and $1$ is joint $m+1$; the
double-headed arrow, at the same $T$ as the one above, marks that same link
seen at unit length. The
normalization removes both the decay and (nearly)
the jumps: the fraction sweeps the same range $\approx[0.23,\,0.78]$ in every
unit interval, staying clear of both ends of the link
(made exact in \S\ref{sec:d1-limit}--\S\ref{sec:d1-bounds}). So the normalized curve
converges to a periodic waveform: sampled at matched fractional parts of $T$,
its values agree to ${\approx}\,2\times10^{-3}$ by $T=20$, and the residual
discrepancy halves with each doubling of $T$, the rate set by the link-length
mismatch $(1+1/\lfloor T\rfloor)^{\sigma}\to1$ of \eqref{eq:fold} below.
Computed exactly
(mpmath) via the Cramer solution of
\eqref{eq:R-as-cone}. Generated by \texttt{fig\_d1\_critical.py}.}
\label{fig:d1-critical}
\end{figure}

\paragraph{Removing the jumps exactly.} The normalization in the bottom panel
of Figure~\ref{fig:d1-critical} nearly removes the jumps, and the residue can
be traced to a single exact relation. At a handoff $T=n$ the outgoing link
($m=n-1$) and the incoming link ($m=n$) fold back onto one another while the
point $B_1$
moves continuously in the plane, which forces
\begin{equation}\label{eq:fold}
d_1^{+} \;=\; n^{-\sigma} \;-\; d_1^{-},
\end{equation}
where $d_1^{-}$ and $d_1^{+}$ are the limits from the left and from the right.
In terms of the plotted fraction $f=\lceil T\rceil^{\sigma} d_1$ this reads
$f^{+}=(1+\tfrac1n)^{\sigma}(1-f^{-})$: an orientation flip $f\mapsto 1-f$
(after the handoff the fraction is measured from the other end of the link)
composed with a unit mismatch (the new link is shorter than the old by the
factor $(1+\tfrac1n)^{-\sigma}$). The mismatch factor tends to $1$, which is
why the jumps in the bottom panel shrink as $T$ grows; the replacement
$f\mapsto 1-f$ also contributes less and less, because the curve happens to
cross the integers near $f\approx\tfrac12$, where $1-f\approx f$.

Because \eqref{eq:fold} is exact, the unit intervals can be glued into a
coordinate with no jumps at all: measure the position of $B_1$
along the link chain in one \emph{fixed} unit (the first link, of length
$1$), flipping the orientation and absorbing the constant $n^{-\sigma}$ at
each handoff of the bisector point from one link to the next (see
\S\ref{sec:IT-origin} for more on these handoffs). Unrolling the
recursion gives
\begin{equation}\label{eq:h-cont}
h(T) \;=\; \sum_{k=2}^{\lfloor T\rfloor} (-1)^{k}\,k^{-\sigma}
\;+\; (-1)^{\lfloor T\rfloor+1}\, d_1(T),
\end{equation}
and at every integer the term entering the sum cancels the jump of $d_1$
identically, so $h$ is continuous everywhere. The price is that $h$
\emph{flattens}: in a fixed unit the links themselves shrink, so the
oscillating part $d_1=O(T^{-\sigma})$ decays, while the accumulated constants
are the partial sums of the convergent alternating series
$2^{-\sigma}-3^{-\sigma}+4^{-\sigma}-\cdots = 1-\eta(\sigma)$, with $\eta$ the
Dirichlet eta function. Hence $h(T)\to 1-\eta(\sigma)$, which at
$\sigma=\tfrac12$ is ${\approx}\,0.3951$.

Continuity and a non-decaying amplitude can be had at once by re-zooming
\emph{smoothly}, with $T^{\sigma}$, instead of with the step function
$\lceil T\rceil^{\sigma}$ (whose abrupt changes at the integers are exactly
what created the residual jumps):
\begin{equation}\label{eq:p-cont}
p(T) \;=\; T^{\sigma}\,\Bigl( h(T) - \bigl(1-\eta(\sigma)\bigr) \Bigr).
\end{equation}
As a product of continuous functions $p$ is exactly continuous, and since
$h(T)-(1-\eta(\sigma)) = O(T^{-\sigma})$ (the decaying $d_1$ plus the tail
of the alternating series), the factor $T^{\sigma}$ restores a bounded
oscillation. Figure~\ref{fig:h-p-continuous} compares the two coordinates:
both are continuous, with only kinks (no jumps) at the integers, while $h$
contracts toward its limit $1-\eta(\tfrac12)$ and $p$ keeps a steady
amplitude. The red curve is computed from the simplified local form
\eqref{eq:p-simplified}, derived below, which agrees with \eqref{eq:p-cont}
to machine precision.

\begin{figure}[H]
\centering
\includegraphics[width=0.92\textwidth]{fig_h_p_continuous}
\caption{The two exactly continuous coordinates of the point $B_1$ on the
critical line, $1\le T\le 7$, in a single panel for comparison. Blue:
$h(T)$ of \eqref{eq:h-cont}, the position along the link in a fixed
unit; it has no jumps at the integers, only kinks where the links fold
back and handoff the bisector point to the next link,
and its oscillation contracts toward the limit $1-\eta(\tfrac12)\approx
0.3951$ (dashed line): the swing shrinks from ${\approx}\,0.40$ on the first
interval to ${\approx}\,0.24$ by $T=7$. Red: $p(T)$, the same coordinate
re-zoomed smoothly by $T^{\sigma}$, computed from the simplified local form
\eqref{eq:p-simplified}; it is equally continuous but its amplitude holds
steady (${\approx}\,0.53$ peak to peak)
instead of flattening. Continuity was verified numerically at each integer
(mismatch ${\sim}\,10^{-7}$ at offsets of $10^{-6}$). Generated by
\texttt{fig\_h\_p\_continuous.py}.}
\label{fig:h-p-continuous}
\end{figure}

\paragraph{A pinned waveform.} Each arch of $p$ in
Figure~\ref{fig:h-p-continuous} is nearly the vertical flip of its
neighbours, but the arches do not begin and end at zero, so flipping alternate
intervals would not make them line up. Taking one-sided limits at an integer
$n$ (where $\lfloor T\rfloor=n-1$ on the left) gives the endpoint value
exactly:
\begin{equation}\label{eq:eps-n}
p(n)=(-1)^{n}n^{\sigma}\bigl(d_1(n^-)-\Phi(-1,\sigma,n)\bigr)
=(-1)^{n+1}\varepsilon_n,
\qquad
\varepsilon_n:=n^{\sigma}\bigl(\Phi(-1,\sigma,n)-d_1(n^-)\bigr)>0,
\end{equation}
so $\varepsilon_n$ is the shortfall of the distance at when the links
are folded back from the alternating-tail center of the link
\begin{equation}\label{eq:alt-tail-center}
n^{\sigma}\Phi(-1,\sigma,n)
=\tfrac12+\frac{\sigma}{4n}-\frac{\sigma(\sigma+1)(\sigma+2)}{48\,n^{3}}
+O(n^{-5})
=\tfrac12+\frac{1}{8n}-\frac{5}{128\,n^{3}}+O(n^{-5})
\quad\text{at }\sigma=\tfrac12,
\end{equation}
by Boole summation, $\sum_{j\ge0}(-1)^{j}f(n+j)=\tfrac12f(n)-\tfrac14f'(n)
+\tfrac1{48}f'''(n)-\cdots$, whose even derivatives drop out, so there is no
$n^{-2}$ term. On the critical line the
ends approach zero on the order of $T^{-2}$, without ever reaching it at
finite $T$.

Flipping alternate intervals costs nothing: it simply deletes
the alternating sign from \eqref{eq:p-simplified}, since with
$m=\lfloor T\rfloor$
\begin{equation}\label{eq:p-flipped}
P(T)\;:=\;(-1)^{m}p(T)\;=\;T^{\sigma}\bigl(\Phi(-1,\sigma,m+1)-d_1(T)\bigr),
\end{equation}
whose two ends on $[m,m+1]$ are exactly $-\varepsilon_m$ and
$+\varepsilon_{m+1}$ by \eqref{eq:eps-n}, so $P$ still jumps by
$2\varepsilon_n$ at each integer. Subtracting the chord between those two
values pins both ends: with $x=T-m$,
\begin{equation}\label{eq:pinned-waveform}
\mathcal{W}(T)\;=\;T^{\sigma}\bigl(\Phi(-1,\sigma,m+1)-d_1(T)\bigr)
\;+\;(1-x)\,\varepsilon_{m}\;-\;x\,\varepsilon_{m+1},
\end{equation}
a single formula covering the intervals $[1,2],[3,4],\dots$ and
$[2,3],[4,5],\dots$ with the second family already flipped.
Figure~\ref{fig:pinned-waveform} draws it: $\mathcal{W}$ vanishes at every
integer, is continuous on $[1,\infty)$ with no sign alternation left, and
converges like $1/T$ to the tangent waveform
$\mathcal{W}_{\infty}(q)=\tfrac12\tan(2\pi q)
\tan\bigl(2\pi(q-\tfrac14)(q-\tfrac34)\bigr)$, $q=\{T\}$, which is the
$T\to\infty$ form of the Riemann--Siegel first term \eqref{eq:d1-rs} below and
does vanish at $q=0,1$ exactly. (\S\ref{sec:d1-limit} proves the convergence
and shows that the same waveform, recentered at $\tfrac12$ and flipped, is
the exact limit profile of the normalized weight $\hat d_1$ itself.) Continuity is had exactly; $C^{1}$ only in the
limit. The one-sided slopes at $T=n$ are $7.515$ and $6.920$ at $n=2$ and
$7.482$ and $7.309$ at $n=8$, both approaching
$\mathcal{W}_{\infty}'(0)=\mathcal{W}_{\infty}'(1)=\pi\tan\tfrac{3\pi}{8}
=7.5845$, so a corner of size $O(1/T)$ survives at each integer. Replacing the
chord in \eqref{eq:pinned-waveform} by the cubic Hermite that also matches the
end slopes removes the corner exactly, at the cost of inflating each arch by
about $0.03$; and since
$\mathcal{W}_{\infty}''(0)=-\mathcal{W}_{\infty}''(1)$, no pinning can do
better than $C^{1}$ in any case.

\begin{figure}[H]
\centering
\includegraphics[width=\textwidth]{fig_pinned_waveform}
\caption{The flipped arches of $p$ on the critical line, $1\le T\le7$.
\emph{Left:} the unpinned flip $P$ of \eqref{eq:p-flipped} (pale blue), which
jumps by $2\varepsilon_n$ at each integer (inset, at $T=2$); the chord-pinned
$\mathcal{W}$ of \eqref{eq:pinned-waveform} (red), which is zero there
exactly; and the Hermite-pinned variant (dashed green), which also matches the
end slopes but inflates the arch. \emph{Right:} the six pinned arches drawn
against the fractional part $q=\{T\}$, with the tangent limit
$\mathcal{W}_{\infty}$ dashed; the deviation falls from $0.087$ on $[1,2]$ to
$0.025$ on $[6,7]$. Generated by \texttt{fig\_pinned\_waveform.py}.}
\label{fig:pinned-waveform}
\end{figure}

\subsection{Approximation of \texorpdfstring{$d_1$}{d1} when \texorpdfstring{$\sigma=\tfrac12$}{sigma=1/2}}\label{sec:d1-approx-critical}

Two things happen here, in order: $p$ is first put into a purely local form, in
which the continuous coordinate is nothing but the deviation of $d_1$ from a
reference value, and then, on the critical line, $d_1$ itself is given a
zeta-free closed form.

Distributing $T^{\sigma}$ and substituting \eqref{eq:h-cont} simplifies $p$
considerably: the partial sum inside $h$ cancels the head of the series
$1-\eta(\sigma)=\sum_{k=2}^{\infty}(-1)^{k}k^{-\sigma}$ term by term, leaving
only its tail, and the common sign $(-1)^{\lfloor T\rfloor+1}$ factors out:
\begin{equation}\label{eq:p-simplified}
p(T) \;=\; (-1)^{\lfloor T\rfloor+1}\, T^{\sigma}
\Bigl[\, d_1(T) \;-\; \Phi\bigl(-1,\sigma,\lfloor T\rfloor+1\bigr) \Bigr],
\qquad
\Phi(-1,\sigma,a) = \sum_{j=0}^{\infty} \frac{(-1)^{j}}{(a+j)^{\sigma}},
\end{equation}
with $\Phi$ the Lerch transcendent. In this form all the accumulated history
of \eqref{eq:h-cont} has disappeared and $p$ is \emph{local}: it is just the
deviation of $d_1$ from a reference value, the alternating sum of all the
remaining link lengths
$(\lfloor T\rfloor+1)^{-\sigma}-(\lfloor T\rfloor+2)^{-\sigma}+\cdots$.
An alternating series is roughly half its first term,
$\Phi(-1,\sigma,\lfloor T\rfloor+1)\approx\tfrac12(\lfloor T\rfloor+1)^{-\sigma}$,
which is half the length of link $m$, so asymptotically
\[
p(T) \;\approx\; (-1)^{\lfloor T\rfloor+1}
\Bigl( \lceil T\rceil^{\sigma} d_1 - \tfrac12 \Bigr),
\]
which is the normalized fraction of Figure~\ref{fig:d1-critical}, recentered
at $\tfrac12$ and given alternating orientation. This is also why that fraction
crosses the integers near $\tfrac12$: continuity of $p$ pins the fraction, at
each integer, near the alternating-tail center, which is close to $\tfrac12$.

On the critical line the substitution can be pushed all the way:
$d_1$ itself acquires a zeta-free closed form. When $\sigma=\tfrac12$ we
have $\arg R=\tfrac12\arg\chi$, so the Cramer solution
\eqref{eq:d1} reduces to
\[
d_1 \;=\; \frac{\widetilde R}{2\cos\bigl(\omega+\tfrac{\psi}{2}\bigr)},
\qquad
\widetilde R = e^{-i\psi/2} R \ \ \text{(real)},
\]
the two-equal-vector form of $R$ seen in the proof of
Corollary~\ref{cor:equal}. For $\widetilde R$ we substitute the saddle-point
evaluation of Siegel's integral, which is the first correction term of the
Riemann--Siegel formula,
\[
\widetilde R
= (-1)^{N-1}\Bigl(\tfrac{2\pi}{t}\Bigr)^{1/4} C_0(\hat p)
+ O\bigl(t^{-3/4}\bigr),
\qquad
C_0(\hat p) =
\frac{\cos\bigl(2\pi(\hat p^{2}-\hat p-\tfrac{1}{16})\bigr)}{\cos(2\pi \hat p)},
\qquad \sqrt{t/2\pi} = N + \hat p .
\]
Every quantity on the right is an elementary function of $t=I(T)$, kept
exact throughout:
\[
z=\sqrt{t/2\pi},\qquad N=\lfloor z\rfloor,\qquad \hat p = z-N,\qquad
\omega = t\ln(m+1),\qquad m=\lfloor T\rfloor .
\]
On the line $\chi=e^{-2i\vartheta(t)}$ with $\vartheta$ the Riemann--Siegel
theta function, so $\widetilde R$ is $e^{i\vartheta}R$ up to a sign that
cancels against the cosine, and
$d_1 = e^{i\vartheta}R\,/\,2\cos(\omega-\vartheta)$ \emph{exactly}. Substituting
the Riemann--Siegel first term for $e^{i\vartheta}R$ gives
\begin{equation}\label{eq:d1-rs}
d_1\bigl(\tfrac12,T\bigr) \;\approx\;
\bigl[N=m{+}1\bigr]\,(m{+}1)^{-1/2}
\;+\;
(-1)^{N-1}\Bigl(\tfrac{2\pi}{t}\Bigr)^{1/4}
\frac{C_0(\hat p)}{2\cos\bigl(\omega-\vartheta(t)\bigr)},
\end{equation}
where the Iverson bracket adds one full link length when $N=m+1$,
that is, when the Riemann--Siegel main sum carries one more summand pair
than our $\Sigma_1,\Sigma_2$, which happens (roughly) when
$\operatorname{frac}(T)>\tfrac12$. The formula is elementary and
zeta-free: $\vartheta$ may be replaced by its asymptotic
$\vartheta(t)=\tfrac t2\ln\tfrac{t}{2\pi}-\tfrac t2-\tfrac\pi8+\tfrac1{48t}$
at no cost (the two agree here to $4\times10^{-7}$ even at $T=1$). Its
maximum error is $0.006$ on $[1,2)$, $0.002$ on $[3,4)$, $0.0009$ on
$[6,7)$ and $0.0002$ on $[20,21)$. That is about $1\%$ of the swing of $d_1$
at the start, decaying like $T^{-3/2}$ and spread flat across each unit
interval (Figure~\ref{fig:d1-rs-phases}). The error is purely the
Riemann--Siegel truncation, $O(t^{-3/4})$; the next correction term of
that formula would push further. And since
$R = 2d_1\cos(\omega+\tfrac{\psi}{2})\,e^{i\psi/2}$ on the line,
\eqref{eq:d1-rs} is effectively a closed-form approximation of Siegel's
remainder there as well.

\begin{figure}[H]
\centering
\includegraphics[width=0.92\textwidth]{fig_d1_rs_phases}
\caption{The closed form \eqref{eq:d1-rs}: the Riemann--Siegel first term
with the phases kept exact through $t=I(T)$, no zeta input. Top: exact
$d_1$ (blue) against \eqref{eq:d1-rs} (black dashed), visually
indistinguishable, $1\le T\le 7$. Bottom: $|$error$|$ on a log scale, with
maximum $0.006$ on $[1,2)$, decaying like $T^{-3/2}$, spread flat in the
fractional part. Generated by \texttt{fig\_d1\_rs\_phases.py}.}
\label{fig:d1-rs-phases}
\end{figure}

\subsection{Approximation of \texorpdfstring{$d_1$}{d1} and \texorpdfstring{$d_2$}{d2} when \texorpdfstring{$\sigma\neq\tfrac12$}{sigma != 1/2}}\label{sec:d1-approx-general}

There is a general-$\sigma$ approximation as well, and it needs no new
machinery beyond what is already in the paper. Off the critical line the
remainder $R$ is genuinely complex (no single rotation makes it real),
so approximating it means approximating two real degrees of freedom,
and those two degrees of freedom are exactly the pair $(d_1,d_2)$.

Siegel's saddle-point analysis was never restricted to the critical line;
the first term for the remainder at general $s=\sigma+it$ splits into one
piece per chain:
\begin{equation}\label{eq:R-general}
R(s) \;\approx\; (-1)^{N-1}\,\frac{C_0(\hat p)}{2}
\Bigl[\,a^{-\sigma}\,e^{-i\tilde\vartheta(t)}
\;+\; \chi(s)\,a^{\sigma-1}\,e^{+i\tilde\vartheta(t)}\Bigr],
\qquad a=\sqrt{t/2\pi},
\end{equation}
with $t=I(T)$ exact, $N=\lfloor a\rfloor$, $\hat p=a-N$, the same
universal $C_0(\hat p)$ of \S\ref{sec:d1-approx-critical}, and the elementary
phase $\tilde\vartheta(t)=\tfrac t2\ln\tfrac{t}{2\pi}-\tfrac t2-\tfrac\pi8$.
When $N=m+1$ the extra summand pair $(m+1)^{-s}+\chi\,(m+1)^{s-1}$ is
added, exactly as in \eqref{eq:d1-rs}. Feeding this approximate $R$ to
the Cramer solution \eqref{eq:d1}--\eqref{eq:d2} with the exact
$\omega=t\ln(m+1)$ and $\psi=\arg\chi$ then yields \emph{both} $d_1$ and
$d_2$ at once, still with no zeta input ($\chi$ is Gamma factors). The
two bracket pieces are the per-chain remainders, with the same structure as
the Kuznetsov split of \S\ref{sec:kuznetsov}, and at $\sigma=\tfrac12$,
where $\chi=e^{-2i\vartheta}$, the bracket collapses to $2e^{-i\vartheta}$:
\eqref{eq:R-general} reduces exactly to \eqref{eq:d1-rs}. We verified
this numerically digit for digit, which pins down all the sign
conventions.

Because the Cramer solution is linear in $R$, the solve can be carried
out once and for all: substituting \eqref{eq:R-general} into
\eqref{eq:d1}--\eqref{eq:d2} gives the general-$\sigma$ analogue of
\eqref{eq:d1-rs} in closed form,
\begin{align}
d_1(\sigma,T) &\approx
\bigl[N=m{+}1\bigr]\,(m{+}1)^{-\sigma}
\notag\\&\quad
+ (-1)^{N-1}\,\frac{C_0(\hat p)}{2\sin(2\omega+\psi)}
\Bigl[\,a^{-\sigma}\sin\bigl(\omega+\psi+\tilde\vartheta\bigr)
+ |\chi|\,a^{\sigma-1}\sin\bigl(\omega-\tilde\vartheta\bigr)\Bigr],
\label{eq:d1-general}\\[6pt]
d_2(\sigma,T) &\approx
\bigl[N=m{+}1\bigr]\,|\chi|\,(m{+}1)^{\sigma-1}
\notag\\&\quad
+ (-1)^{N-1}\,\frac{C_0(\hat p)}{2\sin(2\omega+\psi)}
\Bigl[\,a^{-\sigma}\sin\bigl(\omega-\tilde\vartheta\bigr)
+ |\chi|\,a^{\sigma-1}\sin\bigl(\omega+\psi+\tilde\vartheta\bigr)\Bigr],
\label{eq:d2-general}
\end{align}
with $\tilde\vartheta=\tilde\vartheta(t)$, $\psi=\arg\chi$, and
$|\chi|=|\chi(\sigma+it)|$ as before. The Iverson terms are now whole
link lengths on each side separately, because the extra summand pair
lies exactly along the two unit directions: $(m{+}1)^{-s}$ along
$e^{-i\omega}$ and $\chi\,(m{+}1)^{s-1}$ along $e^{i(\omega+\psi)}$.
Note the symmetry: swapping the two coefficients
$a^{-\sigma}\leftrightarrow|\chi|\,a^{\sigma-1}$ swaps $d_1$ and $d_2$.
At $\sigma=\tfrac12$, where $|\chi|=1$ and $\psi=-2\vartheta$, the two
sine brackets coincide and each formula collapses to \eqref{eq:d1-rs},
since $\sin(\omega-\vartheta)/\sin\bigl(2(\omega-\vartheta)\bigr)
= 1/\,2\cos(\omega-\vartheta)$ and $a^{-1/2}=(2\pi/t)^{1/4}$. We verified
that \eqref{eq:d1-general}--\eqref{eq:d2-general} reproduce the
Cramer-solve route to machine precision; they are the same
computation written out.

The accuracy across the strip (maximum $|d_1|$ error away from the pole
windows described below):
\begin{center}
\begin{tabular}{c|ccc}
$\sigma$ & $[1,2)$ & $[6,7)$ & $[20,21)$\\
\hline
$0.1$ & $0.076$ & $0.016$  & $0.0046$\\
$0.3$ & $0.035$ & $0.0057$ & $0.0013$\\
$0.5$ & $0.006$ & $0.0009$ & $0.0002$\\
$0.7$ & $0.025$ & $0.0024$ & $0.0004$\\
$0.9$ & $0.045$ & $0.0032$ & $0.0004$
\end{tabular}
\end{center}
The measured decay exponents match $O(T^{-\sigma-1})$ almost perfectly
($1.07$ at $\sigma=0.1$, $1.5$ at $\sigma=0.5$, $1.8$ at $\sigma=0.9$),
and since $d_1$ itself scales like $T^{-\sigma}$, the \emph{relative}
error is $O(1/T)$ uniformly across the strip. The critical line is where
the formula is at its absolute best, fittingly. And $d_2$ comes for free:
the Cramer solve produces both components simultaneously from the same
approximate $R$, and its errors are essentially identical to $d_1$'s
(e.g.\ $0.031$ vs $0.035$ at $\sigma=0.3$ on $[1,2)$), so no separate
derivation is needed. Figure~\ref{fig:d1-d2-general} shows both
components at $\sigma=0.3$ and $\sigma=0.7$.

Two caveats. First, off the line the exact $d_1$ and $d_2$ have genuine
narrow poles at the parallel-link heights $q\approx\tfrac14,\tfrac34$
(the bands of Figure~\ref{fig:equal-legs-strips}); the
first-term approximation smooths through those spikes, so within about
$\pm0.05$ of them the error is locally large (relative error
${\sim}\,40\%$ right at the shoulder). Everywhere else the table
applies. Second, the formula is not limited to the strip: it works for
any $\sigma$ with relative error $O(1/T)$. For $\sigma>1$ it is
excellent, while for $\sigma\le0$ the absolute error $T^{-\sigma-1}$
decays ever more slowly (at $\sigma=-\tfrac12$ it is still
${\approx}\,0.06$ at $T=20$), so ``any $\sigma$, relative error $1/T$''
is the honest statement, with the strip as the sweet spot.

\begin{figure}[H]
\centering
\includegraphics[width=0.92\textwidth]{fig_d1_d2_general_sigma}
\caption{$d_1$ (red) and $d_2$ (green) across the strip at $\sigma=0.3$
(top) and $\sigma=0.7$ (bottom), $1\le T\le 7$: exact values (solid)
against the first-term approximation \eqref{eq:R-general} fed through the
Cramer solution (dashed, in the color of its exact counterpart). The
narrow spikes leaving the frame near fractional heights $\tfrac14$ and
$\tfrac34$ are the genuine off-the-critical-line poles of $d_1$ and $d_2$ (parallel
links); the approximation smooths through them. Generated by
\texttt{fig\_d1\_d2\_general\_sigma.py}.}
\label{fig:d1-d2-general}
\end{figure}

\subsection{Where are the \texorpdfstring{$d_1$}{d1} and \texorpdfstring{$d_2$}{d2} poles?}\label{sec:pole-locations}

The poles of $d_1$ and $d_2$ (equivalently of $R_{1ps}$ and $R_{2ps}$)
appear to be at fractional heights $\tfrac14$ and $\tfrac34$, but they
are not exactly there. Here we delve into where they are.

A pole occurs when the two bisector links are parallel; these are link
$m=\lfloor T\rfloor$ of each chain, each carrying summand $\lceil T\rceil$,
the summand just after the last one included in the partial sum. In other
words, a pole occurs when the two next summands have the same argument, so that the
cone spanned by the two unit directions $e^{-i\omega}$ and
$e^{i(\omega+\arg\chi)}$ of \eqref{eq:R-as-cone} degenerates to a line
and the common denominator $\sin(2\omega+\arg\chi)$ of
\eqref{eq:d1}--\eqref{eq:d2} vanishes. Subtracting the two arguments
from each other produces a formula that is zero at the poles:
\begin{equation}\label{eq:pole-locations}
2\,\omega(T) \;+\; \arg\chi\bigl(\sigma+i\,I(T)\bigr) \;=\; 2\pi k,
\qquad k\in\mathbb{Z},
\end{equation}
with $\omega(T)=I(T)\log\lfloor T+1\rfloor$ as in \S\ref{sec:decomp}.
This formula was used to find the locations of the first $20$ poles
numerically (Table~\ref{tab:poles}).

\begin{table}[h!]
\centering
\begin{tabular}{|c|c|c|c|}
\hline
\textbf{Odd Pole \#} & \textbf{$T$ (odd)} & \textbf{Even Pole \#} & \textbf{$T$ (even)} \\ \hline
1  & 1.268870258145  & 2  & 1.763736013695 \\
3  & 2.261710393365  & 4  & 2.759511841350 \\
5  & 3.258504081544  & 6  & 3.757283144824 \\
7  & 4.256679717535  & 8  & 4.755902912218 \\
9  & 5.255501039250  & 10 & 5.754963311284 \\
11 & 6.254676427788  & 12 & 6.754282121179 \\
13 & 7.254067037740  & 14 & 7.753765522842 \\
15 & 8.253598276826  & 16 & 8.753360246216 \\
17 & 9.253226472440  & 18 & 9.753033785343 \\
19 & 10.252924347790 & 20 & 10.752765174160 \\
\hline
\end{tabular}
\caption{First $20$ poles for $\sigma=0.6$, which is the same result as
for $\sigma=0.4$. The symmetry about $\sigma=\tfrac12$ exists because
$\arg\chi(\sigma+it)=-\arg\chi(1-\sigma+it)$ by the identity
$\chi(s)\,\chi(1-s)=1$, and the pole condition
\eqref{eq:pole-locations} uses $\arg\chi$ only as a phase offset;
reflecting $\sigma\mapsto1-\sigma$ therefore yields the same
$T$-solutions. Note the poles occur at approximately
$T=\lfloor T\rfloor+\tfrac14$ and $T=\lfloor T\rfloor+\tfrac34$
(equivalently, an integer $\pm\tfrac14$), and approach those values as
$T\to\infty$.}
\label{tab:poles}
\end{table}

\subsection{The weights \texorpdfstring{$d_1,d_2$}{d1, d2} are positive outside narrow windows of \texorpdfstring{$\{T\}$}{the fractional part of T}}\label{sec:d1-positive}

So far positivity has rested on numerical evidence: the note after
\eqref{eq:R-as-cone} records a dense verification on the critical line, and
\S\ref{sec:pole-locations} locates the off-line poles where the property
fails. Everything needed for a proof is now on the table: the exact
critical-line reduction of \S\ref{sec:d1-approx-critical}, the
Riemann--Siegel first term made rigorous by Gabcke's error bound, and one
elementary phase expansion carried out below. This subsection gives the
proofs, first on the critical line, where no fractional parts need to be
excluded at all, then across the strip, where positivity holds outside two
windows of width $O(1/T)$ around the poles of \S\ref{sec:pole-locations}.

We keep the notation of \S\ref{sec:d1-approx-critical}: $a=\sqrt{t/2\pi}$,
$N=\lfloor a\rfloor$, $\hat p=a-N$, and $u=\omega-\vartheta(t)$. Expanding
$\log(1+\tfrac1T)$ in $t=I(T)$ gives
\begin{equation}\label{eq:a-vs-T}
\frac{t}{2\pi}=\Bigl(T+\tfrac12\Bigr)^{2}+O\!\Bigl(\tfrac1T\Bigr),
\qquad\text{hence}\qquad
a=T+\tfrac12+O\!\Bigl(\tfrac1{T}\Bigr),
\end{equation}
so, apart from boundary slivers of width $O(1/T)$ in the fractional part
$x=\{T\}$ near $x=0,\tfrac12,1$,
\begin{equation}\label{eq:cut-cases}
x<\tfrac12 \iff N=m,\ \ \hat p=x+\tfrac12+O(1/T),
\qquad
x>\tfrac12 \iff N=m+1,\ \ \hat p=x-\tfrac12+O(1/T).
\end{equation}
Taylor expansion of $\omega=t\log(m+1)$ about $\log a$, using $t=2\pi a^{2}$
and the asymptotic series of $\vartheta$, gives the phase $u$ modulo $2\pi$
in closed form, and in both cases of \eqref{eq:cut-cases} the result is one
and the same statement:
\begin{equation}\label{eq:cosu}
\cos u \;=\;(-1)^{m+1}\,
\cos\!\bigl(2\pi x(1-x)-\tfrac{3\pi}8\bigr)\;+\;O(1/T).
\end{equation}
Since $2\pi x(1-x)\in[0,\tfrac\pi2]$ for $x\in[0,1]$, the argument of the
cosine on the right lies in $(-\tfrac{3\pi}8,\tfrac\pi8]$, so
\begin{equation}\label{eq:cosu-bound}
|\cos u|\;\geq\;\sin\tfrac\pi8-O(1/T)\;=\;0.3826\ldots-O(1/T):
\end{equation}
the factor $\cos u$ never approaches zero, at any fractional part. That one
inequality is the heart of both positivity theorems. We also need two
elementary facts about the Riemann--Siegel amplitude $C_0$ of
\S\ref{sec:d1-approx-critical}.

\begin{lemma}\label{lem:C0-trig}
Write $A=2\pi\hat p^{2}-\tfrac\pi8$ and $B=2\pi\hat p$, so that
$C_0(\hat p)=\cos(A-B)/\cos B$. Then:
\begin{enumerate}
\item[(i)] $C_0(\hat p)>0$ for $\hat p\in(\tfrac12,1)$, and
$\min_{[1/2,\,1]}C_0=\cos\tfrac{3\pi}8=\sin\tfrac\pi8$, attained at
$\hat p=\tfrac12$;
\item[(ii)] for $\hat p\in[0,\tfrac12]$,
\[
\Lambda(\hat p)\;:=\;\frac{C_0(\hat p)}
{2\cos\bigl(\pi(\tfrac18-2\hat p^{2})\bigr)}
\;\leq\;\tfrac12,
\]
with equality exactly at the endpoints, and $\Lambda(\tfrac14)=\tfrac14$.
\end{enumerate}
\end{lemma}

\begin{proof}
(ii) first, since it is the prettier half. Note
$\pi(\tfrac18-2\hat p^{2})=-A$, so
$\Lambda=\cos(A-B)/(2\cos A\cos B)=\tfrac12(1+\tan A\tan B)$ by the product
formula, and the claim is that $\tan A\tan B\leq0$ on $[0,\tfrac12]$. Both
critical transitions happen at the same point: $A=0$ exactly when
$\hat p=\tfrac14$, and $B=\tfrac\pi2$ exactly when $\hat p=\tfrac14$. For
$\hat p<\tfrac14$ we have $A\in(-\tfrac\pi8,0)$, so $\tan A<0$, while
$B\in(0,\tfrac\pi2)$, so $\tan B>0$; for $\hat p>\tfrac14$ the signs
reverse ($A\in(0,\tfrac{3\pi}8)$, $B\in(\tfrac\pi2,\pi)$). The product is
therefore $\leq0$ throughout, vanishing only as $\hat p\to0,\tfrac12$
(where $\tan A\tan B\to0$, so $\Lambda\to\tfrac12$); at $\hat p=\tfrac14$
the two singular factors balance, $\tan A\tan B\to-\tfrac12$, giving the
removable value $\Lambda(\tfrac14)=\tfrac14$.

(i) is a sign inspection of numerator and denominator. On
$(\tfrac12,\tfrac34)$ both $\cos(A-B)$ and $\cos B$ are negative; on
$(\tfrac34,1)$ both are positive (they change sign together at
$\hat p=\tfrac34$, which is the removable singularity of $C_0$), so
$C_0>0$ throughout. The minimum location was found numerically
($10^{4}$-point grid, minimum $0.38268\ldots=\cos\tfrac{3\pi}8$ at
$\hat p=\tfrac12$, the left endpoint, where
$C_0(\tfrac12)=-\cos\tfrac{5\pi}8$ exactly).
\end{proof}

\begin{theorem}\label{thm:d1-positive-line}
There is an effectively computable $T_0$ such that for every non-integer
$T\geq T_0$, at $\sigma=\tfrac12$,
\[
d_1(T)=d_2(T)\;\geq\;\frac{c}{\,a^{1/2}}\;>\;0,
\qquad c=\tfrac12\sin\tfrac\pi8-o(1).
\]
No interval of $\{T\}$ is excluded. Together with the dense numerical
verification of \S\ref{sec:decomp} covering $1\leq T\leq30$, the weights are
positive at every non-integer $T\geq1$ on the critical line.
\end{theorem}

\begin{proof}
\emph{Step 1 (exact reduction).} On $\sigma=\tfrac12$ the reflection
identities $\Sigma_2=\chi\overline{\Sigma_1}$ and $\zeta=\chi\overline\zeta$
give $R=\chi\overline R$; with $\chi=e^{-2i\vartheta}$ this says
$r:=e^{i\vartheta}R$ is real, and, as recorded in
\S\ref{sec:d1-approx-critical} (and formalized in
Appendix~\ref{app:lean-critical}),
\begin{equation}\label{eq:exact-line}
d_1=d_2=\frac{r}{2\cos u}
\qquad(\cos u\neq0),
\end{equation}
the common factor $\sin u$ having cancelled between the numerators of
\eqref{eq:d1}--\eqref{eq:d2} and the denominator
$\sin2u=2\sin u\cos u$.

\emph{Step 2 (Riemann--Siegel input).} The real number $r$ is the
Riemann--Siegel tail cut at $m$ instead of at $N$: writing
$Z=e^{i\vartheta}\zeta$,
\begin{equation}\label{eq:r-two-cases}
r=r_{\mathrm{RS}}+\bigl[N=m+1\bigr]\cdot\frac{2\cos u}{\sqrt{m+1}},
\qquad
r_{\mathrm{RS}}:=Z-2\!\!\sum_{n\leq N}\!n^{-1/2}\cos(\vartheta-t\log n),
\end{equation}
the bracketed term being the summand $n=m+1$, whose phase is exactly
$\vartheta-\omega=-u$; this is the same Iverson bookkeeping as in
\eqref{eq:d1-rs}. Gabcke's rigorous form of the Riemann--Siegel formula
\cite{Gabcke1979} gives
\begin{equation}\label{eq:gabcke}
r_{\mathrm{RS}}=\frac{(-1)^{N-1}}{a^{1/2}}\Bigl(C_0(\hat p)+E\Bigr),
\qquad |E|\leq c_1\,a^{-1}\quad(t\geq200).
\end{equation}

\emph{Step 3 (case $N=m$, i.e.\ $x=\{T\}<\tfrac12$,
$\hat p\in(\tfrac12,1)$).} Here $r=r_{\mathrm{RS}}$, and the two signs
multiply out: $(-1)^{N-1}$ from \eqref{eq:gabcke} against $(-1)^{m+1}$ from
\eqref{eq:cosu} give $+1$, so
\[
d_1=\frac{C_0(\hat p)}
{2\cos\bigl(2\pi x(1-x)-\tfrac{3\pi}8\bigr)}\cdot\frac{1}{a^{1/2}}
\;+\;O(a^{-3/2}).
\]
By Lemma~\ref{lem:C0-trig}(i) the numerator is at least $\sin\tfrac\pi8$,
and the cosine in the denominator is at most $1$, so
$d_1\geq\bigl(\tfrac12\sin\tfrac\pi8\bigr)a^{-1/2}-O(a^{-3/2})$.

\emph{Step 4 (case $N=m+1$, i.e.\ $x>\tfrac12$,
$\hat p\in(0,\tfrac12)$).} Now \eqref{eq:r-two-cases} has the extra term,
\[
d_1=\frac1{\sqrt{m+1}}+\frac{r_{\mathrm{RS}}}{2\cos u},
\]
and this time the two signs multiply to $-1$: the correction is
\emph{subtracted}, with magnitude $\Lambda(\hat p)\,a^{-1/2}+O(a^{-3/2})$
(here $\cos(2\pi x(1-x)-\tfrac{3\pi}8)=\cos(\pi(\tfrac18-2\hat p^{2}))$
under $x=\hat p+\tfrac12$). Lemma~\ref{lem:C0-trig}(ii) caps
$\Lambda\leq\tfrac12$, and $a=m+1+\hat p\geq m+1$, so
\[
d_1\;\geq\;\frac1{\sqrt{m+1}}-\frac{1}{2a^{1/2}}-O(a^{-3/2})
\;\geq\;\frac{1}{2\,a^{1/2}}-O(a^{-3/2}).
\]

In both cases $d_1\geq(\tfrac12\sin\tfrac\pi8-o(1))a^{-1/2}$ once $a$ is
large enough that the explicit $O(1/a)$ errors of \eqref{eq:gabcke} and of
the phase expansion \eqref{eq:cosu} are dominated. The boundary slivers of
\eqref{eq:cut-cases} are covered because the two case formulas agree to
leading order at the seams: as $x\to\tfrac12$ from either side, and across
each integer, both give $d_1\approx\tfrac12a^{-1/2}$. The range $T<T_0$ is
closed by the verification of \S\ref{sec:decomp} ($30{,}562$ samples on
$1\le T\le30$ with magnified windows at fractional parts
${\approx}\,\tfrac14,\tfrac34$; \texttt{check\_d1\_positive\_critical.py}).
\end{proof}

\begin{remark}[why the critical line has no poles of $d_1$]
The denominator $\sin(2\omega+\psi)=\sin2u$ vanishes twice per unit
interval of $T$, near $\{T\}=\tfrac14$ and $\tfrac34$, exactly where
\S\ref{sec:pole-locations} finds the off-line poles. On the line, both
vanishings belong to the factor $\sin u$, which cancels identically against
the numerator in \eqref{eq:exact-line}; the dangerous factor $\cos u$ stays
bounded away from zero by \eqref{eq:cosu-bound}. Off the line the
cancellation is spoiled, and the same two zeros become genuine poles. The
next theorem quantifies how much of each unit interval they actually
poison.
\end{remark}

\begin{theorem}\label{thm:d1-positive-offline}
Fix $\sigma\in(0,1)$, $\sigma\neq\tfrac12$. There are constants $C(\sigma)$
and $T_0(\sigma)$ such that for all $T\geq T_0$, both $d_1(\sigma,T)>0$ and
$d_2(\sigma,T)>0$ whenever $\{T\}$ lies outside two intervals of width at
most $C(\sigma)/T$, centered at points that converge to $\tfrac14$ and
$\tfrac34$ at rate $O(1/T)$. Inside each window one of the two weights is
negative, so the excluded intervals cannot be removed. In particular, for
every fixed $\delta>0$ there is $T_1(\sigma,\delta)$ with $d_1>0$ and
$d_2>0$ whenever $T\geq T_1$ and
$\{T\}\notin(\tfrac14-\delta,\tfrac14+\delta)\cup
(\tfrac34-\delta,\tfrac34+\delta)$.
\end{theorem}

\begin{proof}
Write $u'=\omega+\psi/2$ and let $\tau=\arg R-\psi/2\pmod\pi$, folded to
$[-\tfrac\pi2,\tfrac\pi2)$, measure the tilt of $R$ away from the bisector
direction. Up to a common positive factor and a common sign,
\eqref{eq:d1}--\eqref{eq:d2} read
\begin{equation}\label{eq:offline-form}
d_1\propto\frac{\sin(u'-\tau)}{\sin 2u'},
\qquad
d_2\propto\frac{\sin(u'+\tau)}{\sin 2u'}.
\end{equation}
Two inputs control the picture.

\emph{(i) The tilt $\tau$ is small.} The first-term approximation
\eqref{eq:R-general} has argument $\psi/2$ modulo $\pi$ up to
$O_\sigma(1/t)$; the deviation of the exact $R$ from that argument comes
from the next term of Siegel's expansion, giving
$|\tau|\leq C_2(\sigma)\,t^{-1/2}$. (Numerically at $\sigma=0.3$ the
product $|\tau|\,t^{1/2}$ stays between about $0.03$ and $0.09$ over
$t\in[200,4\times10^{4}]$.)

\emph{(ii) The phase $u'$ has the same skeleton as on the line.} Since
$\psi(\sigma+it)=-2\vartheta(t)+O_\sigma(1/t)$ uniformly for $\sigma$ in
compacta, $u'=u+O_\sigma(1/t)$, and \eqref{eq:cosu}--\eqref{eq:cosu-bound}
apply: $|\cos u'|\geq\sin\tfrac\pi8-o(1)$, so the only sign changes of
$\sin2u'=2\sin u'\cos u'$ within a unit interval are the two zeros of
$\sin u'$, at fractional parts converging to $\tfrac14$ and $\tfrac34$,
crossed with slope $|du'/d\{T\}|=\pi+O(1/T)$.

Away from those zeros (say $|\sin u'|\geq2|\tau|$) the ratios
$\sin(u'\mp\tau)/\sin u'$ are positive, so \eqref{eq:offline-form} carries
the same signs as on the line, and the Step-3/Step-4 analysis of
Theorem~\ref{thm:d1-positive-line}, with the $O_\sigma$ errors absorbed
into the constants, gives positivity of both weights. Near a zero of
$\sin u'$ the numerator of $d_1$ crosses zero at $u'=\tau$ while the
denominator crosses at $u'=0$: between the two crossings, an interval of
$u'$-length exactly $|\tau|$, the signs disagree and $d_1<0$;
mirror-symmetrically $d_2<0$ on the interval of length $|\tau|$ on the
other side. Converting $u'$-length to $\{T\}$-length by the slope
$\pi+O(1/T)$ bounds each window by
$C_2(\sigma)\,t^{-1/2}/\pi\asymp C(\sigma)/T$, since
$t=I(T)\asymp2\pi T^{2}$, centered $O(1/T)$ from the denominator zero. On the
critical line $\tau\equiv0$ exactly (Step~1 of
Theorem~\ref{thm:d1-positive-line}), the two crossings coincide, and the
windows are empty, consistent with no poles there.
\end{proof}

\begin{remark}
The mechanism reproduces Remark~\ref{rem:d-ratio}: just outside a window,
$d_1$ and $d_2$ have opposite-signed numerators of size ${\approx}\,|\tau|$
against a denominator passing through zero, which is why values as large as
$d_1=-d_2\approx-4\times10^{36}$ are observed adjacent to a pole while the
window itself is only $O(1/T)$ wide.
\end{remark}

\begin{remark}[numerical corroboration]
At $\sigma=0.3$, scanning $\{T\}\in[0.22,0.28]\cup[0.72,0.78]$ with step
$2\times10^{-4}$ (\texttt{check\_positivity\_windows.py}):
\begin{center}
\begin{tabular}{ccccc}
$N$ & window near $\tfrac14$ & center & window near $\tfrac34$ & center\\
\hline
5  & $0.0034$ & $0.2556$ & $0.0010$ & $0.7550$\\
10 & $0.0020$ & $0.2529$ & $0.0006$ & $0.7528$\\
20 & $0.0008$ & $0.2515$ & $0.0004$ & $0.7515$\\
40 & $0.0006$ & $0.2508$ & $0.0002$ & $0.7508$
\end{tabular}
\end{center}
Widths scale like $1/T$ and centers drift to $\tfrac14,\tfrac34$ like
$1/T$, as Theorem~\ref{thm:d1-positive-offline} predicts; on the line the
same scan finds no negative values at all. The leading-order formula of
Theorem~\ref{thm:d1-positive-line} reproduces the exact $d_1$ to relative
error $0.028$ on $T\in[20,21]$ and $0.014$ on $T\in[40,41]$, consistent
with the $O(1/a)$ error terms, and its lower bound $0.19\,a^{-1/2}$ lies
below the observed minima ($0.0498$ and $0.0355$ respectively).
\end{remark}

\subsection{The limit profile of the fractional amount: the waveform in closed form}\label{sec:d1-limit}

The bottom panel of Figure~\ref{fig:d1-critical} showed the normalized
fraction $\hat d_1=\lceil T\rceil^{1/2}d_1$ sweeping nearly the same arc in
every unit interval, and \S\ref{sec:d1-approx-critical} met the tangent
waveform $\mathcal{W}_{\infty}$ as the limit of the pinned coordinate
$\mathcal{W}$. The machinery of Theorem~\ref{thm:d1-positive-line} upgrades
both observations to a theorem: at fixed fractional part the fraction
converges, and its limit is exactly $\tfrac12$ minus the tangent waveform.
Note that it is the hatted weight that has a nontrivial limit: $d_1$ itself
decays, $d_1\asymp a^{-1/2}$ by Theorem~\ref{thm:d1-positive-line}, and the
normalization \eqref{eq:frac-vs-weight} is what removes the decay.

\begin{theorem}\label{thm:d1-limit}
Let $\sigma=\tfrac12$ and fix $x\in(0,1)$. Then
\[
\lim_{\substack{T\to\infty\\ \{T\}=x}}\hat d_1(T)
=\lim_{N\to\infty}\sqrt{N+1}\;d_1(N+x)
=d(x),
\qquad
d(x)\;=\;\tfrac12-\mathcal{W}_{\infty}(x)
\;=\;\tfrac12-\tfrac12\tan(2\pi x)\,
\tan\!\bigl(2\pi(x-\tfrac14)(x-\tfrac34)\bigr),
\]
with the removable points $x=\tfrac14,\tfrac34$ filled in by continuity
($d(\tfrac14)=\tfrac14$, $d(\tfrac34)=\tfrac34$). The same limit holds for
$\hat d_2$, since $d_2=d_1$ on the line. The convergence rate is $O(1/T)$.
\end{theorem}

\begin{proof}
Fix $x$ and let $T=N+x\to\infty$, $m=\lfloor T\rfloor=N$. Refining
\eqref{eq:a-vs-T} one order, $a=T+\tfrac12-\tfrac1{24T}+O(T^{-2})$, so the
offset $\hat p$ converges to its case value in \eqref{eq:cut-cases} with
rate $1/T$; moreover $(m+1)/a\to1$ at the same rate, and all error terms in
Steps 2--4 of Theorem~\ref{thm:d1-positive-line} are $O(1/a)$. Multiplying
the two case formulas of that proof by $\sqrt{m+1}$ therefore gives
\[
\hat d_1\longrightarrow
\frac{C_0(x+\tfrac12)}
{2\cos\bigl(2\pi x(1-x)-\tfrac{3\pi}8\bigr)}
\quad(x<\tfrac12),
\qquad
\hat d_1\longrightarrow 1-\Lambda\bigl(x-\tfrac12\bigr)
\quad(x>\tfrac12).
\]
It remains to collapse both branches to the tangent form. For
$x<\tfrac12$, substituting $\hat p=x+\tfrac12$ into $C_0$ gives
$\hat p^{2}-\hat p-\tfrac1{16}=x^{2}-\tfrac5{16}$ and
$\cos(2\pi\hat p)=-\cos(2\pi x)$, so with $\alpha=2\pi x$ and
$\beta=2\pi(x-\tfrac14)(x-\tfrac34)=2\pi(x^{2}-x)+\tfrac{3\pi}8$ the limit
is
\[
\frac{-\cos\bigl(2\pi x^{2}-\tfrac{5\pi}8\bigr)}
{2\cos(2\pi x)\cos\bigl(2\pi x(1-x)-\tfrac{3\pi}8\bigr)}
=\frac{\cos(\alpha+\beta)}{2\cos\alpha\,\cos\beta}
=\tfrac12\bigl(1-\tan\alpha\tan\beta\bigr)
=\tfrac12-\mathcal{W}_{\infty}(x),
\]
using $\alpha+\beta=2\pi x^{2}+\tfrac{3\pi}8$,
$\cos(\alpha+\beta)=-\cos(2\pi x^{2}-\tfrac{5\pi}8)$, the evenness of the
cosine for the $\beta$ factor, and the product formula. For $x>\tfrac12$,
the same substitution algebra applied to $\Lambda(x-\tfrac12)$ (as in
Lemma~\ref{lem:C0-trig}(ii), with $\hat p=x-\tfrac12$) yields
$\Lambda(x-\tfrac12)=\tfrac12-\mathcal{W}_{\infty}(1-x)$, and
$\mathcal{W}_{\infty}(1-x)=-\mathcal{W}_{\infty}(x)$ because
$\tan(2\pi(1-x))=-\tan(2\pi x)$ while the quadratic factor is symmetric
under $x\mapsto1-x$; hence
$1-\Lambda(x-\tfrac12)=\tfrac12-\mathcal{W}_{\infty}(x)$ as well.
\end{proof}

So the same elementary function that governs the pinned waveform of
\S\ref{sec:d1-approx-critical} governs the fraction itself, recentered at
$\tfrac12$ and flipped. This is the closed-form waveform promised at the
start of the section.

\begin{remark}[properties of the profile]\label{rem:d-profile}
Since $\mathcal{W}_{\infty}(1-x)=-\mathcal{W}_{\infty}(x)$, the profile
satisfies the functional equation
\[
d(x)+d(1-x)=1,
\]
so its graph is antisymmetric about the point $(\tfrac12,\tfrac12)$. Three
exact values follow:
\[
d\bigl(\tfrac14\bigr)=\tfrac14,\qquad
d\bigl(\tfrac12\bigr)=\tfrac12,\qquad
d\bigl(\tfrac34\bigr)=\tfrac34:
\]
at the two abscissas where the off-line weights have their poles, the
limiting fraction of the next summand equals the abscissa itself, and the
middle value is forced by the functional equation. Furthermore
$\mathcal{W}_{\infty}(0)=\mathcal{W}_{\infty}(1)=0$ gives
$d(0^{+})=d(1^{-})=\tfrac12$: although $\hat d_1$ jumps at every integer
$T$ by the relation \eqref{eq:fold}, both sides of the jump tend to
$\tfrac12$, and the limit profile extends to a \emph{continuous periodic
function} on $\mathbb{R}/\mathbb{Z}$ once $d(0):=\tfrac12$. Its range is
$[m_0,\,1-m_0]$ with
\[
m_0=d(x^{*})=0.2268951\ldots,
\qquad x^{*}=0.1629962\ldots
\]
(the maximum $1-m_0=0.7731048\ldots$ at the mirror point $1-x^{*}$):
asymptotically, the split always uses between $22.7\%$ and $77.3\%$ of the
$(m+1)$st summand, which is the range
${\approx}\,[0.23,\,0.78]$ measured in Figure~\ref{fig:d1-critical}.
\S\ref{sec:d1-bounds} upgrades this to uniform bounds valid at every finite
$T$.
\end{remark}

\begin{figure}[H]
\centering
\includegraphics[width=\textwidth]{fig_d_limit}
\caption{The closed-form profile $d(x)$ of Theorem~\ref{thm:d1-limit}
(blue) against the exact $\sqrt{m+1}\,d_1$ sampled at $T=10+x$, $50+x$,
$400+x$ on the critical line. Already at $T=10+x$ the samples sit on the
curve; the marked squares are the exact values
$(\tfrac14,\tfrac14)$, $(\tfrac12,\tfrac12)$, $(\tfrac34,\tfrac34)$, and
the dashed red lines are the extremes $m_0=0.2268951\ldots$ and
$1-m_0=0.7731048\ldots$ of the profile. Generated by
\texttt{fig\_d\_limit.py}.}
\label{fig:d-limit}
\end{figure}

\begin{remark}[numerical corroboration]
With $\hat d_1$ computed from the exact remainder
(\texttt{check\_d\_limit.py}), the error $|\hat d_1(N+x)-d(x)|$ at
$N=50,200,800$ decreases by a factor $4.0$ at each quadrupling of $N$,
matching the $O(1/T)$ rate; at $x=0.35$, for instance, the errors are
$3.6\times10^{-4}$, $9.2\times10^{-5}$, $2.3\times10^{-5}$.
\end{remark}

\begin{remark}[off the critical line the profile is universal]
For fixed $\sigma\in(0,1)$ the same argument, run through
Theorem~\ref{thm:d1-positive-offline}, gives the \emph{identical} limit for
both fractional amounts of \eqref{eq:frac-vs-weight},
\[
\lim_{\substack{T\to\infty\\ \{T\}=x}}\hat d_1(\sigma,T)
=\lim_{\substack{T\to\infty\\ \{T\}=x}}\hat d_2(\sigma,T)
= d(x)
\qquad\bigl(x\neq\tfrac14,\tfrac34\bigr):
\]
the leading-order amplitudes agree to $O(1/t)$ and the tilt
$\tau=O(t^{-1/2})$ vanishes in the limit, so the $\sigma$-dependence
survives only inside the two shrinking pole windows. Numerically at
$\sigma=0.3$ both normalized weights match $d(x)$ with errors falling like
$1/N$ (at $x=0.35$: errors $9.7\times10^{-5}$ and $5.1\times10^{-5}$ at
$N=800$). At $x=\tfrac14,\tfrac34$ exactly, the finite-$T$ poles prevent a
uniform statement off the line.
\end{remark}

\subsection{Uniform bounds: the fraction stays between \texorpdfstring{$\tfrac15$ and $\tfrac45$}{1/5 and 4/5}}\label{sec:d1-bounds}

Theorem~\ref{thm:d1-limit} confines the fraction to $[m_0,1-m_0]$ in the
limit. At finite $T$ the excursions are slightly larger, and round
constants cover every height at once: the fractional summand never uses
less than a fifth, or more than four fifths, of the next term.

\begin{theorem}\label{thm:d1-bounds}
Let $\sigma=\tfrac12$. For every non-integer $T>1$,
\[
\frac15\;<\;\hat d_1(T)=\hat d_2(T)\;<\;\frac45 .
\]
The constants are nearly sharp: no better $T$-independent ones exist than
\[
\inf_{T>1}\hat d_1=m_0=0.2268951\ldots,
\qquad
\sup_{T>1}\hat d_1=0.78499047437\ldots,
\]
the infimum approached but not attained as $T\to\infty$ along
$\{T\}\to x^{*}=0.1629962\ldots$, the supremum attained at
$T=1.85307336433\ldots$.
\end{theorem}

\begin{proof}
On the line $d_2=d_1$ and $|\chi|=1$, so $\hat d_2=\hat d_1$ by
\eqref{eq:frac-vs-weight} and it suffices to bound
$\hat d_1=\sqrt{m+1}\,d_1$.

\emph{Step 1 (exact finite-$T$ form).} Multiplying the two case formulas
in the proof of Theorem~\ref{thm:d1-positive-line} by $\sqrt{m+1}$ and
writing $\rho=\sqrt{(m+1)/a}$,
\begin{align*}
N=m\ (\{T\}<\tfrac12):&\qquad
\hat d_1=\rho\cdot
\frac{C_0(\hat p)+E}{2\cos\bigl(2\pi x(1-x)-\tfrac{3\pi}8+\delta\bigr)},
&&\rho=\sqrt{1+\tfrac{1-\hat p}{a}}
\in\Bigl[1,\sqrt{1+\tfrac1{2a}}\,\Bigr],\\[2pt]
N=m+1\ (\{T\}>\tfrac12):&\qquad
\hat d_1=1-\rho\cdot
\frac{C_0(\hat p)+E}{2\cos\bigl(\pi(\tfrac18-2\hat p^{2})+\delta\bigr)},
&&\rho=\sqrt{1-\tfrac{\hat p}{a}}
\in\Bigl[\sqrt{1-\tfrac1{2a}},1\Bigr],
\end{align*}
with $|E|\leq c_1a^{-1}$ from \eqref{eq:gabcke} and $|\delta|\leq
c_2a^{-1}$ from the phase expansion \eqref{eq:cosu}. At $\rho=1$,
$E=\delta=0$ these are exactly the two branches of
Theorem~\ref{thm:d1-limit}.

\emph{Step 2 (range of the main terms).} On $(0,\tfrac12)$ the branch
function $d(x)$ descends from $\tfrac12$ to its single interior minimum
$m_0=d(x^{*})$ and climbs back to $\tfrac12$ (one sign change of $d'$; the
critical point $x^{*}$ is the root of an explicit elementary equation).
Hence the first main term ranges over $[m_0,\tfrac12]$ and the second over
$[\tfrac12,\,1-\rho\,m_0]$.

\emph{Step 3 (all $T\geq30$).} The total deviation of $\hat d_1$ from its
main term is at most $(\rho-1)\cdot\tfrac12$ plus a term collecting $E$
and $\delta$, with $\rho-1\leq\tfrac1{4a}$; the collected constant is
validated by the convergence-rate measurements of
Theorem~\ref{thm:d1-limit} (observed absolute error ${\leq}\,0.008$ at
$T\approx20$, decreasing like $1/T$), giving a total deviation below
$0.02$ for $T\geq30$. Therefore, for $T\geq30$,
\[
\{T\}<\tfrac12:\quad
\hat d_1\in\bigl(m_0-0.02,\ \tfrac12(1+\tfrac1{120})+0.02\bigr)
\subset(0.206,\,0.525),
\qquad
\{T\}>\tfrac12:\quad
\hat d_1\in(0.48,\,0.795),
\]
and both windows lie inside $(\tfrac15,\tfrac45)$.

\emph{Step 4 ($1<T\leq30$, computational).} A dense evaluation of the
exact $\hat d_1$ (grid step $5\times10^{-4}$ on $(1,3]$, $10^{-3}$ on
$(3,12]$, $2\times10^{-3}$ on $(12,30]$;
\texttt{check\_uniform\_bounds.py}) gives the observed range
$[0.2274,\ 0.78499]$, with margin at least $0.015$ to both $\tfrac15$ and
$\tfrac45$ and sampled variation between adjacent grid points below
$0.004$, a factor of nearly four inside the margin. This covers the
remaining range and, refined near its endpoints, produces the sharp values
in the statement: the maximum $0.78499047437\ldots$ at
$T=1.85307336433\ldots$ (on the $m=1$ interval, where $\rho<1$ pushes the
second branch above its limiting value), and interval minima decreasing
toward $m_0$ from above as $T$ grows ($0.227369$ at $T=59.164$, for
instance).
\end{proof}

\begin{figure}[H]
\centering
\includegraphics[width=0.78\textwidth]{fig_interval_range}
\caption{Range of the exact $\hat d_1$ over each unit interval $(N,N+1)$,
$N=1,\dots,29$, on the critical line (blue bands; grid step
$5\times10^{-4}$ for $N<3$, $10^{-3}$ after). Dashed black lines: the
uniform bounds $\tfrac15$ and $\tfrac45$ of
Theorem~\ref{thm:d1-bounds}; the red right-hand scale marks the extremes
$m_0=0.2268951\ldots$ and $1-m_0=0.7731048\ldots$ of the limit profile of
Theorem~\ref{thm:d1-limit}. The interval maxima decrease toward the
$1-m_0$ mark, starting from the supremum $0.78499\ldots$ attained on the
first interval at $T=1.8530\ldots$; the interval minima decrease toward
$m_0$ from above without attaining it. Generated by
\texttt{fig\_interval\_range.py}.}
\label{fig:interval-range}
\end{figure}

\begin{remark}
For the unnormalized weight the theorem reads
$\tfrac15(m+1)^{-1/2}<d_1<\tfrac45(m+1)^{-1/2}$; in particular
$d_1\leq0.78499\ldots/\sqrt2=0.55508\ldots$, attained at the same
$T=1.853\ldots$, and no positive $T$-independent lower bound exists for
$d_1$ itself.
\end{remark}

\begin{remark}[what is cited rather than proved]\label{rem:d1-provenance}
For the record, the proofs of
\S\ref{sec:d1-positive}--\S\ref{sec:d1-bounds} take four inputs on faith
or on computation. (1)~Gabcke's explicit error bound \eqref{eq:gabcke} for
the first Riemann--Siegel correction \cite{Gabcke1979}. (2)~The
first-order Siegel asymptotics of $R$ at fixed $\sigma\in(0,1)$ behind
Theorem~\ref{thm:d1-positive-offline}(i), from Siegel's account of the
Nachlass \cite{Siegel1932}; its error term is classical, but the constant
$C_2(\sigma)$ is not tracked explicitly here, which is why
$T_0(\sigma)$ is not effective as stated. (3)~The numerically located
minimum in Lemma~\ref{lem:C0-trig}(i) and the single-minimum property of
the branch in Step~2 of Theorem~\ref{thm:d1-bounds}. (4)~
Theorem~\ref{thm:d1-bounds} is computer-assisted: Step~3's error constant
is validated numerically rather than derived with explicit Gabcke and
Taylor constants, and Step~4 is a dense grid evaluation with margin, not
interval arithmetic.
\end{remark}

\clearpage
% ======================================================================
\section{The geometry behind the result (experimental mathematics)}\label{sec:geometry}
% ======================================================================

The decomposition $R=R_{1ps}+R_{2ps}$ was not first derived; it was first
\emph{seen}. Over roughly nine years we built and used an interactive
visualization tool, Zest,\footnote{Source code and standalone executables (Mac
and PC) at \url{https://github.com/paul-stahura/zest}.} which draws the
partial-sum chains of $\zeta$ in the complex plane and lets us sweep $\sigma$
and $T$ continuously while measuring what appears on the screen. This section
reconstructs the series of observations made with that tool. It led, in order,
to a distinguished link and point of the spiral (the \emph{bisector link} and
\emph{bisector point}), and from their motion to the mapping $I(T)$ of
\S\ref{sec:IT}. The final observation in that sequence, the \emph{yin} and
\emph{yang} curves from which the weights $d_1,d_2$ of \S\ref{sec:decomp}
were read off, is developed at length in \S\ref{sec:yinyang}.

Throughout we use the link and joint numbering of
\S\ref{sec:matrix-product}: joint $k$ of a chain is the value of its partial
sum after $k$ summands, and link $k$ runs from joint $k$ to joint $k+1$, so
that link $k$ represents the $(k+1)$st summand and each remainder is the
fractional link $m$. The \emph{forward} chain ($\Sigma_1$) is
anchored at joint $0$, the origin. The \emph{reverse} chain ($\Sigma_2$) is
anchored at $\zeta$ and traversed backwards, so its joint $m$ lands at
$\zeta-\Sigma_2=\Sigma_1+R$. In the pictures the chains are drawn well past
joint $m=\lfloor T\rfloor$: the sum, though ultimately divergent, first winds
into an Euler-like spiral whose center lies near $\zeta$; it seemingly
converges before it diverges. Figure~\ref{fig:full-spirals} shows the picture:
the forward chain winds from the origin into a spiral centered near $\zeta$,
while the reverse chain winds from $\zeta$ into a mirror spiral centered near
the origin, and the two link chains cross. If more links were drawn, the links would
continue to spiral outward around $\zeta$ and cover the entire figure, because
the zeta sum diverges on the critical line, unlike the Dirichlet
$L$-functions of \S\ref{sec:IT-Lfunctions}, whose characters have bounded partial sums and
so those chains close in on their limit instead of covering the plane.\footnote{At $t=176.13$ on the critical line, the partial
sums of $L(s,\chi_3)$ come within $0.021$, $0.0021$ and $0.00019$ of the value
by $N=10^{3}$, $10^{5}$ and $3\times10^{6}$, a drift inward like $N^{-1/2}$,
while zeta's are off by $0.18$, $1.80$ and $9.83$, a drift outward at
$\sqrt{N}/t$. Convergence in $\sigma>0$ is special to a non-principal
character: the principal character returns $\zeta$ itself, and $L$-functions
of higher degree have coefficients too large to converge anywhere in the
critical strip.} We stop drawing links at link number
\begin{equation}\label{eq:LN}
L_N(T,S_n)=\left\lfloor\frac{I(T)}{\pi\,(2S_n+1)}\right\rfloor,
\end{equation}
where $S_n$ is the spiral number, spiral number $0$ being the last one (the
biggest, at $\zeta$). Note that $\zeta$ itself is not typically \emph{on} this
last link; the last link is the link closest to $\zeta$. The same formula
locates the other spirals and pseudo-spirals too: with $S_n=1$, for example, it yields the link
number of the link most in the center of the second-to-last spiral.

\begin{figure}[p]
\centering
\includegraphics[height=0.82\textheight,keepaspectratio]{fig_full_spirals}
\caption{The forward spiral (blue, anchored at the origin $O$) and the reverse
spiral (green, anchored at $\zeta$) for $T=14.3085$
(so $t=I(T)\approx 1377.33$ and $m=14$), on the critical line
$\sigma=\tfrac12$ (top, $\zeta\approx-1.163+0.870i$) and off it at
$\sigma=0.673$ (bottom, $\zeta\approx 0.092+1.434i$). Each spiral is drawn out
to link $\lfloor I(T)/\pi\rfloor=438$, the link nearest its spiral center, for
a total of $439$ links per spiral (the count is independent of $\sigma$). The
red dot is the bisector point $B_1=\Sigma_1+R_{1ps}$, where the two chains
cross at their bisector links (\S\ref{sec:bisector}); the dotted line is the
bisector line through the bisector point and $\zeta/2$ (open circle). On the
critical line (top) the whole configuration is symmetric about the bisector
line; at $\sigma=0.673$ (bottom) the symmetry is lost. Generated by
\texttt{fig\_full\_spirals.py}.}
\label{fig:full-spirals}
\end{figure}

\begin{figure}[p]
\centering
\includegraphics[height=0.8\textheight,keepaspectratio]{fig_last_spiral_zoom}
\caption{Zooming in on the last spiral, the spiral centered on
$\zeta$, of the forward chain at $T=41.73$ (so
$t=I(T)\approx 11204.74$, $\sigma=\tfrac12$,
$\zeta\approx-4.796-7.853i$). The chain carries
$\lfloor I(T)/\pi\rfloor+1=3567$ links; the last link is number $3566$
by \eqref{eq:LN} with $S_n=0$, and it is highlighted in black in every
panel.
Each panel is more zoomed in than the last (panels $1$--$6$, left to
right from the top). Near the
spiral center each link turns by nearly $\pi$, so the links sweep back
and forth across $\zeta$ and the joints settle onto a tight ring of
radius about half a link length: at coarse scale the spiral renders as
a dense annulus (panels $1$--$3$), and the near-diameter links envelope
a lens-shaped clearing around $\zeta$ (panels $4$--$5$). Observe in
panel $6$ that $\zeta$ is \textbf{\emph{not on}} the last link: the last link
is only the link closest to $\zeta$. Generated by
\texttt{fig\_last\_spiral\_zoom.py}.}
\label{fig:last-spiral-zoom}
\end{figure}

\clearpage
\subsection{The bisector link and the bisector point}\label{sec:bisector}

Drawn at height $T$, each chain exhibits spirals and pseudo-spirals, and (as
others have also observed~\cite{Erickson2005,Lander2018,Nickel2013,Kapitonets2019}) there is a one-to-one correspondence between the
joints and the spiral centers, in a type of symmetry; Nickel puts it as the
center-to-center distances and angles of the large linked spirals reproducing
the initial steps exactly. The link
at index $m=\lfloor T\rfloor$ sits at the boundary between the two
regimes, and we call it the \emph{bisector link}. Each chain has one: the \emph{forward bisector link} (link $m$ of the
$\Sigma_1$ chain, running from $\Sigma_1$ along the $(m+1)$st summand
direction) and the \emph{reverse bisector link} (link $m$ of the $\Sigma_2$
chain, running from $\Sigma_1+R$ along its $(m+1)$st summand direction).

The two spirals intersect at their bisector links, and we call the point where
the forward and reverse bisector links cross the \emph{bisector point}. One
qualification: the crossing can momentarily fail. Twice per unit of $T$, near
fractional parts ${\approx}\,0.25$ and ${\approx}\,0.75$, the two bisector
links become parallel and, off the critical line, do not intersect; these are
the isolated instants at which $\sin(2\omega+\arg\chi)$, the denominator of
\eqref{eq:d1}--\eqref{eq:d2}, vanishes. Those are exactly the places where
$d_1$ and $d_2$ have poles: the parallel-link condition and the pole
condition are the same equation, which is why the precise locations were
worked out in \S\ref{sec:pole-locations}. On the critical line the two links are then collinear rather than
merely parallel, so the bisector point survives there.
In the notation of the earlier sections the bisector point is exactly
\begin{equation}\label{eq:bisector-point}
B_1 \;=\; \Sigma_1 + R_{1ps} \;=\; \Sigma_1 + R - R_{2ps};
\end{equation}
it is visible as the joint between the red and orange fractional
links in Figures~\ref{fig:spiral-summands} and~\ref{fig:remainder-zoom}.

\subsection{Legs}\label{sec:legs}

When $\sigma=\tfrac12$ the spiral object
has a type of symmetry about the line through the bisector point and
$\tfrac{\zeta}{2}$, so the bisector point sits at the apex of an isosceles
triangle whose base runs from the origin to $\zeta$ and whose equal sides are
the two legs $B_1$ and $\zeta-B_1$ (this is Corollary~\ref{cor:equal-legs}
below, proved in Appendix~\ref{app:lean-critical} from the reflection
$B_2=\chi\overline{B_1}$; equivalently, the perpendicular projection of the
bisector point onto the base $\mathbb{R}e^{\,i\psi/2}$ carrying $\zeta$, where
$\psi=\arg\chi$, is exactly $\tfrac{\zeta}{2}$, which is
Corollary~\ref{cor:bisector-proj}). Off the line we keep the same
terminology even though the triangle's legs are then (usually) unequal.

\begin{figure}[H]
\centering
\includegraphics[width=\textwidth]{fig_legs}
\caption{The two legs, drawn as thick dark-yellow segments over the spirals of
Figure~\ref{fig:full-spirals}: one from the origin $O$ to the bisector point
($B_1$), the other from the bisector point to $\zeta$ ($\zeta-B_1$). On the
critical line (left, $\sigma=\tfrac12$) the legs have equal length
($|B_1|=|\zeta-B_1|\approx 3.577$): the bisector point is the apex of an
isosceles triangle whose base runs from $O$ to $\zeta$, and the dotted
bisector line is its axis of symmetry. Off the critical line (right,
$\sigma=0.673$) the legs are unequal ($\approx 2.899$ and $\approx 1.797$) at
this $T$. Both panels also mark the two angles defined below: $\vartheta_1$,
at $O$ from the positive real axis to Leg~1, and $\vartheta_2$, at the
bisector point from the extension of Leg~1 (dashed) to Leg~2.
Generated by \texttt{fig\_legs.py}.}
\label{fig:legs}
\end{figure}

On the critical line the two fractional summands are not independent: each is
the other's complex conjugate, rotated by the functional-equation factor.
Writing $R_{1ps}=d_1e^{-i\omega}$ and $R_{2ps}=d_2e^{\,i(\omega+\arg\chi)}$ as in
\S\ref{sec:matrix-product} and using that $|\chi|=1$ when $\sigma=\tfrac12$, the
equality $d_1=d_2$ proved in \S\ref{sec:d1-exact}
(Corollary~\ref{cor:equal}) is exactly
equivalent to the compact relation
\begin{equation}\label{eq:R2-conj-R1}
R_{2ps}=\chi\,\overline{R_{1ps}}\qquad\bigl(\sigma=\tfrac12\bigr).
\end{equation}
Geometrically, \eqref{eq:R2-conj-R1} says the second fractional summand is the
mirror image of the first (conjugation across the real axis) followed by
the rotation $\chi$. Because $\chi$ has unit modulus on the line this rotation
is an isometry, so $|R_{1ps}|=|R_{2ps}|$ immediately. This is the
functional-equation reflection $s\leftrightarrow 1-s$ acting as a symmetry of
the two remainders, and it is the source of the equal-leg phenomenon we now
develop.

The two legs $B_1=\Sigma_1+R_{1ps}$ and $B_2=\Sigma_2+R_{2ps}$ of
\S\ref{sec:matrix-product} meet at a ``bisector point'': Leg~1 runs from the
origin to $\Sigma_1+R_{1ps}$, and Leg~2 from there to $\zeta$. Their lengths are
\begin{align}
L_{1}(\sigma,T) &= \bigl|\,\Sigma_{1}+R_{1ps}\,\bigr|,\\
L_{2}(\sigma,T) &= \bigl|\,\Sigma_{2}+R_{2ps}\,\bigr|,
\end{align}
while $\vartheta_1$ is the angle Leg~1 makes with the real axis and $\vartheta_2$
is the angle from Leg~1 to Leg~2 (oriented so that a positive $\vartheta_2$
rotates Leg~1 toward Leg~2):
\begin{align}
\vartheta_{1}(\sigma,T) &= \arg\!\bigl(\Sigma_{1}+R_{1ps}\bigr),\\[2pt]
\vartheta_{2}(\sigma,T) &= \arg\!\left(\frac{\Sigma_{2}+R_{2ps}}{\Sigma_{1}+R_{1ps}}\right).
\end{align}
With these, $\zeta$ is simply the sum of the two legs,
\begin{equation}\label{eq:zeta-legs-polar}
\zeta(\sigma,T)
= L_{1}\,e^{\,i\vartheta_{1}}
\;+\;
L_{2}\,e^{\,i(\vartheta_{1}+\vartheta_{2})},
\end{equation}
or equivalently, in matrix form,
\begin{equation}\label{eq:zeta-matrix}
\zeta(\sigma,T)
=
\begin{pmatrix}1 & i\end{pmatrix}
\begin{pmatrix}
\cos\vartheta_{1} & \cos(\vartheta_{1}+\vartheta_{2})\\[2pt]
\sin\vartheta_{1} & \sin(\vartheta_{1}+\vartheta_{2})
\end{pmatrix}
\begin{pmatrix}
L_{1}\\[2pt]
L_{2}
\end{pmatrix},
\end{equation}
where the row $\begin{pmatrix}1 & i\end{pmatrix}$ turns the resulting real
$(x,y)$ components into the complex value $x+iy$.

\begin{corollary}\label{cor:equal-legs}
When $\sigma=\tfrac12$ the two legs are always equal in length: $L_1=L_2$.
\end{corollary}

This is immediate from Corollary~\ref{cor:equal} of \S\ref{sec:d1-exact}:
on the critical line every summand of $\Sigma_2$ is $\chi$ times the
conjugate of the corresponding summand of $\Sigma_1$, and by
\eqref{eq:R2-conj-R1}, which is equivalent to $d_1=d_2$, the same holds for the
fractional summands. Hence
$\Sigma_2+R_{2ps} = \chi\,\overline{\Sigma_1+R_{1ps}}$, and since $|\chi|=1$
on the line, the two legs have the same length. A Lean proof is recorded in
Appendix~\ref{app:lean-critical}.

The same reflection locates the apex exactly, which is what the dotted bisector
line of Figures~\ref{fig:full-spirals} and~\ref{fig:legs} draws.

\begin{corollary}\label{cor:bisector-proj}
When $\sigma=\tfrac12$ the perpendicular projection of the bisector point $B_1$
onto the line through the origin carrying $\zeta$ is exactly $\tfrac{\zeta}{2}$:
writing $\psi=\arg\chi$,
\begin{equation}\label{eq:bisector-proj}
\Re\!\bigl(e^{-i\psi/2}B_1\bigr)\,e^{\,i\psi/2}\;=\;\frac{\zeta}{2}.
\end{equation}
\end{corollary}

Indeed $B_2=\chi\overline{B_1}$ gives
$e^{-i\psi/2}\zeta=e^{-i\psi/2}B_1+\overline{e^{-i\psi/2}B_1}
=2\Re(e^{-i\psi/2}B_1)$, and $e^{\,i\psi/2}$ is the direction of $\zeta$ (this is
the reality of the Riemann--Siegel $Z$ function). So the isosceles triangle with
base from $O$ to $\zeta$ has its apex over the midpoint of that base, and the
equal-leg statement $L_1=L_2$ is the same fact read as a distance. This is also
proved in Appendix~\ref{app:lean-critical}.

\paragraph{Consequences for zeta zeros.} A zero of $\zeta$ in the critical strip
requires \emph{both} of the following:
\[
\frac{L_{1}(\sigma,T)}{L_{2}(\sigma,T)} = 1
\qquad\text{and}\qquad
\vartheta_2(\sigma,T) = \pi,
\]
i.e.\ the two legs must be equal in length \emph{and} Leg~2 must fold back onto
Leg~1.\footnote{Tolkien's subtitle for \emph{The Hobbit}, \emph{There and Back
Again}, is the whole of the requirement. A hobbit leaves the Shire (the origin)
for Mordor (the bisector point) and his journey must end where it began, and no
choice of direction for the return trip from Mordor can bring the hobbit back to
the Shire unless the way back is exactly as long as the way out. Equal legs is that necessity; $\vartheta_2=\pi$ is the further
demand that the way back be the way he came.} By Corollary~\ref{cor:equal-legs} the length condition holds automatically on
the critical line; there, a zero occurs exactly when $\vartheta_2(\tfrac12,T) =
\pi$. Away from the critical line, the two legs are unequal except at
exceptional points, so most of the strip is immediately ruled out for zeros. The leg-length ratio is symmetric across the
line,
\begin{equation}
\frac{L_{1}(\sigma,T)}{L_{2}(\sigma,T)}
=
\frac{L_{2}(1-\sigma,T)}{L_{1}(1-\sigma,T)} .
\end{equation}
Numerically we have never found a point with $\sigma\neq\tfrac12$ where the
ratio is $1$ and $\vartheta_2=\pi$ simultaneously; such a point would disprove
the Riemann Hypothesis.

\begin{figure}[p]
\centering
\includegraphics[width=0.9\textwidth]{fig_equal_legs_strips}
\caption{Four windows of the critical strip ($0\le\sigma\le1$, index $T$
vertical): $2\le T\le3$, $3\le T\le4$, $4\le T\le5$, and $5\le T\le6$. The
blue vertical line is the critical line $\sigma=\tfrac12$, the black dots are
the zeta zeros, and the small blue points are the equal-leg points: points
$(\sigma,T)$ where $L_1=L_2$, i.e.\ where the bisector point
$B_1=\Sigma_1+R_{1ps}$ is equidistant from the origin and from $\zeta$. By
Corollary~\ref{cor:equal-legs} the entire critical line consists of equal-leg
points; off the line the equal-leg locus forms isolated ovals and thin
horizontal bands, and every zero sits on the critical line where the length
condition is automatic. The horizontal bands, near fractional parts
${\approx}\,0.25$ and ${\approx}\,0.75$ of $T$, are the instants where the two
bisector links become parallel and fail to intersect (\S\ref{sec:bisector}):
the bisector point has a pole there ($d_1$ and $d_2$ blow up), and as $T$
passes through it the equal-leg point sweeps across the full width of the
strip. The sweep gets faster as $T$ grows, so the fixed-step sampling leaves
fewer dots and the bands appear fainter. Data from the Zest tool point set
\emph{Zps Equal Leg Lengths}; generated by
\texttt{fig\_equal\_legs\_strips.py}.}
\label{fig:equal-legs-strips}
\end{figure}

\begin{figure}[p]
\centering
\includegraphics[width=0.9\textwidth]{fig_theta2_strips}
\caption{The same four windows of the critical strip as in
Figure~\ref{fig:equal-legs-strips}, now showing the \emph{folded-leg} points
(red): points $(\sigma,T)$ where $\vartheta_2=\pi$, i.e.\ where Leg~2 folds
straight back onto Leg~1 and a zeta zero is therefore possible. The blue
vertical line is the critical line and the black dots are the zeta
zeros. The folded-leg locus forms nearly horizontal curves sweeping the
strip; a zero requires \emph{both} $\vartheta_2=\pi$ and $L_1=L_2$, and since
the length condition holds automatically on the critical line
(Corollary~\ref{cor:equal-legs}), every crossing of a red curve with the
critical line is a zeta zero. Data from the Zest tool point set \emph{Zps Leg Angle =
PI}; generated by \texttt{fig\_theta2\_strips.py}.}
\label{fig:theta2-strips}
\end{figure}

\begin{figure}[p]
\centering
\includegraphics[width=0.9\textwidth]{fig_combined_strips}
\caption{Figures~\ref{fig:equal-legs-strips} and~\ref{fig:theta2-strips}
overlaid: the equal-leg locus $L_1=L_2$ (blue), the folded-leg locus
$\vartheta_2=\pi$ (red), the critical line (blue, since it belongs to the
equal-leg locus), and the zeta zeros (black dots). A zero requires both
conditions at once, so zeros occur exactly where
red meets blue. On the critical line the blue locus is the line itself
(Corollary~\ref{cor:equal-legs}), and every red curve crossing it produces a
zero;
away from the line, the red curves and the blue ovals interleave without
touching. A red/blue crossing off the critical line would be a
counterexample to the Riemann Hypothesis. Generated by
\texttt{fig\_combined\_strips.py}.}
\label{fig:combined-strips}
\end{figure}

\clearpage

Figure~\ref{fig:combined-strips} makes the stakes plain: wherever a region of
the critical strip contains no red and no blue dots, a zero of $\zeta$ is
impossible. On the critical line the two legs are always equal in length, so
whenever $\vartheta_2=\pi$ there, a zero occurs. Off the line the legs are
seldom equal in length, but sometimes they are (the blue ovals and bands).
If one could prove that, off the critical line, in the rare times when
the legs are equal, that $\vartheta_2$ is simultaneously never $\pi$, one would have proven
the Riemann Hypothesis true.

\subsection{The strip lines at ${\approx}\,0.25$ and ${\approx}\,0.75$ are not flat}\label{sec:strip-lines-not-flat}

\paragraph{The pole heights are not exactly flat.}
Viewed across the strip, the pole heights of \S\ref{sec:pole-locations}
are nearly invariant (``flat'') in $\sigma$, and symmetric about
$\sigma=\tfrac12$ (Table~\ref{tab:poles} and its caption). But they are
\emph{not exactly} flat: there is a slight, symmetric curvature in
$\sigma$, visible in Figure~\ref{fig:pole-heights-zoom}.

\begin{figure}[H]
\centering
\includegraphics[width=0.5\textwidth]{fig_pole_heights_zoom}
\caption{Two magnifications of the critical strip in the style of
Figure~\ref{fig:combined-strips}: equal-leg locus $L_1=L_2$ (blue) and
folded-leg locus $\vartheta_2=\pi$ (red), each recomputed at $N=400$
values of $\sigma$ across $(0,1)$; critical line (blue vertical) and
zeta zeros (black). Samples where the PS split is singular (tiny
determinant) are omitted from the red locus, so the $d_1$ pole is not
mislabeled as legs folded back. Left: $2.75\le T\le2.76$; the blue equal-leg
band sits near $T\approx2.757$ and the red folded band near
$T\approx2.759$, with a gap at $\sigma=\tfrac12$ (no zero in this
window). Right: $5.25\le T\le5.26$; blue near $T\approx5.255$ and red
near $T\approx5.2596$, meeting the critical line at the zeta zero in
the window (not at the pole band). The dashed horizontal guide marks
$\lfloor T\rfloor\pm\tfrac14$; the blue strip lines bow slightly across
$\sigma$, symmetric about $\sigma=\tfrac12$. At this resolution the
red and blue loci do not cross off the critical line. Generated by
\texttt{fig\_pole\_heights\_zoom.py}.}
\label{fig:pole-heights-zoom}
\end{figure}

\clearpage
\subsection{PS, AK and RS Legs and Angles}\label{sec:ps-ak-r2}

The equal-leg and folded-leg loci above were drawn with the exact
bisector point $B_1=\Sigma_1+R_{1ps}$. The same two conditions
$L_1=L_2$ and $\vartheta_2=\pi$ can be restated for any other choice of
bisector point. Two natural alternatives come from Kuznetsov's split
\eqref{eq:Rak-split}: take $B_1=\Sigma_1+R_{1ak}$, or take the midpoint
remainder $B_1=\Sigma_1+(R_{1ak}+R_{2ak})/2$, which is the geometric
analogue of Siegel's half-remainder $\tfrac12 R$ once
$R_{1ak}+R_{2ak}\approx R$.

Figure~\ref{fig:ps-ak-r2} magnifies a short window of the strip,
$4.65\le T\le4.80$, and overlays the three choices side by side. On the
left, equal-leg loci: PS in blue, AK in green, and $R/2$ in purple. The
three families form nested ovals about the critical line, but the
nesting order is not fixed: already in this short window the lower
oval has PS outermost and AK innermost, while the upper oval has AK
outermost and $R/2$ innermost, and across $3\le T\le20$ all six
orderings of PS, AK and $R/2$ occur with no simple alternating
pattern. On the right, folded-leg loci $\vartheta_2=\pi$: PS in red,
AK in orange, and $R/2$ in purple (hollow). The zeros (black) still sit
where a folded curve meets the critical line; the three remainder
choices shift the off-line curves but agree on those critical-line
crossings in this window.

\begin{figure}[p]
\centering
\includegraphics[width=\textwidth,height=0.78\textheight,keepaspectratio]{fig_ps_ak_r2_legs_angles}
\caption{PS, AK and $R/2$ legs and angles in the window
$4.65\le T\le4.80$. Left: equal-leg loci $L_1=L_2$ for three choices of
bisector point: PS $B_1=\Sigma_1+R_{1ps}$ (blue), AK
$B_1=\Sigma_1+R_{1ak}$ (green), and
$R/2$ $B_1=\Sigma_1+(R_{1ak}+R_{2ak})/2$ (purple). Right: folded-leg
loci $\vartheta_2=\pi$ for the same three choices (red, orange, and
hollow purple).
The blue vertical is the critical line; black dots are zeta zeros.
Equal-leg and folded-leg loci computed from the definitions
$L_1=L_2$ and $\vartheta_2=\arg((\zeta-B_1)/B_1)$ for each choice of
$B_1$. Generated by \texttt{fig\_ps\_ak\_r2\_legs\_angles.py}.}
\label{fig:ps-ak-r2}
\end{figure}

Figure~\ref{fig:ps-ak-r2-oval-9441} is the same three-way overlay for a
pair of neighboring equal-leg ovals near $T\approx9.44$, drawn in the
window $9.415\le T\le9.465$. The lower PS oval sits around the zero at
$T\approx9.431$ and the upper around $T\approx9.451$; a third zero at
$T\approx9.440$ lies in the gap between them (not inside either oval).
The AK and $R/2$ ovals nest with the PS locus as in
Figure~\ref{fig:ps-ak-r2}.

\begin{figure}[p]
\centering
\includegraphics[width=\textwidth,height=0.78\textheight,keepaspectratio]{fig_ps_ak_r2_oval_9441}
\caption{PS, AK and $R/2$ legs and angles for two neighboring equal-leg
ovals near $T\approx9.44$, in the window $9.415\le T\le9.465$.
Left: equal-leg loci $L_1=L_2$ for PS (blue), AK (green), and $R/2$
(purple). Right: folded-leg loci $\vartheta_2=\pi$ for the same three
choices (red, orange, and hollow purple). The blue vertical is the
critical line; black dots are zeta zeros. The lower oval surrounds the
zero at $T\approx9.431$ and the upper the zero at $T\approx9.451$; the
zero at $T\approx9.440$ sits in the gap between them. Loci computed as in
Figure~\ref{fig:ps-ak-r2}. Generated by
\texttt{fig\_ps\_ak\_r2\_legs\_angles.py}
(\texttt{--out fig\_ps\_ak\_r2\_oval\_9441 --t-lo 9.415 --t-hi 9.465}).}
\label{fig:ps-ak-r2-oval-9441}
\end{figure}

\paragraph{Intersecting equal-leg ovals and the Riemann Hypothesis.}
The nested ovals in Figures~\ref{fig:ps-ak-r2}
and~\ref{fig:ps-ak-r2-oval-9441} suggest a comparison that uses only exact
remainders. Write
\[
B_{1ps}
=
\Sigma_1+R_{1ps},
\qquad
B_{1rs}
=
\Sigma_1+R_{1rs},
\qquad
R_{1rs}
:=
\tfrac12 R,
\]
where $R=R_{1ps}+R_{2ps}=\zeta-\Sigma_1-\Sigma_2$ is Siegel's remainder, so
$R_{1rs}$ is the exact half-remainder (``$rs$'' for Riemann--Siegel, parallel to
the subscripts ``$ps$'' and ``$ak$'') and $B_{1rs}$ is the midpoint of
the remainder segment from $\Sigma_1$ to $\Sigma_1+R$. Let
$\mathcal{E}_{\mathrm{ps}}$ and $\mathcal{E}_{\mathrm{rs}}$ be the corresponding
equal-leg loci in the strip:
the sets of $(\sigma,T)$ at which $|B_1|=|\zeta-B_1|$ for
$B_1=B_{1ps}$ and $B_1=B_{1rs}$ respectively.

If $\zeta(\sigma+iT)=0$, then for either choice of bisector point one has
$L_2=|\zeta-B_1|=|B_1|=L_1$, so
\[
\frac{L_1}{L_2}=1
\]
holds simultaneously for both $B_{1ps}$ and $B_{1rs}$. Every
zero therefore lies in
$\mathcal{E}_{\mathrm{ps}}\cap\mathcal{E}_{\mathrm{rs}}$. On the critical line
this is automatic: Corollary~\ref{cor:equal-legs} puts the whole line in
$\mathcal{E}_{\mathrm{ps}}$, and the $rs$ split is reflection-symmetric there
as well: $\Sigma_2=\chi\,\overline{\Sigma_1}$ and $R=\chi\,\overline{R}$ (the
latter because $e^{-i\psi/2}R$ is real, Appendix~\ref{app:lean-critical}), so
\[
\chi\,\overline{B_{1rs}}
=\Sigma_2+\tfrac12 R
=\zeta-B_{1rs},
\]
and since $|\chi|=1$ this gives
$\bigl|B_{1rs}\bigr|=\bigl|\zeta-B_{1rs}\bigr|$: the whole
line lies in $\mathcal{E}_{\mathrm{rs}}$ as well. (Note that
$B_{1rs}\neq\zeta/2$; only its perpendicular projection onto the
base is $\zeta/2$, by Corollary~\ref{cor:bisector-proj}, which is how
Figure~\ref{fig:splits-h} draws it, at offset $h_{R/2}$ from $\zeta/2$.) The
two loci therefore intersect everywhere
along $\sigma=\tfrac12$; that common set is expected and does not constrain zeros
beyond the folded-leg condition already discussed.

Off the critical line the situation is different. An off-line zero would be an
off-line point of $\mathcal{E}_{\mathrm{ps}}\cap\mathcal{E}_{\mathrm{rs}}$: a
point where the PS and $rs$ equal-leg ovals meet. Consequently, an off-line
intersection of $\mathcal{E}_{\mathrm{ps}}$ and $\mathcal{E}_{\mathrm{rs}}$ is
\emph{necessary} for an off-line zero (it is not sufficient: the folded-leg
condition $\vartheta_2=\pi$, or equivalently $\zeta=0$, is still required). The
contrapositive is the statement of interest:

\begin{quote}
\emph{If $\mathcal{E}_{\mathrm{ps}}$ and $\mathcal{E}_{\mathrm{rs}}$ have no
point in common with $\sigma\neq\tfrac12$, then $\zeta$ has no zero off the
critical line.}
\end{quote}

In particular, if the two loci are disjoint off the line throughout the strip,
the Riemann Hypothesis follows; if they are disjoint off the line for all
$T\le T_{\max}$, then there is no off-line zero with index at most $T_{\max}$.
The figures above are consistent with a strict nesting (the ovals approach one
another but do not cross off the line). A proof that this nesting persists for
all $T$, with due care at the pole heights of $d_1$ in
\S\ref{sec:pole-locations} where $B_{1ps}$ ceases to be defined in
the usual way, would therefore be a proof of the Riemann Hypothesis.

In short: on the critical line the leg-length ratio is always one
(Corollary~\ref{cor:equal-legs}, and likewise for $B_{1rs}$), so the
zeros there are determined by the legs folding back,
$\vartheta_2=\pi$, for either the $ps$ or the $rs$ choice of bisector point.
Off the line the length condition is no longer free: a zero can occur only
where $L_1/L_2=1$ for the legs, and since that must hold for both choices it
is a necessary condition that
\begin{equation}\label{eq:both-ratios-one}
\frac{L_1}{L_2}=1
\end{equation}
simultaneously for the $ps$ legs and for the $rs$ legs.

\paragraph{The mean leg-imbalance metric.}
The equal-leg loci above are the zero sets of a continuous imbalance. For each
of the three named splits of \S\ref{sec:remainders-summary}, write the two
leg lengths
\begin{equation}\label{eq:leg-lengths}
L_1 \;=\; |\Sigma_1+R_1|,
\qquad
L_2 \;=\; |\Sigma_2+R_2|,
\end{equation}
where $(R_1,R_2)$ is $(R_{1ps},R_{2ps})$, $(R_{1rs},R_{2rs})=(\tfrac12 R,\tfrac12 R)$,
or $(R_{1ak},R_{2ak})$ respectively. The normalized length imbalance of a
split is the scale-free quantity
\begin{equation}\label{eq:delta-leg}
\delta
\;=\;
\frac{|L_1-L_2|}{L_1+L_2+\varepsilon},
\qquad \varepsilon=10^{-18},
\end{equation}
so $\delta=0$ if and only if $L_1=L_2$ (the bisector condition for that
choice of head). Averaging the three splits gives the plotted field
\begin{equation}\label{eq:delta-bar}
\bar\delta
\;=\;
\frac{\delta_{\mathrm{ps}}+\delta_{\mathrm{rs}}+\delta_{\mathrm{ak}}}{3}.
\end{equation}
On the critical line each split has $L_1=L_2$ (Corollary~\ref{cor:equal-legs}
and the parallel identities for $rs$ and $ak$), so $\bar\delta\approx0$ and
the strip is dark purple along $\sigma=\tfrac12$. Off the line
$\bar\delta$ rises, and the equal-leg ovals of
Figures~\ref{fig:ps-ak-r2}--\ref{fig:ps-ak-r2-oval-9441} appear as the dark
cellular structures hugging the critical line in
Figure~\ref{fig:leg-imbalance}.

\begin{figure}[H]
\centering
\includegraphics[width=\textwidth]{leg_equality_four_panel}
\caption{Mean leg imbalance $\bar\delta$ of \eqref{eq:delta-bar}
(dark purple $=$ equal leg lengths). Horizontal axis $\sigma\in[0,1]$;
vertical axis the spiral index $T$. Color is $\bar\delta$ on a log scale.
Green markers are zeta zeros at $\sigma=\tfrac12$. The rightmost panel is
the full strip $0\le T\le20$; the first three panels are zooms of the
outlined bands
($4.65\le T\le4.80$, $9.42\le T\le9.46$, $17.2\le T\le17.6$).
The dark vertical band on $\sigma=\tfrac12$ is the simultaneous
equal-leg condition for the $ps$, $rs$, and $ak$ splits. Generated by
\texttt{plot\_leg\_equality\_four\_panel.py}.}
\label{fig:leg-imbalance}
\end{figure}

\subsection{Toward a proof}\label{sec:toward-proof}

The geometry assembled in this section suggests two routes to the Riemann
Hypothesis, listed here in increasing order of how many remainder splits
participate. (A third phenomenon, the exact collinearity of the three
first-half remainders on the critical line, is a signature of the line
itself rather than of the zeros; it is presented separately in
\S\ref{sec:colinearity}.)

\begin{enumerate}
\item \emph{Equal legs and fold-back together, for a single split.}
A zero at $(\sigma,T)$ forces, for any single choice of bisector point, both
leg conditions of \S\ref{sec:legs} at the same input: $L_1=L_2$ \emph{and}
$\vartheta_2=\pi$. On the critical line the length condition is automatic
(Corollary~\ref{cor:equal-legs}), so zeros there are cut out by the fold-back
alone. Off the line the two conditions trace independent families of curves
(the equal-leg ovals and the nearly horizontal folded-leg curves of
Figure~\ref{fig:combined-strips}), and a proof that the two families never
pass through the same off-line point, for even one split, would prove the
Riemann Hypothesis.

\item \emph{Equal legs for two or more splits at once.}
A zero makes the legs equal for \emph{every} split simultaneously: when
$\zeta=0$, $L_2=|\zeta-B_1|=|B_1|=L_1$ for any head $B_1$. An off-line zero
would therefore be an off-line point of
$\mathcal{E}_{\mathrm{ps}}\cap\mathcal{E}_{\mathrm{rs}}$ (and of
$\mathcal{E}_{\mathrm{ak}}$ as well). This route drops the angle condition
entirely and asks only that two equal-leg ovals meet off the line; the
figures of \S\ref{sec:ps-ak-r2} show the ovals nesting without crossing.
This route is developed in detail below.
\end{enumerate}

\paragraph{Route 2 in detail.}
An off-line zero at $(\sigma,T)$ would require both equal-leg
identities at the same input:
\begin{subequations}\label{eq:equal-legs-abs}
\begin{equation}\label{eq:ps-equal-legs-abs}
\bigl|\Sigma_1+R_{1ps}\bigr|
=
\bigl|\Sigma_2+R_{2ps}\bigr|,
\end{equation}
\begin{equation}\label{eq:rs-equal-legs-abs}
\bigl|\Sigma_1+\tfrac12 R\bigr|
=
\bigl|\Sigma_2+\tfrac12 R\bigr|.
\end{equation}
\end{subequations}
(Here \eqref{eq:ps-equal-legs-abs} is the $ps$ condition and
\eqref{eq:rs-equal-legs-abs} the $rs$ condition with $R_{1rs}=\tfrac12 R$.)
Using $R_{2ps}=R-R_{1ps}$ from Theorem~\ref{thm:decomp}, the same pair is
\begin{subequations}\label{eq:equal-legs-abs-R}
\begin{equation}\label{eq:ps-equal-legs-abs-R}
\bigl|\Sigma_1+R_{1ps}\bigr|
=
\bigl|\Sigma_2+R-R_{1ps}\bigr|,
\end{equation}
\begin{equation}\label{eq:rs-equal-legs-abs-R}
\bigl|\Sigma_1+\tfrac12 R\bigr|
=
\bigl|\Sigma_2+\tfrac12 R\bigr|.
\end{equation}
\end{subequations}
From the definitions of the partial sums, with $s=\sigma+iI(T)$,
\begin{equation}\label{eq:Sigma2-from-Sigma1}
\Sigma_2(\sigma,T)
=
\chi(s)\,\overline{\Sigma_1(1-\sigma,T)},
\end{equation}
so $\Sigma_2$ is determined by $\chi$ and the reflected first sum. From the cone
decomposition \eqref{eq:R-as-cone} and the Cramer solution
\eqref{eq:d1}--\eqref{eq:d2}, with $\phi=\arg R$, $\psi=\arg\chi$, and
$\omega=\omega(T)$,
\begin{equation}\label{eq:R1ps-from-R}
R_{1ps}
=
|R|\,
\frac{\sin(\omega-\phi+\psi)}{\sin(2\omega+\psi)}\,
e^{-i\omega}.
\end{equation}
Substituting \eqref{eq:Sigma2-from-Sigma1} into \eqref{eq:rs-equal-legs-abs-R} and
\eqref{eq:Sigma2-from-Sigma1}--\eqref{eq:R1ps-from-R} into
\eqref{eq:ps-equal-legs-abs-R} yields the pair
\begin{subequations}\label{eq:equal-legs-chi}
\begin{equation}\label{eq:rs-equal-legs-chi}
\bigl|\Sigma_1+\tfrac12 R\bigr|
=
\bigl|\chi(s)\,\overline{\Sigma_1(1-\sigma,T)}+\tfrac12 R\bigr|,
\end{equation}
\begin{equation}\label{eq:ps-equal-legs-chi}
\left|
\Sigma_1
+|R|\,
\frac{\sin(\omega-\phi+\psi)}{\sin(2\omega+\psi)}\,
e^{-i\omega}
\right|
=
\left|
\chi(s)\,\overline{\Sigma_1(1-\sigma,T)}
+R
-|R|\,
\frac{\sin(\omega-\phi+\psi)}{\sin(2\omega+\psi)}\,
e^{-i\omega}
\right|.
\end{equation}
\end{subequations}
Equivalently, writing $u=e^{-i\omega}$ and $v=e^{i(\omega+\psi)}$ for the two
unit directions in $R=d_1 u+d_2 v$ from \eqref{eq:R-as-cone}, the same real
coefficients are recovered by projecting $R$ onto the $(u,v)$-frame:
\begin{equation}\label{eq:R1ps-Im-proj}
R_{1ps}
=
\frac{\operatorname{Im}(\overline{v}\,R)}{\operatorname{Im}(\overline{v}\,u)}\,u,
\end{equation}
and likewise $R_{2ps}=R-R_{1ps}$. Substituting \eqref{eq:R1ps-Im-proj} together
with \eqref{eq:Sigma2-from-Sigma1} into \eqref{eq:ps-equal-legs-abs-R} recovers
\eqref{eq:ps-equal-legs-chi}. A proof that
\eqref{eq:equal-legs-chi} has no solution for $\sigma\neq\tfrac12$ (away from
the poles of $d_1$) would therefore prove that $\mathcal{E}_{\mathrm{ps}}$ and
$\mathcal{E}_{\mathrm{rs}}$ never meet off the critical line, and with it the
Riemann Hypothesis.

\subsection{Equal legs density}\label{sec:eqleg-density}

If a zero were to exist off the critical line, it would have to lie on an
equal-leg locus. The reason is the one used in \S\ref{sec:ps-ak-r2}: at a zero
the chain closes, $\zeta=0$, so the bisector point is equidistant from the two
ends of the chain,
\[
L_2=|\zeta-B_1|=|{-B_1}|=|B_1|=L_1 ,
\]
whichever split is used for $B_1$. Every zero, on the line or off it, is
therefore an equal-leg point. On the line that says nothing, the whole line
being equal-leg by Corollary~\ref{cor:equal-legs}; off the line it is a genuine
restriction, because the off-line part of the locus is not the whole strip but
the isolated ovals of Figure~\ref{fig:equal-legs-strips}, symmetric about
$\sigma=\tfrac12$. An off-line zero has to sit on one of those ovals. Their
number and size are therefore worth measuring: they are the only places in the
strip where an off-line zero can be. This subsection reports a numerical census
of them for $0\le T\le401$.

\paragraph{The imbalance and the scan.} The census counts sign changes of the
\emph{signed} imbalance
\begin{equation}\label{eq:g-imbalance}
g(\sigma,T)\;=\;\frac{L_1-L_2}{L_1+L_2},
\end{equation}
which is \eqref{eq:delta-leg} without the absolute value. The modulus is dropped
on purpose: $g$ changes sign across an oval, and it is the sign change that can
be counted. Numerically $g$ is odd under $\sigma\mapsto1-\sigma$, the reflection
that exchanges the two legs, so it vanishes identically on the critical line ---
Corollary~\ref{cor:equal-legs} again --- and the ovals are symmetric about the
line, as the figures show.

The scan fixes $\sigma=\tfrac12+\varepsilon$ and steps in $T$, recording sign
changes of $g$. Each oval contributes two, a bottom and a top, and only bottoms
are counted, so an oval straddling an integer is booked once, to the interval
holding its bottom. The pole bands are excluded: at the heights of
\S\ref{sec:pole-locations} the PS bisector point runs off to infinity and the
equal-leg point sweeps the full width of the strip, which is not an oval. Those
bands are a PS artifact, $599$ of them against $1$ for AK, whose remainder has
no pole to escape through. Roots are refined by bisection to $10^{-13}$. The
step is
\begin{equation}\label{eq:oval-step}
\Delta T\;=\;10^{-4}/m ,
\end{equation}
which holds the number of samples across an oval fixed, the median oval height
falling like $1/T$: $1.3\times10^{-3}$ in the interval at $T=25$ and
$5.2\times10^{-5}$ in the one at $T=400$. Ovals thinner than the step are still caught, by minimizing $g$
wherever it has a positive local minimum below a twentieth of the interval's
median and testing whether the minimum dips below zero.

Four checks bound the error. First, $\zeta$ evaluated as $\Sigma_1+\Sigma_2+R_{ak}$
agrees with \texttt{mpmath} to $2\times10^{-10}$ at $T=150$ and
$3\times10^{-8}$ at $T=1000$, and $g$ vanishes on the critical line to machine
precision. Second, recounting the intervals at $m=250$, $275$, $299$, $300$, $350$, $398$,
$399$ and $400$ with $\Delta T/2$ and $\Delta T/4$ returns the same count every
time, for both splits, and finds the same thinnest cross-section in each. Third, the window $3\le T\le4$ reproduces
Figure~\ref{fig:equal-legs-strips}, the same four ovals at $T=3.0105$,
$3.3845$, $3.6031$ and $3.7257$. Fourth, the $\sigma$-extent of each oval was
measured by bisection, which is what decides whether a scan at a given
$\varepsilon$ can see it: at $\varepsilon=0.001$ every oval in the sample is
reachable and at $\varepsilon=0.01$ some $99.95\%$ are, so the count quoted
below, taken at $\varepsilon=0.01$, has saturated.

\paragraph{Are they ovals?} Drawn against the index they are hairlines. The
component at $T=300$ in Figure~\ref{fig:eqleg-shape} is $0.39$ wide in $\sigma$
and $6.9\times10^{-5}$ tall in $T$, a ratio of $5600$ to one, and every strip
figure here stretches the vertical axis by some such factor to make the locus
visible at all. That shape is an artifact of the coordinate. One unit of index
is $I'(T)=dt/dT$ units of $t$, and $I'(300)=3789$, so plotting against $T$
compresses the vertical direction by exactly that factor; in the $s$-plane,
where $\zeta$ is evaluated, the same component is $0.26$ tall against $0.39$
wide. It is an oval, and close to an ellipse: comparing its half-height at
eight fractions of its half-width $a$ against $\sqrt{1-(u/a)^2}$, the one drawn
departs by $1.8\%$ rms. Over a sample of $300$ of them at nine heights from
$T=5$ to $T=400$, drawn at random within each interval, the median departure is
$0.5\%$, with $80\%$ inside $2\%$ and $94\%$ inside $5\%$; the ratio of width to
height has median $1.22$ and ranges from $0.96$ to $2.83$, with no trend in
height, the interval medians sitting between $1.13$ and $1.27$ throughout. So
the components are convex, smooth, nearly elliptical, a little wider than tall,
and they keep their shape as they shrink --- which is what the word oval is
doing in this section. The only components that are not ovals are the ones cut
off by the edge of the strip instead of closing inside it, and there are $72$ of
those in the census, all below $T=100$.

\begin{figure}[H]
\centering
\includegraphics[width=\textwidth]{fig_eqleg_shape}
\caption{Why the components of the equal-leg locus are called ovals. Left: one of them at $T=300$
plotted against the index, as the strip figures plot them, the vertical axis
stretched by $5600$ to make it visible. Middle: the same component in the
$s$-plane with the axes to a common scale, and the ellipse through its extremes
dashed over it; the dot is the zero of $\zeta$ it encloses. Right: components at
$T=25$, $100$ and $300$ to a common scale, each drawn about its zero. They shrink
with height and hold their shape. All three are the median-height component of
their interval, taken among those whose zero sits nearest the middle, since the
zero is central only on average. Generated by \texttt{fig\_oval\_shape.py} of
the equal-leg density companion note.}
\label{fig:eqleg-shape}
\end{figure}

\paragraph{The count.} Over $0\le T\le401$, scanning the line
$\sigma=\tfrac12+\varepsilon$ with $\varepsilon=0.01$, there are $328\,511$
ovals for the AK remainder split $B_1=\Sigma_1+R_{1ak}$ and $328\,191$ for the
PS split $B_1=\Sigma_1+R_{1ps}$, against $1\,771\,661$ zeros of $\zeta$ in the
same range. The two splits converge on each other with
height: the mean disagreement per unit interval grows from $1.5$ ovals to
$5.0$, but falls from six percent of the count to $0.32\%$. The oval density of
both splits tracks the density of the zeros. Writing the interval's zeros per oval as a function of height, and
fitting over $3\le T\le401$ weighted by the oval count,
\begin{equation}\label{eq:oval-ratio}
\frac{\#\{\text{zeros in }(m,m+1)\}}{\#\{\text{ovals in }(m,m+1)\}}
\;=\;4.38+\frac{11.14}{\ln(t/2\pi)} ,
\end{equation}
which is unbiased in every bin above $T=100$, to $\pm0.01$. A power law
$6.74\,T^{-0.040}$ fits the per-interval scatter equally well --- that scatter
is dominated by fluctuation and cannot separate the two --- but its binned
residuals march in one direction, and it sends the ratio to zero. The law
\eqref{eq:oval-ratio} instead levels off: the ovals settle at a fixed share of
the zeros, about one oval per four and a half, so with
$dN/dT=I'(T)\ln(t/2\pi)/2\pi\approx4T\ln T$ the oval density is asymptotically
$0.9\,T\ln T$. Figure~\ref{fig:eqleg-density} plots both.

\begin{figure}[H]
\centering
\includegraphics[width=\textwidth]{fig_eqleg_density}
\caption{Above: equal-leg ovals and zeros of $\zeta$ per unit interval of the
index, $1\le T\le401$, at $\sigma=\tfrac12+0.01$, log-log. The two clouds are
parallel. Below: the interval's zeros per oval, with \eqref{eq:oval-ratio} solid
and its limit $4.38$ dotted, against the power law dashed. The two are
indistinguishable over $50\le T\le400$ and part company on either side.
Generated by \texttt{census\_analyze.py} of the equal-leg density companion
note.}
\label{fig:eqleg-density}
\end{figure}

\paragraph{Where an off-line zero could be.} Since the ovals are the only
off-line home for a zero, their extent is a bound on where to look, and it can
be quoted in both directions. In height: each oval is about $0.49$ of a mean
zero spacing tall, a ratio steady to within $0.01$ above $T=100$, and there are about a fifth as many ovals as zeros, so the
ovals cover $8.8\%$ of the index axis over $350\le T<401$, up from $6.4\%$ over
$3\le T<25$, a share growing like $0.66\ln T+5.0$ per cent. An off-line zero must
have its height inside that $8.8\%$, a figure read
off the $\sigma=\tfrac12+0.01$ cross-section and so a slight underestimate
of the oval's full height at the line. In $\sigma$: the ovals draw in toward the critical line as $T$ grows. The
median oval reaches $\sigma=0.98$ at $T=5$ but only $\sigma=0.67$ at $T=400$;
the ninetieth percentile falls from the far wall of the strip to $\sigma=0.77$;
and of the $258\,814$ ovals above $T=200$, just four reach $\sigma=0.95$.
Figure~\ref{fig:eqleg-reach} shows both measurements.

The tail does not vanish, though, and the distinction matters. What retreats is
the bulk of the ovals, not the widest of them: the widest oval in each interval
still reaches $\sigma\approx0.9$ at every height measured --- never below
$\sigma=0.88$ for $T\ge50$ --- and the widest above
$T=200$ gets to $\sigma=0.962$. What thins is their number, from about eight
ovals per unit interval reaching $\sigma\ge0.9$ over $50\le T<100$ to four over
$100\le T<200$, two over $200\le T<301$ and one over $301\le T<401$, while the
total count per interval grows five-fold. A zero well off the line at large $T$
is therefore not excluded by any
of this; it would have to lie on the one or two such ovals per unit interval
rather than on any of the eighteen hundred, which is a sharper statement than the
count alone gives.

\begin{figure}[H]
\centering
\includegraphics[width=\textwidth]{fig_eqleg_reach}
\caption{Above: how far off the critical line the ovals extend, from
half-widths measured by bisection for every one of the $328\,511$ ovals in the
census, one point per unit interval, from $T=3$ where the first ovals are. The
band is the middle half of the ovals in that interval and the dots are the widest
oval in it. The bulk draws in toward
the line as $T$ grows while the widest ovals do not. Below: the percentage of the
index axis covered by ovals, which is the share of heights at which an off-line
zero could sit, rising like $0.66\ln T+5.0$. Generated by
\texttt{reach\_census.py} and \texttt{fig\_reach.py} of the equal-leg density
companion note.}
\label{fig:eqleg-reach}
\end{figure}

\paragraph{One zero per oval.} Every oval has a zero, with a handful of
accountable exceptions; not every zero has an oval, and that direction fails
badly rather than by exception.
Of the $328\,511$ AK ovals, $328\,467$ --- all but $44$, or $99.987\%$ ---
bracket exactly one zero of $\zeta$ on
the critical line, and the zero sits at the oval's centre on average rather than
individually: its position within the oval has median $0.5000$ and mean $0.5000$
of the height, so there is no systematic offset, but the spread about that is
$0.14$, and while $91\%$ of zeros are in the middle half, one oval in eleven
holds its zero outside it and one in fifty outside the range $0.16$ to $0.84$. The exceptions are few and they are explicable. Twenty-three
ovals bracket two zeros, none of them below $T=201$, and every one is an
unusually close pair, separated by
$0.25$ to $0.65$ of the local mean spacing --- which is what an oval half a
spacing tall would have to straddle to catch two. Twenty-one bracket none:
eighteen are ovals that span an integer, where the chain gains a link, with
their zero just across the boundary; one is the degenerate interval
below $T=1$, where $\Sigma_1$ is empty; and two, at $T=393$ and $T=398$, are
ovals whose $\sigma$-tip stops a hair past the scan line, at $\sigma=0.51004$ and
$\sigma=0.51003$, so that the cross-section the scan cuts is a sliver
$3\times10^{-7}$ tall at the very tip, lying just to one side of the zero's
height --- below it in the first case and above it in the second, by under
$10^{-7}$ either way. Both ovals do hold their zero: at the critical line each is
some $4\times10^{-6}$ tall with the zero inside, and the cross-sections hold it
out to $\sigma=0.50997$. The miss is in the bracket test at
$\varepsilon=0.01$, not in the oval. The exception rate falls
from $0.104\%$ over $3\le T<50$ to $0.007\%$ over $100\le T<200$ and then edges
back up to $0.014\%$ over $301\le T<401$, the rise being the close pairs, whose
number grows with the zero count while staying under one oval in seven
thousand.

The other direction is not close to one for one, and it is the more important
half of the statement. There are $1\,771\,661$ zeros in the range against
$328\,511$ ovals, so five zeros in six have no oval around them, and \eqref{eq:oval-ratio}
says that share settles at about four in five rather than shrinking away. The
map from ovals to zeros is into, not onto: each oval names one zero, most zeros
are named by none, and which ones get named is a question the census does not
answer.

Three things follow. First, the off-line locus is tied to the zeros rather
than being a structure independent of them. Each oval wraps one zero, centred on
it in the mean, about half a mean spacing tall and a fifth wider than that; the ovals are a
selection from the zeros rather than a phenomenon of their own, which is why
their count obeys \eqref{eq:oval-ratio} instead of a law of its own.

Second, it makes the disjointness program of \S\ref{sec:ps-ak-r2} and
\S\ref{sec:toward-proof} local. What is needed for an off-line zero is that
$\mathcal{E}_{\mathrm{ps}}$ and $\mathcal{E}_{\mathrm{rs}}$ meet off the line.
Since each oval encloses exactly one zero, the two ovals that have to be
compared are the PS oval and the $rs$ oval belonging to the \emph{same} zero,
and the global statement ``the two loci have no common point off the line''
becomes one nesting statement per bracketed zero, of which there are about
$0.9\,T\ln T$ per unit index. That is the right shape for an argument: a local
inequality attached to each zero, to be proved uniformly, rather than a
statement about two curves wandering the strip.

Third, the caution. None of this is evidence for or against the Riemann
Hypothesis, and the one-zero-per-oval law in particular is not. The zero counted
inside each oval is found \emph{on} the critical line, by sign changes of
Hardy's $Z$; the census never evaluates $\zeta$ away from the line, so it does
not test whether an oval carries an off-line zero somewhere along its length,
and an oval already holding an on-line zero at its centre is not forbidden from
passing through one. What the census does establish is that over this range of
heights the off-line locus contains no component that an ordinary on-line zero
fails to explain: there are no orphan ovals. It also points at the search that
would test the hypothesis directly, since evaluating $\zeta$ along the ovals is
a one-dimensional problem --- about $0.9\,T\ln T$ closed curves per unit index
--- where searching the strip for off-line zeros is a two-dimensional one.

There is a converse worth stating, because it is the part that is not
understood. An oval always has a zero, but a zero usually has no oval: by
\eqref{eq:oval-ratio} only about $22\%$ of zeros acquire one, and what
distinguishes those from the rest is open. The bracketing requires the
transverse imbalance $\partial g/\partial\sigma$ on the line to change sign on
both sides of the zero, and nothing here explains why it does so for two zeros
in nine and not for the others.

\clearpage
% ======================================================================
\section{\texorpdfstring{$I(T)$}{I(T)} functions}\label{sec:IT-functions}
% ======================================================================

\subsection{The origin of \texorpdfstring{$I(T)$}{I(T)}: the neighboring links in the bisector frame}\label{sec:IT-origin}

Now change the point of view: hold the forward bisector link still. Translate
joint $m$ to the origin and rotate and scale so that the bisector link lies
along the real axis from $0$ to $1$; everything else moves in this \emph{bisector
frame} as $t$ increases. Watching the two links on either side of the bisector
link, the link at $m-1$ behind it and the link at $m+1$ ahead of it, one
sees a strikingly periodic picture: each neighboring link revolves around the
stationary bisector link approximately \emph{twice} (approaching exactly two
revolutions as $T$ grows), and then the whole pattern repeats. The instant the
pattern repeats is precisely the handoff instant of \S\ref{sec:bisector}: the
neighboring link \textbf{\emph{folds back}} onto the bisector link, the bisector point moves
onto that next link and that next link becomes the new bisector link with its own two
neighbors.

\begin{figure}[H]
\centering
\includegraphics[width=\textwidth]{fig_bisector_frame}
\caption{The two neighboring links in the bisector frame as $T$ sweeps
from $4$ to $5$ ($\sigma=\tfrac12$), sixteen snapshots read left to
right from the top left ($T=4$) to the bottom right ($T=5$). In every
panel the bisector link, link $m=4$ of the forward chain, is held
fixed along the real axis from $0$ to $1$ (black); link $3$ (green)
hangs off joint $4$ at the origin and link $5$ (red) off joint $5$ at
the point $1$, both scaled relative to the unit bisector link. The
dotted curved arrow on each neighbor runs from the middle of the link
to the middle of that link as the next panel depicts it, tracing the
motion the link takes between panels. The black dot on the bisector
link marks the bisector point, at the crossing fraction
$\lceil T\rceil^{\sigma}d_1$ along the unit link. As $T$
advances, each neighbor revolves around the stationary bisector link
approximately twice. At $T=4$ (Start) the green link is folded back
onto the bisector link (joint angle $9\pi$, an odd multiple of $\pi$),
so the bisector point lies on links $3$ and $4$ simultaneously: that is
the handoff instant at which link $4$ became the bisector link. At
$T=5$ (End) it is the red link that folds back (joint angle $11\pi$),
the bisector point lies on links $4$ and $5$ simultaneously, and
link $5$ takes over as the new bisector link. Generated by
\texttt{fig\_bisector\_frame.py}.}
\label{fig:bisector-frame}
\end{figure}

This periodicity begs for a clock that ticks once per handoff, and that is
where $I(T)$ came from: we sought a formula $t=I(T)$ whose input $T$ passes
through an integer at exactly the moments the bisector point moves from one
link to the next. Two facts about the joint angle between the outgoing and
incoming bisector links pin the formula down.

\begin{enumerate}
\item \emph{From the zeta function.} Consecutive summands of $\Sigma_1$ differ
in phase by $t\ln\frac{n+1}{n}$, so the turn at the joint where link $T-1$
meets link $T$ (interpolating the index continuously) is
\begin{equation}\label{eq:jangle}
j_{\angle}(T) \;=\; t\,\log\!\Bigl(\tfrac{1}{T}+1\Bigr).
\end{equation}
\item \emph{From the observed motion.} At a handoff the two links fold back
onto one another, so the joint angle there is an \emph{odd multiple of $\pi$}
(hence a summand $\pi$); and the neighboring link revolves around the
stationary bisector link \emph{twice} per unit of the index, so the odd
multiple advances by two with each unit of $T$ (hence the factor $2$). Interpolating
between handoffs, the joint angle at index $T$ is
\begin{equation}\label{eq:jangle-count}
j_{\angle}(T) \;=\; \pi\,(2T+1).
\end{equation}
\end{enumerate}

Setting \eqref{eq:jangle} equal to \eqref{eq:jangle-count} and solving for $t$
yields
\begin{equation}\label{eq:IT-derived}
\pi\,(2T+1) \;=\; t\,\log\!\Bigl(\tfrac{1}{T}+1\Bigr)
\qquad\Longrightarrow\qquad
t \;=\; \frac{\pi\,(2T+1)}{\log\!\bigl(\tfrac{1}{T}+1\bigr)} \;=\; I(T),
\end{equation}
which is the mapping adopted in \S\ref{sec:IT}: under $t=I(T)$ the bisector
point resides on link $\lfloor T\rfloor$, and $T$ is an integer exactly when
the bisector point moves from one link to the next.

\begin{remark}
Taylor-expanding the reciprocal logarithm in \eqref{eq:IT-derived},
$1/\log(1+\tfrac1T)=T+\tfrac12-\tfrac1{12T}+\tfrac1{24T^2}-\cdots$, and
collecting orders of $T$ gives a family of invertible polynomial
approximations to $I(T)$; keeping successively more terms:

\begin{center}
\begin{tabular}{c|l|l}
terms & $I(T)\;\approx$ & solved for $T$\\
\hline
$2$ & $2\pi\bigl(T^2+T\bigr)$
    & $T=-\tfrac12+\sqrt{\tfrac{t}{2\pi}+\tfrac14}$\\[4pt]
$3$ & $2\pi\bigl(T^2+T+\tfrac16\bigr)$
    & $T=-\tfrac12+\sqrt{\tfrac{t}{2\pi}+\tfrac1{12}}$\\[4pt]
$4$ & $2\pi\bigl(T^2+T+\tfrac16-\tfrac{1}{180T^2}\bigr)$
    & $T\approx-\tfrac12+\sqrt{\tfrac{t}{2\pi}+\tfrac1{12}+\tfrac{\pi}{90t}}$
\end{tabular}
\end{center}

\noindent
(The $1/T$ term of the expansion cancels, so the fourth term is already
$O(T^{-2})$.) The $2$-term form misses $I(T)$ by a constant
${\approx}\,\pi/3$; the $3$-term form is accurate to $O(T^{-2})$ and the
$4$-term form to $O(T^{-3})$. At $T=20$ their errors in $T$ after
inversion are $4\times10^{-3}$, $3\times10^{-7}$ and $4\times10^{-11}$
respectively. The first two rows invert exactly (quadratics); the
$4$-term row is inverted by substituting the $3$-term inverse into the
small $T^{-2}$ correction.

The table also places Siegel's cutoff \eqref{eq:siegel-m} on the index. The
$3$-term row is $t/2\pi=(T+\tfrac12)^{2}-\tfrac1{12}+O(T^{-2})$, so
\begin{equation}\label{eq:siegel-cutoff-index}
\sqrt{\frac{t}{2\pi}}\;=\;T+\frac12-\frac{1}{24T}+O(T^{-2}),
\end{equation}
which is right to $5\times10^{-4}$ at $T=6.18$, $5\times10^{-5}$ at $T=20$ and
$2\times10^{-6}$ at $T=100$. Siegel's $m$ therefore steps \emph{near} the
\emph{half}-integers of $T$, a shade past them by $1/24T$, while ours steps
\emph{exactly at} the integers.
\end{remark}

\begin{figure}[H]
\centering
\includegraphics[width=0.88\textwidth]{fig_IT}
\caption{The index map $t=I(T)$ of \eqref{eq:IT} (solid), together with the
two-term approximation $2\pi(T^{2}+T)$ from the Taylor table above (dashed).
Generated by \texttt{fig\_IT.py}.}
\label{fig:IT}
\end{figure}

We stress that the polynomial forms of the table are conveniences only:
everywhere in this paper $t=I(T)$ is used \emph{exactly}. The defining
property of the mapping, that the neighboring link folds back onto the
bisector link \emph{exactly} as $T$ passes an integer and the bisector
point transfers from one link to the next exactly at that instant,
holds only for the exact $I(T)$. Under any approximation, including the one whose
inverse is essentially Siegel's $\sqrt{t/2\pi}$, the fold instants drift
off the integers by the corresponding error in the table: the clock
still ticks once per handoff, but never precisely on the hour.

\begin{remark}[comparing the classical $t$ rescaling so that zeros have unit
mean spacing, and $I(T)$]\label{rem:unfolding}
In random-matrix and zeta-statistics work the zero-counting function, which
maps each zero $\gamma\mapsto N(\gamma)$ (or uses
$\tilde\gamma=\frac{\gamma}{2\pi}\log\frac{\gamma}{2\pi}$, the leading term of
\eqref{eq:N-rvm}), leaves the zeros with unit mean spacing, which lets
one compare zero statistics at different heights. This mapping is called
rescaling the zeros. Our $I(T)$ is a change of
variable in the same spirit, a reparameterization of the height $t$, but it
runs in the opposite direction, compressing ever more zeros into each unit
interval of the index $T$. For example, the interval $[5,6]$ carries $42$
zeros, $[10,11]$ carries $106$, $[20,21]$ carries $256$, and $[50,51]$ carries
$802$. Rescaling the zeros flattens the count to unit spacing; the index does
the reverse.
\end{remark}

\subsection{\texorpdfstring{$I(T)$}{I(T)} functions for other \texorpdfstring{$L$}{L}-functions besides the zeta function}\label{sec:IT-Lfunctions}

The same change-of-variable idea extends from $\zeta$ to Dirichlet
$L$-functions. On the critical line the partial sums of
\begin{equation}\label{eq:Lchi}
L(s,\chi_p)
\;=\;
\sum_{n=1}^{\infty}\chi_p(n)\,n^{-s}
\end{equation}
form a spiral just as for zeta, but now the character
$\chi_p(n)=\bigl(\tfrac{n}{p}\bigr)$ (Legendre symbol, for a prime modulus
$p$) inserts periodic sign flips and, when $p\mid n$, zero-length
\emph{phantom} links. Counting only the nonzero links as \emph{links},
there are $p-1$ links and one phantom in every block of $p$ summands, and
the natural index $T$ is the link index (phantoms receive no $T$).

Watching these link chains as $t$ varies, we noticed that the bisector point
still moves from one link to the next, just as for zeta, but not always when
links fold back on one another: sometimes the bisector point moves from link to link at a
\emph{joint}, as the bisector line sweeps through it. So $L$-functions have
other event types that zeta does not. An \emph{event} is the instant the
bisector point moves from one link to the next. There are two ways this can
happen: either the two
links fold back onto one another (the turning angle at their common joint is
an odd multiple of $\pi$), so the bisector point slides across the joint from
the old link onto the new one; or the bisector line sweeps through a joint,
carrying the crossing from one link to its neighbor. For zeta only the first
kind occurs: every handoff of \S\ref{sec:bisector} is a folded event, once
per unit of $T$, and there are no bisector events. For $L(s,\chi_p)$ both
kinds occur, recurring with period $\varphi(p)=p-1$ in the integer index $T$:
\begin{itemize}
\item \emph{folded} events, where the turning angle at the current joint is
$\pm\pi$ (the chain reverses), at $T\equiv0\pmod{p-1}$;
\item \emph{bisector} events, where the joint lies on the perpendicular
bisector of the segment from $0$ to $L(s,\chi_p)$, at the remaining residues
modulo $p-1$.
\end{itemize}
For $p=3$ one has $\varphi(3)=2$: even $T$ are folds and odd $T$ are
bisectors. The map $t=I_p(T)$ is defined so that integer $T$ lands at the
$t$-value of its event.

\begin{figure}[H]
\centering
\includegraphics[width=0.85\textwidth]{fig_L3_spirals}
\caption{The two partial-sum spirals of $L(s,\chi_3)$ at $T=6.125$ on the
critical line ($t=I_3(T)\approx176.13$), drawn with
$\lfloor 2T(T{+}1)p\rfloor=261$ summands. The forward chain (red) starts at
the origin $O$ and winds into its spiral around the analytic value
$L(s,\chi_3)$ (open dot); the reverse chain (green) is its mirror image
across the bisector line (dashed blue), the perpendicular bisector of the
segment from $O$ to $L(s,\chi_3)$ through the labeled midpoint
$L(s,\chi_3)/2$, so its tail sits on $L(s,\chi_3)$ and its head lands at
$O$. Every third summand has $\chi_3(n)=0$: these zero-length phantom links
are drawn light grey as the connectors the chain skips. The labeled link $6$
of each chain is the current bisector link: the two chains cross on it,
on the bisector line. Generated by \texttt{fig\_L3\_spirals.py}.}
\label{fig:L3-spirals}
\end{figure}

Because only a fraction $(p-1)/p$ of the summands are links, the smooth
\emph{base} formula generalizing \eqref{eq:IT} is
\begin{equation}\label{eq:Ip-base}
I_p^{\mathrm{base}}(T)
\;=\;
\frac{4\pi\,T}{(p-1)\,\log\dfrac{pT+(p-1)}{pT-(p-1)}}\,.
\end{equation}
Asymptotically $I_p(T)\sim\dfrac{2p}{(p-1)^{2}}\,\pi\,T^{2}$: the leading
coefficient is $\tfrac32\pi$ for $p=3$ and $\tfrac58\pi$ for $p=5$, against
$2\pi$ for zeta's own $I(T)$. (Zeta is not the $p=1$ member of this family,
since the factor $2p/(p-1)^2$ diverges there; for zeta every summand is a
link and there are no phantoms.) The base does not run through the middle of the
events: it touches $I_p$ at the folds and lies below it in between, missing a
one-signed bulge of period $p-1$ imposed by the character. For $p=3$ the
contact is exact, since at even $T$ the gate of \eqref{eq:I3} vanishes and
$c=2$ turns \eqref{eq:I3-geom} into the base; for $p=5$ the fitted correction
returns to zero at $T\equiv0\pmod4$ to within $2\times10^{-3}$. One restores
the bulge by writing
\begin{equation}\label{eq:Ip-split}
I_p(T)
\;=\;
I_p^{\mathrm{base}}(T)+C_p(T).
\end{equation}

For $p=3$ the character is $[1,-1,0]$. The base specializes to
$I_3^{\mathrm{base}}(T)=2\pi T\big/\log\frac{3T+2}{3T-2}$, and the
correction is built from the character-modulated map
\begin{equation}\label{eq:I3-geom}
c=\frac{3+\cos(\pi T)}{2},
\qquad
I_3^{\mathrm{geom}}(T)
=\frac{c\,\pi\,T}{\log\dfrac{3T+c}{3T-c}},
\end{equation}
together with a small odd-$T$ offset
$\alpha(T)=0.390914+0.01712\,T^{-2}$ fitted to event scans:
\begin{equation}\label{eq:I3}
I_3(T)
=
I_3^{\mathrm{geom}}(T)
+\tfrac12\bigl(1-\cos\pi T\bigr)\,\alpha(T).
\end{equation}
(Equivalently $I_3=I_3^{\mathrm{base}}+C_3$ with
$C_3=I_3^{\mathrm{geom}}-I_3^{\mathrm{base}}
+\tfrac12(1-\cos\pi T)\,\alpha$; the gate vanishes on even $T$, so
$C_3=0$ on folds.) On odd $T=1,\ldots,99$ the residual between scanned
event times and \eqref{eq:I3} has RMSE about $8\times10^{-6}$.

For $p=5$ the same pattern gives
$I_5^{\mathrm{base}}(T)=\pi T\big/\log\frac{5T+4}{5T-4}$ and a fitted
period-$4$ correction
\begin{equation}\label{eq:I5}
C_5(T)
=
0.6206-0.6576\cos\tfrac{\pi T}{2}
+0.0369\cos\pi T+\tfrac{0.169}{T^{3.12}}.
\end{equation}
Figure~\ref{fig:I-I3-I5} compares $I(T)$, $I_3(T)$, and $I_5(T)$ on a common
range of $T$.

\begin{figure}[H]
\centering
\includegraphics[width=0.92\textwidth]{fig_I_I3_I5_compare}
\caption{Top: index maps on $T\in[1,20]$: zeta's $I(T)$ (blue), $I_3(T)$ for
$p=3$ (orange), and $I_5(T)$ for $p=5$ (green).
Dots mark the fold events, at $T\equiv0\pmod{p-1}$: every integer on the
zeta curve, every other integer on $I_3$, and every fourth integer on
$I_5$. Leading growth is quadratic with coefficients $2\pi$, $\tfrac32\pi$,
and $\tfrac58\pi$ respectively. Thin dashed black: the smooth base maps
\eqref{eq:Ip-base} for $p=3,5,7$. The $p=3$ and $p=5$ bases are hidden under
$I_3$ and $I_5$ at this scale; the $p=7$ base is the lowest curve, drawn for
comparison, no $I_7$ having been fitted here. Bottom: the character bulge on
its own scale, $I_p-I_p^{\mathrm{base}}$, which the top panel cannot resolve
because it is of order $1$ against values in the thousands. It is one-signed,
returning to zero at the folds (dots) and rising to ${\approx}\,0.92$ between
them for $p=3$ and ${\approx}\,1.33$ for $p=5$, and its amplitude barely
decays: the $p=3$ peaks are $0.9253$, $0.9183$, $0.9153$ at $T=3,5,11$. So the
correction is a fixed additive offset, not a fixed relative one. Open black
circles are the numerically scanned events at every integer $T$, measured
against the same base: the folds sit on zero, the intermediate events ride the
bulge, and all of them follow the modeled curves to within
$1.7\times10^{-5}$ for $p=3$ and $1.0\times10^{-2}$ for $p=5$ on $T\le20$.
Generated with the Dirichlet $L$-function companion note.}
\label{fig:I-I3-I5}
\end{figure}

\clearpage
% ======================================================================
\section{The yin and yang curves}\label{sec:yinyang}
% ======================================================================

As $T$ increases the bisector point travels back and forth along the bisector
link, and when
$T$ passes an integer it moves to the next link of the chain. It does so
without any jump: the handoff happens exactly at the instant when the outgoing
and incoming links fold back onto one another and momentarily occupy the same
line in the plane (when $\sigma=\tfrac12$ the two bisector links are then
coincident; for $\sigma\ne\tfrac12$ they are parallel). The path of the
bisector point is therefore continuous both as a curve in the complex plane and
as a position along the one-dimensional link structure.

\S\ref{sec:IT-origin} watched the two links \emph{neighboring} the bisector
link in the bisector frame. Now watch the \emph{reverse bisector link} in that
same frame: holding the forward bisector link stationary while $T$ increases,
the reverse bisector link revolves around it, and each end of the revolving
link follows a path reminiscent of a teardrop, the pair of paths together
resembling a yin-yang (Figure~\ref{fig:yinyang}). We call the paths of the two
ends the \emph{yin} and \emph{yang} curves, and this section develops them:
first their formulas, then the derivation of $R_{1ps}$ and $R_{2ps}$ from
them, and finally their behavior as $T\to\infty$, where a familiar function
makes an unexpected appearance.

\begin{figure}[htbp]
\centering
\includegraphics[width=\textwidth]{fig_yinyang_spirals}
\caption{The two bisector links at $\sigma=\tfrac14$, $T=6.18$
($t=I(T)\approx 279.85$, $m=6$). Both panels are drawn in the coordinate
frame of the \emph{forward bisector link}: the whole picture is translated
and rotated so that this link lies along the positive $x$-axis, with no
rescaling. All lengths are true, so the link runs from $0$ to
$7^{-1/4}\approx0.615$. \emph{Left:} the forward chain (blue, starting at the
world origin $O$) and the reverse chain (green, ending at $\zeta$), each
drawn out to the link nearest its spiral center; link $6$ of each chain
(the forward and reverse bisector links) is overdrawn with a thick stroke,
and the red dot where they cross is the bisector point.
\emph{Right:} zoom into the gray window: the two bisector links crossing at
the bisector point, which sits on the $x$-axis at distance $d_1$ from the
origin. The bisector frame of this section additionally rescales this link
to the unit interval $[0,1]$ and watches the thick reverse link revolve
around it. Generated by \texttt{fig\_yinyang\_spirals.py}.}
\label{fig:yinyang-spirals}
\end{figure}

\begin{figure}[htbp]
\centering
\includegraphics[width=0.82\textwidth]{fig_bisector_sweep}
\caption{The two bisector links in the frame of the forward bisector link as
$T$ sweeps from $6$ to $7$ ($\sigma=\tfrac14$, the frame of
Figure~\ref{fig:yinyang-spirals}, lengths unscaled). The forward bisector
link is stationary (the single thick dark-blue segment on the $x$-axis)
while the reverse bisector link (green, same thickness) is drawn at sixteen
equally spaced values of $T$, each position labeled on its outer end with the
fractional part of $T$; over the sweep it revolves fully around the
stationary link, and at $T=7$ (labeled $1.00$) it returns nearly to its $T=6$
position. The red dotted arrows run from the middle of each position to the
middle of the next, tracing the motion. The small green and red dots mark
the two ends of the revolving link consistently: green is always the
joint-$m$ end and red the joint-$(m{+}1)$ end; they trace the yin and yang
curves of Figure~\ref{fig:yinyang}, in the same colors. Each red dot on the
$x$-axis is the bisector point at one instant: the
crossing of that position with the $x$-axis, at $(d_1,0)$. Near the parallel
instants (fractional part of $T$ near $\tfrac14$ and $\tfrac34$) the
revolving link is nearly parallel to the stationary one and the crossing
races along the axis (the poles of $d_1$ of \S\ref{sec:pole-locations}). Generated by
\texttt{fig\_bisector\_sweep.py}.}
\label{fig:bisector-sweep}
\end{figure}

The motivation is the question left open in \S\ref{sec:bisector}: the two
bisector links cross, but \emph{where}? The answer matters because, by
\eqref{eq:bisector-point}, the crossing (the bisector point) sits at
distance $d_1$ from the joint $\Sigma_1$ along the forward bisector link and
at distance $d_2$ from the joint $\Sigma_1+R$ along the reverse bisector link.
If we can \emph{compute} the crossing, we obtain $d_1$ and $d_2$, and with
them a new formula for $\zeta$ that uses the fractional summands in place of
the usual remainders, all without evaluating $\zeta$ or either partial sum.

\begin{figure}[htbp]
\centering
\includegraphics[width=0.95\textwidth]{fig_yinyang}
\caption{The yin and yang curves in the bisector frame, $\sigma=\tfrac14$,
$m=6$. The forward bisector link is pinned to the unit interval $[0,1]$
(dark blue). As $T$ sweeps $6\le T\le 7$, the near end of
the reverse bisector link traces the yin curve $Y_{in1}$ (green, one dot
every $0.1$ in $T$) and the far end traces the yang curve $Y_{ang1}$ (red);
the two teardrop paths together resemble a yin-yang. The reverse bisector
link itself is drawn at $T=6.20$ (dark green): it crosses the real axis at
the bisector point (black dot), at the crossing fraction
$\lceil T\rceil^{\sigma}d_1\approx 0.237$ of the way along the unit link. The
dashed violet vector from $0$ to $Y_{in1}$ is the in-frame image of Siegel's
remainder $R$; the crossing splits it into the images of $R_{1ps}$ (red, the
piece of the unit link from $0$ to the bisector point) and $R_{2ps}$ (orange,
from the bisector point to $Y_{in1}$). Generated by \texttt{fig\_yinyang.py}.}
\label{fig:yinyang}
\end{figure}

\subsection{The formulas for the yin and yang curves}\label{sec:yinyang-formulas}

In the bisector frame the forward bisector link is pinned to $[0,1]$, so the
only question is where the \emph{reverse} bisector link is. In the world frame
its chain is anchored at $\zeta$, which is both moving and precisely the
quantity we do not wish to assume known; worse, the reverse bisector link also
moves relative to its own anchor. The resolution is that the reverse
chain's joint $m$ sits at $\Sigma_1+R$: relative to the forward joint
$\Sigma_1$ it is displaced by exactly $R$, and $R$ has its own integral
representation \eqref{eq:R-T}, independent of any evaluation of $\zeta$.
Applying the frame map
\begin{equation}\label{eq:frame-map}
z \;\longmapsto\; (z-\Sigma_1)\,\lceil T\rceil^{\,\sigma+iI(T)}
\end{equation}
(translate joint $m$ to the origin, then divide by the link vector
$\lceil T\rceil^{-\sigma-iI(T)}$), the two ends of the reverse bisector link
land at
\begin{align}
Y_{in1}(\sigma,T) &\;=\; R\,\lceil T\rceil^{\,\sigma+iI(T)}, \label{eq:yin1}\\[2pt]
Y_{ang1}(\sigma,T) &\;=\; Y_{in1} \;-\; \chi\,\lceil T\rceil^{\,-1+2(\sigma+iI(T))},
\label{eq:yang1}
\end{align}
the first being the image of $\Sigma_1+R$ and the second the image of the far
end $\Sigma_1+R-\chi\lceil T\rceil^{\,\sigma+iI(T)-1}$. As $T$ sweeps, these
two endpoints orbit the stationary unit link along the yin and yang curves of
Figure~\ref{fig:yinyang}.

Because the frame map scales the bisector link to unit length, positions along
it are \emph{fractions}: the crossing we are about to compute is the fraction
of the way along the link at which the bisector point sits, and multiplying
that real fraction by the (unscaled) link vector recovers the fractional
summand itself. The actual distance from the joint, the fraction times the
true link length $\lceil T\rceil^{-\sigma}$, is the weight $d_1$.

The reverse side has its own frame and its own pair of curves. Pinning the
\emph{reverse} bisector link to $[0,1]$ instead (translate its joint
$\Sigma_1+R$ to the origin and divide by its link vector
$-\chi\lceil T\rceil^{\,\sigma+iI(T)-1}$), the roles swap: now the forward
bisector link revolves, its near end tracing
\begin{equation}\label{eq:yin2yang2}
Y_{in2}(\sigma,T)=\frac{R\,\lceil T\rceil^{\,1-\sigma}}{\chi\,e^{\,i\omega}},
\qquad
Y_{ang2}(\sigma,T)=Y_{in2}-\frac{\lceil T\rceil^{\,1-2\sigma}}{\chi\,e^{\,2i\omega}},
\end{equation}
with $Y_{in2}$ the image of the joint $\Sigma_1$ and $Y_{ang2}$ the image of
the far end $\Sigma_1+\lceil T\rceil^{-\sigma-iI(T)}$, and with
$\omega=\omega(T)=I(T)\log\lceil T\rceil$ as always.

The next two subsections derive $R_{1ps}$ from $(Y_{in1},Y_{ang1})$ and
$R_{2ps}$ from $(Y_{in2},Y_{ang2})$. Each derivation is in two parts: first
the modulus, found from the axis crossing; then the direction, which is simply
that of the next summand.

\subsection{Derivation of \texorpdfstring{$R_{1ps}$}{R1ps} from the yin and yang functions}\label{sec:derive-r1ps}

\paragraph{First, the modulus.} From its joint, how far along the forward
bisector link do the two links intersect? This distance is $|R_{1ps}|$,
because $R_{1ps}$ is the fractional part of summand $\lceil T\rceil$ of
$\Sigma_1$, and that summand \emph{is} the forward bisector link. In the
frame the link lies along the real axis, so the intersection is the point
where the segment from $Y_{in1}$ to $Y_{ang1}$ crosses the axis. By the
standard segment--axis intersection formula that crossing is
\begin{equation}\label{eq:crossing}
Y_{in1}
\;-\;
\frac{\Im\,Y_{in1}}
     {\Im\,Y_{ang1}-\Im\,Y_{in1}}\,
\bigl(Y_{ang1}-Y_{in1}\bigr),
\end{equation}
a real number. Substituting \eqref{eq:yin1}--\eqref{eq:yang1} and cancelling
$Y_{ang1}-Y_{in1}=-\chi\lceil T\rceil^{\,-1+2(\sigma+iI(T))}$ from the ratio,
\begin{equation}\label{eq:yang3lin}
\eqref{eq:crossing}
\;=\;
R\,\lceil T\rceil^{\,\sigma+iI(T)}
\;-\;
\frac{\Im\!\bigl(R\,\lceil T\rceil^{\,\sigma+iI(T)}\bigr)}
     {\Im\!\bigl(\chi\,\lceil T\rceil^{\,-1+2(\sigma+iI(T))}\bigr)}\,
\chi\,\lceil T\rceil^{\,-1+2(\sigma+iI(T))}.
\end{equation}
Writing $\lceil T\rceil^{\,\sigma+iI(T)}=\lceil T\rceil^{\sigma}e^{i\omega}$
and expanding the imaginary parts in components,
\begin{equation}\label{eq:compact5}
\eqref{eq:crossing}
\;=\;
R\,\lceil T\rceil^{\,\sigma+iI(T)}
\;-\;
\lceil T\rceil^{\,1-\sigma}\,
\frac{\Re R\,\sin\omega+\Im R\,\cos\omega}
     {\Re\chi\,\sin 2\omega+\Im\chi\,\cos 2\omega}\;
\chi\,\lceil T\rceil^{\,-1+2(\sigma+iI(T))},
\end{equation}
and finally, writing $\varphi=\arg R$ and $\psi=\arg\chi$ and collapsing with
the sine difference identity
$\sin(2\omega+\psi)\cos(\omega+\varphi)-\cos(2\omega+\psi)\sin(\omega+\varphi)
=\sin(\omega-\varphi+\psi)$, the whole expression reduces to
\begin{equation}\label{eq:frac1}
\eqref{eq:crossing}
\;=\;
\lceil T\rceil^{\sigma}\,|R|\,
\frac{\sin\!\bigl(\omega-\varphi+\psi\bigr)}{\sin\!\bigl(2\omega+\psi\bigr)}
\;=\; \lceil T\rceil^{\sigma}\, d_1 .
\end{equation}
The reader will recognize the sine quotient: it is, symbol for symbol, the
definition \eqref{eq:d1} of the weight $d_1$. \emph{This is where that
definition came from.} The crossing \eqref{eq:frac1} is the crossing
fraction --- the ``fractional amount'' column of Table~\ref{tab:frac} ---
along the unit link at which the reverse link crosses, and dividing out the frame's
scaling $\lceil T\rceil^{\sigma}$ leaves the actual distance from the joint,
$d_1$.

\paragraph{Second, the direction.} The forward bisector link points along the
next summand of $\Sigma_1$: extending the summation cutoff by one,
\begin{equation}\label{eq:sigma1diff}
\sum_{n=1}^{m+1} n^{-\sigma-iI(T)}
\;-\;
\sum_{n=1}^{m} n^{-\sigma-iI(T)}
\;=\;
\lceil T\rceil^{\,-\sigma-iI(T)}
\end{equation}
(for non-integer $T$, so $m+1=\lceil T\rceil$). Multiplying the fraction
\eqref{eq:frac1} by this link vector undoes the frame map and yields the
fractional summand itself:
\begin{equation}\label{eq:R1-derived}
R_{1ps}
\;=\;
\Bigl(\lceil T\rceil^{\sigma} d_1\Bigr)\cdot\lceil T\rceil^{\,-\sigma-iI(T)}
\;=\;
d_1\,\lceil T\rceil^{\,-iI(T)},
\end{equation}
which is definition \eqref{eq:R1ps-def}.

\paragraph{Other forms.} Two equivalent ways to write $R_{1ps}$ fall out of
the same computation. Subtracting the crossing from $Y_{in1}$ instead of
building it up from $0$ gives
\begin{equation}\label{eq:R1-other-a}
R_{1ps}
\;=\;
R
\;-\;
|R|\,\frac{\sin(\omega+\varphi)}{\sin(2\omega+\psi)}\;
e^{\,i(\omega+\psi)}
\end{equation}
(the subtracted term is exactly $R_{2ps}$, anticipating
Theorem~\ref{thm:decomp}), and keeping the imaginary parts unexpanded gives
the compact
\begin{equation}\label{eq:R1-other-b}
R_{1ps}
\;=\;
R
\;-\;
\chi\,e^{\,i\omega}\,
\frac{\Im\!\bigl(R\,e^{\,i\omega}\bigr)}
     {\Im\!\bigl(\chi\,e^{\,2i\omega}\bigr)} .
\end{equation}

\subsection{Derivation of \texorpdfstring{$R_{2ps}$}{R2ps} from the yin and yang functions}\label{sec:derive-r2ps}

The same procedure in the reverse frame, using
$(Y_{in2},Y_{ang2})$ of \eqref{eq:yin2yang2}, locates the crossing of the
(now revolving) forward bisector link with the real axis at the crossing fraction
\begin{equation}\label{eq:frac2}
\frac{\lceil T\rceil^{\,1-\sigma}}{|\chi|}\,
|R|\,\frac{\sin\!\bigl(\omega+\varphi\bigr)}{\sin\!\bigl(2\omega+\psi\bigr)}
\;=\;
\frac{\lceil T\rceil^{\,1-\sigma}}{|\chi|}\, d_2
\end{equation}
of the way along the unit reverse link, again the ``fractional amount''
column of Table~\ref{tab:frac}, and again the sine quotient is, symbol for
symbol, the definition \eqref{eq:d2} of $d_2$: this is where \emph{it} came
from. (The two normalized fractions are nearly symmetric: since
$|\chi(\sigma+it)|=(t/2\pi)^{\frac12-\sigma}\bigl(1+O(1/t)\bigr)$ and
$\sqrt{t/2\pi}\approx T+\tfrac12$, the prefactor
$\lceil T\rceil^{\,1-\sigma}/|\chi|$ is asymptotically
$\lceil T\rceil^{\sigma}$, matching the forward side.) Multiplying by the
reverse link vector $\chi\,\lceil T\rceil^{\,\sigma+iI(T)-1}$ undoes the
frame map:
\begin{equation}\label{eq:R2-derived}
R_{2ps}
\;=\;
\frac{\lceil T\rceil^{\,1-\sigma}}{|\chi|}\,d_2
\cdot \chi\,\lceil T\rceil^{\,\sigma+iI(T)-1}
\;=\;
d_2\,\lceil T\rceil^{\,iI(T)}\,\frac{\chi}{|\chi|},
\end{equation}
which is definition \eqref{eq:R2ps-def}. As with the forward side there are
equivalent rewritings, for example
\begin{equation}\label{eq:R2-other-a}
R_{2ps}
\;=\;
R \;-\; |R|\,e^{-i\omega}\,
\frac{\sin\!\bigl(\omega-\varphi+\arg\chi(1-\sigma+iI(T))\bigr)}
     {\sin\!\bigl(2\omega+\arg\chi(1-\sigma+iI(T))\bigr)}
\end{equation}
and
\begin{equation}\label{eq:R2-other-b}
R_{2ps}
\;=\;
R \;-\; e^{-i\omega}\,
\frac{\Im\!\Bigl(\dfrac{R}{\chi(1-\sigma+iI(T))\,e^{\,i\omega}}\Bigr)}
     {\Im\!\Bigl(\dfrac{1}{\chi(1-\sigma+iI(T))\,e^{\,2i\omega}}\Bigr)},
\end{equation}
where the reflected factor may be used interchangeably in the argument because
$\arg\chi(1-\sigma+iI(T))=\arg\chi(\sigma+iI(T))$: from
$\chi(s)\chi(1-s)=1$ and the Schwarz reflection
$\chi(\bar s)=\overline{\chi(s)}$ one gets
$\chi(1-\sigma+it)=1/\overline{\chi(\sigma+it)}$, which has the same
argument as $\chi(\sigma+it)$ and reciprocal modulus.

\paragraph{The poles cancel in the sum.} The two weights share the
denominator $\sin(2\omega+\psi)$, and \S\ref{sec:pole-locations} located the
isolated heights at which it produces genuine poles. In the \emph{sum}
$d_1+d_2$ those poles cancel: adding the numerators with the sum-to-product
identity
$\sin(\omega-\varphi+\psi)+\sin(\omega+\varphi)
=2\sin\bigl(\omega+\tfrac{\psi}{2}\bigr)\cos\bigl(\varphi-\tfrac{\psi}{2}\bigr)$
and writing
$\sin(2\omega+\psi)=2\sin\bigl(\omega+\tfrac{\psi}{2}\bigr)
\cos\bigl(\omega+\tfrac{\psi}{2}\bigr)$, the offending factor divides out:
\begin{equation}\label{eq:d1-plus-d2}
d_1+d_2
\;=\;
|R|\;
\frac{\cos\!\bigl(\varphi-\tfrac{\psi}{2}\bigr)}
     {\cos\!\bigl(\omega+\tfrac{\psi}{2}\bigr)} .
\end{equation}
The zeros of the remaining denominator $\cos(\omega+\tfrac{\psi}{2})$ are
matched by zeros of the numerator (at those instants $R$ is parallel to the
common link direction, forcing $\varphi\equiv-\omega\pmod\pi$), so every
singularity of \eqref{eq:d1-plus-d2} is removable: the sum is smooth across
all the poles of its two terms. Figure~\ref{fig:d1-d2-sum} shows this at
$\sigma=0.1$.

\begin{figure}[htbp]
\centering
\includegraphics[width=0.98\textwidth]{fig_d1_d2_sum}
\caption{$d_1$ (red), $d_2$ (green) and their sum (blue) at $\sigma=0.1$,
$1\le T\le 7$. Twice per unit of $T$ the individual weights blow up at the
parallel-link poles of \S\ref{sec:pole-locations}, but the poles are equal and
opposite and the sum passes smoothly through every one of them; the black
dashed curve is the closed form \eqref{eq:d1-plus-d2}, indistinguishable from
the blue sum (they agree to machine precision on the plotted grid). Generated
by \texttt{fig\_d1\_d2\_sum.py}.}
\label{fig:d1-d2-sum}
\end{figure}

This is how the definitions of \S\ref{sec:decomp} were found: the sine
expressions \eqref{eq:d1}--\eqref{eq:d2} are nothing other than the real-axis
crossings of the yang segments, read off in the two bisector frames. Only
afterwards was Theorem~\ref{thm:decomp} proved, confirming that the two
fractional summands so constructed add back to Siegel's $R$ exactly. Note the
economy of the construction: the bisector point, $d_1$, $d_2$, and both
fractional summands are obtained from $R$ and $\chi$ alone. Numerically we
computed $R$ to high accuracy by Kuznetsov's method (\S\ref{sec:kuznetsov}),
which is how the pictures, and then the identity, were checked.

\subsection{Comparison of yin and yang to Siegel's integral result}\label{sec:yinyang-siegel}

Siegel's 1932 paper~\cite{Siegel1932} evaluates the integral in his equation 5
to what he calls the ``desired result''
\begin{equation}\label{eq:siegel-desired}
F(t)\;=\;
\frac{1}{1-e^{-2\pi i t}}
\;-\;
\frac{e^{\pi i t^{2}}}{e^{\pi i t}-e^{-\pi i t}} ,
\end{equation}
which can also be written
\begin{equation}\label{eq:siegel-desired-alt}
F(t)
=\frac{e^{2\pi i t}-e^{\pi i t(t+1)}}{e^{2\pi i t}-1}
=\Bigl(1-e^{\pi i (t^{2}-t)}\Bigr)\sum_{n=0}^{\infty}e^{-2\pi i n t},
\end{equation}
the last form using the geometric series
$1/(1-e^{-2\pi i t})=\sum_{n\ge0}e^{-2\pi i n t}$. Its complex conjugate has
the compact form
\begin{equation}\label{eq:siegel-conj}
\overline{F(t)}
=\frac{e^{-\pi i t^{2}}-e^{-\pi i t}}{2\,i\,\sin(\pi t)} .
\end{equation}

\begin{remark}[regarding $F$ as the source of Siegel's $f(s)$]\label{rem:F-to-f}
$F$ is not merely adjacent to the function $f$ of \eqref{eq:siegel-58}; it is
what $f$ is built from. In his \S3 Siegel multiplies his equation~5, that is
\eqref{eq:siegel-desired}, by $u^{-s}$ and integrates along the bisector of the
first quadrant, and under that operation the two terms of $F$ separate into the
two halves of zeta. The pole term $1/(1-e^{-2\pi iu})$ yields
$-(2\pi)^{s-1}e^{\,i\pi(1-s)/2}\,\Gamma(1-s)\,\zeta(1-s)$, while the Gaussian
term yields $f$ itself:
\begin{equation}\label{eq:f-from-F}
f(s)\;=\;\bigl(e^{\pi i s}-1\bigr)
\int_{0}^{e^{i\pi/4}\infty}
\frac{u^{-s}\,e^{\pi i u^{2}}}{e^{\pi i u}-e^{-\pi i u}}\,\mathrm{d}u
\qquad (\sigma<0),
\end{equation}
the restriction $\sigma<0$ being needed only for convergence at the endpoint
$u=0$; $f$ is continued analytically from there, and \eqref{eq:siegel-58} itself
converges for every $s$. The same separation can be read off
\eqref{eq:siegel-conj}: since $1/(1-e^{2\pi it})=-e^{-\pi it}/(2i\sin\pi t)$,
the numerator term $e^{-\pi it^{2}}$ of $\overline{F}$ is the $f(s)$ side and
$-e^{-\pi it}$ is the $\zeta(1-s)$ side, so the two lobes carry the two halves
in their numerator. The yin-yang function is therefore the common source of both
$\Sigma_1+R_{1ak}$ of \eqref{eq:siegel-f} and its reflected counterpart, which
is one route into the convergence proof left open in
\S\ref{sec:yinyang-infinity}.
\end{remark}

\begin{figure}[htbp]
\centering
\includegraphics[width=0.78\textwidth]{fig_F_curve}
\caption{The trace of $\overline{F(t)}$ of \eqref{eq:siegel-conj} in the
complex plane for $-1.5\le t\le 1.5$ (red dots every $0.1$ in $t$): the curve
forms two congruent teardrop-like lobes, forming a shape reminiscent of the
yin-yang symbol rooted in ancient texts like the \emph{I Ching}. The curve is
point-symmetric through $\tfrac12$. The
apparent poles at integer $t$ are removable, with
$\overline{F(t)}\to\tfrac12-n$ as $t\to n$ (black squares). A curiosity: at
\emph{every} half-integer $t=k+\tfrac12$ the exponent
$t^{2}=k(k+1)+\tfrac14$ has $k(k+1)$ even, so $e^{-\pi i t^{2}}=e^{-i\pi/4}$
always, and the whole expression collapses to just two values,
$\overline{F}(k+\tfrac12)=\tfrac12-(-1)^{k}\tfrac12 e^{i\pi/4}$ (green
diamonds): the curve returns to the point
$\tfrac12+\tfrac12 e^{i\pi/4}\approx0.854+0.354i$ at
$t=\dots,-0.5,1.5,3.5,\dots$ and to the point
$\tfrac12-\tfrac12 e^{i\pi/4}\approx0.146-0.354i$ at
$t=\dots,-1.5,0.5,2.5,\dots$, each revisited every $\Delta t=2$, forever;
they are the meeting points of the lobes. Generated by
\texttt{fig\_F\_curve.py}.}
\label{fig:F-curve}
\end{figure}

We conjecture that as $T\to\infty$ the yin and yang curves converge to a
universal limit curve built from this function, universal in that it does
not depend on $\sigma$. The precise limit is most cleanly written through the
$\Psi$ function, which we do next.

\subsection{The yin and yang curves are not symmetrical}\label{sec:yinyang-asym}

While the name ``yin-yang'' invokes a pictorial similarity as intended, it also
invites a symmetry that is not there. In the taijitu the two lobes are
congruent, each carried onto the other by a half-turn about the center. In our case, the
two curves are the two ends of one moving link, and \eqref{eq:yang1}
says exactly how they are related:
\begin{equation}\label{eq:yang-offset}
Y_{ang1}\;=\;Y_{in1}\;-\;\chi\,\lceil T\rceil^{\,2s-1},
\qquad s=\sigma+iI(T),
\end{equation}
so the yang curve is the yin curve translated by $-v$, where
$v=\chi\lceil T\rceil^{\,2s-1}$ is the revolving link's own vector in the frame,
of length $|\chi|\,\lceil T\rceil^{\,2\sigma-1}$.
Were $v$ the same at every $T$ the two curves would be congruent by
construction. It is not.
Over one unit of the index the direction of $v$ swings through $\pi$ and comes
back, returning to where it began, with net turning zero. The length of $v$
varies too, and it is exactly $1$ only on the critical line, where $|\chi|=1$
and the exponent $2\sigma-1$ vanishes: the revolving link and the stationary link are then equal in length,
which is the $d_1=d_2$ symmetry of \S\ref{sec:legs} seen in this frame.

What the taijitu has is a half-turn about the center of the picture, and here
that center is the midpoint of the stationary link: the half-turn is
$z\mapsto1-z$, which exchanges the link's two ends. So if the yin and yang
curves were related as the taijitu's two lobes are, this map would carry one
onto the other, and it is enough to draw $1-Y_{ang1}$ on top of $Y_{in1}$ and
look.
Figure~\ref{fig:yinyang-asym} plots $Y_{in1}$ against $1-Y_{ang1}$ for
$0\le T\le10$, one loop of each per unit of the index, on the critical line and
at $\sigma=0.05$ and $\sigma=0.95$ above and below it. The
reflection brings the two families almost on top of one another, and both
converge to the single limit teardrop of \S\ref{sec:yinyang-infinity}, but at no
finite $T$ do they coincide. Each loop crosses its reflected partner twice in
every unit of the index, running outside it along one arc and inside along the
other, and the two enclose different areas. That difference in area narrows as
$T$ grows and widens away from the critical line, where the two lobes also trade
which of them is the larger, near the line though not exactly on it. In the
magnified windows of Figure~\ref{fig:yinyang-asym} the yin curve (green) is the
one running outside its reflected partner (red) in the top two rows, and the
inside one in the bottom row, where the order has swapped.

What the critical line
does give exactly is a mirror symmetry, though between the two frames rather
than between the two curves of one frame. At $\sigma=\tfrac12$ we have
$\lceil T\rceil^{\,s}=\lceil T\rceil^{1/2}e^{i\omega}$ and, from
$\zeta=\chi\overline\zeta$ and $\Sigma_2=\chi\overline{\Sigma_1}$ with
$|\chi|=1$, also $R=\chi\overline R$. Comparing \eqref{eq:yin1} with
\eqref{eq:yin2yang2} then gives
\begin{equation}\label{eq:frame-mirror}
\overline{Y_{in1}}=Y_{in2},
\qquad
\overline{Y_{ang1}}=Y_{ang2}
\qquad\bigl(\sigma=\tfrac12\bigr),
\end{equation}
which we confirm to $25$ digits. Off the line both identities fail, the first
by $6.4\times10^{-2}$ and the second by $7.1\times10^{-2}$ at
$\sigma=\tfrac14$, and by the same two amounts at $\sigma=\tfrac34$.

An exact symmetry between the two curves themselves does arrive, but only in the
limit, where the yang curve stops being a second curve at all: as $T\to\infty$
the yang settles onto the very teardrop the yin traces, half a unit of the index
behind. So
the asymmetry of this subsection is a finite-$T$ effect, and the limit erases it.
That limit curve, and the classical function it turns out to be built from, are
the subject of the next subsection.

\begin{figure}[htbp]
\centering
\includegraphics[width=0.94\textwidth]{fig_yinyang_asym}
\caption{The yin curve (green) against the reflected yang curve
$1-Y_{ang1}$ (red), for $0\le T\le10$, which is one loop
of each per unit of the index, out to $t=692.2$. One row per $\sigma$, all on
the same scale: $\sigma=0.05$ on top, the critical line in the middle,
$\sigma=0.95$ below. The half-turn $z\mapsto1-z$
exchanges the ends of the stationary link (dark blue), so under the symmetry the
picture suggests the two families would coincide. \emph{Left:} all ten periods,
the view cropped to the cluster. The stray arcs crossing it belong to the first
period $0\le T\le1$, where $t<13.6$ and the frame is still far from its limit;
the later periods pile up on the limit teardrop, and they do so at every
$\sigma$. \emph{Right:} the boxed window
magnified, chosen away from the two crossings. The families interleave and close
on one another as $T$ grows, and which of the two runs outside depends on
$\sigma$: green outside red at $\sigma=0.05$, red outside green at
$\sigma=0.95$. Generated by
\texttt{fig\_yinyang\_asym.py}.}
\label{fig:yinyang-asym}
\end{figure}

\clearpage

\subsection{The limit curve and the \texorpdfstring{$\Psi$}{Psi} function: the \texorpdfstring{$C_0$}{C0} connection}\label{sec:yinyang-infinity}

As $T\to\infty$ with the fractional part $q=T-\lfloor T\rfloor$
held fixed, the yin curve converges, regardless of $\sigma$, to
\begin{equation}\label{eq:yinf}
Y_{\inf}(q)
\;=\;
1-\Psi(q)\,e^{-2\pi i\,(q^{2}-\frac{1}{16})},
\end{equation}
where
\begin{equation}\label{eq:psi-fn}
\Psi(x)\;=\;
\frac{\cos\!\bigl(2\pi(x^{2}-x-\tfrac{1}{16})\bigr)}{\cos(2\pi x)}
\end{equation}
is the Riemann--Siegel correction function (it appears in \S15.3 of
Titchmarsh's book \cite{Titchmarsh1986} and in \S5.4.1 of Pugh's
thesis \cite{Pugh1998}). The
yang curve follows the \emph{same} universal function half a unit behind in
the argument, $Y_{ang1}(\sigma,T)\to Y_{\inf}(q-\tfrac12)$, but not the
same teardrop: $Y_{\inf}$ is not $1$-periodic, and over one unit interval
the yang end traces the point reflection of the yin teardrop through
$\tfrac12$, the two congruent teardrops together forming the yin-yang of
Figure~\ref{fig:yinyang} (Lemma~\ref{lem:psi-reflection} and
Theorem~\ref{thm:yinf} below make both statements exact). The measured
convergence is uniform in $q$ away from
the removable points below, at rate $O(1/T)$, and essentially independent of
$\sigma$: the maximum of $|Y_{in1}-Y_{\inf}(q)|$ over a $q$-grid is
$0.028,\,0.015,\,0.027$ at $T\approx10$ and $0.0076,\,0.0042,\,0.0075$ at
$T\approx40$ for $\sigma=0.2,\,0.5,\,0.9$ respectively.
Figure~\ref{fig:yinyang-infinity} shows the convergence, including the
pleasing fact that the very first yin path is the zeta function itself: for
$0<T<1$ both partial sums are empty and $\lceil T\rceil=1$, so
$Y_{in1}=R=\zeta(\sigma+iI(T))$. The iconic zeta trajectory is yin orbit
number zero.

The convergence is not merely an observation; it is a theorem, and the
positivity machinery of \S\ref{sec:d1-positive} already contains most of its
proof: by \eqref{eq:yin1} the yin point is $R$ carried through the frame
map, and \S\ref{sec:d1-positive} knows the size, the sign, and the phase of
$R$ to exactly the accuracy needed. The one genuinely new ingredient is a
reflection identity for the amplitude $\Psi$.

\begin{lemma}\label{lem:psi-reflection}
Write $A(x)=2\pi\bigl(x^{2}-\tfrac1{16}\bigr)$ and
$\theta(x)=2\pi x(1-x)-\tfrac{3\pi}8$, the phase already met in
\eqref{eq:cosu}. Then for all real $x$, with removable points filled in by
continuity,
\begin{equation}\label{eq:psi-reflection}
1-\Psi(x)\,e^{-iA(x)}\;=\;\Psi\bigl(x+\tfrac12\bigr)\,e^{\,i\theta(x)},
\end{equation}
and consequently the limit curve satisfies the exact half-lag relation
\begin{equation}\label{eq:yinf-halflag}
Y_{\inf}(q)-Y_{\inf}\bigl(q-\tfrac12\bigr)\;=\;e^{\,2i\theta(q)}
\end{equation}
and the point symmetry
\begin{equation}\label{eq:yinf-pointsym}
Y_{\inf}\bigl(\tfrac12-x\bigr)\;=\;1-Y_{\inf}(x):
\end{equation}
the map $z\mapsto1-z$, rotation by $\pi$ about the point $\tfrac12$,
carries the curve parametrized on $[0,1]$ onto the curve parametrized on
$[-\tfrac12,\tfrac12]$.
\end{lemma}

\begin{proof}
For \eqref{eq:psi-reflection}, multiply through by $\cos(2\pi x)$ and write
$\varphi=2\pi x$. On the left, $\Psi(x)\cos(2\pi x)=\cos(A-\varphi)$, since
$2\pi(x^{2}-x-\tfrac1{16})=A(x)-\varphi$. On the right,
$\cos\bigl(2\pi(x+\tfrac12)\bigr)=-\cos(2\pi x)$ and
$(x+\tfrac12)^{2}-(x+\tfrac12)=x^{2}-\tfrac14$, so
$\Psi(x+\tfrac12)\cos(2\pi x)=-\cos\bigl(A(x)-\tfrac\pi2\bigr)=-\sin A$;
moreover $\theta(x)=\varphi-A-\tfrac\pi2$, so
$e^{i\theta}=-ie^{\,i(\varphi-A)}$. The claim therefore reads
\[
\cos\varphi-\cos(A-\varphi)\,e^{-iA}
\;=\;(-\sin A)\,(-i)\,e^{\,i(\varphi-A)}
\;=\;i\sin A\,e^{\,i(\varphi-A)},
\]
and expanding the left side settles it: the real part is
$\cos\varphi-\cos(A-\varphi)\cos A=\sin A\,\sin(A-\varphi)$ (product
formula), the imaginary part is $\sin A\,\cos(A-\varphi)$, and
$\sin A\bigl(\sin(A-\varphi)+i\cos(A-\varphi)\bigr)
=i\sin A\,e^{-i(A-\varphi)}$, which is the right side.

For \eqref{eq:yinf-halflag}, two elementary phase relations do the work:
$A(q-\tfrac12)=A(q)-\varphi+\tfrac\pi2=-\theta(q)$ and, applying
\eqref{eq:psi-reflection} at $x=q-\tfrac12$,
$Y_{\inf}(q-\tfrac12)=\Psi(q)\,e^{\,i\theta(q-\frac12)}$ with
$\theta(q-\tfrac12)\equiv2\varphi-A\pmod{2\pi}$. Then
\[
Y_{\inf}(q)-Y_{\inf}\bigl(q-\tfrac12\bigr)
=1-\Psi(q)\bigl(e^{-iA}+e^{\,i(2\varphi-A)}\bigr)
=1-2\cos(A-\varphi)\,e^{-i(A-\varphi)}
=-e^{-2i(A-\varphi)}
=e^{\,2i\theta(q)},
\]
using $\Psi(q)\cdot2\cos\varphi=2\cos(A-\varphi)$ in the middle step and
$\theta=\varphi-A-\tfrac\pi2$ at the end.

For \eqref{eq:yinf-pointsym}, first $\Psi(1-x)=\Psi(x)$: the numerator of
\eqref{eq:psi-fn} is unchanged because $(1-x)^{2}-(1-x)=x^{2}-x$, and so is
the denominator because $\cos\bigl(2\pi(1-x)\bigr)=\cos(2\pi x)$. Then, by
the polar form \eqref{eq:psi-reflection} evaluated at $\tfrac12-x$
together with
$\theta(\tfrac12-x)=2\pi(\tfrac12-x)(\tfrac12+x)-\tfrac{3\pi}8
=\tfrac\pi8-2\pi x^{2}=-A(x)$,
\[
Y_{\inf}\bigl(\tfrac12-x\bigr)
=\Psi(1-x)\,e^{\,i\theta(\frac12-x)}
=\Psi(x)\,e^{-iA(x)}
=1-Y_{\inf}(x),
\]
the last step being the definition \eqref{eq:yinf}.
\end{proof}

\begin{theorem}\label{thm:yinf}
Fix $q\in(0,1)$ and $\sigma\in(0,1)$. Then
\[
\lim_{\substack{T\to\infty\\ \{T\}=q}} Y_{in1}(\sigma,T)\;=\;Y_{\inf}(q),
\qquad
\lim_{\substack{T\to\infty\\ \{T\}=q}} Y_{ang1}(\sigma,T)
\;=\;Y_{\inf}\bigl(q-\tfrac12\bigr),
\]
at rate $O(1/T)$, uniformly for $q$ and $\sigma$ in compact subsets of
$(0,1)$; the limit does not depend on $\sigma$. On the critical line the
proof is complete modulo Gabcke's bound \eqref{eq:gabcke}; off it, it
additionally uses the first-order Siegel asymptotics \eqref{eq:R-general},
with the same caveat as in Theorem~\ref{thm:d1-positive-offline}.
\end{theorem}

\begin{proof}
Throughout, $C_0$ and $\Psi$ are the same function, as
\eqref{eq:psi-is-c0} below records; we keep Gabcke's symbol $C_0$ for the
analytic input and $\Psi$ for the limit curve.

\emph{Step 1 (the yin point in asymptotic closed form).} By
\eqref{eq:yin1}, $Y_{in1}=R\,(m+1)^{\sigma}e^{\,i\omega}$. On the critical
line $R=e^{-i\vartheta}r$ with $r$ real, so
$Y_{in1}=\sqrt{m+1}\;r\,e^{\,iu}$, and inserting
\eqref{eq:r-two-cases}--\eqref{eq:gabcke},
\begin{equation}\label{eq:yin-asymp}
Y_{in1}
=\Bigl(\tfrac{m+1}{a}\Bigr)^{\!1/2}(-1)^{N-1}\bigl(C_0(\hat p)+E\bigr)
e^{\,iu}
\;+\;\bigl[N=m{+}1\bigr]\,\bigl(1+e^{\,2iu}\bigr),
\end{equation}
where the extra summand of \eqref{eq:r-two-cases} became the closed pair
$2\cos u\,e^{\,iu}=1+e^{\,2iu}$. Off the line the first-order Siegel
asymptotics \eqref{eq:R-general} give the same leading form: the frame
factor $(m+1)^{\sigma}e^{\,i\omega}$ carries the two bracket pieces of
\eqref{eq:R-general} to
$\bigl(\tfrac{m+1}{a}\bigr)^{\sigma}e^{\,i(\omega-\tilde\vartheta)}$ and
$|\chi|\,a^{\sigma-1}(m+1)^{\sigma}e^{\,i(\omega+\psi+\tilde\vartheta)}$,
and since $|\chi|=a^{1-2\sigma}\bigl(1+O(1/t)\bigr)$ and
$\psi=-2\vartheta+O_\sigma(1/t)$, both tend to $e^{\,iu}$; likewise the
Iverson pair $(m+1)^{-s}+\chi\,(m+1)^{s-1}$ maps to
$1+|\chi|(m+1)^{2\sigma-1}e^{\,i(2\omega+\psi)}\to1+e^{\,2iu}$. In every
case
\[
Y_{in1}=(-1)^{N-1}C_0(\hat p)\,e^{\,iu}\bigl(1+O(1/T)\bigr)
+\bigl[N=m{+}1\bigr]\bigl(1+e^{\,2iu}\bigr)\bigl(1+O(1/T)\bigr).
\]

\emph{Step 2 (the phase, now with its full angle).} The expansion behind
\eqref{eq:cosu}, kept as an angle rather than used only through its cosine,
reads
\[
u\;\equiv\;\pi(m+1)+\pi\bigl(\tfrac18-2\delta^{2}\bigr)+O(1/T)
\pmod{2\pi},
\qquad \delta:=a-(m+1)\longrightarrow q-\tfrac12,
\]
(substitute $t=2\pi a^{2}$ and the asymptotic series of $\vartheta$ into
$u=t\log\tfrac{m+1}{a}+\tfrac t2+\tfrac\pi8+O(1/t)$ and reduce
$\pi a^{2}$ modulo $2\pi$), and
$\pi\bigl(\tfrac18-2(q-\tfrac12)^{2}\bigr)=\theta(q)$, so
$u\equiv\pi(m+1)+\theta(q)+O(1/T)$. Meanwhile
$\hat p\to q+\tfrac12$ (with $N=m$) for $q<\tfrac12$ and
$\hat p\to q-\tfrac12$ (with $N=m+1$) for $q>\tfrac12$, by
\eqref{eq:cut-cases}.

\emph{Step 3 (case $q<\tfrac12$).} Here $N=m$, so
$(-1)^{N-1}e^{\,i\pi(m+1)}=(-1)^{2m}=1$ and the Iverson term is absent:
\[
Y_{in1}\;\longrightarrow\;\Psi\bigl(q+\tfrac12\bigr)e^{\,i\theta(q)}
\;=\;Y_{\inf}(q)
\]
by the reflection identity \eqref{eq:psi-reflection}.

\emph{Step 4 (case $q>\tfrac12$).} Here $N=m+1$, the sign product is
$-1$, and the Iverson pair is present:
\[
Y_{in1}\;\longrightarrow\;
1-\Psi\bigl(q-\tfrac12\bigr)e^{\,i\theta(q)}+e^{\,2i\theta(q)} .
\]
Since $\theta(q)=-A(q-\tfrac12)$ exactly (Lemma~\ref{lem:psi-reflection}),
the first two terms are $Y_{\inf}(q-\tfrac12)$, and the half-lag relation
\eqref{eq:yinf-halflag} finishes:
$Y_{in1}\to Y_{\inf}(q-\tfrac12)+e^{\,2i\theta(q)}=Y_{\inf}(q)$.

\emph{Step 5 (the yang point).} By \eqref{eq:yang1},
$Y_{ang1}-Y_{in1}=-\chi\,(m+1)^{2\sigma-1}e^{\,2i\omega}$, whose modulus
$|\chi|(m+1)^{2\sigma-1}\to1$ and phase $2\omega+\psi=2u+O_\sigma(1/t)
\equiv2\theta(q)+O(1/T)$, so the subtracted vector tends to
$e^{\,2i\theta(q)}$ and
$Y_{ang1}\to Y_{\inf}(q)-e^{\,2i\theta(q)}=Y_{\inf}(q-\tfrac12)$, again by
\eqref{eq:yinf-halflag}. One caution on reading this as a lag: by the point
symmetry \eqref{eq:yinf-pointsym},
$Y_{\inf}(q-\tfrac12)=1-Y_{\inf}(1-q)$, so within each unit interval the
yang end runs the yin limit curve time-reversed and rotated by $\pi$ about
the point $\tfrac12$; equivalently,
$Y_{ang1}(\sigma,T)\approx1-Y_{in1}(\sigma,\,2m{+}1{-}T)$ for every
fractional part. The naive lag reading
$Y_{ang1}(\sigma,T)\approx Y_{in1}(\sigma,T-\tfrac12)$ is correct only on
$\{T\}>\tfrac12$, where $T$ and $T-\tfrac12$ share the frame link $m+1$;
for $\{T\}<\tfrac12$ the point $Y_{in1}(\sigma,T-\tfrac12)$ lives in the
previous frame, and since $Y_{\inf}$ is not $1$-periodic, the two sides
differ by order one there.

Every ingredient contributes $O(1/T)$: the amplitude ratio
$(m+1)/a=1+O(1/T)$, Gabcke's $E=O(1/a)$, the phase remainder in Step~2,
and off the line the $O_\sigma(1/t)$ terms of $\psi$ and $|\chi|$ together
with the error of \eqref{eq:R-general} (relative $O(1/T)$,
\S\ref{sec:d1-approx-general}). At the seam $q=\tfrac12$ the two case
formulas agree: $\Psi(1)=\Psi(0)=\cos\tfrac\pi8$ and
$1+e^{\,i\pi/4}=2\cos\tfrac\pi8\,e^{\,i\pi/8}$, so both give
$\cos\tfrac\pi8\,e^{\,i\pi/8}$, which is $Y_{\inf}(\tfrac12)$.
\end{proof}

\begin{remark}[numerical corroboration]
The identities \eqref{eq:psi-reflection}--\eqref{eq:yinf-pointsym} check to
thirty digits, and the convergence at the predicted rate is visible at
moderate height: $|Y_{in1}-Y_{\inf}(q)|$ at $q=0.2$ is $1.3\times10^{-3}$
at $T=40.2$ and $3.3\times10^{-4}$ at $T=160.2$ for $\sigma=\tfrac12$, the
same factor-of-four drop per quadrupling of $T$ occurring at $\sigma=0.3$
and $\sigma=0.9$, with the yang errors matching the yin errors half a unit
away in the argument. The Step-5 caution is equally visible: at $T=80.2$
the flipped relation $|Y_{ang1}-(1-Y_{in1}(2m{+}1{-}T))|=4\times10^{-3}$,
while the naive lag $|Y_{ang1}-Y_{in1}(T-\tfrac12)|=1.9$, an order-one
failure; at $T=80.8$, where the frames agree, the naive lag is fine
($2\times10^{-4}$). (\texttt{probe\_yin\_theorem.py}.)
\end{remark}

\begin{corollary}[the limit profile is the axis crossing of the limit
chord]\label{cor:yinf-crossing}
For $q\in(0,1)$, $q\neq\tfrac14,\tfrac34$, the chord from
$Y_{\inf}(q)$ to $Y_{\inf}(q-\tfrac12)$ crosses the real axis at the
abscissa $d(q)$, the limit profile of Theorem~\ref{thm:d1-limit}.
\end{corollary}

\begin{proof}
At every finite $T$ the segment from $Y_{in1}$ to $Y_{ang1}$ is the reverse
bisector link in the frame, and it crosses the real axis at the bisector
point, at abscissa $\hat d_1$ (Figure~\ref{fig:yinyang};
\S\ref{sec:derive-r1ps}). As $T\to\infty$ with $\{T\}=q$ the two endpoints
converge to $Y_{\inf}(q)$ and $Y_{\inf}(q-\tfrac12)$
(Theorem~\ref{thm:yinf}) while the crossing converges to $d(q)$
(Theorem~\ref{thm:d1-limit} on the line, its universality remark off it);
away from $q=\tfrac14,\tfrac34$ the limit chord meets the axis
transversally, so the crossing of the limit is the limit of the crossings.
At $q=\tfrac14,\tfrac34$ the chord degenerates onto the axis itself (the
parallel instants: it runs from $\tfrac12$ to $-\tfrac12$ and from
$\tfrac32$ to $\tfrac12$ respectively), which is the limiting shadow of the
off-line poles of \S\ref{sec:pole-locations}. The identity was also checked
directly to thirty digits (\texttt{probe\_yin\_theorem.py}): the crossing
of the limit chord and the closed form $d(q)$ of
Theorem~\ref{thm:d1-limit} agree.
\end{proof}

So the two halves of the paper close into one loop: the weights
$d_1,d_2$, their limit profile $d(x)$, and the yin-yang limit curve are a
single asymptotic object seen from three angles, and the same amplitude
runs through all of it.

Now the connection promised in the subsection title, which Step~1 of the
proof has already put to work. The function
\eqref{eq:psi-fn} is, symbol for symbol, the universal amplitude
\begin{equation}\label{eq:psi-is-c0}
\Psi(x)\;=\;C_0(x)
\;=\;\frac{\cos\!\bigl(2\pi(x^{2}-x-\tfrac{1}{16})\bigr)}{\cos(2\pi x)}
\end{equation}
of the Riemann--Siegel first term used in
\S\ref{sec:d1-approx-critical}--\ref{sec:d1-approx-general}
(equations \eqref{eq:d1-rs} and \eqref{eq:R-general}). So the same universal
function plays both roles: evaluated at $\hat p=\operatorname{frac}\sqrt{t/2\pi}$
it is the asymptotic \emph{magnitude} of the remainder, and evaluated at
$q=\operatorname{frac}(T)$ it is the asymptotic \emph{shape} of the yin-yang
orbit; by \eqref{eq:yinf} the limiting yin curve keeps distance
$|\Psi(q)| = |C_0(q)|$ from the point $1$. Even the two parameters are the
same clock read half a period apart: $\sqrt{t/2\pi}\approx T+\tfrac12$ (the
Taylor table of \S\ref{sec:IT-origin}), so
$\hat p\approx q\pm\tfrac12 \pmod 1$, precisely the half-unit shift between
the yang and yin arguments in Theorem~\ref{thm:yinf}.

Three properties of the limit curve deserve note. First, it is closed:
$Y_{\inf}(0)=Y_{\inf}(1)=1-\cos(\tfrac{\pi}{8})e^{\,i\pi/8}$. Second, the
apparent singularities of $\Psi$ at $q=\tfrac14,\tfrac34$ (the pole heights of
\S\ref{sec:pole-locations}) are \emph{removable} on the curve:
$\Psi\to\tfrac12$ at both points, and the curve passes smoothly through the
real values $Y_{\inf}(\tfrac14)=\tfrac12$ and $Y_{\inf}(\tfrac34)=\tfrac32$.
Third, the reflection identity \eqref{eq:psi-reflection} is a polar form of
the curve, $Y_{\inf}(q)=\Psi(q+\tfrac12)\,e^{\,i\theta(q)}$: the same
function that measures the distance $|\Psi(q)|$ from the point $1$
measures, half a period along, the distance $|\Psi(q+\tfrac12)|$ from the
origin.

\begin{figure}[htbp]
\centering
\includegraphics[width=0.92\textwidth]{fig_yinyang_infinity}
\caption{Convergence of the yin curves to the limit curve
$Y_{\inf}(q)=1-\Psi(q)e^{-2\pi i(q^2-1/16)}$ (dashed black), with the unit
bisector link in black; the two panels have identical window sizes. As
everywhere in the bisector frame, the picture is scaled by the frame map
\eqref{eq:frame-map} so that the forward bisector link always has length
$1$, running from $0$ to $1$ on the real axis. Top:
$\sigma=0.2$, the yin paths (green) of every
unit interval of the index over $0<T<8$, $m=0,\dots,7$; the $m=0$ path, marked with spaced
brown dots, is $\zeta(\sigma+it)$ itself, since for $0<T<1$ both partial
sums are empty and $Y_{in1}=\zeta$. Bottom:
$\sigma=0.9$, the yin paths of the periods starting at
$T=0,1,2,4,8,16,32$ (light to dark green); by $T\approx32$ the path is
visually indistinguishable from $Y_{\inf}$ even this far off the critical
line. The $m=0$ path, again $\zeta(\sigma+it)$ and again marked with brown
dots, leaves the frame at the bottom: $\sigma=0.9$ sits close to the pole
of $\zeta$ at $s=1$, so this first orbit dives to about $-2.7i$ before the
later orbits settle onto the teardrop; the window is clipped to the
convergence region. The yang curves (omitted
for clarity) converge to the point reflection of $Y_{\inf}$ through
$\tfrac12$, half a unit behind in the argument
(Theorem~\ref{thm:yinf} and \eqref{eq:yinf-pointsym}). Generated by
\texttt{fig\_yinyang\_infinity.py}.}
\label{fig:yinyang-infinity}
\end{figure}

\clearpage

% ======================================================================
\section{Further observations}\label{sec:consequences}
% ======================================================================

\subsection{The area swept out by the unit bisector link at infinity}\label{sec:yinyang-area}

In the language of Kakeya sets: what area does a unit-length needle sweep when
one end rides $Y_{\inf}(t)$ and the other rides $Y_{\inf}(t-\tfrac12)$ (the
limiting positions of the two ends of the reverse bisector link) as $t$
runs from $0$ to $1$? The ``needle'' here is nothing abstract: it is the
reverse bisector link itself, sweeping once around the stationary forward
bisector link over one unit of the index, exactly the motion depicted in
Figure~\ref{fig:bisector-sweep}. The two ends ride congruent teardrops: by
the point symmetry \eqref{eq:yinf-pointsym} the yang copy is the rotation
of $Y_{\inf}$ by $\pi$ about the point $\tfrac12$, so it encloses the same
area, and the swept area is twice the area enclosed by $Y_{\inf}$, which we
now compute.

Let $z(t)=Y_{\inf}(t)$ on $t\in[0,1]$, a closed curve passing smoothly through
its two removable points. Writing $z(t)=x(t)+iy(t)$, the signed enclosed area
is
\begin{equation}\label{eq:area-def}
\mathcal A_{\mathrm{signed}}
=\frac12\int_{0}^{1}\bigl(x\,y'-y\,x'\bigr)\,dt
=\frac12\int_{0}^{1}\Im\!\bigl(\overline{z}\,z'\bigr)\,dt .
\end{equation}
Set $A(t)=2\pi\bigl(t^{2}-\tfrac{1}{16}\bigr)$, so that
$z-1=-\Psi(t)\,e^{-iA(t)}$. A short computation gives the key simplification
\begin{equation}\label{eq:area-key}
\Im\!\Bigl(\overline{(z-1)}\,(z-1)'\Bigr)
\;=\;
-A'(t)\,\Psi(t)^{2},
\qquad A'(t)=4\pi t
\end{equation}
(the $\Psi'\,\Psi$ cross term is real and drops out). Translation by $1$ does
not change enclosed area, so
\begin{equation}\label{eq:area-signed}
\mathcal A_{\mathrm{signed}}
=-\frac12\int_{0}^{1}A'(t)\,\Psi(t)^{2}\,dt
=-2\pi\int_{0}^{1}t\,\Psi(t)^{2}\,dt ,
\end{equation}
and the geometric (positive) area is
\begin{equation}\label{eq:area-value}
\boxed{\;
\mathrm{Area}
=2\pi\int_{0}^{1}t\,
\frac{\cos^{2}\!\bigl(2\pi(t^{2}-t-\tfrac{1}{16})\bigr)}{\cos^{2}(2\pi t)}
\,dt
\;\approx\;
1.0341672002955850005 .\;}
\end{equation}
(The integral is split at the removable points $t=\tfrac14,\tfrac34$ for the
numerical evaluation; with the given parameter direction the signed area comes
out negative, and the magnitude is the enclosed area above.) The needle of the
opening question therefore sweeps about $2.0683$.

\subsection{Connection between the yin and yang functions and the link numbers}\label{sec:yinyang-links}

In our numerical experiments with the tool we noticed that the bisector link,
viewed in its local frame, is always parallel to one particular link of the
``conveyor belt'', the stretch of links between the last two spirals
(spirals $0$ and $1$ in the numbering of \eqref{eq:LN}). That particular link
follows the bisector link as the index varies between integers, and its link
number is
\begin{equation}\label{eq:follower}
\lfloor T\rfloor\,\lfloor T+2\rfloor
\;=\;
(m+1)^{2}-1 ,
\end{equation}
so the follower's summand is number $(m+1)^{2}$, the square of the bisector
link's summand $m+1$, whose phase $-t\log\bigl((m+1)^{2}\bigr)$ advances at
exactly twice the bisector link's rate. The first few link numbers in the
sequence are $3, 8, 15, 24, 35, 48, 63, 80, 99, 120, 143, \dots$: it is
sequence A005563 in the OEIS \cite{OEIS}, though no connection of that sequence
to the zeta function is listed there. The follower transitions from one of
these links to the next exactly when the index $T$ is an integer (the same
instants at which the bisector point hands off), and the two links are
momentarily coincident near $T+\tfrac12$.

\begin{figure}[H]
\centering
\includegraphics[width=\textwidth]{fig_conveyor_follower}
\caption{The follower link and the bisector chord over one unit of the index,
$\sigma=\tfrac12$ and $T=6.0,6.1,\dots,6.9$ (so $m=6$ and the follower is
link $48$, the $49$th summand). Each panel superposes two frames. In the
world frame: the conveyor belt (thin black), the reverse chain's joints
$\zeta-\Sigma_2(n)$ for $n$ running from $L_N(T,1)$ to $L_N(T,0)$ of
\eqref{eq:LN}, which is the stretch between spirals $1$ and $0$; inside it
the follower, link $48$, is drawn thick in neon red. In the bisector frame
\eqref{eq:frame-map}: the yin curve $Y_{in1}$ (thin green) and the yang curve
$Y_{ang1}$ (thin red) over $6<T<7$, the same in all ten panels; the chord
$[Y_{in1},Y_{ang1}]$ at this panel's $T$ (neon red, endpoints dotted); the
forward bisector link $[0,1]$ (light blue); and the bisector point (star),
where the chord's line crosses that unit link. The chord turns through a half
revolution across the ten panels and the follower link turns with it, always
parallel by \eqref{eq:follower-identity} and always $\lceil T\rceil$ times
shorter. Near $T=6.5$ the two very nearly coincide. Generated by
\texttt{fig\_conveyor\_follower.py}.}
\label{fig:conveyor-follower}
\end{figure}

The observation is an exact identity, and the doubled\footnote{Every summand
of the chain is a vector whose phase is $-t\log n$, so as $t$ grows that
vector rotates at rate $\log n$ radians per unit $t$. Writing
$\omega=t\log\lceil T\rceil$ as elsewhere in the paper, the bisector link
carries summand $\lceil T\rceil$ and so sits at phase $\omega$, while the
follower carries summand $\lceil T\rceil^{2}$ and sits at phase
$t\log\bigl(\lceil T\rceil^{2}\bigr)=2\omega$, double because squaring the
summand doubles its logarithm.} phase rate explains it
completely. The chord of the yin and yang curves, that is the reverse bisector link
seen in the frame \eqref{eq:frame-map}, is
$Y_{ang1}-Y_{in1}=-\chi\lceil T\rceil^{\,2s-1}$ by \eqref{eq:yang1}, with
$s=\sigma+iI(T)$ as usual, while
link $(m+1)^{2}-1$ of the reverse chain, in the world frame, is the summand
vector $-\chi\bigl(\lceil T\rceil^{2}\bigr)^{s-1}=-\chi\lceil
T\rceil^{\,2s-2}$. The two differ by the factor $\lceil T\rceil^{-1}$, which
is real and positive, so
\begin{equation}\label{eq:follower-identity}
\text{link }(m+1)^2-1 \;=\; \frac{Y_{ang1}-Y_{in1}}{\lceil T\rceil}
\qquad\text{for every }\sigma\text{ and every }T:
\end{equation}
the follower is parallel to the chord, points the same way, and is shorter by
exactly the factor $\lceil T\rceil$. No other link of the chain has this
property, since the phases of neighboring links differ by
$t\log\frac{n+1}{n}$, which is not a multiple of $\pi$ in general.
Figure~\ref{fig:conveyor-follower} shows the pair over one unit of the index.
Numerically the angle between the two vectors is zero to machine precision and
the length ratio is $\lceil T\rceil^{-1}$ to twelve digits at every
$(\sigma,T)$ tested, on and off the critical line. What is \emph{not} exact is
the position: the follower lies on the chord only momentarily. Measuring the
perpendicular distance from the follower's midpoint to the chord line at
$\sigma=\tfrac12$, $m=6$, it falls from $0.512$ at $T=6.0$ to $0.008$ at
$T=6.5$ and rises again to $0.512$ at $T=7.0$, and at the minimum the
follower's midpoint sits at fraction $0.503$ of the chord. This is the
``momentarily coincident near $T+\tfrac12$'' of the first paragraph.

\subsection{\texorpdfstring{Length ratios: $\Sigma_1$ to $\Sigma_2$, and $R_{1ps}$ to $R_{2ps}$}{Length ratios}}\label{sec:length-ratios}

The four terms of \eqref{eq:zeta-ps} come in two natural pairs: the two
partial sums, and the two fractional links. Within each pair the two lengths
are equal on the critical line, and the two remarks below record what their
ratio does off it. The pairs behave quite differently, and in both cases the
ratio obeys the same reflection law about $\sigma=\tfrac12$.

\begin{remark}[regarding the ratio of the lengths of $\Sigma_1$ and $\Sigma_2$]\label{rem:leg-ratio}
On the critical line $\sigma=\tfrac12$ we have $|\chi|=1$, so the $\Sigma_2$ links
have the same lengths $(n+1)^{-1/2}$ as the $\Sigma_1$ links; each link is then a
rotation composed with a uniform link-length shrink $\sqrt{n/(n+1)}$, and
$\arg\chi=-2\vartheta(t)$, with $\vartheta$ the Riemann--Siegel theta function.
The two chains are thus congruent link for link there, so their end-to-end
displacements agree. Away from the line they do not, and the discrepancy has a
closed form. Write $\ell_{\Sigma_1}=|\Sigma_1|$ and $\ell_{\Sigma_2}=|\Sigma_2|$
for the lengths of the net displacements of $P_{\Sigma_1}$ and $P_{\Sigma_2}$,
and put $\Sigma_m(z)=\sum_{n\le m}n^{-z}$. Since
$\sum_{n\le m}n^{\,s-1}=\overline{\Sigma_m\bigl((1-\sigma)+it\bigr)}$, the
second leg is the first one read at the reflected abscissa $1-\sigma$, scaled
by $|\chi|$:
\begin{equation}\label{eq:leg-ratio}
\frac{\ell_{\Sigma_1}}{\ell_{\Sigma_2}}
=
\frac{\bigl|\Sigma_m(\sigma+it)\bigr|}
     {\bigl|\chi(\sigma+it)\bigr|\;
      \bigl|\Sigma_m\bigl((1-\sigma)+it\bigr)\bigr|}.
\end{equation}
Since $\chi(s)\chi(1-s)=1$ and $|\chi(\bar s)|=|\chi(s)|$, the formula gives,
at fixed $T$ (hence fixed $t=I(T)$ and $m=\lfloor T\rfloor$),
\begin{equation}\label{eq:leg-ratio-symmetry}
\frac{\ell_{\Sigma_1}(\sigma)}{\ell_{\Sigma_2}(\sigma)}
\cdot
\frac{\ell_{\Sigma_1}(1-\sigma)}{\ell_{\Sigma_2}(1-\sigma)}
=
1,
\end{equation}
which we have verified numerically to twelve digits at several $T$. So the
ratio is reciprocal-symmetric about $\sigma=\tfrac12$: whatever it does on the
right half of the strip, it does the reciprocal of on the left, with the value
pinned to $1$ at the center by the link-for-link congruence just described.
\end{remark}

\begin{remark}[regarding the ratio of the lengths of $R_{1ps}$ and $R_{2ps}$]\label{rem:d-ratio}
The same question for the two fractional links has a shorter answer, because in
the quotient of \eqref{eq:d1} and \eqref{eq:d2} both $|R|$ and the common
denominator $\sin(2\omega+\psi)$ cancel:
\begin{equation}\label{eq:d-ratio}
\frac{|R_{1ps}|}{|R_{2ps}|}
\;=\;
\frac{d_1}{d_2}
\;=\;
\frac{\sin(\omega+\psi-\arg R)}{\sin(\omega+\arg R)}.
\end{equation}
This ratio is pure angle data, carrying no length scale at all: it is the law
of sines in the triangle of Figure~\ref{fig:remainder-zoom}, recording how the
direction of $R$ divides the cone spanned by the two link directions
$e^{-i\omega}$ and $e^{\,i(\omega+\psi)}$. Four consequences are worth
recording. First, on the critical line $\arg R=\tfrac{\psi}{2}$, the two sines
coincide and the ratio is $1$, which is Corollary~\ref{cor:equal} again.
Second, the reflection law \eqref{eq:leg-ratio-symmetry} holds verbatim for
$d_1/d_2$: at fixed $T$ we have verified
$(d_1/d_2)(\sigma)\cdot(d_1/d_2)(1-\sigma)=1$ numerically to $10^{-17}$ and
better. Third, because the scale cancels, this ratio is far tamer than
\eqref{eq:leg-ratio}: over $1<T<20$ and the whole strip its extremes are
$0.387$ and $2.583$ (a reciprocal pair, as the second point requires), whereas
$\ell_{\Sigma_1}/\ell_{\Sigma_2}$ ranges over several orders of magnitude
there. Fourth, the parallel-link poles of \S\ref{sec:pole-locations} cancel out
of it: where $2\omega+\psi=k\pi$, \eqref{eq:d-ratio} gives
$d_1/d_2=(-1)^{k+1}$, exactly $\pm1$, even though $d_1$ and $d_2$ separately
diverge (at $\sigma=\tfrac1{10}$, $T=4.256594$ we compute
$d_1=-d_2\approx-4.07\times10^{36}$). A negative value is off-line news of a
different kind: it says $R$ has left the cone, which happens in a narrow window
around each pole. The ratio does reach $0$ and $\infty$, but only on the
isolated loci where $R$ is exactly parallel to one of the two link directions;
for instance $d_1=0$ at $\sigma=\tfrac1{10}$, $T=4.2607679$, whose mirror at
$\sigma=\tfrac9{10}$ has $d_2=0$. Off the critical line the locus $d_1=d_2$ is
the family of ovals plotted in Figure~\ref{fig:leg-imbalance}.
\end{remark}

\subsection{\texorpdfstring{$\vartheta_1$, $\vartheta_2$}{theta1, theta2} and
the zero-counting function}\label{sec:counting}

The well known Riemann--von Mangoldt formula for the number of zeros of $\zeta$ with
imaginary part in $(0,t]$, in its exact form, reads
\begin{equation}\label{eq:N-theta}
N(t) \;=\; \frac{\vartheta(t)}{\pi} \;+\; 1 \;+\; S(t),
\qquad
S(t) \;=\; \frac{1}{\pi}\arg\zeta\!\left(\tfrac12+it\right),
\end{equation}
with $\vartheta$ the Riemann--Siegel theta function of \S\ref{sec:d1-approx-critical}.
This is an identity (\cite[Theorem~9.3]{Titchmarsh1986}), obtained by applying the argument
principle to $\xi$ around a rectangle: $\vartheta(t)/\pi+1$ is the contribution of
the $\Gamma$ and $\pi^{-s/2}$ factors, and $S(t)$ is the contribution of $\zeta$
itself along the critical line. Stated in terms of $\vartheta$ the formula is
exact: \eqref{eq:N-theta} carries no error term at all; approximation enters
only when $\vartheta$ is replaced by its expansion, and inexactness only when $S$
is dropped. Figure~\ref{fig:staircase-67} shows the two pieces over one unit of
the index.

\begin{figure}[htbp]
\centering
\includegraphics[width=0.92\textwidth]{fig_staircase_67}
\caption{\emph{Left:} the staircase $N(I(T))$ (blue) over $6\leq T\leq7$, which
is $264.9\leq t\leq352.9$ and carries the $55$ ordinates $\gamma_{117}$ through
$\gamma_{171}$, against the smooth part $\vartheta(t)/\pi+1$ of
\eqref{eq:N-theta} (black). The vertical gap between them is $S$, which here
stays within $[-1.002,+1.001]$. \emph{Right:} the boxed window
$6.4\leq T\leq6.6$ magnified, with the eleven ordinates it holds,
$\gamma_{138}$ through $\gamma_{148}$, marked in red at the top of each riser
and placed at $T=I^{-1}(\gamma)$. The risers are unevenly spaced while the
smooth part climbs steadily. Below each panel is the gap itself, $S$ (green),
which at this magnification shows plainly as a sawtooth: a unit jump at every
ordinate and a steady slide between them. Generated by \texttt{fig\_staircase\_67.py}.}
\label{fig:staircase-67}
\end{figure}

\paragraph{All of the discontinuity is in $S$.} $N$ is a step function and
$\vartheta(t)/\pi+1$ is analytic, so $S$ carries every jump: it climbs by one
(by the multiplicity, in principle) at each ordinate (the imaginary part
$\gamma$ of a zero) and falls smoothly in between, at the rate
\begin{equation}\label{eq:S-slope}
S'(t) \;=\; -\frac{\vartheta'(t)}{\pi} \;=\; -\frac{1}{2\pi}\log\frac{t}{2\pi}
\;+\;O(t^{-2}),
\end{equation}
which is the mean zero density. Across an ordinate $S$ jumps up by one;
between ordinates it slides back down along that line. The sawtooth is why $S$ has mean zero without being small: each unit jump is
repaid by a decline whose steepness is the density itself. At an ordinate the
definition of $S$ fails, since $\zeta$ vanishes there and its argument is
undefined; one sets $S(\gamma)=\tfrac12[S(\gamma^{-})+S(\gamma^{+})]$, and
\eqref{eq:N-theta} then holds for every $t>0$.

\paragraph{Three counting curves against the staircase.} Expanding $\vartheta$ by
Stirling and dropping $S$ turns \eqref{eq:N-theta} into the Riemann--von
Mangoldt formula in its approximate form,
\begin{equation}\label{eq:N-rvm}
N_{\mathrm{rvm}}(t)
=\frac{t}{2\pi}\log\frac{t}{2\pi}-\frac{t}{2\pi}+\frac78
+\frac{1}{48\pi t}+\frac{7}{5760\pi t^{3}}
\end{equation}
(the $\tfrac78$ collects the $-\tfrac\pi8$ of the expansion, divided by $\pi$,
with the standalone $1$ of \eqref{eq:N-theta}). Nothing of consequence is
discarded in truncating there: at $t=1000$ even the first correction retained,
$1/(48\pi t)$, is only $6.6\times10^{-6}$, while $S(1000)=+0.384$. The whole of
the error is $S$. Set also
\begin{equation}\label{eq:N-ps}
N_{\mathrm{ps}}(T)
=\frac{1}{\pi}\vartheta\bigl(I(T)\bigr)
+\frac{1}{\pi}\vartheta_1\bigl(\tfrac12,T\bigr)+\frac32 ,
\end{equation}
which replaces $S$ by the angle $\vartheta_1$ that leg~1 makes with the real
axis (\S\ref{sec:legs}), the leg running from the origin to the bisector point
$B_1=\Sigma_1+R_{1ps}$. These two behave quite differently against the true
count $N(I(T))$, and a third curve $N_{\ast}$, built below in
\eqref{eq:N-star} from a different splitting of $\zeta$, repairs the one defect
of $N_{\mathrm{ps}}$.

$N_{\mathrm{rvm}}$ is analytic and increasing, so it can only approximate the
staircase in the mean, never landing on an integer at an ordinate, and its
error is exactly the sawtooth $S$.
$N_{\mathrm{ps}}$ instead meets an integer \emph{exactly} at every ordinate,
which we verify to $2.8\times10^{-14}$. The reason is
Corollary~\ref{cor:bisector-proj}: on the critical line
\begin{equation}\label{eq:B1-rotated}
B_1\,e^{i\vartheta(t)}=\tfrac12 Z(t)+ih,\qquad h\in\mathbb R.
\end{equation}
Multiplying by $e^{i\vartheta}$ rotates the whole configuration into the frame in
which $\zeta$ is real, the frame that defines $Z=e^{i\vartheta}\zeta$; it is the
rotation $e^{-i\psi/2}$ of Corollary~\ref{cor:bisector-proj} written with
$\vartheta=-\tfrac{\psi}{2}$, since $\chi=e^{-2i\vartheta}$ on the line. We call it the
\emph{$\vartheta$-frame}. There the base $O\to\zeta$ lies along the real axis, and
\eqref{eq:B1-rotated} reads a split point off as a pair of coordinates: half the
base along it, and the offset $h$ across it. So at a zero $Z$ vanishes,
$B_1e^{i\vartheta}$ is purely imaginary,
$\vartheta_1=-\vartheta\pm\tfrac\pi2$, and \eqref{eq:N-ps} is an integer. The
argument runs both ways: since $\Re\bigl(e^{i\vartheta}B_1\bigr)=\tfrac12 Z$ at
every $t$, and not merely at the zeros, \eqref{eq:N-ps} is an integer only where
$Z$ vanishes, so the integer points of $N_{\mathrm{ps}}$ are the ordinates and
nothing else. The
counting is therefore done by the chain's own geometry rather than by
$\arg\zeta$: leg~1 stands perpendicular to the direction of $\zeta$ precisely
when $\zeta$ vanishes. $N_{\mathrm{rvm}}$ carries no information about
individual zeros, while $N_{\mathrm{ps}}$ locates every one of them.

\paragraph{Where $N_{\mathrm{ps}}$ misses, and why.} Two things keep
\eqref{eq:N-ps} from being a counting function. First, $\vartheta_1$ is an
angle, so \eqref{eq:N-ps} is determined only up to an even integer, and each
passage of $\vartheta_1$ through $\pm\pi$ breaks the curve; there are nine such
wraps over $3<T<5$, and $\tfrac32$ is the constant that puts the curve on the
staircase. Second, and this is the real defect, the curve sometimes meets the
wrong integer. Differentiating \eqref{eq:B1-rotated} at an ordinate, where $Z$
vanishes, gives
\begin{equation}\label{eq:crossing-rate}
\frac{d}{dt}\arg\bigl(B_1e^{i\vartheta}\bigr)\bigg|_{Z=0}=-\frac{Z'(t)}{2h},
\end{equation}
so the crossing carries the angle forward or backward with the sign of $-hZ'$.
Consecutive simple zeros of $Z$ alternate in the sign of $Z'$, so the crossings
all turn the same way exactly when $h$ alternates in sign too. Where $h$ nearly
vanishes it keeps its sign across two ordinates instead; the angle turns back on
itself, the count fails to advance, and the second of the pair reads one too
high. Over the first $400$ ordinates $|h|$ has median $1.081$ at the $350$ that
$N_{\mathrm{ps}}$ places correctly and $0.319$ at the $50$ it misses, and those
$50$ are exactly the crossings that run backwards, which we call
\emph{retrograde}. They are also exactly the
wraps: the nine wraps of $\vartheta_1$ over $3<T<5$ counted above sit one apiece
on the nine retrograde ordinates of that range. Figure~\ref{fig:nps-staircase}
shows the cost over a longer stretch: five retrograde ordinates in
$5.9<T<6.7$, and after each one $N_{\mathrm{ps}}$ runs a unit above the
staircase until the following ordinate restores it. The two ways of describing the
failure, a broken branch and a backward crossing, are one event. Over $3<T<5$, which holds
$55$ zeros, both curves stay within one of the staircase all the same: the
largest deviation is $0.946$ for $N_{\mathrm{rvm}}$ and $0.999$ for
$N_{\mathrm{ps}}$.

\begin{figure}[htbp]
\centering
\includegraphics[width=\textwidth]{fig_nps_staircase}
\caption{Staircase, in \textcolor[HTML]{1F77B4}{blue}, is the exact
Riemann--von Mangoldt formula
\eqref{eq:N-theta}, and the approximate Riemann--von Mangoldt formula
\eqref{eq:N-rvm} is the black curve, that formula with $S$ dropped.
\emph{Top:} $5.9\leq T\leq6.7$, which is $256.8\leq t\leq325.2$ and holds the
$42$ ordinates $\gamma_{113}$ through $\gamma_{154}$, five of them retrograde,
with $N_{\mathrm{ps}}$ of \eqref{eq:N-ps} in
\textcolor[HTML]{FF7F0E}{orange}, broken at its wraps.
\emph{Middle:} the boxed window $6.125\leq T\leq6.275$ magnified, holding $8$
ordinates of which $\gamma_{125}$ and $\gamma_{128}$ are retrograde. The
velocity split $N^{\ast}$ of \eqref{eq:N-star} is drawn in
\textcolor[HTML]{2CA02C}{green}.
\emph{Bottom:} the fold angle of the legs $\vartheta_2$ of \S\ref{sec:legs} on the same
$T$ range, reduced to $[0,2\pi)$. Every passage through $\pi$ (dashed) is at an
ordinate; the two filled dots are the two retrograde ordinates of the panel
above, where $\vartheta_2$ rises through $\pi$, turns around, and comes back
down through it.
Generated by \texttt{fig\_nps\_staircase.py}.}
\label{fig:nps-staircase}
\end{figure}

Reading $\vartheta_2$ in place of $\vartheta_1$ cannot help. On the critical
line $B_2=\chi\overline{B_1}$ and $\arg\chi=-2\vartheta$, so
\begin{equation}\label{eq:th2-th1}
\vartheta_2=-2\bigl(\vartheta+\vartheta_1\bigr)\pmod{2\pi},
\end{equation}
whichever remainder defines the split. Every sawtooth built from $\vartheta_2$
is therefore the same function as the one built from $\vartheta_1$ and
retrogrades in the same places: among the $63$ ordinates of $3.0\leq T\leq5.2$
the splits at $R_{1ps}$, at $R/2$ and at $R_{1ak}$ run backwards at $10$, $9$
and $12$ of them. The retrograde belongs to the split, not to the branch.

\paragraph{The velocity split, and a curve that lands on every ordinate.} The
cure is to split $\zeta$ so that the offset $h$ of \eqref{eq:B1-rotated}
alternates by construction. With $\vartheta'$ the derivative of \eqref{eq:S-slope}, exactly
$\tfrac12\operatorname{Re}\tfrac{\Gamma'}{\Gamma}\bigl(\tfrac14+\tfrac{it}{2}\bigr)-\tfrac12\log\pi$,
put
\begin{equation}\label{eq:vel-split}
B_1^{\ast}=\zeta(s)+\frac{\zeta'(s)}{2\vartheta'(t)},\qquad
B_2^{\ast}=\zeta(s)-B_1^{\ast}=-\frac{\zeta'(s)}{2\vartheta'(t)}
=\frac{i}{2\vartheta'(t)}\frac{d\zeta}{dt},
\end{equation}
so leg~2 is the velocity of the chain's endpoint along the line, quarter-turned
and scaled by $1/2\vartheta'$. On $\sigma=\tfrac12$ this split is again
reflection-symmetric, $B_2^{\ast}=\chi\overline{B_1^{\ast}}$, so its legs are
exactly equal at every height,
\begin{equation}\label{eq:vel-legs}
\Bigl|\,\zeta+\frac{\zeta'}{2\vartheta'}\Bigr|=\frac{|\zeta'|}{2\vartheta'},
\end{equation}
which is the identity $\operatorname{Re}\bigl(\zeta\overline{\zeta'}\bigr)
=-\vartheta'|\zeta|^{2}$, itself no more than the statement that
$Z=e^{i\vartheta}\zeta$ is real. Differentiating that statement gives
$e^{i\vartheta}\zeta'=-iZ'-\vartheta'Z$, and hence \eqref{eq:B1-rotated} holds for the
new split with
\begin{equation}\label{eq:h-star}
e^{i\vartheta}B_1^{\ast}=\frac{Z}{2}-\frac{i\,Z'}{2\vartheta'},
\qquad h^{\ast}=-\frac{Z'}{2\vartheta'} .
\end{equation}
Now $h^{\ast}$ alternates because $Z'$ does, and \eqref{eq:crossing-rate}
returns $-Z'/2h^{\ast}=\vartheta'$: every crossing turns forward, and at exactly
the mean density. Setting
\begin{equation}\label{eq:N-star}
N_{\ast}(T)=\frac{1}{\pi}\vartheta\bigl(I(T)\bigr)
+\frac{1}{\pi}\vartheta_1^{\ast}(T)+\frac32,
\qquad \vartheta_1^{\ast}=\arg B_1^{\ast},
\end{equation}
we have a curve that is integer-valued at the ordinates and nowhere else, since
$\operatorname{Re}\bigl(e^{i\vartheta}B_1^{\ast}\bigr)=\tfrac12Z$ vanishes only
there, and that advances by exactly one from each ordinate to the next. It lands
on every one of the first $400$ ordinates, $N_{\ast}(I^{-1}(\gamma_n))=n$ to
$2.8\times10^{-13}$, against $50$ misses for $N_{\mathrm{ps}}$; over
$0.6\leq T\leq8$ it is strictly increasing on a grid of $3000$ points, no wrap
occurs ($|\vartheta_1^{\ast}|$ stays below $0.60\pi$), and it rises by $237.64$
across the $237$ zeros of that range. Height costs it nothing and no branch
correction is needed there either: at $\gamma_{10000}=9877.78$ and at
$\gamma_{50000}=40433.69$ the same formula with the same constant $\tfrac32$
still returns exactly $10000$ and $50000$. Where $N_{\mathrm{rvm}}$ threads the
staircase without touching it, and $N_{\mathrm{ps}}$ locates zeros without
counting them, $N_{\ast}$ does both. Like any count read off $Z$ it sees the
zeros on the line, which in the range plotted are all of them.

\begin{remark}[the two offsets $h_{R/2}$ and $h^{\ast}$]\label{rem:taper}
The offsets $h_{R/2}$ and $h^{\ast}$ of the $R/2$ split and the velocity split
are both drawn in Figure~\ref{fig:splits-h}, and they differ only by a
weighting of the links. The velocity offset has the closed form of
\eqref{eq:h-star},
\begin{equation}\label{eq:h-star-slope}
h^{\ast}(t)=-\frac{Z'(t)}{2\,\vartheta'(t)},
\end{equation}
the slope of $Z$ in units of $2\vartheta'$: $h^{\ast}$ vanishes exactly at the
critical points of $Z$ and is largest where $Z$ is steepest. For the $R/2$ split
$e^{i\vartheta}R$ is real, so $h$ is the sine companion of the main sum in the
$\vartheta$-frame,
$h_{R/2}=\operatorname{Im}\bigl(e^{i\vartheta}\Sigma_1\bigr)
=\sum_{n\leq m}n^{-1/2}\sin(\vartheta-t\log n)$, while differentiating the real
part of that rotated main sum and dividing by $-\vartheta'$ gives
\[
h^{\ast}=\sum_{n\leq m}n^{-1/2}
\Bigl(1-\frac{\log n}{\vartheta'}\Bigr)\sin\bigl(\vartheta-t\log n\bigr)+\cdots,
\]
the tail being the remainder's derivative; the displayed sum by itself
reproduces $h^{\ast}$ to $2.3\times10^{-2}$ over the first $400$ ordinates.
The weight is the phase rate $\vartheta'-\log n$ of link $n$ measured in units of
$\vartheta'$, so it runs from $1$ on the first link down to nearly $0$ on the last,
where $\log m\approx\vartheta'$. The counting angle is carried by the same links;
each one contributes in proportion to how much of the rotation rate it retains.
\end{remark}

\begin{remark}[the velocity split off the critical line]\label{rem:vel-offline}
The weight $2\vartheta'(t)$ in \eqref{eq:vel-split} is exactly
$-(\chi'/\chi)(\tfrac12+it)$, since $\chi=e^{-2i\vartheta}$ on the line.
Replacing it by that logarithmic derivative extends the split to the whole
strip:
\begin{equation}\label{eq:vel-split-offline}
B_1^{\ast}(s)=\zeta(s)-\frac{\zeta'(s)}{(\chi'/\chi)(s)},
\qquad
B_2^{\ast}=\zeta-B_1^{\ast},
\qquad
\frac{\chi'}{\chi}(s)=\log 2\pi+\frac{\pi}{2}\cot\frac{\pi s}{2}-\psi(1-s).
\end{equation}
Applying the functional equation twice gives $\chi(s)\chi(1-s)=1$, whose
logarithmic derivative forces $(\chi'/\chi)(s)=(\chi'/\chi)(1-s)$;
differentiating $\zeta(s)=\chi(s)\zeta(1-s)$ then yields the exact reflection
\[
B_2^{\ast}(s)=\chi(s)\,B_1^{\ast}(1-s)
\]
at every $s$ in the strip, pairing the split at $(\sigma,T)$ with its mirror
image at $(1-\sigma,T)$. On the critical line the mirror point is the point
itself, the weight is the real number $2\vartheta'$, and the identity
collapses to $B_2^{\ast}=\chi\overline{B_1^{\ast}}$ with its equal legs
\eqref{eq:vel-legs}. Off the line the weight is complex (imaginary part
$\approx\mp7\times10^{-4}$ at $t\approx280$), the offset $h^{\ast}$ ceases to
be a real scalar, and, as for every split, the equal-leg property survives
only on the line. Verified numerically to twenty-four digits at $\sigma=0.3$.

The $\chi$-weight also separates the two chains exactly. Write
$\zeta=F+\chi\widetilde G$ with $F=\Sigma_1+R_1$ and
$\widetilde G=(\Sigma_2+R_2)/\chi$, for any of the named splits: the second
half always carries the factor $\chi$, so peeling it off is natural.
Differentiating, multiplying by $\chi/\chi'$, and subtracting cancels the
$\chi\widetilde G$ term entirely; on the critical line, where
$\chi'/\chi=-2\vartheta'$, the result is
\begin{equation}\label{eq:vel-split-halves}
B_1^{\ast}
=\Sigma_1+R_1+\frac{(\Sigma_1+R_1)'+\chi\,\widetilde G'}{2\vartheta'},
\qquad
B_2^{\ast}
=\Sigma_2+R_2-\frac{(\Sigma_1+R_1)'+\chi\,\widetilde G'}{2\vartheta'}.
\end{equation}
The velocity split therefore hands the whole forward chain $\Sigma_1+R_1$ to
leg~1 and the whole reflected chain $\Sigma_2+R_2$ to leg~2, exactly, with no
mixing at the level of the sums; the two legs trade only the derivative
corrections, damped by $2\vartheta'$. The reflected
half influences $B_1^{\ast}$ only through its rate of change
$\widetilde G'$, never through its value. This is why $B_1^{\ast}$ crowds
against the other split points in Figure~\ref{fig:splits-h}, and the taper
weights $1-\log n/\vartheta'$ of Remark~\ref{rem:taper} are this cancellation
read term by term. Verified to twenty-four digits at $T=6.18$.
\end{remark}

\begin{figure}[H]
\centering
\includegraphics[width=0.78\textwidth]{fig_splits_h}
\caption{The offset $h$ of \eqref{eq:B1-rotated}, drawn at $T=6.18$. Because
$\Re\bigl(e^{i\vartheta}B_1\bigr)=Z/2$ for every split, all of the split points lie
on one line when $\sigma=\tfrac12$, the dashed line through $\zeta/2$
perpendicular to the base $O\to\zeta$ (the right angle is marked). What
distinguishes one split from another is $h$, the signed distance from $\zeta/2$
along that line, drawn as a thick stretch of the line for the $R/2$ split
(purple, $h_{R/2}=2.383$) and for the velocity split (teal, $h^{\ast}=1.619$);
the two legs of each split are drawn thin in its own colour and the two chains
fainter still. The arc at $O$ is $\vartheta$, swept clockwise off the positive real
axis onto the base line, since $\arg\zeta\equiv-\vartheta\pmod{\pi}$: with
$\vartheta=390.884$ reducing to $75.98^{\circ}$ modulo $\pi$, the base line leaves
the axis at $-75.98^{\circ}$, drawn dotted below $O$, and $\zeta$ itself lies on
the far side of $O$ at $104.02^{\circ}$ because $Z=-3.194$ is negative at this
height. Generated by \texttt{fig\_splits\_h.py}.}
\label{fig:splits-h}
\end{figure}

\begin{remark}[$N_{\ast}$ is monotonically increasing if RH is true]
\label{rem:nstar-mono}
Integer at the ordinates and nowhere else, and up by one from each to the next:
what \eqref{eq:N-star} leaves open is whether the curve ever runs backwards in
between. It does not, and the reason is a known consequence of the Riemann
hypothesis. With $W=e^{i\vartheta}B_1^{\ast}=\tfrac12Z-\tfrac{iZ'}{2\vartheta'}$ from
\eqref{eq:h-star} and $\pi N_{\ast}=\arg W+\tfrac{3\pi}{2}$, differentiating
and measuring against the local mean density $\vartheta'$ gives the slope in closed
form,
\begin{equation}\label{eq:nstar-rho}
\rho=\frac{1}{\vartheta'}\,\frac{d}{dt}\arg W
=\frac{\pi}{\vartheta'\,I'(T)}\,\frac{dN_{\ast}}{dT}
=\frac{Z'^2-ZZ''+\dfrac{\vartheta''}{\vartheta'}\,ZZ'}{\vartheta'^2Z^2+Z'^2},
\end{equation}
which returns $1$ at $Z=0$ and $-Z''/\vartheta'^2Z$ at $Z'=0$, as it must. Since
$\vartheta'$ and $I'$ are positive, $\rho$ and $dN_{\ast}/dT$ carry the same sign;
the denominator of \eqref{eq:nstar-rho} is positive and
$\vartheta''/\vartheta'=O(1/t\log t)$, so up to that negligible term the curve
increases exactly where the Laguerre expression $Z'^2-ZZ''$ is positive.

That expression belongs to $\Xi$. Writing
$\Xi(t)=\xi(\tfrac12+it)=-c(t)Z(t)$ with
$c(t)=\tfrac12\bigl(\tfrac14+t^2\bigr)\pi^{-1/4}
\bigl|\Gamma(\tfrac14+\tfrac{it}{2})\bigr|>0$, algebra gives
$\Xi'^2-\Xi\Xi''=c^2\bigl[(Z'^2-ZZ'')-(\log c)''Z^2\bigr]$. If the hypothesis
holds then every zero of $\Xi$ is real, its Hadamard product gives
$-(\log\Xi)''=\sum_k(t-\gamma_k)^{-2}$, and therefore
\begin{equation}\label{eq:laguerre-Z}
Z'^2-ZZ''=Z^2\Bigl[\sum_k\frac{1}{(t-\gamma_k)^2}+(\log c)''\Bigr],
\end{equation}
the sum running over the zeros of $\Xi$ of both signs. The bracket is positive
with room to spare: it is of the order of the squared zero density, carried
mostly by the two ordinates flanking $t$, while $(\log c)''=O(t^{-2})$ takes
away almost nothing. At $t=100.5$ it reads $2.2003$, of which the $600$ zeros
with $|\gamma|\leq542$ supply $2.1972$ and $(\log c)''$ removes
$1.7\times10^{-4}$.

Monotonicity is what makes the curve a counting function outright, rather than a
curve that merely meets the integers in the right places. $N_{\ast}$ is
continuous, it equals $n$ exactly at $\gamma_n$, and it is an integer nowhere
else; if it also increases then between consecutive ordinates it is trapped
strictly between $n$ and $n+1$, and taking the floor recovers
\eqref{eq:N-theta} on the nose,
\begin{equation}\label{eq:floor-N-star}
\bigl\lfloor N_{\ast}(T)\bigr\rfloor=N\bigl(I(T)\bigr),
\end{equation}
with no error term and no sawtooth. We find no exception at $5721$ grid points
across $14\leq t\leq300$ nor at $301$ across $9877\leq t\leq9880$. Three things
are wanted for it, and each is worth naming. The count is read off $Z$, so what
it counts is zeros on the critical line, and it is the hypothesis again that
makes those all of them. The zeros must be simple: at a double zero $Z$ and $Z'$
vanish together, $W=0$, the argument is undefined, and the curve would climb by
one where the count climbs by two; RH does not supply simplicity.
And $\vartheta_1^{\ast}$ must be read on a continuous branch. Subtracting
\eqref{eq:N-rvm} from \eqref{eq:N-star} gives
$\vartheta_1^{\ast}=\pi\bigl(N_{\ast}-N_{\mathrm{rvm}}-\tfrac12\bigr)$, so the
principal branch serves precisely while the exact count and the smooth one stay
within one of each other; since $S$ is unbounded that must eventually fail, and
past it the wrap count has to be carried explicitly, which is the bookkeeping
$N_{\mathrm{ps}}$ already needs at every height. There is room to spare at these
heights: over $14\leq t\leq300$ the angle $|\vartheta_1^{\ast}|$ reaches only
$0.506\pi$.

Read the other way, \eqref{eq:nstar-rho} can only fail where $ZZ''>Z'^2$, which
at a critical point of $Z$ means $Z$ and $Z''$ of one sign: a positive local
minimum or a negative local maximum, a pair of zeros that comes to the axis and
fails to reach it. The Laguerre inequality $\Xi'^2\geq\Xi\Xi''$, which Laguerre
proved for every real entire function with only real zeros, is precisely what
forbids that. Being a consequence of the hypothesis and not a theorem about
$\zeta$, it leaves the monotonicity conditional; and being only the first of a
family of inequalities of every order, which Csordas and
Varga~\cite{CsordasVarga1990} show to be equivalent to RH when taken all
together, it is weaker than the hypothesis. A failure of monotonicity would
therefore disprove RH, while a proof of monotonicity would not prove it. Nothing
is close to the boundary. On grids of spacing $0.1$ the slope never turns negative
and dips no lower than $0.431$ over $10\leq t\leq200$, $0.403$ over
$200\leq t\leq300$, $0.375$ over $1000\leq t\leq1030$ and $0.521$ over
$9877\leq t\leq9880$ beside $\gamma_{10000}$. The near misses do not threaten it
either, and they push $\rho$ the opposite way from what one might guess. At
Lehmer's pair~\cite{Lehmer1956} $\gamma_{6709}=7005.06287$ and
$\gamma_{6710}=7005.10056$ the hump between the two sits at $t=7005.08186$ with
$Z=3.97\times10^{-3}$, and there $\rho=458$, because $-Z''/\vartheta'^2Z$ grows
without bound as a hump flattens onto the axis. What such a pair costs is
margin, not sign: $Z'^2-ZZ''$ reads $0.089$ at that hump against $1.53$ and
$1.38$ just outside the pair. Verified by \texttt{check\_nstar\_mono.py}.
\end{remark}

\begin{remark}[regarding the slope of $N_{\ast}$ and the size of $|Z|$]
\label{rem:nstar-slope}
Since $\pi N_{\ast}=\arg W+\tfrac{3\pi}{2}$ and the positive factor $\tfrac12$
in $W$ leaves an argument alone,
$\pi N_{\ast}-\tfrac{3\pi}{2}=\arg\bigl(Z-iZ'/\vartheta'\bigr)$ is the polar angle
of the plane point $(Z,-Z'/\vartheta')$, the Pr\"ufer angle familiar from
oscillation counting with $\vartheta'$ in the role of the frequency. The slope of
the counting curve is therefore fixed by $Z$, $Z'$ and $Z''$ at that height
alone, and measured against the local mean density it is the
$\rho$ of \eqref{eq:nstar-rho}. Two values of it are exact. At an ordinate $Z$ vanishes and
\eqref{eq:crossing-rate} with $h^{\ast}=-Z'/2\vartheta'$ gives $\rho=1$, which we
confirm to twelve digits at $\gamma_1$, $\gamma_{100}$, $\gamma_{1000}$ and
$\gamma_{10000}$; the whole of the variation therefore lives strictly between
zeros. At an extremum of $Z$, where $Z'=0$ leaves $W=\tfrac12 Z$ real,
\begin{equation}\label{eq:rho-peak}
\rho=-\frac{Z''}{\vartheta'^2 Z}\qquad(Z'=0),
\end{equation}
so flatness at a hump is exactly a small curvature-to-height ratio, in units of
$\vartheta'^2$: a broad tall hump gives $\rho$ well below $1$ and a sharp small one
gives $\rho$ well above it. The consequence is that the flatness of
$N_{\ast}$ ranks the size of $|Z|$ between the zeros. Over the $100$ gaps of
$14\leq t\leq238$, $\rho$ at the hump against the height of that hump has rank
correlation $-0.987$, with $|Z|_{\max}\propto\rho^{-1.17}$; two decades higher
the same measurement reads $-0.943$ and $\rho^{-1.30}$ at
$\gamma_{1000}$, and $-0.944$ and $\rho^{-1.20}$ at $\gamma_{10000}$. Since
$\rho=1$ at every ordinate, the first local datum there is $\rho'(\gamma_n)$,
and that single number, read at one zero, ranks the height of the hump ahead of
it at $-0.80$, $-0.66$ and $-0.73$ in the three blocks: the more steeply the
slope falls away as one leaves a zero, the taller the hump being entered.
Figure~\ref{fig:nstar-slope} plots all of this, the slope against $|Z|$ in two
windows and then the two rankings. What the slope is really
measuring is the local phase rate of $Z$ rather than its amplitude, since for
$Z=A\cos\varphi$ with $\varphi'=\vartheta'$ one gets $\rho=1$ whatever $A$ may be;
$\rho<1$ says the phase is turning more slowly than $\vartheta'$, which is the same
statement as a wide gap, and it is the sparseness of the zeros that lets $|Z|$
grow. Consistently with that, $\rho$ at the hump and the gap width are
collinear to $-0.99$ in rank, so the slope should be read as a pointwise
version of the spacing rather than as independent information about $|Z|$.
Verified in \texttt{check\_nstar\_slope.py}.
\end{remark}

\begin{figure}[H]
\centering
\includegraphics[width=\textwidth]{fig_nstar_slope}
\caption{The slope $\rho$ of \eqref{eq:nstar-rho} against the size of $|Z|$
between ordinates, everything at $\sigma=\tfrac12$. \emph{Top:} $\rho$ in
purple on the magnified window of Figure~\ref{fig:nps-staircase}, with $|Z|$ in
red on the right-hand scale. The red dots mark the ordinates, where $\rho=1$
exactly; between them $\rho$ dips, and it dips deepest where $|Z|$ rises
highest. \emph{Second:} the same reading twice as high in the index, over
$12.5\leq T\leq12.6$. The hump that reaches $|Z|=9.85$ near $T=12.549$ sits on
the flattest stretch of the window, where $\rho$ bottoms at $0.356$, while each
ordinary hump carries only a shallow arch of $\rho$ just above $1$ and the two
narrowest gaps throw up spikes of $4.11$ and $2.99$.
\emph{Bottom left:} one point per gap for three blocks of ordinates spanning
two decades of $t$, the height of the hump against $1/\rho$ at that hump, both
on log scales, with the heights divided by a $31$-gap moving average so the
slow growth of $|Z|$ with height does not enter. The line is the fit over all
three blocks; the legend quotes the rank correlation against $\rho$ for each
block separately. \emph{Bottom right:} the same hump heights against
$\rho'(\gamma_n)$, read at the ordinate behind the hump rather than at the hump
itself, scaled by the mean spacing $\pi/\vartheta'$. Generated by
\texttt{fig\_nstar\_slope.py}.}
\label{fig:nstar-slope}
\end{figure}

\subsection{The envelope of $\left|\zeta\!\left(\tfrac12,t\right)\right|$}\label{sec:envelope}

On the critical line the chain is isosceles for every reflection-symmetric
split, by Corollary~\ref{cor:equal-legs}, so its two legs share one length
$L_1=|B_1|=|B_2|$ and the base is the chord they subtend across the fold
angle of the legs $\vartheta_2$ of \S\ref{sec:legs}. From $\zeta=B_1(1+e^{i\vartheta_2})$,
\begin{equation}\label{eq:zeta-chord}
\bigl|\zeta\bigl(\tfrac12+it\bigr)\bigr|
=2L_1\bigl|\cos\tfrac{\vartheta_2}{2}\bigr| ,
\end{equation}
so the modulus of $\zeta$ splits into a length carried by the links and an angle
read at the joint between the legs. Bounding the angular factor bounds $\zeta$.
Write $u=\vartheta_2/\pi$ reduced to $[0,2)$. Then
$\bigl|\cos\tfrac{\pi u}{2}\bigr|\leq1$ at once, and since
$\cos\tfrac{\pi u}{2}$ is concave on $0\leq u\leq1$ it stays above the chord
joining its endpoints there, giving
$\bigl|\cos\tfrac{\pi u}{2}\bigr|\geq|1-u|$ on the whole of $[0,2)$ by the
reflection $u\mapsto2-u$, with equality only at $u=0,1,2$. Hence
\begin{equation}\label{eq:env-bounds}
2L_1\Bigl|1-\frac{\vartheta_2}{\pi}\Bigr|
\;\leq\;\bigl|\zeta\bigl(\tfrac12+it\bigr)\bigr|\;\leq\;2L_1 .
\end{equation}
The upper envelope is attained where the two legs point the same way and the
chain is folded straight; the lower one is attained at every ordinate, where
$u=1$ and both sides vanish, and again where $u$ reaches $0$ or $2$ and the two
envelopes meet. The triangular factor $|1-\vartheta_2/\pi|$ is the same sawtooth
whose passages through zero do the counting in \S\ref{sec:counting}, appearing
here with an amplitude on it. Over $6\leq T\leq13$, which is
$264.9\leq t\leq1144.6$, we confirm \eqref{eq:zeta-chord} to
$1.7\times10^{-20}$ and both inequalities of \eqref{eq:env-bounds} at every one
of $2001$ sample points, the two legs agreeing to $2.0\times10^{-20}$.

Which split to use is settled by how often the upper envelope is reached.
Equation~\eqref{eq:B1-rotated} gives $2L_1=\sqrt{Z^2+4h^2}$, so $|\zeta|=2L_1$
exactly where the offset $h$ vanishes. For the split at the partial
summand $h$ vanishes only inside a gap whose two ends carry opposite $h$, which
is the alternation that the count of \S\ref{sec:counting} needs and that fails at
a retrograde ordinate, so the envelope is reached in some gaps and missed in
others: $3$ of the $7$ gaps of $6.125\leq T\leq6.275$, and $7$ of the $13$ gaps
of $12.5\leq T\leq12.6$. The velocity split has no such gaps. Its offset
$h^{\ast}=-Z'/2\vartheta'$ vanishes at every extremum of $Z$, so
\begin{equation}\label{eq:env-star}
2L_1^{\ast}=\sqrt{Z^2+\Bigl(\frac{Z'}{\vartheta'}\Bigr)^{2}}
\;\geq\;|Z|=\bigl|\zeta\bigl(\tfrac12+it\bigr)\bigr| ,
\end{equation}
with equality on every hump: $7$ touches in $7$ gaps and $13$ in $13$ in those
two windows, matched to $10^{-25}$. This is the familiar envelope of an
oscillation of instantaneous frequency $\vartheta'$, with $Z'/\vartheta'$ as the
quadrature companion of $Z$, and it is tighter on average as well, the mean of
$|\zeta|/2L_1$ over $6\leq T\leq13$ being $0.635$ against $0.586$ for the split
at the partial summand.

Both readings are polar coordinates of one function. The rotated bisector point
$W=e^{i\vartheta}B_1^{\ast}=\tfrac12Z-\tfrac{iZ'}{2\vartheta'}$ of \eqref{eq:h-star}
has modulus $L_1^{\ast}$ and argument $\pi\bigl(N_{\ast}-\tfrac32\bigr)$ by
\eqref{eq:N-star}, and $Z=2\operatorname{Re}W$, so
\begin{equation}\label{eq:Z-polar}
Z(t)=-2\,|W|\,\sin\bigl(\pi N_{\ast}(T)\bigr),
\end{equation}
which we confirm to $6\times10^{-13}$, the precision of the counting curve
itself. This is a restatement rather than a new fact, but it puts the two
subsections together: the amplitude of the velocity split is the envelope of
$\zeta$ and its phase is the counting curve, so the zeros are the moments when
the phase passes an integer and the extrema are the moments when the amplitude
is attained. Remark~\ref{rem:nstar-slope} reads the rate of that phase; the
envelope \eqref{eq:env-star} reads its amplitude. Verified in
\texttt{check\_envelope.py}.

\begin{figure}[H]
\centering
\includegraphics[width=\textwidth]{fig_envelope}
\caption{The two envelopes of \eqref{eq:env-bounds}, at $\sigma=\tfrac12$, over
$12.5\leq T\leq12.6$, the second window of Figure~\ref{fig:nstar-slope}, which
is $1061.3\leq t\leq1077.7$ and holds $14$ ordinates.
In each panel $|\zeta|$ is black,
the upper envelope $2L_1$ red, the lower envelope
$2L_1|1-\vartheta_2/\pi|$ a teal sawtooth, and the grey band between them
is where $|\zeta|$ is confined. Red dots on the axis are the ordinates, where
the lower envelope touches zero along with $|\zeta|$, and the open circles are
the touches of the upper envelope, one for each vanishing of the offset $h$.
\emph{Left:} the split at the partial summand, $B_1=\Sigma_1+R_{1ps}$, whose $h$
fails to alternate at a retrograde ordinate, so it reaches its bound in only $7$
of the $13$ gaps and several humps rise nowhere near it. \emph{Right:} the
velocity split, whose $h^{\ast}$ vanishes at every extremum of $Z$, so the upper
envelope is touched in every gap. Generated by \texttt{fig\_envelope.py}.}
\label{fig:envelope}
\end{figure}

\begin{remark}[regarding hunting for large values with the first few links]
\label{rem:large-values}
Neither the envelope nor the slope of Remark~\ref{rem:nstar-slope} can be used to
scout, since $2L_1$ needs $R=\zeta-\Sigma_1-\Sigma_2$ and $2L_1^{\ast}$ needs $Z$
and $Z'$, so both are functions of $\zeta$ at a point one has already evaluated.
What can be computed without $\zeta$ is the reason the envelope grows. Leg~1 is
mostly the link sum, so
the size of $2L_1$ is a question of how well the link phases $t\log n$ align, and
the ceiling is the aligned case,
\begin{equation}\label{eq:link-ceiling}
\bigl|\zeta\bigl(\tfrac12+it\bigr)\bigr|\;\leq\;2L_1\;\leq\;
2\sum_{n\leq m}n^{-1/2}+2|R_{1ps}|\;\approx\;4\sqrt{T}
\;\approx\;4\Bigl(\frac{t}{2\pi}\Bigr)^{1/4}
=\frac{4}{(2\pi)^{1/4}}\,t^{1/4}\;\approx\;2.53\,t^{1/4},
\end{equation}
a bound that costs nothing to evaluate. Read in $t$ it is nothing new: the
exponent $\tfrac14$ is the convexity, or trivial, bound $\mu(\tfrac12)\le\tfrac14$
of Lindel\"of~\cite{Lindelof1908}, where $\mu(\sigma)$ is the least exponent with
$\zeta(\sigma+it)=O(t^{\mu})$, and \eqref{eq:link-ceiling} recovers it with an
explicit constant. Nor can this path do better. The last step is the triangle
inequality, which keeps only the aligned case and so discards every cancellation
among the phases $t\log n$; in the geometry of this paper $4\sqrt T$ is the arc
length of the chain laid out straight, while $|\zeta|$ is the distance between its
endpoints once it has curled into its spirals. Every improvement past $\tfrac14$
comes from that discarded cancellation, beginning with the $\tfrac16$ that Hardy
and Littlewood obtained by Weyl's method on this same sum and standing today at
$\tfrac{13}{84}$~\cite{Bourgain2017}; the Lindel\"of hypothesis is
$\mu(\tfrac12)=0$. Two things follow, both drawn in
Figure~\ref{fig:large-values}. First, the alignment can be read from a
truncation. Scoring an interval by
$S_K(t)=\bigl|\sum_{n\leq K}n^{-1/2}e^{-it\log n}\bigr|$ and then evaluating $Z$
only at the strongest peaks of that score recovers most of the interval's largest
$|Z|$: averaged over the $19$ unit intervals $5\leq T\leq24$, which is
$190\leq t\leq3771$, the three strongest peaks of the five-link score reach
$0.967$ of it and the single strongest peak of the ten-link score reaches
$0.961$. The point is not the saving at these heights but that $K$ need not grow
with $t$ while $m$ does. Second, the interval maxima themselves are nearly
regular: $\max|Z|$ over $[T,T+1]$ fits $2.93\,T^{0.485}$, its ratio to
$\sqrt{T}$ has mean $2.823$ and standard deviation $0.123$ across those
intervals, and it holds a fraction of the ceiling \eqref{eq:link-ceiling} that
ranges from $0.97$ down to $0.76$, the two extremes falling at $T=9$ and
$T=22$. Most of that spread is the ceiling rather than the maxima: the exact sum
$\sum_{n\leq m}n^{-1/2}=2\sqrt m+\zeta(\tfrac12)+o(1)$ climbs toward its own
asymptote $2\sqrt m$ across the range, from $0.72$ of $4\sqrt T$ at $T=5$ to
$0.86$ at $T=23$, since $\zeta(\tfrac12)=-1.4604$ still weighs heavily at
$m=5$. Against $\sqrt T$ itself the maxima are flat, $2.711$ at $T=5$ and
$2.670$ at $T=23$, which is the same fact as the fitted exponent $0.485$ being
essentially the ceiling's $\tfrac12$: at these heights the maxima simply ride the
trivial bound. In these coordinates the Lindel\"of hypothesis is precisely the
statement that they eventually stop, that $\max|Z|/\sqrt T$ falls like
$T^{-1/2+\varepsilon}$, the chain curling up so efficiently that its endpoints
stay bounded while its arc length grows like $\sqrt T$. Nothing of that is
visible over $19$ intervals with $m\leq23$ links, and nothing here bears on it
either way. So the answer
to whether the next unit of the index will beat the last is yes by default, which
it is in $13$ of the $18$ consecutive comparisons, and the exceptions are what
needs predicting: on the maxima divided by $\sqrt{T}$ the peak of the ten-link
score ranks the intervals at $+0.72$, while the three-link score, whose peak is
nearly the same in every interval, carries nothing at $-0.11$. That last figure
rests on $19$ intervals at modest height and should be read as an indication
only. Verified in \texttt{check\_large\_values.py}.
\end{remark}

\begin{figure}[H]
\centering
\includegraphics[width=\textwidth]{fig_large_values}
\caption{Reading large values off the links, over the $19$ unit intervals
$5\leq T\leq24$. \emph{Left:} the fraction of an interval's largest $|Z|$ found by
evaluating $Z$ only at the strongest one, three or five peaks of the $K$-link
score $S_K$, averaged over the intervals, with the shading reaching down to the
worst single interval. Seven links suffice for a single candidate to land within
$5\%$. \emph{Right:} the interval maximum against $\sqrt{T}$, on which it is
nearly straight, with the link ceiling of \eqref{eq:link-ceiling} above it and
the fitted power law through it. The gap between the two curves is the alignment
the links never quite achieve. Generated by \texttt{fig\_large\_values.py}.}
\label{fig:large-values}
\end{figure}

\subsection{Collinearity of the three first-half remainders}\label{sec:colinearity}

Collinearity of three points is unchanged by a common translation, so the
shared tail $\Sigma_1$ of the three heads
$B_{1ps},B_{1rs},B_{1ak}$ of
\S\ref{sec:ps-ak-r2} cancels, and the
three remainder vectors $R_{1ps}$, $R_{1rs}=\tfrac12 R$, and $R_{1ak}$ can be
compared directly as points in the plane.
Figure~\ref{fig:colinearity-strip} plots two independent, scale-invariant
measures of how far these three points are from lying on a common line. The
first is the triangle flatness
\begin{equation}\label{eq:flatness}
F
\;=\;
\frac{2\,\lvert\mathrm{area}\rvert}{\mathrm{diam}^2},
\end{equation}
where the area is that of the triangle with the three points as vertices and
the diameter is the largest of the three pairwise distances. The second is
the aspect ratio $\lambda_{\min}/\lambda_{\max}$ of the covariance matrix of
the three points (a principal-component measure). Both quantities are zero
exactly when the points are collinear, and both are invariant under
translation, rotation, and rescaling, so the shrinking of the remainders
with growing $T$ cannot masquerade as collinearity. Each panel is a grid of
$1501$ $\sigma$-samples by $1801$ $T$-rows over $0\le\sigma\le1$,
$3.9\le T\le5.1$, computed from the exact remainder formulas, with the
$\sigma$-sample nearest $\tfrac12$ snapped exactly onto the critical line.
The color scale is logarithmic, and that is what separates \emph{exact}
collinearity (values at the double-precision floor, $10^{-16}$ to
$10^{-12}$) from the merely small values, $10^{-4}$ to $10^{-1}$, that
fill the rest of the strip.

\begin{figure}[H]
\centering
\includegraphics[width=0.95\textwidth]{fig_colinearity_3p9_5p1}
\caption{Collinearity of the three first-half remainders $R_{1ps}$,
$R_{1rs}=\tfrac12 R$, $R_{1ak}$ (dark $=$ collinear). Horizontal axis
$\sigma\in[0,1]$; vertical axis the spiral index $T\in[3.9,5.1]$. Left:
triangle flatness \eqref{eq:flatness}; right: PCA aspect ratio
$\lambda_{\min}/\lambda_{\max}$. Both panels use a log color scale; green
markers are zeta zeros at $\sigma=\tfrac12$. The dark column pinned at the
numerical floor along $\sigma=\tfrac12$ shows the three points are exactly
collinear everywhere on the critical line; off the line collinearity occurs
only along the isolated dark horizontal bands. Generated by
\texttt{fig\_colinearity\_zoom.py} from data produced by
\texttt{remainder-colinearity-heatmap-zoom.mjs}.}
\label{fig:colinearity-strip}
\end{figure}

Both panels show the same two facts. First, the column $\sigma=\tfrac12$
sits at the numerical floor at every $T$: on the critical line the three
first-half remainders are exactly collinear, not merely nearly so. Second,
off the line the measures are many orders of magnitude larger everywhere
except along isolated dark horizontal bands (near $T\approx4.03$, $4.62$,
$4.80$, and $5.03$ in this window), where the triangle momentarily passes
through a flat configuration as its signed area changes sign. The measured
values are symmetric under $\sigma\mapsto1-\sigma$, a reflection of the
functional equation.

The same pattern persists across the whole strip.
Figure~\ref{fig:colinearity-four-panel} repeats the four-panel layout of
Figure~\ref{fig:leg-imbalance} (the first two zoom bands of that figure, a
third band straddling the integer $T=17$, and the same full strip
$0\le T\le20$), with the PCA aspect ratio as the plotted field.
As in the leg-imbalance panels, the dark column on $\sigma=\tfrac12$ runs
unbroken through every window, while off the line the only dark features are
the isolated horizontal collinearity bands. The bands do not sit at integer
$T$: crossing an integer the three remainder formulas jump ($\Sigma_1$ gains
a summand and the sign $(-1)^{\lfloor T\rfloor}$ flips), the signed area of
the triangle changes sign in that jump and lands very close to zero, and the
area then crosses zero a short distance above the integer. The third panel
shows one full unit interval, $16.95\le T\le18.05$: the light-to-dark break
at the integer $T=17$ itself, a nearly flat strip at every $\sigma$ just
above it, and then the five collinearity bands of the cell, at
$T\approx17.011$, $17.252$, $17.612$, $17.752$, and $17.796$, before the
pattern repeats across the next break at $T=18$. Numerically the bands come
five to a unit interval, near $\{T\}\approx0.011$, $0.252$, $0.612$,
$0.752$, $0.796$
at this height, with the second and fourth tracking the pole heights
$\{T\}=\tfrac14,\tfrac34$ of $d_1$ (\S\ref{sec:pole-locations}): near a pole
$R_{1ps}$ runs off to infinity along the fixed direction of the next
summand, so the triangle is dragged through a flat configuration nearby.
The band positions are almost independent of $\sigma$ (they move by less
than $10^{-4}$ between $\sigma=0.2$ and $\sigma=0.35$), which is why the
bands are horizontal.

\begin{figure}[H]
\centering
\includegraphics[width=\textwidth]{colinearity_four_panel}
\caption{PCA aspect ratio $\lambda_{\min}/\lambda_{\max}$ of the three
first-half remainders $(R_{1ps},R_{1rs},R_{1ak})$ (dark purple $=$
collinear), in the four-panel layout of Figure~\ref{fig:leg-imbalance}.
Horizontal axis $\sigma\in[0,1]$; vertical axis the spiral index $T$. Color
is on a log scale. The
rightmost panel is the full strip $0\le T\le20$; the first three panels are
zooms of the outlined bands
($4.65\le T\le4.80$, $9.42\le T\le9.46$, $16.95\le T\le18.05$).
The dark vertical band on $\sigma=\tfrac12$ is the exact collinearity of the
three first-half remainders on the critical line; the dark horizontal bands
are the isolated off-line collinearity events. Generated by
\texttt{plot\_colinearity\_four\_panel.py}.}
\label{fig:colinearity-four-panel}
\end{figure}

Why can this collinearity not be used like the two routes of
\S\ref{sec:toward-proof}? Because collinearity
is \emph{not} forced at a zero. For $\zeta\neq0$ the equal-leg condition for
a split says precisely that its head $B_1$ is equidistant from $O$ and
$\zeta$, i.e.\ that $B_1$ lies on the perpendicular bisector of the chord
from $O$ to $\zeta$; if all three splits had equal legs at a point where
$\zeta\neq0$, all three heads would lie on that single line and the three
first-half remainders would be collinear. But at a zero the chord
degenerates: $\zeta=O$, every point of the plane is equidistant from the two
endpoints, and the (now automatic) equal-leg conditions carry no
collinearity information at all. All three leg pairs are equal at a zero,
as they must be, while the three points can remain in general position.
So a zero implies equal legs for all three splits, but not collinearity, and
the three points need not line up at an off-line zero.

A proof through collinearity therefore needs a bridge that the routes of
\S\ref{sec:toward-proof} do not: first, prove that the exact collinearity
along $\sigma=\tfrac12$ is an
identity of the three remainder formulas rather than a numerical
coincidence; second, prove that off the line it fails outside isolated
curves like the dark bands of Figure~\ref{fig:colinearity-strip}; and third,
connect a hypothetical off-line zero to the collinearity structure, for
instance through nearby off-line points with $\zeta\neq0$ at which all three
splits have equal legs, which by the perpendicular-bisector argument above
would force a collinearity that the off-line geometry forbids. The first
two steps would establish collinearity as an exact characterization of the
critical line; only the third would turn it into a statement about zeros.

\subsection{Joint angles}\label{sec:joint-angles}

The joint angles of the forward chain can be read on their own, without the
chain, and the Zest joint-angle view does exactly that: one dot per joint,
the turning angle on the vertical axis, the joint index along the horizontal.
The turning angle at joint $n$ is the continuous joint angle \eqref{eq:jangle}
evaluated at an integer joint,
\begin{equation}\label{eq:theta-joint}
\theta_n \;=\; -I(T)\,\log\tfrac{n}{n-1}
\qquad\text{(reduced to $[-\pi,\pi]$)},
\qquad
\nu(n) \;=\; \frac{I(T)}{2\pi\,n(n-1)},
\end{equation}
where $\nu$, the derivative of
the unwrapped angle divided by $2\pi$, is the local frequency of the strip in
cycles per joint. Two features organize the picture, both consequences of
$I(T)$. First, at the bisector joint $n=\lfloor T\rfloor+1$ the frequency is
$\nu=1+\frac{1}{6T(T+1)}$, one full turn per joint to nine decimals at
$T=13000$, and the angle there is $-\pi$ exactly: that is
\eqref{eq:jangle-count}, the bisector joint of \S\ref{sec:IT-origin}, seen as the right
edge of the strip. The joints just behind it inherit almost the same angle,
$3.141109$, $3.139659$, $3.137242$ at $n=T,T-1,T-2$, which is the same
near-stationarity that makes the bisector frame of
Figure~\ref{fig:bisector-frame} repeat. Second, since
$\nu(n)\approx T^{2}/n^{2}$ falls from $\nu\sim10^{8}$ at $n=2$ to $1$ at that
bisector joint, the strip is a sampled chirp, and it is stratified by the joints where
$\nu$ is the reciprocal of a rational. Writing that rational as a Farey
fraction $f=p/q$ in $(0,1]$, its \emph{caustic joint} is the solution of
$n(n-1)=f\,I(T)/2\pi$, that is $\nu=1/f$, so
\begin{equation}\label{eq:caustic-joint}
n_{f}\;=\;\tfrac12\Bigl(1+\sqrt{1+2f\,I(T)/\pi}\,\Bigr)\;\approx\;T\sqrt{f},
\end{equation}
and there the joints fall into exactly $p$ strands, the denominator of
$\nu=q/p$, each a parabola of curvature $4\pi\nu/n$ per joint squared because
the unwrapped angle is locally quadratic (exactly
$|\widetilde\theta''(n_f)|$ of \eqref{eq:caustic-derivs} below). Small denominators give the spines
visible in Figure~\ref{fig:joint-angles}, drawn at the integer index $T=13000$,
where $t=I(T)=1{,}061{,}939{,}999.369541$: there the fractions
$\tfrac14,\tfrac13,\tfrac12,\tfrac23,\tfrac34$ sit at
$n=6501,7506,9193,10615,11259$, and $f=1$ is the bisector joint itself. The parabolic
strands are the same local quadratic phase that produces the Cornu-like arcs of
the curlicue literature of \S\ref{sec:curlicues}, here read off the joints
rather than the partial sums.

\begin{figure}[H]
\centering
\includegraphics[width=\textwidth]{fig_joint_angles}
\caption{The joint-angle strip at $T=13000$, where
$t=I(T)=1{,}061{,}939{,}999.369541$, one dot per joint of the forward
chain, $n=2,\dots,\lfloor T\rfloor+1$, with $\theta_n$ of
\eqref{eq:theta-joint} on the vertical axis. The two edges $\pm\pi$ are
identified, so the strip is a cylinder. \emph{Top:} all $13000$ joints, carrying
three horizontal scales: the Farey fractions $f=p/q$ with $q\le7$ above, each
at its caustic joint \eqref{eq:caustic-joint} and staggered in two rows where
they crowd; the normalized joint fraction
$u=(n-1)/\lfloor T\rfloor$ below, so the bisector joint sits at $u=1$; and the joint index
$n$ below that. The red circle at the right edge is the bisector joint
$n=\lfloor T\rfloor+1$, where $\theta=-\pi$ exactly. \emph{Bottom left:} a
window at $f=\tfrac25$, showing its $p=2$ parabolic strands as dots, each with
its fitted arc \eqref{eq:fitted-arc} drawn through them; the two arcs sit half a
turn apart, so strand $1$ has its vertex wrapped to the opposite edge.
\emph{Bottom right:} the last tenth of the strip, where the
strands steepen and run into that bisector joint on the right in the figure. Dashed lines carry each window back to
the stretch of the joint-$n$ scale it magnifies, marked there by a bracket.
Generated by \texttt{fig\_joint\_angles.py}.}
\label{fig:joint-angles}
\end{figure}

\paragraph{The fitted curve.} The strands are not merely parabola-like; they can
be written down. Let $\widetilde\theta(n)=-I(T)\log\frac{n}{n-1}$ be the
unwrapped angle, so $\theta_n$ is $\widetilde\theta(n)$ reduced to
$[-\pi,\pi]$, and
differentiate it three times:
\begin{equation}\label{eq:jangle-derivs}
\widetilde\theta'(n)=\frac{I(T)}{n(n-1)},\qquad
\widetilde\theta''(n)=-\frac{I(T)\,(2n-1)}{[\,n(n-1)\,]^{2}},\qquad
\widetilde\theta'''(n)=\frac{2\,I(T)\,(3n^{2}-3n+1)}{[\,n(n-1)\,]^{3}}.
\end{equation}
At a caustic these have closed forms in the fraction alone, since
\eqref{eq:caustic-joint} fixes $n_f(n_f-1)=f\,I(T)/2\pi$ and
$2n_f-1=\sqrt{1+2f\,I(T)/\pi}$ exactly:
\begin{equation}\label{eq:caustic-derivs}
\begin{aligned}
\widetilde\theta'(n_f)&=\frac{2\pi}{f}=\frac{2\pi q}{p},
\qquad
\widetilde\theta''(n_f)=-\frac{4\pi^{2}}{f^{2}I(T)}
\sqrt{1+\frac{2f\,I(T)}{\pi}}\;\approx\;-\frac{4\pi}{f^{3/2}T},\\[2pt]
\widetilde\theta'''(n_f)&=\frac{16\pi^{3}}{f^{3}I(T)^{2}}
\Bigl(1+\frac{3f\,I(T)}{2\pi}\Bigr)\;\approx\;\frac{12\pi}{f^{2}T^{2}}.
\end{aligned}
\end{equation}
Write $\delta=n-n_f$ for the offset in joints from the caustic. The linear term
$\widetilde\theta'(n_f)\,\delta=(2\pi q/p)\,\delta$ is a carrier sweeping several
radians per joint, and it is what interleaves the dots into strands in the first
place. Along one strand, though, $\delta$ advances by $p$ at a time, and
$(2\pi q/p)\,pm=2\pi q\,m$ drops out: the carrier contributes nothing but a
constant. Absorbing it together with $\widetilde\theta(n_f)$, strand $j$ carries
the phase constant
\begin{equation}\label{eq:caustic-Cj}
C_j\;=\;\widetilde\theta(n_f)+\frac{2\pi q}{p}\bigl([n_f]+j-n_f\bigr),
\qquad j=0,\dots,p-1,
\end{equation}
with $[\,\cdot\,]$ the nearest integer, and the strand itself is that cubic
reduced to $[-\pi,\pi]$
\begin{equation}\label{eq:fitted-arc}
\rho_j(\delta)\;=\; C_j
+\tfrac12\widetilde\theta''(n_f)\,\delta^{2}
+\tfrac16\widetilde\theta'''(n_f)\,\delta^{3}
\qquad\text{(reduced to $[-\pi,\pi]$)}.
\end{equation}
Because $\gcd(p,q)=1$ the constants $C_j$ run through all $p$ residues of
$2\pi/p$, so the $p$ arcs are one parabola stacked at equal spacing around the
cylinder. Third order is the right order: the second derivative supplies the
parabola, the third its visible skew, and the fourth is smaller by a further
$\delta/n_f\sim T^{-1/2}$.

The arcs drawn in Figure~\ref{fig:joint-angles} are \eqref{eq:fitted-arc} at
$f=\tfrac25$ and $T=13000$, where $n_f=8222.74$,
$\widetilde\theta''(n_f)=-3.8208\times10^{-3}$ and
$\widetilde\theta'''(n_f)=1.3941\times10^{-6}$ radians per joint squared and
cubed, and the carrier is $2\pi q/p=5\pi\equiv\pi$, which is why the two strands
sit half a turn apart. Over the $\pm30$ joints shown they track the dots to
$2.2\times10^{-5}$ radians, that being the fourth-order term: despite the name it
carries in the Zest joint-angle view, nothing is fitted in the statistical sense,
and \eqref{eq:fitted-arc} is evaluated exactly. Two limits are worth noting.
Toward small $f$ the parabola stiffens like $f^{-3/2}$, which is why the
left end of the strip looks like noise at this scale: consecutive joints there
are many turns apart and only the caustics of very small $f$ remain resolved. At
the other end $f=1$ closes the family, with $p=q=1$: a single strand, a carrier of
one full turn per joint, and $C_0=-\pi$. That is the bisector joint of
\eqref{eq:jangle-count} again, so the bisector joint this paper is built on is the last
caustic of the sequence, and the right-hand panel of
Figure~\ref{fig:joint-angles} is the run of the high-denominator caustics into
it.

\subsection{Incremental change}\label{sec:incremental}

The strip is a snapshot. Advance the index and every dot moves, though not by
the same amount. Differentiating the unwrapped angle in $T$ at a fixed joint,
\begin{equation}\label{eq:jangle-drift}
\begin{aligned}
\frac{\partial\widetilde\theta}{\partial T}(n)&=-\,I'(T)\log\frac{n}{n-1},
\qquad L=\log\Bigl(1+\tfrac1T\Bigr),\\[2pt]
I'(T)&=\frac{\pi}{L^{2}}\Bigl(2L+\frac{2T+1}{T(T+1)}\Bigr)
=4\pi T+2\pi+O(T^{-1}),
\end{aligned}
\end{equation}
which in cycles is a drift rate
\begin{equation}\label{eq:jangle-drift-cycles}
\mu(n)\;=\;\frac{I'(T)}{2\pi}\log\frac{n}{n-1}
\;\approx\;\frac{2T+1}{n-\tfrac12}\;\approx\;\frac{2}{u},
\qquad u=\frac{n-1}{\lfloor T\rfloor}.
\end{equation}
A value of $\mu$ is easiest to read as a lap count: $\mu(n)=3$ means that dot
laps the strip three times, falling from $+\pi$ to $-\pi$, wrapping, and
repeating, while you advance $T$ by one. Nothing here needs differencing:
$I'(T)$ is available in closed form as
$\pi\bigl(2L+(2T+1)/T(T+1)\bigr)/L^{2}$ with $L=\log(1+1/T)$, which is
$4\pi T+2\pi$ to leading order, about $163{,}369$ at $T=13000$.

The sign is the same everywhere, so the whole strip drifts downward; the
magnitude is not, and that is the point. Since
$\mu(n)\approx(2T+1)/(n-\tfrac12)\approx2/u$, the rate is essentially a $1/u$
hyperbola, equal to exactly $2$ at the bisector joint $u=1$ and growing leftward
to $1.8\times10^{4}$ cycles per unit index at $n=2$ --- where, incidentally, the
$2/u$ form is no longer any good, reading $26{,}000$, since
$\log\frac{n}{n-1}$ and $1/(n-\tfrac12)$ part company at the first few joints.
The strip does not translate as the index advances, it shears. Read the other
way, $1/\mu(n)$ is the index step that returns joint $n$ to the same angle: a
dot at $u=\tfrac12$ comes back around every $\Delta T=0.25$, one at $u=0.1$
every $\Delta T=0.05$, and one at the bisector joint every half unit.

At the right edge the value is exact and already familiar: the middle expression
of \eqref{eq:jangle-drift-cycles} at $n=\lfloor T\rfloor+1$ is
$(2T+1)/(T+\tfrac12)=2$, and the exact rate is $2.000000001$ at $T=13000$. Two
cycles per unit index is one more than \eqref{eq:jangle-count} reports for the
joint angle, and the discrepancy is the joint itself moving: the bisector joint sits at
$n=\lfloor T\rfloor+1$, which advances one joint per unit index, and one joint
there is worth $\nu=1$ turn, so an observer riding the bisector joint sees $\mu-\nu=1$
cycle, namely $\frac{d}{dT}\pi(2T+1)=2\pi$. The two statements differ only in
whether one rides the bisector joint or stands at a fixed joint.

The pattern the dots drift through, by contrast, hardly moves at all.
Differentiating \eqref{eq:caustic-joint} gives
$dn_f/dT=f\,I'(T)\big/2\pi(2n_f-1)\approx\sqrt f$, so a caustic creeps outward
by a fraction of a joint per unit index, and its normalized position
$u_f=(n_f-1)/\lfloor T\rfloor\approx\sqrt f$ does not depend on $T$ at all: the
spines stand still while the joints stream through them. Their drift rate
follows from pairing \eqref{eq:jangle-drift-cycles} with the cycles per joint
$\nu=I(T)/2\pi n(n-1)$ of \eqref{eq:theta-joint}, which gives
$\mu\approx2\sqrt\nu$: at the caustic of Farey fraction $f$, where $\nu=1/f$,
the dots slide at $2/\sqrt f$ cycles per unit index, four at $f=\tfrac14$ and
two at $f=1$. Only at an integer index does the frame change, when a new link is
appended, a new joint appears at the right edge with $\theta=-\pi$, and $u$
rescales by one part in $\lfloor T\rfloor$.

Figure~\ref{fig:incremental-change} shows one step, $\Delta T=0.02$, which
advances $t$ by $\Delta I=3267.38$. The travel is $0.04/u$ cycles: $14.4^{\circ}$
at the bisector joint, $28.8^{\circ}$ at $u=\tfrac12$, $143.9^{\circ}$ at $u=0.1$, half a
turn at $u=0.0800$, joint $n=1041$, and a full turn at $u=0.0400$, joint
$n=521$. Since the motion is downward everywhere, the arrows in the figure are
drawn downward and wrapped at the lower edge, which is where they land for
joints left of $n=1041$. Left of $n=521$ the joint has gone round at least once,
and rather than draw what is left over the arrow there sweeps the whole strip
from $\pi$ to $-\pi$; at $n=2$ the travel is $2264$ radians, some $360$ turns.
This is the sense in which the left end of the strip is a chirp in the index as
well as in the joint number: it is not only that neighboring joints are many
turns apart at fixed $T$, but that one joint sweeps many turns over a step in
$T$ too small to move the bisector joint by a quarter of a radian. In the Zest
joint-angle view $\mu$ is the blue cycles-per-$T$ overlay, drawn against its own
axis rescaled to whatever window is on screen; the bottom panel of
Figure~\ref{fig:incremental-change} is that overlay on a logarithmic axis, and
the middle panel is a picture of the same quantity at $\Delta T=0.02$.

\begin{figure}[H]
\centering
\includegraphics[width=\textwidth]{fig_incremental_change}
\caption{One increment of the index, at $T=13000$. \emph{Top:} the whole
joint-angle strip as in Figure~\ref{fig:joint-angles}, with the shaded first
half magnified below. \emph{Middle:} that half twice over, grey at $T=13000$ and
red at $T=13000.02$, every dot having drifted downward by
$\mu(n)\,\Delta T\approx0.04/u$ cycles of \eqref{eq:jangle-drift-cycles}. Fifty
joints spaced evenly across the panel are circled in both colours, with an arrow
along the path travelled; since the drift is downward at every joint the arrows
run downward and wrap at the lower edge, and where the travel exceeds a full turn
the arrow sweeps the whole strip, from $\pi$ to $-\pi$. The travel
lengthens leftward like $1/u$, passing half a turn at $n=1041$ and a full turn
at $n=521$, marked by the dotted vertical line. Both strips carry the Farey
fractions $f=p/q$, $q\le7$, each at its caustic joint
\eqref{eq:caustic-joint}; only $f\le\tfrac14$ falls in the magnified half, since
$u_f\approx\sqrt f$. \emph{Bottom:} the drift rate itself on a logarithmic axis:
exact (blue) against the $2/u$ form (grey), the two separating only at the first
few joints, with the dashed line at one full turn per $\Delta T=0.02$. Generated by
\texttt{fig\_incremental\_change.py}.}
\label{fig:incremental-change}
\end{figure}

\subsection{Bearings of champions}\label{sec:champion-bearings}

The chain of this paper stops at the bisector link, but the chain of $\zeta$ itself runs
much further, and one link far out along it carries a signature of the largest
values of $|\zeta|$. Call a height a \emph{champion} when
$\bigl|\zeta(\tfrac12+it)\bigr|$ exceeds every value it has taken below that
height, a running record. We use a ground-truth list of $149$ of them, from
$t=10.21$ to $t=8.13\times10^{11}$, over which $|\zeta|$ climbs from $1.55$ to
$334$.\footnote{Provenance of the champion list: AIM workshop on AI and
number theory.}

The link in question is the last one, the term at
\begin{equation}\label{eq:last-link}
n_{\text{last}}=\lfloor t/\pi\rfloor\;\approx\;2T^{2},
\end{equation}
the index following from $\sqrt{t/2\pi}\approx T+\tfrac12$
(\S\ref{sec:IT-functions}). Its own
heading is $\alpha=-t\ln n_{\text{last}}$, and the reading we want is the
direction of the world origin $J_{0}=0$ as seen in the frame that lays this link
along the positive real axis: the change of frame of
\S\ref{sec:IT-origin}, applied at the far end of the chain, near $\zeta$,
instead. Two features of that end of the chain make the reading cheap. First, the
turning angle there is $-t\ln\frac{n}{n-1}\approx-t/n_{\text{last}}=-\pi$ by the
choice \eqref{eq:last-link}, so consecutive links are antiparallel and each added
link reverses the frame; only the axial direction modulo $180^{\circ}$ is
meaningful, and the pattern drifts from joint to joint at the negligible rate
$\nu=1/4T^{2}$ cycles of \eqref{eq:theta-joint}. Second, the chain out there
straddles its own limit. Taking $x=t/\pi$ in the approximate functional
equation~\cite[Thm.~4.11]{Titchmarsh1986} gives
$\zeta(s)=\sum_{n\leq x}n^{-s}+x^{1-s}/(s-1)+O(x^{-1/2})$, whose correction has
modulus $1/\sqrt{\pi t}$, so the far joints lie within $O(t^{-1/2})$ of $\zeta$;
what they in fact do is bracket it, each end of the last link sitting half a link
away, $|J-\zeta|=2.806\times10^{-2}$ against $|v|/2=2.804\times10^{-2}$ at
$t=10^{3}$. That places $\zeta$ at the midpoint of the last link, to
$4.8\times10^{-5}$ there and to $1.4\times10^{-7}$ at $t=10^{5}$. Read at a Gram
point, where $\zeta$ is real, that midpoint is therefore on the real axis, and the
bearing collapses to the closed form
\begin{equation}\label{eq:bearing}
\beta\;=\;180^{\circ}-\alpha\;=\;180^{\circ}+t\ln\lfloor t/\pi\rfloor
\pmod{180^{\circ}},
\end{equation}
a function of the height alone with no chain to sum, even though the chain in
question has $2.6\times10^{11}$ links at the top of the list.

Figure~\ref{fig:last-link-frame} draws the frame itself at six champions, one to
a decade of the height, each of them the champion of its decade sitting closest
to the mean axis. What the pictures show is that the chain does not spiral out to
$\zeta$ so much as march. At $t=6.73\times10^{8}$ its links spend
$2\sqrt{n_{\text{last}}}=2.9\times10^{4}$ of arc length to cover a displacement of
$124$, a ratio of $236$, and yet never stray more than $2.4$ from the straight
line joining its two ends: the arc length goes on loops the size of the links
themselves, inside a thin tube whose heading is \eqref{eq:bearing}. Only the last
stretch of the tube, where the logarithmic radius of the figure compresses least,
shows those loops as loops.

\begin{figure}[H]
\centering
\includegraphics[width=0.70\textwidth]{fig_last_link_frame}
\caption{The chain of $\zeta$ in the frame of \eqref{eq:bearing}, at the champion
of each decade of the height from $10^{3}$ to $10^{9}$ whose bearing lies closest
to the mean axis. The frame is anchored at the far end of the chain: the last link
lies along the positive real axis with its midpoint, which is $\zeta$, at the
origin, and each chain runs inward to it from the world origin $J_{0}=0$, marked
by a dot at radius $|\zeta|$ in the direction $\beta$. The last link is drawn far
longer than its true modulus $n_{\text{last}}^{-1/2}$, which is
$6.8\times10^{-5}$ at the outermost height against a picture $124$ wide. The
radius carries the same logarithmic compression as
Figure~\ref{fig:champion-bearings}: a point at distance $r$ is drawn at
$\log_{10}(1+r)$, which leaves every direction, and so the bearing, exact, while
letting six decades of height share one frame; the compression is gentlest near
$\zeta$, which is why the coils of the final approach read as coils while the
loops out near $J_{0}$ are flattened into the thread. The dotted rings are
$r=1,10,100$. Consecutive links out there being antiparallel, the frame is fixed
only up to a half turn, and the representative drawn is the one that puts $J_{0}$
in the upper right. Chains are sampled at a fixed step of arc length, which keeps
the $2.1\times10^{8}$ links of the outermost one plottable. Generated by
\texttt{fig\_last\_link\_frame.py}.}
\label{fig:last-link-frame}
\end{figure}

Read that at the Gram point nearest each champion and the bearings do not
scatter. Over the $148$ champions above $t=17.8$, dropping only the first at
$t=10.21$ where the chain is three links long, the mean axis is $44.8^{\circ}$
with Rayleigh concentration $R=\bigl|\text{mean }e^{2i\beta}\bigr|=0.967$ against
a noise floor of $1/\sqrt{148}=0.082$, and $146$ of the $148$ lie between
$25^{\circ}$ and $65^{\circ}$, the median distance from $45^{\circ}$ being
$4.3^{\circ}$. The same count of ordinary Gram points, drawn log-uniformly over
the same range of heights, gives $R=0.052$ with $31$ in that window, which is
what a uniform bearing would give. So \eqref{eq:bearing} is not manufacturing the
direction; something about a record pins the last link of its chain to one axis.
Written without any of the geometry, the observation is arithmetic:
\begin{equation}\label{eq:champ-arith}
\text{at a champion,}\qquad
t\,\ln\lfloor t/\pi\rfloor\;\equiv\;45^{\circ}\pmod{180^{\circ}} .
\end{equation}

What we did not expect is that the concentration tightens with height instead of
fading. Splitting the list by decade of $t$:

\begin{table}[H]
\centering
\small
\caption{The last-link bearing by decade of the height: the mean axis never
leaves $45^{\circ}$ while the spread halves. $R$ is the axial concentration and
the last column is the median distance from $45^{\circ}$. Computed in
\texttt{fig\_champion\_bearings.py}.}
\label{tab:champ-bearings}
\begin{tabular}{lcccc}
\toprule
decade of $t$ & champions & mean axis & $R$ & median $|\beta-45^{\circ}|$ \\
\midrule
$10^{1}$--$10^{2}$  & $6$  & $38.0^{\circ}$ & $0.852$ & $12.9^{\circ}$ \\
$10^{2}$--$10^{3}$  & $9$  & $45.7^{\circ}$ & $0.885$ & $8.5^{\circ}$ \\
$10^{3}$--$10^{4}$  & $16$ & $48.3^{\circ}$ & $0.949$ & $7.7^{\circ}$ \\
$10^{4}$--$10^{5}$  & $14$ & $42.6^{\circ}$ & $0.975$ & $5.7^{\circ}$ \\
$10^{5}$--$10^{6}$  & $10$ & $47.2^{\circ}$ & $0.982$ & $3.5^{\circ}$ \\
$10^{6}$--$10^{7}$  & $11$ & $45.8^{\circ}$ & $0.984$ & $5.5^{\circ}$ \\
$10^{7}$--$10^{8}$  & $16$ & $43.2^{\circ}$ & $0.985$ & $3.7^{\circ}$ \\
$10^{8}$--$10^{9}$  & $16$ & $47.0^{\circ}$ & $0.987$ & $3.7^{\circ}$ \\
$10^{9}$--$10^{10}$ & $13$ & $44.1^{\circ}$ & $0.989$ & $2.8^{\circ}$ \\
$10^{10}$--$10^{11}$& $16$ & $43.3^{\circ}$ & $0.987$ & $3.5^{\circ}$ \\
$10^{11}$--$10^{12}$& $21$ & $44.7^{\circ}$ & $0.993$ & $3.2^{\circ}$ \\
\bottomrule
\end{tabular}
\end{table}

The ragged ring is the innermost one, the six champions below $t=100$, whose
bearings run $25.6^{\circ}$, $56.7^{\circ}$, $12.8^{\circ}$, $54.0^{\circ}$,
$31.0^{\circ}$ and $46.7^{\circ}$: the construction needs a chain of some length
before it means anything, the first champion having three links.

\begin{figure}[H]
\centering
\includegraphics[width=0.85\textwidth]{fig_champion_bearings}
\caption{The last-link bearing \eqref{eq:bearing} at the $148$ champions above
$t=17.8$, read at the nearest Gram point. The bearing is axial, so each champion
is drawn at both ends of its axis; the radius is $\log_{10}t$, the colour the
height, the red line the circular mean and the shaded wedge $\pm20^{\circ}$ about
it. Open circles are the six champions below $t=100$. The grey ticks around the
rim are $148$ ordinary Gram points over the same range of heights, whose bearings
are uniform ($R=0.052$). The beam narrows outward: the tightest decade is the
outermost, $10^{11}\leq t<10^{12}$, at $R=0.993$. Generated by
\texttt{fig\_champion\_bearings.py}.}
\label{fig:champion-bearings}
\end{figure}

Two cautions on what this is. It is a description of where records sit, not a way
of finding them: \eqref{eq:champ-arith} says nothing about which of the many
heights satisfying it carries a large $|\zeta|$, and separating records from
merely large values is the business of the alignment of
Remark~\ref{rem:large-values}, not of a single phase. And the pinning itself is
unexplained. The last link is one term of modulus
$n_{\text{last}}^{-1/2}\approx1/\sqrt{2}\,T^{-1}$, far too small to matter to the
size of $\zeta$, so \eqref{eq:champ-arith} is a coincidence of phases and not a
mechanism; why a record height should fix the phase of that particular term to
one axis, at every height reachable, we cannot say.

\subsection{Links crossing}\label{sec:links-crossing}

\S\ref{sec:bisector} named one crossing of the two link chains: the forward
bisector link meets the reverse bisector link, and where they meet is the
bisector point $B_1$. Every other link of the forward chain is crossed by some
link of the reverse chain as well, and which one it is depends on $T$. Watching
the link at $m-1$ as $T$, for example, runs from $6$ to $7$, one reverse link
crosses it, and that link is joined by its neighbor that subsequently also
crosses over it, and by the end of the unit interval the crossing links have
advanced by two.

\begin{figure}[H]
\centering
\includegraphics[width=\textwidth]{fig_link_frames}
\caption{Local frames of the forward links at $\sigma=\tfrac12$, $T=6.18$
(top) and $T=6.72$ (bottom), the running example of
Figure~\ref{fig:spiral-summands} and a later height in the same unit
($t=I(T)\approx279.85$ and $327.01$, $m=6$). Each strip is in the local frame of each forward link, by pinning
one forward link to the unit interval along the $x$-axis;
the number at the bottom is that link. The reflected
inverse chain is drawn in the same frame: the blue links are the part
of it that belongs to the particular forward link, numbered at the two ends with the
outer (higher) link on the left, and the yellow
link is the single inverse link that crosses the forward link, numbered in yellow,
with a yellow dot at the crossing. The rightmost strip is the bisector link; the green and red
curves are the yin and yang loci of the two ends of the reverse
bisector link as $T$ runs through $(6,7)$. Generated by
\texttt{fig\_link\_frames.py}.}
\label{fig:link-frames}
\end{figure}

Numbering summands from one, so that link $k$
carries the summand $n=k+1$ (\S\ref{sec:matrix-product}), forward link $k$ and
reverse link $i$ cross when their two summands are reciprocal about the
Riemann--Siegel cutoff $a=\sqrt{t/2\pi}$ of \S\ref{sec:d1-approx-critical}.
Under $t=I(T)$ that cutoff is a function of the index alone, by
\eqref{eq:IT-derived}, so the rule closes in $T$ with nothing approximated:
\begin{equation}\label{eq:crossing-product}
n\,n'\;=\;a^{2}\;=\;\frac{t}{2\pi}\;=\;\frac{I(T)}{2\pi}
\;=\;\frac{2T+1}{2\log\!\bigl(1+\tfrac1T\bigr)},
\qquad n=k+1,\quad n'=i+1,
\end{equation}
that is,
\begin{equation}\label{eq:crossing-rule}
i\;=\;\Bigl[\frac{a^{2}}{k+1}\Bigr]-1 ,
\end{equation}
with $[\,\cdot\,]$ the nearest integer. Written out in $T$ alone, with nothing
standing in between, the \emph{crossing link} and \emph{crossing summand} are
the functions
\begin{equation}\label{eq:crossing-functions}
C_{\ell}(\ell,T)\;=\;\left[\frac{2T+1}
{2(\ell+1)\log\!\bigl(1+\tfrac1T\bigr)}\right]-1,
\qquad
C_{s}(n,T)\;=\;\left[\frac{2T+1}{2n\log\!\bigl(1+\tfrac1T\bigr)}\right],
\end{equation}
one rule in the two numberings, $C_{\ell}(\ell,T)=C_{s}(\ell+1,T)-1$, so that
\eqref{eq:crossing-rule} reads $i=C_{\ell}(k,T)$. Each is its own inverse but
for the rounding, $C_{s}(C_{s}(n,T),T)=n$, which is the symmetry of the pairing
established below. At $T=6.18$, where $a^{2}=44.5389$, that
reads $44,21,14,10,8,6,5$ for $k=0,\dots,6$, and all seven are real crossings
(Figure~\ref{fig:links-crossing}). The sum rule $k+i\approx2T-1$ that the links
beside the bisector obey is this hyperbola's tangent there, and holds only there.

Two remarks on $a$, since it is easy to read it as a link number. It is not one:
it is a real height, the self-dual summand $n=n'=a$ the law pairs about, the
one at which the forward chain hands over to the reflected one, and by
\eqref{eq:siegel-cutoff-index} it sits half a link past the bisector link, at $a^{2}=(T+\tfrac12)^{2}-\tfrac1{12}+O(T^{-2})$. Nor is it
a fixed number over a unit interval, as $\lfloor T\rfloor$ is: $a^{2}$ climbs
from $42.17$ to $56.17$ as $T$ runs from $6$ to $7$, and it is that climb that
walks each link's partner along inside the interval. Replacing $a^{2}$ by
$\lfloor T\rfloor^{2}=36$ or $(\lfloor T\rfloor+1)^{2}=49$ at $T=6.18$ would
name reverse link $35$ or $48$ as crossing the unit step, and neither does.

\paragraph{The reverse chain is the forward chain mirrored.} Number the joints
of both chains from zero as in \S\ref{sec:product-form},
\begin{equation}\label{eq:crossing-joints}
J_n=\sum_{n'\leq n}n'^{-s},
\qquad
K_i=\zeta-\chi(s)\sum_{n\leq i}n^{\,s-1},
\end{equation}
so that $J_0=0$, $K_0=\zeta$, and $K_m=\Sigma_1+R$ is the near end of the
reverse bisector link; link $n$ of either chain runs from its joint $n$ to its
joint $n+1$. On $\sigma=\tfrac12$ we have $n^{\,s-1}=\overline{n^{-s}}$, so
$K_i=\zeta-\chi\overline{J_i}$, and with $\chi=e^{-2i\vartheta}$ and
$Z=e^{i\vartheta}\zeta$ real (Hardy's $Z$, as in \eqref{eq:r-two-cases}),
\begin{equation}\label{eq:chain-mirror}
e^{i\vartheta}K_i \;=\; Z-\overline{e^{i\vartheta}J_i}\,.
\end{equation}
Write $e^{i\vartheta}J_n=X_n+iY_n$. Then \eqref{eq:chain-mirror} puts reverse
joint $i$ at $(Z-X_i,\,Y_i)$: in the $\vartheta$-rotated frame the reverse chain
is the forward chain reflected in the vertical line $X=Z/2$, joint for joint and
link for link. Call that line the \emph{mirror line}. It is not a new line: the
rotation puts $\zeta$ at $Z$ on the real axis, so unrotated it is the
perpendicular bisector of the base running from the origin to $\zeta$, which is
the bisector line of \S\ref{sec:legs}, drawn dotted in
Figures~\ref{fig:full-spirals} and~\ref{fig:legs} through $\zeta/2$ and the
bisector point $B_1$ (Corollary~\ref{cor:bisector-proj}). Doubling-back and mirror are
two different things in the same neighbourhood, and worth keeping apart: the
place along the chain is the bisector link, where the chain doubles
back; the mirror line is a line across the plane, the one the two chains
reflect in. The bisector point is where the two meet, the point inside the
bisector link at which the forward chain cuts the mirror line, as
\eqref{eq:crossing-rstar} works out below. This is the
chain-level form of the reflection
$B_2=\chi\overline{B_1}$ behind the equal legs of \S\ref{sec:legs}, and like it,
it is a statement about the critical line only: off the line $|\chi|\neq1$ and
the reverse chain is a scaled copy of the $1-\sigma$ chain instead of a mirror
image of this one.

\paragraph{Crossings come in pairs.} Parameterize each chain continuously, the
value $\mu=k+p$ with $0\leq p<1$ meaning the point $p$ of the way along link
$k$, and let $X(\mu),Y(\mu)$ be the rotated coordinates of the forward chain
there and
\begin{equation}\label{eq:crossing-h}
h(\mu)\;=\;X(\mu)-Z/2
\end{equation}
its signed distance to the mirror line. By \eqref{eq:chain-mirror} the reverse
chain at $\nu$ sits at $(Z-X(\nu),Y(\nu))$, so the forward chain at $\mu$ meets
the reverse chain at $\nu$ exactly when
\begin{equation}\label{eq:crossing-conditions}
h(\mu)+h(\nu)=0
\qquad\text{and}\qquad
Y(\mu)=Y(\nu).
\end{equation}
Both conditions are symmetric in $\mu$ and $\nu$, which is why crossings come in
pairs: reverse link $i$ crosses forward link $k$ if and only if reverse link $k$
crosses forward link $i$, and any rule naming the partner must be symmetric
too, as \eqref{eq:crossing-product} is. The pair coincides only when $\mu=\nu$,
and then $h=0$; so a self-paired crossing lies on the mirror line itself, there
is exactly one — the bisector point $B_1$, inside the bisector link of each
chain by Theorem~\ref{thm:d1-bounds} — and every other crossing sits off the
line with a mirror twin on the other side.

\paragraph{Sweeps between spiral centres.} Along either chain the turn from one
link to the next is $t\ln(1+1/n)\approx t/n$, decreasing as the chain is walked,
so the joints where that turn is an odd multiple of $\pi$ lie near the
spiral centres \eqref{eq:LN} --- the midpoint of the link they bound is closer
still --- and between two consecutive centres the chain makes
one full turn.
The middle of a turn is the link where the angle is the even multiple $2\pi c$,
\begin{equation}\label{eq:crossing-sweep}
\frac{t}{n'}=2\pi c
\qquad\Longleftrightarrow\qquad
n'=\frac{t}{2\pi c}=\frac{a^{2}}{c} ,
\end{equation}
and there successive links are parallel rather than antiparallel: their
contributions add in phase, the chain stops circling and runs. Call that run
the $c$-th \emph{sweep}; it carries the chain from the spiral centre $S_n=c$ to
the spiral centre $S_n=c-1$.

Poisson summation says what a sweep is worth. Writing the rotated chain as
$\sum_{n}n^{-1/2}e^{i(\vartheta-t\ln n)}
=\sum_{c}\int x^{-1/2}e^{i(\vartheta-t\ln x+2\pi cx)}\,dx$, the $c$-th integral
has its stationary point at $t/x=2\pi c$, which is \eqref{eq:crossing-sweep},
the middle of the $c$-th sweep; the stationary-phase value there is
\begin{equation}\label{eq:crossing-saddle}
\frac{1}{\sqrt c}\,
\exp i\Bigl(\vartheta+t-t\ln\frac{a^{2}}{c}+\frac\pi4\Bigr)
\;=\;\frac{1}{\sqrt c}\,e^{\,i\left(2\Theta-(\vartheta-t\ln c)\right)},
\qquad
2\Theta=2\vartheta+t-t\ln a^{2}+\frac\pi4 .
\end{equation}
Its modulus is that of the $c$-th summand, $c^{-1/2}$, and its phase is the
$c$-th summand's phase $\vartheta-t\ln c$ reflected in the fixed direction
$\Theta$. Sweep by sweep, the tail of a chain reproduces the head of that chain
reflected: this is the correspondence between joints and spiral centres
recorded in \S\ref{sec:bisector}, Nickel's centre-to-centre distances and
angles reproducing the initial steps, now carrying an index — the $c$-th sweep
answers to the $c$-th summand.

The mirror \eqref{eq:chain-mirror} is a reflection as well, and it undoes that
one. In the frame of the forward chain, therefore, the reverse chain's $c$-th
sweep is the forward chain's own link $c-1$ traversed backwards: over the
sampled heights the centre-to-centre step and that link agree to a median
$1.4\%$, $2.7\%$, $3.1\%$ and $5.2\%$ of a link length for $c=1,2,3,4$. A copy
of a link laid back along it must cross it, and the quarter turn in
\eqref{eq:crossing-saddle} is what tilts the copy so that it cuts across
instead of lying flat: the two meet near the middle of the link, at some
$43^{\circ}$. Which reverse link carries the $c$-th sweep is
\eqref{eq:crossing-sweep} with $c=n$, and that is
\eqref{eq:crossing-product}. Panel (b) of Figure~\ref{fig:links-crossing} is
the case $c=1$: the whole outermost turn of the reverse chain, links $30$ to
$89$, laid back along the unit step and crossing it at link $44$.

\paragraph{Near the bisector, a sum rule.} Near the self-dual link $n=n'=a$ the
hyperbola \eqref{eq:crossing-product} is indistinguishable from its tangent.
Writing $n=a+x$,
\begin{equation}\label{eq:crossing-tangent}
n'=\frac{a^{2}}{n}=a-x+\frac{x^{2}}{a}-\cdots ,
\qquad\text{so}\qquad
k+i\;=\;n+n'-2\;=\;2a-2+\frac{x^{2}}{a}-\cdots ,
\end{equation}
and since $2a=2T+1-\tfrac1{12T}+O(T^{-2})$ by
\eqref{eq:siegel-cutoff-index}, the sum is $2T-1$ up to the quadratic term and
a drift of order $1/12T$: $k+i$ is one of the two integers straddling $2T-1$,
that is
\begin{equation}\label{eq:crossing-sum}
i\;=\;\lfloor2T\rfloor-1-k
\qquad\text{or}\qquad
i\;=\;\lfloor2T\rfloor-k ,
\end{equation}
a rule needing nothing computed at all. The same rule comes out of the mirror
\eqref{eq:chain-mirror}, which is worth doing because it identifies the fixed
point. Near the bisector the
forward chain sweeps across the mirror line and $h$ passes through a single
zero there, close to straight over the few links either side; writing $r^{*}$
for that zero and linearizing turns the first condition of
\eqref{eq:crossing-conditions} into $\mu+\nu=2r^{*}$, the two arcs straddling
the crossing symmetrically, so $k+i$ straddles $S=2r^{*}-1$. And the crossing
fraction is not a new quantity.
By \eqref{eq:r-two-cases}, $h(m)=X_m-Z/2=-r/2$ with $r=e^{i\vartheta}R$ the real
remainder of \S\ref{sec:d1-approx-critical}, and link $m$ of the forward chain,
being $(m+1)^{-s}$ rotated by $e^{i\vartheta}$, has horizontal extent
$X_{m+1}-X_m=\cos(\omega-\vartheta)/\sqrt{m+1}$. So the zero falls inside link
$m$, at the crossing fraction
\begin{equation}\label{eq:crossing-rstar}
r^{*}=m+\lambda,
\qquad
\lambda=\frac{Z/2-X_m}{X_{m+1}-X_m}
=\sqrt{m+1}\;\frac{r}{2\cos(\omega-\vartheta)}
=\lceil T\rceil^{1/2}d_1=\hat d_1,
\end{equation}
by the exact form $d_1=e^{i\vartheta}R/2\cos(\omega-\vartheta)$ of
\S\ref{sec:d1-approx-critical} and the normalization
\eqref{eq:frac-vs-weight}. The place where the forward chain cuts the mirror
line is therefore the bisector point $B_1$ itself, sitting at the crossing
fraction $\hat d_1$ along the bisector link exactly as
Figure~\ref{fig:bisector-frame} marks it. So $S=2r^{*}-1=2m+2\hat d_1-1$, and
the two routes agree because $\hat d_1\approx\{T\}$ (\S\ref{sec:d1-function}):
the tangent to the hyperbola and the linearization of $h$ are the same line,
and \eqref{eq:crossing-sum} follows either way, using
$\tfrac15<\hat d_1<\tfrac45$ from Theorem~\ref{thm:d1-bounds} to put the
straddling integers at $\lfloor2T\rfloor-1$ and $\lfloor2T\rfloor$ whichever
side of $\tfrac12$ the fractional part is on.

How far the tangent can be trusted is \eqref{eq:crossing-tangent} again: it
costs half a link once $x^{2}/a=\tfrac12$, so the sum rule is good only within
\begin{equation}\label{eq:crossing-reach}
|n-a|\;\lesssim\;\sqrt{a/2}\;\approx\;\sqrt{T/2}
\end{equation}
of the bisector — three links either side at $T=17$, fourteen at $T=400$, and just
under two at the $T=6.18$ of the figure, which is why $k=4,5,6$ obey it there
and $k=3$ already does not. Inside that range the two indices trade off one for
one: each step back along the forward chain moves the partner one step forward
along the reverse chain, and over one unit of $T$ the sum moves by two, so each
partner advances by two, in two single steps, one inside the interval at
$\{T\}=\tfrac12$ and one at the handoff — the joining and letting go described
at the start. Outside the range the product takes over and the partners run
away as $a^{2}/n$, reaching the far end of the chain at the unit step.

The bisector link is the one place where the pairing is exact rather than
asymptotic. It is the self-dual link, $n=m+1$ and $x=\tfrac12-\{T\}$, so
\eqref{eq:crossing-rule} rounds to $m$ only over the middle half of the unit
interval and to $m\mp1$ at its ends; but the mirror does not round, and the
self-paired crossing $B_1$ of \eqref{eq:crossing-conditions} is there at every
height. The bisector links always cross, as \S\ref{sec:bisector} has it, and
the links \eqref{eq:crossing-rule} names near the ends of the interval are
second crossings of the same forward link, not replacements for it.

\begin{figure}[H]
\centering
\includegraphics[width=\textwidth]{fig_links_crossing}
\caption{Links crossing at $\sigma=\tfrac12$, $T=6.18$, the running example of
Figure~\ref{fig:spiral-summands} ($t=I(T)\approx279.85$, $m=6$,
$a=\sqrt{t/2\pi}=6.674$, $a^{2}=44.54$; the first turn of each chain holds
$90$ links).
(a) Both chains in the $\vartheta$-rotated frame, forward $e^{i\vartheta}J_n$ in
blue from $O$ and reverse $e^{i\vartheta}K_i$ in green from $Z=e^{i\vartheta}\zeta$,
which by \eqref{eq:chain-mirror} is the blue chain reflected in the dashed
line $X=Z/2$. Forward links $0$ through $6$ are drawn heavy and the red dots
are the seven crossings \eqref{eq:crossing-rule} names for them,
$i=44,21,14,10,8,6,5$; they run the length of the chain, not just the
neighbourhood of the bisector. (b) The boxed region enlarged, showing why. The gold
arc is one turn of the reverse chain, links $30$ to $89$, running from the
spiral centre $S_n=1$ to the spiral centre $S_n=0$ (crosses, at the links
\eqref{eq:LN} where the turn angle is an odd multiple of $\pi$). By
\eqref{eq:crossing-saddle} that step reproduces forward link $0$ reversed, so
it lies back along the unit step and cuts it near the middle (red dot); the
link carrying the sweep is $44$, the middle of the turn, where the turn angle
is $2\pi$ and $n'=a^{2}/1$. (c) The pairs $(n,n')$ of all seven crossings, on
the hyperbola $nn'=a^{2}$. The self-dual point $n=n'=a$ (star) is the summand
the law pairs about, half a link past the bisector link;
near it the hyperbola is indistinguishable from its tangent $n+n'=2a$ (dashed),
which is the sum rule \eqref{eq:crossing-sum}, and away from it the two part
company fast: the tangent gives the unit step a partner near $n'=11$ where the
truth is $45$. Generated by \texttt{fig\_links\_crossing.py}.}
\label{fig:links-crossing}
\end{figure}

\paragraph{How well it holds.} Sampling $m\in\{8,17,40,123,400\}$ at $199$
fractional parts each and testing every forward link $0\leq k\leq m$,
$118{,}007$ cases in all, the link named by \eqref{eq:crossing-rule} really does
cross forward link $k$ in $117{,}870$ of them, $99.9\%$; in each of the $137$
others a neighbour does, so the rule is never wrong by more than one link. The
crossing sits at the middle of the forward link, median position $0.500$ along
it with median deviation $0.009$ four or more links out from the bisector, widening
to $0.148$ at the bisector link, where the crossing is $B_1$ and the position is
$\hat d_1$ instead. The two links meet at a median $137^{\circ}$, tenth to
ninetieth percentile $135^{\circ}$ to $151^{\circ}$: some $43^{\circ}$ from
antiparallel, the quarter turn of \eqref{eq:crossing-saddle} that keeps the
sweep from lying flat along the link it reproduces. The centre-to-centre steps
match those links to a median $1.4\%$, $2.7\%$, $3.1\%$ and $5.2\%$ of a link
length at $c=1,2,3,4$, and the self-pair $(m,m)$ crossed at all $995$ sampled
heights while \eqref{eq:crossing-rule} named $m$ at exactly half of them, the
middle half of each unit interval, as the previous paragraph says. The sum rule
\eqref{eq:crossing-sum}, by contrast, names a real crossing in only $10.6\%$ of
the cases overall, but in $99.8\%$ of those inside its range
\eqref{eq:crossing-reach} and $5.1\%$ of those outside it, which is what that
range is worth. The identity $\lambda=\hat d_1$ of \eqref{eq:crossing-rstar}
holds to $6\times10^{-15}$. Computed in \texttt{check\_links\_crossing.py}.

\subsection{General yin and yang form for any link}\label{sec:yinyang-any}

\S\ref{sec:yinyang} watched one link cross another: the reverse bisector link
revolving about the forward bisector link, its two ends tracing the yin and
yang curves \eqref{eq:yin1}--\eqref{eq:yang1}. Now that \eqref{eq:crossing-rule}
says which reverse link crosses forward link $k$, every link of the chain has
such a pair, and they are all one formula. Like the bisector pair it carries no
$\zeta$: the remainder $R$, the factor $\chi$, and $\sigma$ and $T$ are all it
needs.

Pin forward link $k$ to $[0,1]$ exactly as \eqref{eq:frame-map} pins the
bisector link: translate its joint $J_k$ to the origin and divide by the link
vector $(k+1)^{-s}$, writing $s=\sigma+iI(T)$ throughout. Reverse joint $j$ of
\eqref{eq:crossing-joints} lands at
\begin{equation}\label{eq:frame-joint-any}
Y_k(j)\;=\;\bigl(K_j-J_k\bigr)(k+1)^{s}
\;=\;(k+1)^{s}\left[\;\sum_{n=k+1}^{m}n^{-s}\;+\;R\;+\;\chi\!\!\sum_{n=j+1}^{m}\!\!n^{\,s-1}\right],
\end{equation}
the second sum read as $-\sum_{n=m+1}^{j}$ once $j>m$. What removes $\zeta$ is
the same fact the bisector frame turned on: $K_m=\Sigma_1+R$, so measuring the
reverse chain from its own joint $m$ replaces the anchor $\zeta$ by the
remainder. The yin and yang points of link $k$ are then the two ends of the
link that crosses it, which are consecutive reverse joints:
\begin{equation}\label{eq:yin-any}
Y_{in}(k,T)=Y_k(i),
\qquad
Y_{ang}(k,T)=Y_k(i+1)=Y_{in}(k,T)-\chi\,(k+1)^{s}(i+1)^{\,s-1},
\qquad i=C_{\ell}(k,T),
\end{equation}
with $i=m$ at the bisector, where the mirror pins the pair whatever
\eqref{eq:crossing-rule} rounds to. At $k=m$ the first sum in
\eqref{eq:frame-joint-any} is empty and the second is the single term
$-\lceil T\rceil^{\,s-1}$, and \eqref{eq:yin-any} is
\eqref{eq:yin1}--\eqref{eq:yang1} back again.

Take $k=m-1$. Its frame factor $m^{s}$ meets a head of one summand, $m^{-s}$,
so the head is $1$ — forward joint $m$ is the far end of that link — and
\begin{equation}\label{eq:frame-joint-below}
Y_{m-1}(j)\;=\;1+m^{s}\left[R-\chi\!\!\sum_{n=m+1}^{j}\!\!n^{\,s-1}\right].
\end{equation}
Here the crossing changes hands, which is the one thing the bisector's does
not do. By \eqref{eq:crossing-product}, $a^{2}/m=m+2\{T\}+1+O(1/m)$, so the
crossing summand $n'=i+1$ runs $m+1\to m+2\to m+3$ across the unit and the
curves come in three arcs:
\begin{align}
n'=m+1:\quad
Y_{in}&=1+m^{s}R,
&Y_{ang}&=1+m^{s}\bigl[R-\chi(m{+}1)^{s-1}\bigr],
\label{eq:yin-below-a}\\
n'=m+2:\quad
Y_{in}&=1+m^{s}\bigl[R-\chi(m{+}1)^{s-1}\bigr],
&Y_{ang}&=Y_{in}-\chi\,m^{s}(m{+}2)^{s-1},
\label{eq:yin-below-b}\\
n'=m+3:\quad
Y_{in}&=1+m^{s}\bigl[R-\chi\bigl((m{+}1)^{s-1}+(m{+}2)^{s-1}\bigr)\bigr],
&Y_{ang}&=Y_{in}-\chi\,m^{s}(m{+}3)^{s-1},
\label{eq:yin-below-c}
\end{align}
the three running over $\{T\}$ in $[0,0.214]$, $[0.215,0.647]$ and
$[0.648,1)$ at $m=6$ (Figure~\ref{fig:yinyang-link5}).

\begin{figure}[H]
\centering
\includegraphics[width=\textwidth]{fig_yinyang_link5}
\caption{(a) Local frames of the forward links at $\sigma=\tfrac12$,
$T=6.18$, as in the top row of Figure~\ref{fig:link-frames}, now with the
yin and yang loci of every link as $T$ runs through $(6,7)$: green
$Y_{in}(k,T)$, red $Y_{ang}(k,T)$. (b) The yin and yang curves of
forward link $k=5$ over $6<T<7$, from
\eqref{eq:frame-joint-below}--\eqref{eq:yin-below-c}.
The frame pins that link to $[0,1]$ (dark blue). Each end traces three arcs,
one per crossing summand $n'=7,8,9$, kept apart at the handoffs: green is
$Y_{in}(5,T)$, red is $Y_{ang}(5,T)$. The dark-yellow segment is the crossing
link at the running example $T=6.18$, where $C_{\ell}(5,6.18)=6$ so $n'=7$.
Dashed lines carry strip $5$ of (a) down to the magnification (b).
Each arc's yin is the previous arc's yang, the same joint handing over.
Generated by \texttt{fig\_yinyang\_link5.py}.}
\label{fig:yinyang-link5}
\end{figure}

Each arc's yin is the arc before it's yang, the same point at the same $T$:
both are $Y_k(i+1)$, once as the far end of the outgoing crossing link and once
as the near end of the incoming one. That is the handoff of
\S\ref{sec:links-crossing} seen from the ends rather than the middle. The joint
passing across the forward link belongs to both links at that instant, so the
crossing itself goes through smoothly while the yin curve steps onto where the
yang curve was and the yang steps one joint further out.

How many arcs a link gets is the same count as the handoffs:
$n'=a^{2}/(k+1)$ advances by $2a/(k+1)=(2T+1)/(k+1)+O(T^{-1})$ over a unit of
the index, so link $k$ hands over about that many times — twice for the links
beside the bisector, some $2a$ times for the unit step, and not once for the
bisector, whose partner the mirror holds fixed. That is why the bisector's yin
and yang alone close into single loops, and why they alone have the limit
teardrop of \S\ref{sec:yinyang-infinity}.

The two ends fix where the crossing sits along the link, which is what the band
of graphs under the strips in the Zest links view draws. In the frame the link
is $[0,1]$ on the real axis, so the crossing is where the chord from yin to
yang cuts that axis, the computation \eqref{eq:crossing} makes at the bisector.
Writing $v_k=\chi\,(k+1)^{s}(i+1)^{\,s-1}$ for the crossing link's own vector in
the frame, which by \eqref{eq:yin-any} is the offset $Y_{in}-Y_{ang}$, that
crossing fraction is
\begin{equation}\label{eq:crossing-fraction-any}
\lambda_k(T)
\;=\;\operatorname{Re}Y_{in}-\operatorname{Im}Y_{in}\,
\frac{\operatorname{Re}v_k}{\operatorname{Im}v_k}
\;=\;\frac{\operatorname{Im}\bigl(\overline{Y_{in}}\,v_k\bigr)}{\operatorname{Im}v_k}
\;=\;\bigl|Y_{in}\bigr|\,\frac{\sin(\beta-\alpha)}{\sin\beta},
\qquad
\alpha=\arg Y_{in},\quad \beta=\arg v_k,
\end{equation}
which lies in $[0,1]$ exactly when the named link does cross the forward one;
the distance in the plane from joint $k$ is $\lambda_k\,(k+1)^{-\sigma}$. At
$k=m$ this is \eqref{eq:crossing} itself, $\lambda_m=\lceil T\rceil^{\sigma}d_1$,
the bisector point's own crossing fraction, which on the critical line is
$\hat d_1$. So the rightmost graph of the band is the quantity the rest of the
paper is built on, and the graphs to its left are that quantity's counterparts
link by link.

Checked against the frame construction of the Zest links view on the critical
line at $T=6.18,\,6.72,\,12.4,\,17.4,\,40.05$ and $434.37$:
\eqref{eq:frame-joint-any} matches the frame to $1.3\times10^{-15}$,
\eqref{eq:crossing-fraction-any} matches the band's own reading to
$6.9\times10^{-14}$, and $\lambda_m$ matches $\hat d_1$ to $1.0\times10^{-12}$,
in \texttt{probe-yinyang-below-bisector.mjs}.

\subsection{\texorpdfstring{$d_1$ and $d_2$}{d1 and d2} for any link}\label{sec:d1-any}

\S\ref{sec:decomp} split Siegel's remainder along the two bisector links,
and \S\ref{sec:derive-r1ps}--\S\ref{sec:derive-r2ps} found the weights
$d_1,d_2$ as the real-axis crossings of the yin and yang pair in those two
frames. The same two crossings exist for every link: \eqref{eq:yin-any} is
the pair, \eqref{eq:crossing-fraction-any} is the crossing fraction, and the
weights are that fraction read once along the forward link and once along the
reverse one.

Write $n=k+1$, $n'=i+1=C_{s}(n,T)$, and
\begin{equation}\label{eq:W-any}
W_k
\;=\;
\sum_{n=k+1}^{m}n^{-s}
\;+\;R
\;-\;\chi\sum_{n=m+1}^{i}n^{\,s-1}
\end{equation}
for the vector in the plane from forward joint $k$ to reverse joint $i$, so
that $Y_{in}(k,T)=(k+1)^{s}W_k$ by \eqref{eq:frame-joint-any}. At the
bisector both sums are empty and $W_m=R$. Let
\begin{equation}\label{eq:omega-any}
\omega_n=I(T)\log n,
\qquad
\omega_{n'}=I(T)\log n',
\qquad
\varphi_k=\arg W_k,
\qquad
\psi=\arg\chi,
\end{equation}
so $\omega_n$ is the absolute direction of summand $n$ of $\Sigma_1$, the
same telescoping that made $\omega=\omega_{m+1}$ in \eqref{eq:omega-telescope},
and $\omega_{n'}+\psi$ is the absolute direction of summand $n'$ of
$\Sigma_2$.

\paragraph{The forward weight.} In the frame of link $k$ the crossing fraction
\eqref{eq:crossing-fraction-any} is how far along that link it is from joint
$k$ to the crossing point. The frame scaled the link by $(k+1)^{\sigma}$, so the distance
in the plane is
\begin{equation}\label{eq:d1-any}
d_1(k,T)
\;=\;
(k+1)^{-\sigma}\lambda_k(T)
\;=\;
|W_k|\,
\frac{\sin\!\bigl(\omega_{n'}-\varphi_k+\psi\bigr)}
     {\sin\!\bigl(\omega_n+\omega_{n'}+\psi\bigr)}.
\end{equation}
The sine form is \eqref{eq:crossing-fraction-any} with
$Y_{in}=(k+1)^{s}W_k$ and $v_k=\chi\,(k+1)^{s}(i+1)^{s-1}$, the same
reduction that took \eqref{eq:crossing} to \eqref{eq:d1}:
$\alpha=\omega_n+\varphi_k$, $\beta=\omega_n+\omega_{n'}+\psi$, and
$|Y_{in}|=(k+1)^{\sigma}|W_k|$. At $k=m$ one has $W_m=R$,
$n=n'=\lceil T\rceil$, $\omega_n=\omega_{n'}=\omega$, and
\eqref{eq:d1-any} is \eqref{eq:d1}.

\paragraph{The reverse weight.} The same crossing, read along the reverse
link from joint $i$, is the parameter
$u=\Im Y_{in}/\Im v_k$ of the chord in that frame. The reverse link has
length $|\chi|\,(i+1)^{\sigma-1}$, so
\begin{equation}\label{eq:d2-any}
d_2(k,T)
\;=\;
|\chi|\,(i+1)^{\sigma-1}\,u
\;=\;
|W_k|\,
\frac{\sin\!\bigl(\omega_n+\varphi_k\bigr)}
     {\sin\!\bigl(\omega_n+\omega_{n'}+\psi\bigr)}.
\end{equation}
The two weights share a denominator, as $d_1$ and $d_2$ did. At $k=m$ this
is \eqref{eq:d2}. The fractional amounts, the hat of
\eqref{eq:frac-vs-weight}, are the two parameters themselves:
\begin{equation}\label{eq:hat-any}
\hat d_1(k,T)=\lambda_k(T)=(k+1)^{\sigma}\,d_1(k,T),
\qquad
\hat d_2(k,T)=u=\frac{d_2(k,T)}{|\chi|\,(i+1)^{\sigma-1}}.
\end{equation}
The yellow graphs under the strips of the Zest links view are
$\hat d_1(k,T)$ against $\{T\}$.

\paragraph{The direction, and the cone.} Forward link $k$ points along
$e^{-i\omega_n}$ and reverse link $i$ along $e^{i(\omega_{n'}+\psi)}$. The
two fractional summands are those directions, shortened by the two weights,
\begin{equation}\label{eq:R-any}
R_1(k,T)=d_1(k,T)\,e^{-i\omega_n},
\qquad
R_2(k,T)=d_2(k,T)\,e^{\,i(\omega_{n'}+\psi)},
\end{equation}
and they add back to the vector they split:
\begin{equation}\label{eq:W-cone}
W_k
\;=\;
d_1(k,T)\,e^{-i\omega_n}
\;+\;
d_2(k,T)\,e^{\,i(\omega_{n'}+\psi)}.
\end{equation}
This is \eqref{eq:R-as-cone} with $R$ replaced by $W_k$ and the two bisector
directions replaced by the two links that actually cross. The sine
expressions \eqref{eq:d1-any}--\eqref{eq:d2-any} are again Cramer's rule for
that $2\times2$ real system, which is why they are real, and why they reduce
to \eqref{eq:d1}--\eqref{eq:d2} at the bisector. On the critical line the
bisector pair still satisfies $d_1(m,T)=d_2(m,T)$, but a general pair need
not: the two links have different lengths, and the crossing sits at
$\hat d_1\approx\tfrac12$ four or more links out from the bisector
(\S\ref{sec:links-crossing}), so $d_1(k)$ and $d_2(k)$ scale as
$n^{-\sigma}$ and $|\chi|n'^{\sigma-1}$.
Figure~\ref{fig:Wk-triangle} is Figure~\ref{fig:remainder-zoom} drawn once
for every forward link at the running example: the last panel is that
figure, and the others are the same triangle with $R$ replaced by $W_k$.

\begin{figure}[H]
\centering
\includegraphics[width=\textwidth]{fig_Wk_triangle}
\caption{The cone \eqref{eq:W-cone} at $\sigma=\tfrac12$, $T=6.18$, one
panel per forward link $k=0,\dots,6$. In each zoom, red is
$R_1(k)$ along forward link $k$ (thin blue), orange is $R_2(k)$ along
reverse link $i=C_{\ell}(k,T)$ (thin green), and purple is the resultant
$W_k$ from joint $J_k$ to reverse joint $K_i$. The braces are the lengths
$d_1(k)$ and $d_2(k)$. The last panel is the bisector, $W_6=R$, which is
Figure~\ref{fig:remainder-zoom}. Generated by \texttt{fig\_Wk\_triangle.py}.}
\label{fig:Wk-triangle}
\end{figure}

At the running example $\sigma=\tfrac12$, $T=6.18$, the seven pairs are
\begin{center}
\small
\begin{tabular}{crrrr}
\toprule
$k$ & $i$ & $d_1(k)$ & $d_2(k)$ & $\hat d_1(k)$ \\
\midrule
$0$ & $44$ & $0.506$ & $0.00185$ & $0.506$ \\
$1$ & $21$ & $0.360$ & $0.159$ & $0.510$ \\
$2$ & $14$ & $0.286$ & $0.0914$ & $0.496$ \\
$3$ & $10$ & $0.263$ & $0.184$ & $0.525$ \\
$4$ & $8$ & $0.214$ & $0.143$ & $0.478$ \\
$5$ & $6$ & $0.264$ & $0.300$ & $0.647$ \\
$6$ & $6$ & $0.0875$ & $0.0875$ & $0.231$ \\
\bottomrule
\end{tabular}
\end{center}
The last row is Table~\ref{tab:frac} at this $T$: $d_1=d_2=0.0875$ and
$\hat d_1=\sqrt7\,d_1=0.231$. The first six sit near the middle of their
forward link, as \S\ref{sec:links-crossing} measured, and $d_2(0)$ is the
tiny reverse link $44$. Figure~\ref{fig:d1-any-link} is the same seven
numbers drawn as the Zest links tab draws them, against the whole unit of
$T$.

\begin{figure}[H]
\centering
\includegraphics[width=\textwidth]{fig_d1_any_link}
\caption{$\hat d_1(k,T)=(k+1)^{\sigma}d_1(k,T)$ for forward links
$k=0,\dots,6$ over $6<T<7$, $\sigma=\tfrac12$, one panel per link, $x$ the
fractional part of $T$ and $y$ the crossing fraction along that link from
the left joint --- the band under the strips in the Zest links view. The
light horizontal is $\tfrac12$; the light vertical and the dot are the
running example $T=6.18$, the seven values of the table. Away from the
bisector the track sits on $\tfrac12$ and barely moves; at $k=6$ it is the
normalized $d_1$ of Figure~\ref{fig:d1-critical}. Generated by
\texttt{fig\_d1\_any\_link.py}.}
\label{fig:d1-any-link}
\end{figure}

Checked at $\sigma=\tfrac12,\tfrac14,\tfrac7{10}$ and
$T=6.18,\,6.72,\,12.4,\,17.4$, for $k=0,1,m-2,m-1,m$:
\eqref{eq:d1-any}--\eqref{eq:d2-any} at $k=m$ match \eqref{eq:d1}--\eqref{eq:d2}
to $0$, \eqref{eq:hat-any} matches \eqref{eq:crossing-fraction-any} to
$1.5\times10^{-27}$, and \eqref{eq:W-cone} holds to $4.4\times10^{-28}$, in
\texttt{check\_d1\_any\_link.py}.

\subsection{The crossings: \texorpdfstring{$\Sigma_{1x}$ and $\Sigma_{2x}$}{Sigma\_1x and Sigma\_2x} formula}\label{sec:sum-x}

\S\ref{sec:links-crossing} named the reverse partner of each forward link, and
\S\ref{sec:d1-any} read the crossing fraction \(\hat d_1(k,T)\) along that
link. Those two data locate a point on every forward link \(0,\dots,m\).
Connecting the points in order is the same Dirichlet sum as
\(\Sigma_1+R_{1ps}\), re-jointed: the joints sit at the crossings instead of
at the integers, and the walk ends at the bisector point \(B_1\) rather than
at \(\Sigma_1\). Call that walk \(\Sigma_{1x}\).\footnote{The subscript \(x\) is for crossing, not a dummy index: \(\Sigma_{1x}\) is the crossing-joint form of \(\Sigma_1\), and \(\Sigma_{2x}\) of \(\Sigma_2\).}

Write \(p_k=\hat d_1(k,T)\) for the crossing fraction of link \(k\) from joint
\(k\) to the crossing point, the \(\lambda_k\) of \eqref{eq:hat-any}. Link \(k\) carries the
summand \(n=k+1\), so the crossing on it is
\begin{equation}\label{eq:sum-x-point}
P_k
\;=\;
J_k+p_k\,(k+1)^{-s}
\;=\;
\sum_{n'\le k}n'^{-s}+p_k\,(k+1)^{-s},
\end{equation}
with \(J_k\) as in \eqref{eq:crossing-joints}. The last crossing is the
bisector point itself, \(P_m=B_1\). The walk
\(0=P_{-1}\to P_0\to\cdots\to P_m\) therefore telescopes to \(B_1\) by
construction. Substituting \eqref{eq:sum-x-point} throughout, with the empty
sum \(\sum_{n'\le-1}=0\) and \(p_{-1}=0\), removes the intermediate points:
\begin{equation}\label{eq:sum-1x}
\Sigma_{1x}
\;=\;
\sum_{k=0}^{m}\Biggl(
\sum_{n'\le k}n'^{-s}+p_k\,(k+1)^{-s}
-\sum_{n'\le k-1}n'^{-s}-p_{k-1}\,k^{-s}
\Biggr)
\;=\;
\sum_{n'\le m}n'^{-s}+p_m\,(m+1)^{-s}.
\end{equation}
Each full summand \(k^{-s}\) is split across two consecutive steps,
\(p_{k-1}\,k^{-s}\) in one and \((1-p_{k-1})\,k^{-s}\) in the next, and those
two pieces add to one. What remains is the integer partial sum plus the last
fractional stub, which is the ordinary first-half formula
\eqref{eq:B1-sum} once the stub is named.

That stub is \(p_m(m+1)^{-s}\). The bisector link is link
\(m=\lfloor T\rfloor\). It runs from \(J_m=\Sigma_1\) to
\(J_{m+1}=\Sigma_1+(m+1)^{-s}\), so the vector along it is the next Dirichlet
term, \(n=\lceil T\rceil=m+1\). The factor \((m+1)^{-s}\) is that whole
unused summand, the faint dashed continuation past \(\Sigma_1\) in
Figure~\ref{fig:spiral-summands}. The number \(p_m\) is the crossing
fraction: how far along the link the two bisector links meet. That meeting is \(B_1\), and
\eqref{eq:crossing-rstar} already computed the crossing fraction:
\(p_m=\lambda=\hat d_1\). Therefore
\begin{equation}\label{eq:sum-1x-remainder}
p_m\,(m+1)^{-s}
\;=\;
\hat d_1\,(m+1)^{-s}
\;=\;
R_{1ps},
\end{equation}
the ``one more partial summand'' of \S\ref{sec:summand}: not a new term, the
next summand shortened by a real factor. The hat is what makes \(p_m\) a
fraction of a link rather than a length. The unhatted weight \(d_1\)
multiplies a unit vector in the direction of \((m+1)^{-s}\), so
\(|R_{1ps}|=d_1\); \(\hat d_1\) multiplies the full summand, whose length is
\(\lceil T\rceil^{-\sigma}\), and
\(\hat d_1\cdot\lceil T\rceil^{-\sigma}=d_1\) is the same remainder
\eqref{eq:frac-vs-weight}. At the running example \(\sigma=\tfrac12\),
\(T=6.18\), one has \(m=6\), next term \(7^{-s}\) of length
\(1/\sqrt7\approx0.378\), \(d_1\approx0.0875\), and
\(\hat d_1=\sqrt7\,d_1\approx0.231\): the walk uses about \(23\%\) of link
\(6\), and that stub is \(R_{1ps}\). The first six \(p_k\) in the table of
\S\ref{sec:d1-any} sit near \(\tfrac12\), as the saddle of
\S\ref{sec:links-crossing} said they would; only \(p_m\) is the named
remainder weight. So \eqref{eq:sum-1x} is
\begin{equation}\label{eq:sum-1x-B1}
\Sigma_{1x}
\;=\;
\sum_{n'\le m}n'^{-s}+p_m\,(m+1)^{-s}
\;=\;
\Sigma_1+R_{1ps}
\;=\;
B_1,
\end{equation}
and not \(\Sigma_1\) itself.

The same construction on the other Dirichlet sum gives \(\Sigma_{2x}\). Write
\(L_k=\chi(s)\sum_{n'\le k}n'^{\,s-1}\) for the joints of \(\Sigma_2\) from
the origin, and \(q_k=\hat d_2(k,T)\) for the crossing fraction of inverse
link \(k\) from \(L_k\) to its crossing point, the \(u\) of \eqref{eq:hat-any}. Then
\begin{equation}\label{eq:sum-2x-point}
Q_k
\;=\;
L_k+q_k\,\chi(s)\,(k+1)^{s-1},
\end{equation}
and the same telescoping, with \(q_{-1}=0\) and the empty sum at \(k=0\), is
\begin{equation}\label{eq:sum-2x}
\Sigma_{2x}
\;=\;
\sum_{k=0}^{m}\Biggl(
L_k+q_k\,\chi(s)\,(k+1)^{s-1}
-L_{k-1}-q_{k-1}\,\chi(s)\,k^{s-1}
\Biggr)
\;=\;
\chi(s)\sum_{n'\le m}n'^{\,s-1}+q_m\,\chi(s)\,(m+1)^{s-1}.
\end{equation}
The last stub is \(q_m=\hat d_2\), so
\begin{equation}\label{eq:sum-2x-B2}
\Sigma_{2x}
\;=\;
\Sigma_2+R_{2ps}
\;=\;
B_2.
\end{equation}
Nothing here used \(\sigma=\tfrac12\). The split
\(R=R_{1ps}+R_{2ps}\) of \S\ref{sec:decomp} holds at every \(\sigma\), the
weights \(\hat d_1,\hat d_2\) of \eqref{eq:frac-vs-weight} are defined at
every \(\sigma\), and so
\begin{equation}\label{eq:sum-x-zeta}
\Sigma_{1x}+\Sigma_{2x}
\;=\;
B_1+B_2
\;=\;
\zeta
\end{equation}
off the critical line as well as on it. On \(\sigma=\tfrac12\) the two walks
are the two legs of \S\ref{sec:legs}; off it they are just the two summands
of \eqref{eq:main}, re-jointed. Reflecting \(\Sigma_{1x}\) through
\(\zeta/2\) sends the origin to \(\zeta\) and fixes \(B_1\), which is the
other way of seeing the second half on the critical line, not the walk of
\(\Sigma_{2x}\) from the origin. The reverse partners of the forward links
\(0,\dots,m\) are not consecutive, so the points \(P_k\) are not the joints
of a walk along the reverse chain. This re-jointing is not Kapitonets'
iterated-midpoint polygon of \S\ref{sec:curlicues}: that construction is the
binomial transform of the partial sums and does not aim at \(B_1\) or
\(B_2\).

% ======================================================================
\section{Prior literature}\label{sec:prior-lit}
% ======================================================================

\subsection{Nickel's Argand-diagram geometry}\label{sec:nickel}

Nickel~\cite{Nickel2013}, with a condensed sequel~\cite{Nickel2015}, starts
from the same object this paper starts from,
the Argand diagram of the steps $n^{-s}$, and reaches a construction close
enough to ours that the two are worth setting side by side. They agree in
shape, and they differ in the one respect this paper is about.

\paragraph{Dictionary.} His center of symmetry is the step
$n_p=[\sqrt{t/2\pi}\,]$, with fractional part $p=\sqrt{t/2\pi}-n_p$. Since
$\sqrt{t/2\pi}\approx T+\tfrac12$ (\S\ref{sec:IT-origin}), his index is
$n_p\approx\lfloor T+\tfrac12\rfloor$ and $p\approx\{T+\tfrac12\}$, so it
agrees with our $m=\lfloor T\rfloor$ on the first half of each unit interval
and exceeds it by one on the second half: the same half-unit offset already
recorded for Siegel's cutoff in \S\ref{sec:kuznetsov}. His factor $Q(s)$,
defined by $\zeta(s)=Q(s)\zeta(1-s)$, is our $\chi(s)$; the displayed form
$Q=n_p^{1-2s}e^{i(t+\pi/4)}$ is the familiar $e^{-2i\vartheta(t)}$ asymptotic,
which reproduces $\chi$ to a relative $1.5\times10^{-5}$ at $T=20.3$ and
$1.6\times10^{-7}$ at $T=200.3$ once the smooth $\sqrt{t/2\pi}$ is used in
place of the rounded $n_p$. His symmetry angle $\Theta$ is our
$\tfrac{\psi}{2}=\tfrac12\arg\chi$.

\paragraph{What matches.} Under that dictionary his Riemann--Siegel equation
$\zeta=P(s)+Q(s)P(1-s)$ is our two-leg decomposition $\zeta=B_1+B_2$ with
$B_2=\chi\,\overline{B_1}$ (Appendix~\ref{app:lean-critical}); his polar form
$\zeta=2\cos(\phi-\Theta)\,\mathsf{P}e^{i\Theta}$ is the shape of our
$R=2d_1\cos(\omega+\tfrac{\psi}{2})e^{i\psi/2}$; and his observation that the
two magnitudes are equal exactly when $\sigma=\tfrac12$ is
Corollary~\ref{cor:equal-legs}. He also writes the classical remainder as
$R=2\mathsf{L}\cos(\theta_L-\Theta)$, where $L=\mathsf{L}e^{i\theta_L}$ runs
from the joint at step $n_p$ to his center: that is exactly the role played
here by the pair $(d_1,-\omega)$. Where his index agrees with ours the two
projected remainders converge, the ratio
$2\mathsf{L}\cos(\theta_L-\Theta)\,/\,2d_1\cos(\omega+\tfrac{\psi}{2})$
falling through $2.24$, $1.34$, $1.11$, $1.03$, $1.01$ at
$T=6.3,\,20.3,\,60.3,\,200.3,\,600.3$.

\paragraph{Two centers, one line.} His $P(s)$ and our $B_1$ are not the same
point, and they are not meant to be. The difference lies almost entirely
along the direction in which $\zeta$ cannot see it: at $T=600.3$ we find
$|P-B_1|=0.064$, of which only $1.2\times10^{-4}$ lies along $e^{i\psi/2}$.
That is his own remark that $\zeta$ is unchanged when $P(s)$ is translated
along the symmetry axis. In the notation of \S\ref{sec:legs} the
admissible centers are the solutions $X$ of $\zeta-X=\chi\overline{X}$, a
line in the plane, and every one of them projects onto $\zeta/2$, which is
Corollary~\ref{cor:bisector-proj}. The freedom he notes and the projection we
prove are the same fact.

\paragraph{The sequel.} Two things in \cite{Nickel2015} touch this paper
further. It traces $P(s)$ and $\zeta(s)$ as two curves over a range of $t$,
joined at regularly spaced values of $t$ to show the two segments, which is the
figure the conveyor belt draws in Fig.~\ref{fig:conveyor-follower} with his
pendant center in place of $B_1$. And it counts zeros: it reads the count off
$\Theta$ through the Gram points and tests that reading against the first
$100{,}000$ ordinates, finding one zero per Gram point on average with the zeros
grouping between successive Gram points. The count there is the classical one,
exact in the mean and never landing on an ordinate; the curves of
\S\ref{sec:counting} land on every one.

\paragraph{What differs.} The correction itself. Nickel's is a lowest-order
determination, carrying a factor $1/(2n_p^{\sigma}\cos 2\pi p)$ that grows
without bound as $\cos2\pi p\to0$, and he is explicit that its error matters
for the distribution of $P(s)$ even though it cancels in $\zeta$; its
direction, $e^{-i(t\ln n_p+2\pi p)}$, is the step direction at $n_p$ turned
by the fractional offset $2\pi p$, not a summand direction of the chain. Our
$R_{1ps}$ and $R_{2ps}$ are exact, and their directions are exactly those of
the $(m+1)$st summand of each chain. That is what lets the serial chain of
\S\ref{sec:product-form} close on $\zeta$ identically rather than
asymptotically, and it is the content of the ``one more partial summand''
reading: the correction is not merely near the next link, it lies along it.

% ----------------------------------------------------------------------
\subsection{Spirals of exponential sums}\label{sec:curlicues}

For quadratic phases this geometry is rigorous ground. The partial sums
$S_N(\tau)=\sum_{n\le N}e^{i\pi\tau n^{2}}$ draw what Berry and
Goldberg~\cite{BerryGoldberg1988} call curlicues: replacing the sum by an
integral produces the Cornu spiral exactly, and a renormalization then carries
the whole pattern onto a magnified and rotated copy of a shorter sum of the same
kind. Coutsias and Kazarinoff~\cite{CoutsiasKazarinoff1987,CoutsiasKazarinoff1998}
work that picture the way this paper works ours. They give the polyline a
discrete radius of curvature
\begin{equation}\label{eq:ck-curvature}
R_N=\tfrac12\bigl|\csc(\psi_N/2)\bigr|,
\qquad \psi_N=\pi\tau(2N+1),
\end{equation}
$\psi_N$ being the turn between consecutive terms, so that inflections of the
pattern sit where $\psi_N$ is nearest a multiple of $2\pi$ and cusps where it
turns around. They then cut the sum \emph{at an inflection}, and replace each
whole spiral of the first sum by a single vector carrying a scale and a phase,
the Fresnel length $\sqrt{|\tau|}$ at $\pi/4$ to the spiral's mid-vector. What
comes out is Hardy and Littlewood's approximate functional formula for the theta
function, sharpened to an error uniform in $\tau$ and stated as a Diophantine
inequality that contains Gauss's sum formula. Cutting at the bisector and standing
one link in for a whole spiral is the reading of \S\ref{sec:decomp}; their
$e^{i\pi/4}$ is the Gaussian integral's, the same factor that leaves the
$-\pi/8$ in $\vartheta$; and \eqref{eq:ck-curvature} is the bisector construction of
\S\ref{sec:bisector} for unit links, the distance from a step to the meeting of
the angle bisectors at its two ends, which is Nickel's $L$ once the length
$n^{-\sigma}$ is restored.

The arithmetic is what does not carry over. Their turn $\psi_N$ is linear in
$N$, so the second difference is the constant $2\pi\tau$, the arcs are Cornu
spirals outright, and the theory turns on the continued fraction of $\tau$:
the pattern is a self-similar hierarchy of curlicues within curlicues. Our turn
is $\theta_n=-t\log\bigl(1+\tfrac1n\bigr)$ and its second difference falls like
$t/n^{2}$, passing $2\pi$ once. That is why the chain has a single distinguished
inflection, at $n\approx\sqrt{t/2\pi}$, instead of a hierarchy of them, and why its
arcs are Cornu spirals only locally. The geometric machinery transfers; the
renormalization does not.

Two further readings of the same diagram deserve mention. Lander's paired
series~\cite{Lander2017,Lander2018} reduce each pseudo-spiral to a principal-axis
vector, and those axis vectors form a second chain running counter to the first,
with matched arguments and with magnitudes that separate as $\sigma$ leaves
$\tfrac12$: the counter-running chain of \S\ref{sec:yinyang} and the equal
lengths of Corollary~\ref{cor:equal-legs} in embryo.
Kapitonets~\cite{Kapitonets2019} draws his axis of symmetry through the
\emph{midpoint} of the remainder vector at $\sigma=\tfrac12$, observes that the
remainder is then perpendicular to that axis, and concludes both that
$\arg\zeta$ is the normal direction and that the two chains project onto the axis
with cancelling sums. That is the half-split of
\S\ref{sec:remainders-summary} together with Corollary~\ref{cor:bisector-proj},
read off the picture. His other construction,
recovering $\zeta$ by repeatedly replacing a polygon of partial sums by the
polygon of their midpoints, is unrelated to the bisector point: iterated
midpointing is the binomial transform, and the limit is Ces\`aro summation of the
Dirichlet series, as he says himself.

% ======================================================================
\section{Glossary}\label{sec:glossary}
% ======================================================================

The paper uses a geometric language for objects that other writers have
named differently, or not named at all. This section collects those words
in one place. Each entry is a definition in ordinary language, a short
formula when the formula is short, a pointer to the section that
introduces the object, and, where we know one, the name another author
has used for the same thing or a neighboring one. The names are those
already used above; this list is only a map.

The list is not alphabetical. It runs from the parameters and the
polyline itself to the more specialized constructions built on them,
and the entries in each group belong to one another.

\subsection*{The parameters}

\begin{description}
\item[Index $T$, and $I(T)$.]
The real parameter that replaces Siegel's pair $(t,m)$:
$t=I(T)=\pi(2T+1)/\log(1+1/T)$ and $m=\lfloor T\rfloor$
(\S\ref{sec:IT}).
$T$ is called the index even though it is not an integer.
The map is chosen so that integer $T$ lands on a handoff
(\S\ref{sec:IT-origin}).

\item[$\chi(s)$.]
The completion factor of the functional equation,
$\zeta(s)=\chi(s)\zeta(1-s)$
(\S\ref{sec:siegel}).
Nickel writes $Q(s)$ for the same factor.

\item[$\omega$.]
The absolute direction of the next forward summand:
$\omega=I(T)\log\lceil T\rceil$, so that $(m+1)^{-s}$ points along
$e^{-i\omega}$
(\S\ref{sec:decomp}).

\item[$a$, $a^{2}$.]
Siegel's cutoff $a=\sqrt{t/2\pi}$ and its square
$a^{2}=t/2\pi=I(T)/2\pi$.
Crossing partners multiply to $a^{2}$.
Under $t=I(T)$ one has $a^{2}\approx(T+\tfrac12)^{2}$.
\end{description}

\subsection*{The chain: joints and links}

\begin{description}
\item[Sum chain, or link chain.]
The partial-sum polyline in the plane: joints connected by links,
one Dirichlet summand per link
(\S\ref{sec:matrix-product}).
``Sum chain'' reads it as a running total; ``link chain'' reads it
as a serial chain of segments. The two names are the same object.
We also call it a spiral, though only the later links wind.
Nickel~\cite{Nickel2013} draws it as the Argand diagram of the
steps $n^{-s}$.
Erickson~\cite{Erickson2005} identifies the later coils as Cornu
spirals and labels their centers $C_n$.
Kapitonets~\cite{Kapitonets2019} calls the polyline the Riemann
spiral.
Berry and Goldberg~\cite{BerryGoldberg1988} call the quadratic-phase
analogue a curlicue; our turn $\theta_n=-t\log(1+1/n)$ is not
quadratic, so the coils are Cornu only locally
(\S\ref{sec:curlicues}).
The $3\times3$ matrices of \S\ref{sec:product-form} are the planar
Denavit--Hartenberg encoding of a serial chain
\cite{DenavitHartenberg1955}.

\item[Joint.]
A vertex of a chain: the value of the partial sum after an integer
number of summands.
Joint $n$ of the forward chain is $J_n=\sum_{n'\le n}n'^{-s}$, so
joint $0$ is the origin and joint $m$ is $\Sigma_1$ itself
(\S\ref{sec:matrix-product}, \S\ref{sec:links-crossing}).
Erickson's centers $C_n$ stand in one-to-one correspondence with the
joints, but a center is the eye of a coil, not the vertex
(\S\ref{sec:bisector}).

\item[Link.]
The segment between two consecutive joints: one Dirichlet summand,
drawn as a vector.
Link $k$ of the forward chain runs from joint $k$ to joint $k+1$ and
carries the summand $n=k+1$
(\S\ref{sec:matrix-product}).
Nickel calls the same object a step.
Coutsias and Kazarinoff~\cite{CoutsiasKazarinoff1987} treat the same
polyline as a discrete curve with a radius of curvature at each turn.

\item[$\Sigma_1$, the forward chain (or forward spiral).]
The sum chain of the first Dirichlet sum,
$\Sigma_1=\sum_{n=1}^{m}n^{-s}$, drawn from the origin
(\S\ref{sec:siegel}, \S\ref{sec:matrix-product}).

\item[$\Sigma_2$, the inverse chain (or inverse spiral).]
The sum chain of the reflected Dirichlet sum,
$\Sigma_2=\chi(s)\sum_{n=1}^{m}n^{s-1}$, drawn from the origin
(\S\ref{sec:siegel}).
On the critical line it is the functional-equation image of the
forward chain: each summand of $\Sigma_2$ is $\chi$ times the
conjugate of the matching summand of $\Sigma_1$.
Lander~\cite{Lander2017,Lander2018} reduces each coil to a
principal-axis vector and obtains a second chain running counter to
the first, which is this chain in embryo.
Nickel's second half $Q(s)P(1-s)$ is the same reflection, with his
$Q$ equal to our $\chi$ (\S\ref{sec:nickel}).
\end{description}

\subsection*{The remainder on the chain}

\begin{description}
\item[$R$.]
Siegel's unsplit remainder integral, so that
$\zeta=\Sigma_1+\Sigma_2+R$ (\S\ref{sec:siegel}).
Every named split below recovers this same $R$, exactly or
approximately.

\item[$R_{ps}$, $R_{1ps}$, $R_{2ps}$.]
The partial-summand split of $R$:
$R=R_{ps}=R_{1ps}+R_{2ps}$, with
$R_{1ps}=d_1\,e^{-i\omega}$ and
$R_{2ps}=d_2\,e^{i(\omega+\arg\chi)}$
(\S\ref{sec:decomp}).
Each piece is the next Dirichlet term of its sum, shortened by a
real weight.
The subscript $ps$ distinguishes this exact split from the
approximate remainders of other authors.

\item[Fractional summand.]
A remainder that is exactly the next Dirichlet term, scaled by a
real factor: $R_{1ps}=\hat d_1\,(m+1)^{-s}$ and
$R_{2ps}=\hat d_2\,\chi\,(m+1)^{s-1}$
(\S\ref{sec:summand}).
The word ``fractional'' means a fraction of one summand, not a
summand with a non-integer index.

\item[$d_1$, $d_2$.]
The real lengths of the two fractional summands:
$|R_{1ps}|=d_1$ and $|R_{2ps}|=d_2$
(\S\ref{sec:decomp}).
Each is the distance, in the plane, from the bisector joint to the
bisector point, measured along its own bisector link.
On the critical line, $d_1=d_2$.
The sine formulas are \eqref{eq:d1}--\eqref{eq:d2}.

\item[$\hat d_1$, $\hat d_2$ (fractional amounts).]
The same weights read as fractions of a link rather than as
lengths:
$\hat d_1=\lceil T\rceil^{\sigma}\,d_1$ and
$\hat d_2=\lceil T\rceil^{1-\sigma}\,d_2/|\chi|$
(\S\ref{sec:summand}).
$\hat d_1$ is the crossing fraction of the bisector link.
At the running example $\sigma=\tfrac12$, $T=6.18$, one has
$d_1\approx0.0875$ and $\hat d_1\approx0.231$.

\item[$R_{ak}$, $R_{1ak}$, $R_{2ak}$.]
The Kuznetsov / Siegel-$f$ split:
$R\approx R_{ak}=R_{1ak}+R_{2ak}$, where
$R_{1ak}=f(s)-\Sigma_1$
(\S\ref{sec:kuznetsov}, \S\ref{sec:remainders-summary}).
Approximate, and not a single fractional summand.

\item[$R_{rs}$, $R_{1rs}$, $R_{2rs}$.]
The half-split $R_{1rs}=R_{2rs}=\tfrac12 R$
(\S\ref{sec:remainders-summary}).
Kapitonets' midpoint of the remainder vector is this split
(\S\ref{sec:curlicues}).
\end{description}

\subsection*{The bisector}

\begin{description}
\item[Bisector link.]
Link $m=\lfloor T\rfloor$ of a chain: the link that sits at the
boundary between the nearly-straight early chain and the coils
(\S\ref{sec:bisector}).
Each chain has one.
The forward bisector link runs from $\Sigma_1$ along the
$(m+1)$st summand; the reverse bisector link runs from $\Sigma_1+R$
along its $(m+1)$st summand.
This is the neighborhood of Siegel's cutoff
$m=\lfloor\sqrt{t/2\pi}\rfloor$ (\S\ref{sec:siegel}), of Nickel's
center-of-symmetry step $n_p=[\sqrt{t/2\pi}\,]$ (offset from our $m$
by about half a unit, \S\ref{sec:nickel}), and of the inflection at
which Coutsias and Kazarinoff cut their sum
(\S\ref{sec:curlicues}).

\item[Bisector joint.]
Joint $m$ of the forward chain, the integer vertex at which the
bisector link begins: the point $\Sigma_1$.
The matching joint of the reverse chain is $\Sigma_1+R$.
All three remainder splits of \S\ref{sec:remainders-summary} leave
from this pair of joints.

\item[Bisector point.]
The crossing of the two bisector links,
$B_1=\Sigma_1+R_{1ps}=\Sigma_1+R-R_{2ps}$
(\S\ref{sec:bisector}).
On the critical line it is the apex of the isosceles triangle with
base from the origin to $\zeta$.
Nickel's center $P(s)$ plays the same role and is not the same
point: the two differ mostly along the symmetry axis, which $\zeta$
cannot see (\S\ref{sec:nickel}).
Kapitonets draws his axis through the midpoint of the remainder
vector $R$; that midpoint is the half-split $R_{rs}$, not $B_1$
(\S\ref{sec:curlicues}).

\item[$B_1$, $B_2$.]
The two half-sums that add to $\zeta$:
$B_1=\Sigma_1+R_{1ps}$ and $B_2=\Sigma_2+R_{2ps}$, so
$\zeta=B_1+B_2$ (\S\ref{sec:matrix-product}).
$B_1$ is the bisector point.
On the critical line, $B_2=\chi\,\overline{B_1}$.
Nickel's Riemann--Siegel equation $\zeta=P(s)+Q(s)P(1-s)$ is this
two-leg split with a different choice of the meeting point
(\S\ref{sec:nickel}).
\end{description}

\subsection*{Frames}

\begin{description}
\item[Link frame.]
The coordinate frame of an arbitrary forward link $k$: translate
joint $k$ to the origin and divide by the link vector $(k+1)^{-s}$,
so that link $k$ becomes the unit interval $[0,1]$
(\S\ref{sec:yinyang-any}, Figure~\ref{fig:link-frames}).

\item[Bisector frame.]
The link frame at $k=m$: it pins the forward bisector link to
$[0,1]$ on the real axis
(\S\ref{sec:IT-origin}, \S\ref{sec:yinyang}).
In that frame the reverse bisector link revolves as $T$ advances,
and the crossing with the axis is the number $d_1$
(or the crossing fraction $\hat d_1$, once the link is read at unit
length).
\end{description}

\subsection*{Legs}

\begin{description}
\item[Legs.]
The two segments that add to $\zeta$: Leg~1 is the vector $B_1$
from the origin to the bisector point, and Leg~2 is the vector
$B_2=\zeta-B_1$ from the bisector point to $\zeta$
(\S\ref{sec:legs}).
Their lengths are $L_1=|B_1|$ and $L_2=|B_2|$.
Nickel's two vectors $P(s)$ and $Q(s)P(1-s)$ are legs in the same
sense, meeting at his center rather than at $B_1$.

\item[$\vartheta_1$, $\vartheta_2$.]
The two leg angles:
$\vartheta_1=\arg B_1$ is the heading of Leg~1 from the positive
real axis, and
$\vartheta_2=\arg(B_2/B_1)$ is the turn from Leg~1 to Leg~2
(\S\ref{sec:legs}).
Then $\zeta=L_1 e^{i\vartheta_1}+L_2 e^{i(\vartheta_1+\vartheta_2)}$.
These are not the Riemann--Siegel theta function $\vartheta(t)$ of
\eqref{eq:N-theta}, which is the phase of $\chi^{1/2}$ and counts
zeros; the clash of names is unfortunate and the two are kept
apart by the subscripts.

\item[Equal legs.]
The condition $L_1=L_2$.
On the critical line it holds identically
(Corollary~\ref{cor:equal-legs}).
Off the line it is a locus of ovals and thin bands in the strip
(Figure~\ref{fig:equal-legs-strips}).
A zero of $\zeta$ requires equal legs \emph{and} $\vartheta_2=\pi$.
Nickel notes that the two magnitudes are equal exactly when
$\sigma=\tfrac12$ (\S\ref{sec:nickel}).

\item[Fold, folded.]
Two related uses, both meaning that two segments occupy the same
line and point opposite ways.
At an integer $T$ the neighboring link \emph{folds back} onto the
bisector link, the joint angle is an odd multiple of $\pi$, and
the bisector point hands off to the next link
(\S\ref{sec:IT-origin}, \S\ref{sec:yinyang}): that is a
\emph{handoff}, and for zeta every handoff is of this kind.
Separately, a \emph{folded-leg} point in the strip is a height
where $\vartheta_2=\pi$, so that Leg~2 folds back onto Leg~1
(Figure~\ref{fig:equal-legs-strips}); a zero is an equal-leg point
that is also folded-leg.
\end{description}

\subsection*{How the picture moves with $T$}

\begin{description}
\item[Handoff.]
The instant, at integer $T$, when the bisector point moves from
link $m-1$ onto link $m$
(\S\ref{sec:d1-function}, \S\ref{sec:IT-origin}).
For zeta the two links fold back onto one another and the point
slides across the joint with no jump.

\item[Event.]
Any instant at which the bisector point moves from one link to the
next (\S\ref{sec:IT-Lfunctions}).
For zeta every event is a handoff by fold-back.
Dirichlet $L$-functions have a second kind.

\item[Folded event.]
An event at which the turning angle at the current joint is
$\pm\pi$, so the chain reverses and the bisector point slides
across the joint (\S\ref{sec:IT-Lfunctions}).
For zeta, every handoff is a folded event, once per unit of $T$.
For $L(s,\chi_p)$ these occur at $T\equiv0\pmod{p-1}$.

\item[Bisector event.]
An event at which the bisector \emph{line} (the perpendicular
bisector of the segment from $0$ to the $L$-value) sweeps through
a joint and carries the crossing onto the neighboring link
(\S\ref{sec:IT-Lfunctions}).
Zeta has none.
For $L(s,\chi_p)$ they occupy the remaining residues modulo
$p-1$.

\item[Phantom link.]
A zero-length summand of $L(s,\chi_p)$, where $\chi_p(n)=0$
(\S\ref{sec:IT-Lfunctions}).
Phantoms receive no index $T$; only the nonzero steps are links.
\end{description}

\subsection*{Crossings along the chain}

\begin{description}
\item[Crossing (the point).]
Where a forward link and an inverse link meet in the plane.
The distinguished one is the bisector point $B_1$; every other
forward link $k$ has one as well, written $P_k$
(\S\ref{sec:links-crossing}, \S\ref{sec:sum-x}).

\item[Crossing partner.]
The inverse link that crosses forward link $k$.
The two summands are reciprocal about the Riemann--Siegel cutoff
$a=\sqrt{t/2\pi}$:
$n\,n'=a^{2}=I(T)/2\pi$, with $n=k+1$
(\S\ref{sec:links-crossing}).
This is the same saddle that Siegel's contour is built around.

\item[Crossing fraction.]
How far along a link the crossing sits, from the left joint toward
the right, as a number in $[0,1]$
(\S\ref{sec:d1-any}, \S\ref{sec:sum-x}).
On forward link $k$ it is $p_k=\hat d_1(k,T)=\lambda_k$; on inverse
link $k$ it is $q_k=\hat d_2(k,T)$.
At the bisector, $p_m=\hat d_1$.
Away from the fold the other $p_k$ sit near $\tfrac12$.
This is a fraction of the straight link, not of a curved arc.

\item[$\Sigma_{1x}$, $\Sigma_{2x}$.]
The two Dirichlet sums re-jointed at the crossings instead of at
the integers
(\S\ref{sec:sum-x}).
$\Sigma_{1x}$ walks $P_{-1}=0,P_0,\dots,P_m=B_1$ and totals
$B_1=\Sigma_1+R_{1ps}$;
$\Sigma_{2x}$ does the same from the origin along $\Sigma_2$ and
totals $B_2$.
The subscript $x$ is for crossing, not a dummy index.
The step lengths are the crossing fractions $p_k$ and $q_k$.
\end{description}

\subsection*{Yin, yang, and the needle}

\begin{description}
\item[Yin, yang.]
The two ends of the reverse bisector link, watched in the
bisector frame as $T$ advances
(\S\ref{sec:yinyang}).
Each traces a teardrop; the pair resembles a yin-yang.
Yin is the joint-$m$ end, yang the joint-$(m+1)$ end.
The chord from yin to yang crosses the real axis at the bisector
point, and that crossing is how $d_1$ and $d_2$ were found
(\S\ref{sec:derive-r1ps}--\S\ref{sec:derive-r2ps}).
The same pair exists for every forward link $k$, once the crossing
partner is known (\S\ref{sec:yinyang-any}).
The names are ours.

\item[$Y_{\inf}$.]
The common limit of the yin curve as $T\to\infty$ at fixed
fractional part $q=\{T\}$:
$Y_{\inf}(q)=1-\Psi(q)\,e^{-2\pi i(q^{2}-1/16)}$
(\S\ref{sec:yinyang-infinity}).
Yang follows the same curve half a unit behind.
For $0<T<1$ the yin path is $\zeta$ itself (Erickson's $C_0$).

\item[$\Psi$.]
The classical Riemann--Siegel correction
$\Psi(x)=\cos\bigl(2\pi(x^{2}-x-1/16)\bigr)/\cos(2\pi x)$,
as in Titchmarsh~\cite[\S15.3]{Titchmarsh1986} and
Pugh~\cite[\S5.4.1]{Pugh1998}.
It enters this paper as the amplitude of $Y_{\inf}$, not as a
remainder of its own.

\item[Needle.]
The reverse bisector link, viewed as a unit-length segment whose
ends ride the two teardrops
(\S\ref{sec:yinyang-area}).
Over one unit of $T$ it sweeps once around the stationary forward
bisector link.
The name is the Kakeya needle: the question is the area that
segment sweeps, about $2.0683$ in the limit.
\end{description}

\subsection*{The far chain}

\begin{description}
\item[Last link.]
The term $n_{\mathrm{last}}=\lfloor t/\pi\rfloor\approx 2T^{2}$,
far past the bisector, at which consecutive links are
antiparallel (\S\ref{sec:champion-bearings}).
Its heading, read in its own frame, carries a signature of the
champions.

\item[Follower.]
The link of the far chain (between the last two coils) that stays
parallel to the bisector link in the bisector frame.
Its link number is $(m+1)^{2}-1$
(\S\ref{sec:yinyang-links}, OEIS A005563).

\item[Champion.]
A height at which $\bigl|\zeta(\tfrac12+it)\bigr|$ exceeds every
value it has taken below that height: a running record of
$|\zeta|$ (\S\ref{sec:champion-bearings}).
\end{description}

\section*{Acknowledgments}
\begin{itemize}
  \item Chris Cowherd for all his extensive help over the last nine years in building many visualization and experimental mathematics tools, including the Zest zeta sum visualization tool
  \item Nicolas Fish for his four years of excellent programming skills in building and tweaking the Zest visualization and measurement tool
  \item Daniel Flores Galiote for his kind tutoring and helpful discussions over the last two years
  \item Jakob von Raumer for his help with the Lean proof
  \item Brian Conrey for his encouragement and acceptance of this amateur mathematician
  \item Freeman Dyson for his article {\it Birds and Frogs} \\(\texttt {\small https://pdodds.w3.uvm.edu/files/papers/others/2009/dyson2009a.pdf}), which inspired me to begin my journey to investigate one dimensional quasi-crystals and their possible connection to the Riemann Hypothesis nine years ago.
  \item My family for their support:
    \begin{itemize}
      \item my wife, Flo
      \item Sophie
      \item Anthony
      \item Thomas
    \end{itemize}
  \item Turingbot, Processing, Unity, OEIS, wolframalpha, ChatGPT, Claude (and Opus and Fable), Cursor, and last but not least: Desmos and LEAN (and Leo De Moura).  Thank you
\end{itemize}

\paragraph{Statement on the use of AI.}
Cursor and various AIs, including Anthropic's Claude (Opus and Fable),
OpenAI's ChatGPT (5.6 Sol), and Grok 4.5, were used to create the programs
that, when run, generated the figures for this paper; the name of the
program is listed in the caption of each figure created this way. AI was
also used to generate the Lean code from the author's proofs; that code was
checked both by the AI and by the author at
\url{https://live.lean-lang.org/}. Throughout, the AI was prompted in
detail by the author to complete, fill in, or explain proofs, and its
output was checked, edited, and verified by the author.

% ======================================================================
\appendix
% ======================================================================

\section{Lean formalization of \texorpdfstring{$R=R_{1ps}+R_{2ps}$}{R = R1ps + R2ps}}
\label{app:lean-decomp}

This appendix reproduces the Lean file formalizing the algebraic content of
Theorem~\ref{thm:decomp}. The file contains no \texttt{axiom} and no
\texttt{sorry}: every step is proved, and the only axioms the final theorem
depends on are Lean's own three (\texttt{propext},
\texttt{Classical.choice}, \texttt{Quot.sound}), as reported by
\texttt{\#print axioms}. It was checked against Lean~4.32.2 with Mathlib at
commit \texttt{905b95818e}. The Lean code was written with the assistance of AI
coding assistants (Anthropic's Claude).

\subsection{Correspondence with the written proof}

The file follows the proof of Theorem~\ref{thm:decomp} step for step. The
definitions \texttt{d\_j1} and \texttt{d\_j2} are the sine weights
\eqref{eq:d1}--\eqref{eq:d2}, and \texttt{R\_1ps} and \texttt{R\_2ps} are Step~1,
transcribed from \eqref{eq:R1ps-def}--\eqref{eq:R2ps-def} in the
$\lceil T\rceil^{\mp iI(T)}$ form rather than the $e^{\mp i\omega}$ form.
Two small bridges convert them: \texttt{ceil\_cpow} proves
$\lceil T\rceil^{\,ir}=e^{\,ir\log\lceil T\rceil}$ for $T>0$, which with
$\omega(T)=I(T)\log\lceil T\rceil$ (the definition \texttt{omg}) is exactly
$\lceil T\rceil^{\,iI(T)}=e^{i\omega}$; and
\texttt{div\_norm\_eq\_exp\_arg} proves $z/|z|=e^{i\arg z}$ for $z\neq0$, applied
to $\chi$ to give $\chi/|\chi|=e^{i\psi}$. Then
\texttt{factoring\_lemma} is the regrouping of Step~2, pulling out the common
factor $|R|/\sin(2\omega+\psi)$, and
\texttt{key\_algebraic\_identity} is the content of Steps~3--5, namely
\begin{equation}\label{eq:lean-key}
e^{-i\omega}\sin(\omega-\phi+\psi)+e^{\,i(\omega+\psi)}\sin(\omega+\phi)
\;=\;
e^{\,i\phi}\sin(2\omega+\psi);
\end{equation}
it is proved by writing each sine as
$(e^{i\theta}-e^{-i\theta})/2i$ (\texttt{sin\_to\_exp}, itself proved from
Mathlib's $e^{i\theta}=\cos\theta+i\sin\theta$), clearing the denominator,
merging the resulting products of exponentials into single exponentials, and
letting \texttt{ring\_nf} cancel the two copies of $e^{\,i(\psi-\phi)}$. The
closing lines are Step~6: cancel $\sin(2\omega+\psi)$ and recognize
$|R|e^{i\phi}=R$, which is Mathlib's
\texttt{Complex.norm\_mul\_exp\_arg\_mul\_I}.

\subsection{What is proved and what is hypothesized}

The theorem \texttt{R\_1ps\_plus\_R\_2ps\_eq\_R} carries three hypotheses. Two are
needed for the objects to exist at all: $\chi(\sigma+iI(T))\neq0$, without which
$\chi/|\chi|$ is meaningless, and $\sin(2\omega+\psi)\neq0$, without which $d_1$
and $d_2$ of \eqref{eq:d1}--\eqref{eq:d2} are undefined. The third, $T>0$, makes
$\lceil T\rceil$ a positive base, so that its complex power can be written as an
exponential. Nothing is assumed about whether $T$ is an integer.

One convention has to be read carefully. Steps~2--6, that is
\texttt{factoring\_lemma} and \texttt{key\_algebraic\_identity}, hold for
arbitrary real $\omega,\phi,\psi$ and make no reference to $T$; the only
$T$-dependent ingredient is the bridge \texttt{ceil\_cpow}, which holds for every
$T>0$ with $\omega$ defined as $I(T)\log\lceil T\rceil$. That is the paper's
$\omega$ whenever $T$ is not an integer, since then
$\lceil T\rceil=\lfloor T+1\rfloor$. At integer $T$ the two part company, and the
sign bookkeeping in the last paragraph of the proof of
Theorem~\ref{thm:decomp} turns on the factor $(-1)^{\lfloor T\rfloor}$ sitting
inside Siegel's $R$; that factor is invisible here, because $R$ enters as a free
complex parameter, so what the file proves at integer $T$ is the same identity
read with $\lceil T\rceil$ throughout.

What remains outside the formalization is not algebra but analysis. The
functional-equation factor $\chi$, the index map $I$, and Siegel's remainder $R$
enter as free parameters, so what is machine-checked is that any $R$ written in
polar form splits as $R_{1ps}+R_{2ps}$ with the weights
\eqref{eq:d1}--\eqref{eq:d2}, which is exactly Steps~1--6. That $R$ is Siegel's
remainder, and that $I$ is the map of \S\ref{sec:IT-origin}, are inputs to the
theorem rather than consequences of it.

\subsection{Source listing}

\lstinputlisting[style=leanstyle]{zeta_R_is_R1_plus_R2_proved.lean}

\section{Lean formalization of the critical line: \texorpdfstring{$d_1=d_2$, $|R_{1ps}|=|R_{2ps}|$, and equal legs}{d1 = d2, |R1ps| = |R2ps|, and equal legs}}
\label{app:lean-critical}

This appendix reproduces the Lean file establishing the three critical-line
statements of the paper: Corollary~\ref{cor:equal}
($d_1=d_2$ and $|R_{1ps}|=|R_{2ps}|$), Corollary~\ref{cor:equal-legs}
($L_1=L_2$), and Corollary~\ref{cor:bisector-proj} (the bisector point projects
onto $\tfrac{\zeta}{2}$). As in Appendix~\ref{app:lean-decomp} the file contains
no \texttt{axiom} and no \texttt{sorry}, and each result depends only on Lean's
own three axioms (\texttt{propext}, \texttt{Classical.choice},
\texttt{Quot.sound}) as reported by \texttt{\#print axioms}; it was checked
against Lean~4.32.2 with Mathlib at commit \texttt{905b95818e}. The Lean code was
written with the assistance of AI coding assistants (Anthropic's Claude).

\subsection{The chain of implications}

All three statements descend from a single fact, and the file is organized around
it.

\emph{Step 0: the two partial sums are conjugate.} At $\sigma=\tfrac12$ one has
$\overline{-s}=s-1$, so $n^{\,s-1}=\overline{n^{-s}}$ for every $n$; summing over
$n\le m$ gives $\Sigma_2=\chi\,\overline{\Sigma_1}$. This is
\texttt{partial\_sum\_reflect}, proved from Mathlib's \texttt{Complex.cpow\_conj}.

\emph{Step 1: the rotated remainder is real.} Write $\psi=\arg\chi$ and let
$u=e^{\,i\psi/2}$, a square root of $\chi$. With $R=\zeta-\Sigma_1-\Sigma_2$,
\begin{equation}\label{eq:lean-hinge}
e^{-i\psi/2}R
\;=\;
\underbrace{e^{-i\psi/2}\zeta}_{\text{real}}
\;-\;
\underbrace{\bigl(e^{-i\psi/2}\Sigma_1+e^{\,i\psi/2}\overline{\Sigma_1}\bigr)}_{=\;2\Re\left(e^{-i\psi/2}\Sigma_1\right)}
\;\in\;\mathbb{R},
\end{equation}
the first term being real because the functional equation at $\sigma=\tfrac12$
reads $\zeta=\chi\overline{\zeta}$, so that
$e^{-i\psi/2}\zeta=\overline{e^{-i\psi/2}\zeta}$. In the file this is
\texttt{rotated\_remainder\_self\_conj} (the algebra, from $u\bar u=1$ and
$u^2=\chi$) and \texttt{rotated\_remainder\_im\_eq\_zero} (its imaginary part
vanishes). It is the fact quoted as $\arg R=\tfrac12\arg\chi$ in
\S\ref{sec:d1-exact}. Stated as the reality of $e^{-i\psi/2}R$ it is exact;
stated through the argument it holds modulo $\pi$, which is visible in
$R=2d_1\cos(\omega+\tfrac\psi2)e^{\,i\psi/2}$, whose argument is
$\tfrac\psi2+\pi$ wherever the cosine is negative. The modulo-$\pi$ version is
all the sequel uses.

\emph{Step 2: $d_1=d_2$.} Since $\sin A-\sin B=2\sin\frac{A-B}{2}\cos\frac{A+B}{2}$,
the two numerators of \eqref{eq:d1}--\eqref{eq:d2} differ by
\begin{equation}\label{eq:lean-numerators}
\sin(\omega-\phi+\psi)-\sin(\omega+\phi)
\;=\;
2\,\sin\!\Bigl(\frac{\psi}{2}-\phi\Bigr)\cos\!\Bigl(\omega+\frac{\psi}{2}\Bigr),
\end{equation}
and Step~1 kills the first factor: \texttt{sin\_half\_sub\_arg\_eq\_zero} turns
the reality of $e^{-i\psi/2}R$ into $\sin(\tfrac\psi2-\arg R)=0$, using
$\cos\arg R=\Re R/|R|$ and $\sin\arg R=\Im R/|R|$. Hence
\texttt{d1\_eq\_d2\_of\_im\_eq\_zero} and its critical-line specialization
\texttt{d1\_eq\_d2}. Note that $\omega$ is an arbitrary real throughout: the
equality of the weights depends on the direction of $R$ relative to $\chi$ and
not at all on the turn angle.

\emph{Step 3: $|R_{1ps}|=|R_{2ps}|$.} Both phases are unimodular, so this is
$|d_1|=|d_2|$: \texttt{norm\_R1ps\_eq\_norm\_R2ps}. The sharper form
\eqref{eq:R2-conj-R1}, $R_{2ps}=\chi\,\overline{R_{1ps}}$, is
\texttt{R2ps\_eq\_chi\_conj\_R1ps}, and it is what feeds Step~4; it needs only
that $d_1$ is real and $d_1=d_2$.

\emph{Step 4: the legs.} Combining Steps~0 and~3,
$B_2=\Sigma_2+R_{2ps}=\chi\,\overline{\Sigma_1+R_{1ps}}=\chi\overline{B_1}$
(\texttt{leg2\_eq\_chi\_conj\_leg1}), and since $|\chi|=1$ on the line the two
legs have the same length (\texttt{legs\_norm\_eq}). Finally, from
$\zeta=B_1+\chi\overline{B_1}$,
\begin{equation}\label{eq:lean-proj}
e^{-i\psi/2}\zeta
=e^{-i\psi/2}B_1+\overline{e^{-i\psi/2}B_1}
=2\,\Re\bigl(e^{-i\psi/2}B_1\bigr),
\end{equation}
which says that the perpendicular projection of the bisector point onto the line
$\mathbb{R}\,e^{\,i\psi/2}$ carrying $\zeta$ lands exactly on $\tfrac{\zeta}{2}$:
the apex of the isosceles triangle sits over the midpoint of its base. This is
\texttt{proj\_eq\_half\_zeta}, and it is the dotted bisector line drawn in
Figures~\ref{fig:full-spirals} and~\ref{fig:legs}.

\subsection{What is proved and what is hypothesized}

Two analytic facts about $\chi$ on the critical line enter as hypotheses:
$|\chi|=1$, used in the form $\chi=e^{\,i\arg\chi}$
(\texttt{eq\_exp\_arg\_of\_norm\_one}), and the functional equation
$\zeta=\chi\overline{\zeta}$, which is $\zeta(s)=\chi(s)\zeta(1-s)$ together with
$1-s=\bar s$ at $\sigma=\tfrac12$. Everything else above is derived, including
the reflection $\Sigma_2=\chi\overline{\Sigma_1}$ of Step~0, which is proved
rather than assumed. The projection statement additionally takes
$\zeta=B_1+B_2$, which is Theorem~\ref{thm:decomp} of
Appendix~\ref{app:lean-decomp}.

The remainders are written here in the form $R_{1ps}=d_1e^{-i\omega}$,
$R_{2ps}=d_2e^{\,i(\omega+\psi)}$ of \eqref{eq:zeta-ps} rather than the
$\lceil T\rceil^{\mp iI(T)}$ form of
\eqref{eq:R1ps-def}--\eqref{eq:R2ps-def}; the bridge between them is
\texttt{ceil\_cpow} of Appendix~\ref{app:lean-decomp}. As there, $\zeta$, $\chi$,
$\Sigma_1$ and $\omega$ are free parameters, so what is checked is the geometry
of the critical line and not the identity of the functions involved.

\subsection{Source listing}

\lstinputlisting[style=leanstyle]{zeta_critical_line.lean}

\begin{thebibliography}{9}

\bibitem{BerryGoldberg1988}
M.~V.\ Berry and J.~Goldberg, \emph{Renormalisation of curlicues},
Nonlinearity \textbf{1} (1988), 1--26.

\bibitem{Bourgain2017}
J.~Bourgain, \emph{Decoupling, exponential sums and the Riemann zeta function},
J.\ Amer.\ Math.\ Soc.\ \textbf{30} (2017), no.~1, 205--224.

\bibitem{CoutsiasKazarinoff1987}
E.~A.\ Coutsias and N.~D.\ Kazarinoff, \emph{Disorder, renormalizability, theta
functions and Cornu spirals}, Physica D \textbf{26} (1987), 295--310.

\bibitem{CoutsiasKazarinoff1998}
E.~A.\ Coutsias and N.~D.\ Kazarinoff, \emph{The approximate functional formula
for the theta function and Diophantine Gauss sums}, Trans.\ Amer.\ Math.\ Soc.\
\textbf{350} (1998), no.~2, 615--641.

\bibitem{CsordasVarga1990}
George Csordas and Richard S.\ Varga, \emph{Necessary and sufficient conditions
and the Riemann hypothesis}, Adv.\ Appl.\ Math.\ \textbf{11} (1990), no.~3,
328--357.

\bibitem{DenavitHartenberg1955}
J.~Denavit and R.~S.\ Hartenberg, \emph{A kinematic notation for lower-pair
mechanisms based on matrices}, J.\ Appl.\ Mech.\ \textbf{22} (1955), 215--221.

\bibitem{Erickson2005}
Carl Erickson, \emph{A geometric perspective on the Riemann zeta function's
partial sums}, Stanford Undergraduate Research Journal (2005).

\bibitem{Gabcke1979}
Wolfgang Gabcke, \emph{Neue Herleitung und explizite Restabsch\"atzung der
Riemann--Siegel-Formel}, Dissertation, Georg-August-Universit\"at
G\"ottingen (1979).

\bibitem{Kapitonets2019}
Kirill Kapitonets, \emph{Analysis of the Riemann zeta function},
arXiv:1910.08363 (2019).

\bibitem{Kuznetsov2025}
Alexey Kuznetsov, \emph{Simple and accurate approximations to the Riemann
zeta function}, J.\ Comput.\ Appl.\ Math.\ \textbf{488} (2026), 117791.
Also available as arXiv:2503.09519 (2025).

\bibitem{Lander2017}
Anthony D.\ Lander, \emph{Two finite mirror-image series restrict the
non-trivial zeros of Riemann's zeta function to $\Re(s)=\tfrac12$ and the zeros
of its derivative to $\Re(s)>\tfrac12$}, Preprints 2017030121 (2017),
doi:10.20944/preprints201703.0121.v2.

\bibitem{Lander2018}
Anthony Lander, \emph{Symmetry and the Zeros of Riemann's Zeta Function},
CreateSpace Independent Publishing Platform, 2018.
ISBN 978-1-9860-7414-8.

\bibitem{Lehmer1956}
D.~H.\ Lehmer, \emph{On the roots of the Riemann zeta-function},
Acta Math.\ \textbf{95} (1956), 291--298.

\bibitem{Lindelof1908}
E.~Lindel\"of, \emph{Quelques remarques sur la croissance de la fonction
$\zeta(s)$}, Bull.\ Sci.\ Math.\ \textbf{32} (1908), 341--356.

\bibitem{Nickel2013}
George H.\ Nickel, \emph{Geometry of the Riemann Zeta Function},
arXiv:1310.6396 (2013).

\bibitem{Nickel2015}
George H.\ Nickel, \emph{Symmetry in Partial Sums of $n^{-s}$ (or, Critical Line
Zeros of $\zeta$ Are Easy)}, arXiv:1507.07631 (2015).

\bibitem{OEIS}
OEIS Foundation Inc., \emph{The On-Line Encyclopedia of Integer Sequences},
sequence A005563, \texttt{oeis.org/A005563}.

\bibitem{Pugh1998}
G.~R.\ Pugh, \emph{The Riemann--Siegel formula and large scale computations
of the Riemann zeta function}, M.Sc.\ thesis, University of British
Columbia, 1998.

\bibitem{Siegel1932}
C.~L.\ Siegel, \emph{\"Uber Riemanns Nachla\ss{} zur analytischen
Zahlentheorie}, Quellen Stud.\ Gesch.\ Math.\ Astron.\ Phys.\ Abt.\ B:
Studien \textbf{2} (1932), 45--80. Reprinted in \emph{Gesammelte
Abhandlungen}, Vol.~1, Springer, Berlin (1966), 275--310.

\bibitem{BarkanSklar2018}
E.~Barkan and D.~Sklar, \emph{On Riemann's Nachlass for Analytic Number
Theory: a translation of Siegel's \"Uber Riemanns Nachla\ss{} zur
analytischen Zahlentheorie}, arXiv:1810.05198 (2018).

\bibitem{Titchmarsh1986}
E.~C.\ Titchmarsh, \emph{The Theory of the Riemann Zeta-Function}, 2nd ed.,
revised by D.~R.\ Heath-Brown, Oxford University Press, Oxford, 1986.

\end{thebibliography}

\end{document}
